% Generated by build_chapter.py from the GY1967 files. Do not edit by hand; edit the files and rebuild.
\chapter{The academic years 1967--68 and 1968--69}\label{ch:1967}
These two academic years comprise 43 catalogue records and 320 mathematical
occurrences. The entries follow the notebook rather than numerical record
order: Oleinik's record 164 precedes Kazhdan's 163, and Gelfand's 170 precedes
Davis's 169. Catalogue page ranges are retained for identification; each
occurrence has its own narrower source attribution. The supplied scans cover
P000208--P000284, but the explicit dates on the last two pages begin the next
academic year. Their contents are inventoried in the report and are not
assigned to Buchstaber's April continuation.

\begin{center}
\begin{tabular}{@{}clrr@{}}
\toprule
Class & Interpretation & Items & Share \\
\midrule
A & Speaker/Gelfand research published soon afterward & 42 & 13.1\% \\
B & Subsequent development by others & 8 & 2.5\% \\
C & Partial results; a stated residual question & 3 & 0.9\% \\
D & No development established after an adequate search & 0 & 0.0\% \\
E & Known material or a result already available & 239 & 74.7\% \\
F & Exact assertion or attribution not identifiable & 28 & 8.8\% \\
\bottomrule
\end{tabular}
\end{center}

The main research sequences concern classifying spaces for noncompact Lie
groups, character formulas, symbolic dynamics, and differential relations.
Fuchs's three lectures are closely matched by the Gelfand--Fuchs papers of
1967--1968. Their unconventional topology and restricted category of bundles
matter: replacing them by the modern numerable-bundle classification would
lose the actual question. Kirillov's compact character calculation and kernel
trace lemma have exact counterparts in his 1968 paper; the notebook's generic
orbit formulas do not acquire unrestricted validity from those special cases.
Sinai's geometric Markov partitions connect to his 1968 publications, whereas
the two-torus entropy classification and its particular recoding were already
proved by Adler and Weiss in 1967 \citep{gelfandfuchs1967topology,
kirillov1968characters,sinai1968markov,adlerweiss1967entropy}.

A substantial part of the second year is sustained exposition. Bernstein and
Geronimus develop Grothendieck sites and Mumford's moduli viewpoint; Palamodov
explains jets and Spencer sequences. Kazhdan treats Ihara and Dwork; Gromov
passes from formal differential relations to global solutions under essential
openness, invariance and open-manifold hypotheses. Oleinik's presentation of
Radkevich, Bleher on spectral asymptotics, Davis on operator shells, and
Buchstaber on homotopy tori bring analytic and topological methods into the
same seminar. A programme or bare name remains a programme or bare name:
records 141, 148, 158 and 165 contain no separately identifiable mathematical
occurrences.

The critical findings frequently change the truth or scope without changing
the publication match. The fundamental-function interpolation criterion
reported by Krein is undermined by Sharpley's counterexamples; Boyd's dilation
indices supply the relevant distinction. The unextended semisimple
Gelfand--Kirillov conjecture is false in general, despite the positive
finite-centre-extension theorem. The torus proof needs
\(\operatorname{Wh}(\mathbb Z^n)=0\), not \(K_1=0\). The spectral-function
notes confuse a zero normalized leading term with a zero kernel, and omit
normalization factors in spherical multiplicities and averaging. These are
mathematical emendations when the scan confirms the wording, not alleged
transcription mistakes \citep{sharpley1975interpolation,premet2010modular,
basshellerswan1964whitehead,hormander1968spectral}.

The eleven problem entries include resolved historical questions. Two exact
recurrences retain earlier identifiers: GP-1966-01, algebraic classification
of Anosov diffeomorphisms, and GP-1965-08, automorphisms of the first Weyl
division ring. The three C occurrences have different scopes. The
Anosov and Weyl questions remain unresolved in the audited literature; the
hyperbolic-circle question is a qualified reconstruction requiring all small
spectral terms to be subtracted and allowing an epsilon loss. The moment-representation request
is F: its intended category is not uniquely specified, and partial
representation theorems do not certify a particular open conjecture. It is
therefore excluded from the problem index. No D classification is warranted.

One bounded development is included: a complete rederivation of the known
centralizer and rational-coordinate-field stabilizer in the first Weyl
division ring, including explicit inverses and the exact polynomial
restriction. It is not presented as new research or a classification of the
whole group. An independent algebraic review and 522 exact symbolic checks
support its formulas. The other C items do not yield a comparably justified
bounded contribution in this run. Throughout, ``verified'' in the bibliography
means that an identifying source was opened; the separate item field records
whether the result, metadata, or only an indirect account was actually read.

\GStalk{131}{1967-09-18}{A. A. Kirillov}{P000208-P000210}{Local orbit character formula and the rotation group}
The actual Kirillov heading occupies P000208 and the beginning of P000209. It announces an orbit formula for irreducible characters, explicitly without proof, and tests it on \(SO(3)\). The notebook also lists future talks on logic, the continuum hypothesis, Lidskii on eigenvalue asymptotics, and work of Kirillov, Kazhdan, and Graev. The mathematical continuation is Kirillov's 1968 character paper, received on 18 October 1967; its proved cases must be distinguished from the unrestricted announcement.

\begin{GSitem}{1967-131.1}{A}{Local orbit formula for characters}
\GSfield{Kind; source}{method; \GSpage{P000208}}
\GSfield{The notes}{The heading says \ru{без доказательства или противоречащих примеров}. Near the identity it proposes \(\chi(x)=p(x)^{-1/2}\int_\Omega e^{i\langle F,x\rangle}\beta(F)\) for irreducible representations of Lie groups.}
\GSfield{Statement}{For a connected simply connected nilpotent Lie group, the irreducible unitary representation associated with a coadjoint orbit has distribution character equal, in exponential coordinates and with compatible Fourier conventions, to the Fourier transform of the canonical orbital measure. For compact connected groups the corresponding local formula uses the orbit through the highest weight plus \(\rho\), and the square root of the exponential Jacobian. Neither statement is an unrestricted formula for every irreducible representation of every Lie group.}
\GSfield{Reconstruction}{The two precise proved settings replace the notebook's unrestricted scope. No general orbit-method conjecture is silently converted into a theorem.}
\GSfield{Fate}{\citet[Introduction and the nilpotent and compact cases]{kirillov1968characters} develops the announced formula; the manuscript was received 18 October 1967 and published in 1968. The nilpotent orbit correspondence was already available in \citet{kirillov1962unitary}. The publication match is partial because the notebook suppresses the admissibility, normalization, and representation-class restrictions.}
\GSfield{Present status}{As of November 2026, these nilpotent and compact character formulas are established. The universal formulation in the notes is too broad to have a single affirmative status.}
\GSfield{Assessment}{literal statement: correct with missing hypotheses; publication match: partial; verification: read the result; development: annotated; confidence: high}
\end{GSitem}
\begin{GSitem}{1967-131.2}{E}{Distribution character and the missing trace}
\GSfield{Kind; source}{definition; \GSpage{P000208}}
\GSfield{The notes}{For a finite-dimensional representation the notes give \(\chi(g)=\operatorname{tr}T(g)\). In the general case they write \(\langle\chi,\varphi\rangle=\int_G\varphi(g)T(g)\,dg\), adding that the latter already has a trace.}
\GSfield{Statement}{For a representation for which \(T(\varphi)=\int_G\varphi(g)T(g)\,dg\) is trace class for every \(\varphi\in C_c^\infty(G)\), its distribution character is \(\Theta_T(\varphi)=\operatorname{tr}T(\varphi)\). The integral itself is an operator, not a scalar distribution value; trace-class admissibility is an additional hypothesis.}
\GSfield{Reconstruction}{The trace-class condition and the outer trace are explicit editorial repairs.}
\GSfield{Fate}{This is the character convention used in \citet[Introduction]{kirillov1968characters}. The notebook has omitted the outer trace while preserving an explanatory phrase that makes the intended correction clear.}
\GSfield{Present status}{As of November 2026, the qualified statement above is established; the distinctions from the notebook wording remain essential.}
\GSfield{Assessment}{literal statement: false as written; publication match: partial; verification: read the result; development: annotated; confidence: high}
\end{GSitem}
\begin{GSitem}{1967-131.3}{E}{Canonical form and measure on a coadjoint orbit}
\GSfield{Kind; source}{construction; \GSpage{P000208-P000209}}
\GSfield{The notes}{The form \(\beta\) is said to be constructed from the structure constants; the final line mentions integrals of integral forms over cycles of the orbit.}
\GSfield{Statement}{At \(F\in\mathfrak g^*\), the Kirillov form on the coadjoint orbit is \(\omega_F(\operatorname{ad}_X^*F,\operatorname{ad}_Y^*F)=\langle F,[X,Y]\rangle\), up to the sign convention for the coadjoint action. It is closed and nondegenerate. On an orbit of dimension \(2d\), the corresponding Liouville measure is a fixed Fourier-normalized multiple of \(\omega^d/d!\).}
\GSfield{Reconstruction}{The symplectic formula makes precise the otherwise unspecified construction from structure constants; the sign and Fourier normalization are deliberately not inferred from the notebook.}
\GSfield{Fate}{The canonical orbital measure is part of the framework in \citet[Introduction]{kirillov1968characters}, extending the earlier nilpotent setting of \citet{kirillov1962unitary}. The notebook's cycle phrase does not give a separate computable formula for \(\beta\).}
\GSfield{Present status}{As of November 2026, the qualified statement above is established; the distinctions from the notebook wording remain essential.}
\GSfield{Assessment}{literal statement: correct with missing hypotheses; publication match: exact; verification: read the result; development: annotated; confidence: high}
\end{GSitem}
\begin{GSitem}{1967-131.4}{E}{Integrality and the shifted compact orbit}
\GSfield{Kind; source}{remark; \GSpage{P000208-P000209}}
\GSfield{The notes}{Only certain orbits satisfying integrality conditions are to occur. The rotation example selects \(\lambda=2n+1\).}
\GSfield{Statement}{Prequantization of a symplectic orbit requires \([\omega/(2\pi)]\) to be integral. Equivariance under a specified group can impose a stronger lattice condition. In the compact character formula the orbit is shifted by \(\rho\); therefore an unshifted assertion that every integral orbit directly labels a character is insufficient. For \(SO(3)\), the representations have integral spin \(n\geq0\), while the character formula below contains \(2n+1\).}
\GSfield{Reconstruction}{Prequantization, group equivariance, and the \(\rho\)-shift are separated here rather than merged into the ambiguous word integrality.}
\GSfield{Fate}{\citet[compact-group discussion]{kirillov1968characters} supplies the compact character setting. The notes announce an integrality requirement but do not state which lattice or shift is intended.}
\GSfield{Present status}{As of November 2026, the qualified statement above is established; the distinctions from the notebook wording remain essential.}
\GSfield{Assessment}{literal statement: correct with missing hypotheses; publication match: partial; verification: read the result; development: annotated; confidence: high}
\end{GSitem}
\begin{GSitem}{1967-131.5}{E}{Exponential Haar density for the rotation example}
\GSfield{Kind; source}{claim; \GSpage{P000208}}
\GSfield{The notes}{The notebook visibly writes \(p(x)=\sin\tau/\tau\), without a square, and then applies \(p^{-1/2}\).}
\GSfield{Statement}{If \(\tau\) denotes half the rotation angle, the exponential-coordinate Haar Jacobian, normalized at zero, is \(j(x)=(\sin\tau/\tau)^2\). If \(\theta\) is the full angle, it is \((\sin(\theta/2)/(\theta/2))^2\). Its square root, rather than the density itself, is \(\sin\tau/\tau\) near zero.}
\GSfield{Reconstruction}{The half-angle convention is chosen to make the subsequent character formula consistent; the notebook does not explicitly specify this normalization.}
\GSfield{Fate}{The determinant factor in the compact case of \citet{kirillov1968characters} and the elementary eigenvalues \(0,\pm 2i\tau\) of \(\operatorname{ad}_x\) identify the missing square. The displayed formula is faithfully transcribed, not an OCR mistake.}
\GSfield{Present status}{As of November 2026, the qualified statement above is established; the distinctions from the notebook wording remain essential.}
\GSfield{Assessment}{literal statement: false as written; publication match: partial; verification: read the result; development: annotated; confidence: high}
\end{GSitem}
\begin{GSitem}{1967-131.6}{E}{Fourier transform of the spherical orbital measure}
\GSfield{Kind; source}{claim; \GSpage{P000208-P000209}}
\GSfield{The notes}{The notes give \((\delta_{S_\lambda})^\vee=\sin(\lambda\tau)/\tau\), then call \(\beta\) ordinary surface area.}
\GSfield{Statement}{With Fourier kernel \(e^{i\langle F,x\rangle}\) and Euclidean area \(dS\) on the sphere of radius \(\lambda>0\), \(\int_{S_\lambda}e^{i\langle F,x\rangle}dS=4\pi\lambda\sin(\lambda|x|)/|x|\). The measure \(dS/(4\pi\lambda)\) gives the notebook's transform and has mass \(\lambda\).}
\GSfield{Reconstruction}{The calculation is an editorial normalization check, not a new character theorem.}
\GSfield{Fate}{This direct spherical-coordinate calculation fixes the normalization required by the compact orbit formula of \citet{kirillov1968characters}. The words ordinary surface cannot be read literally together with the displayed equality under this Fourier convention.}
\GSfield{Present status}{As of November 2026, the qualified statement above is established; the distinctions from the notebook wording remain essential.}
\GSfield{Assessment}{literal statement: false as written; publication match: partial; verification: read the result; development: annotated; confidence: high}
\end{GSitem}
\begin{GSitem}{1967-131.7}{E}{The irreducible characters of SO(3)}
\GSfield{Kind; source}{theorem; \GSpage{P000209}}
\GSfield{The notes}{The answer is \(\chi(x)=\sin(\lambda\tau)/\sin\tau\), for \(\lambda=2n+1\).}
\GSfield{Statement}{For the irreducible \((2n+1)\)-dimensional representation of \(SO(3)\), a rotation through angle \(\theta\) has character \(\chi_n(\theta)=\sin((n+\tfrac12)\theta)/\sin(\theta/2)\), interpreted continuously at multiples of \(2\pi\). Setting \(\tau=\theta/2\) gives exactly the recorded expression and \(\chi_n(e)=2n+1\).}
\GSfield{Reconstruction}{The parameter is explicitly the half-angle, and the removable singularities are filled in.}
\GSfield{Fate}{This is the rank-one compact character formula, used as the test case for \citet[compact case]{kirillov1968characters}. It was known before the lecture, although its orbit interpretation is the point of the new presentation.}
\GSfield{Present status}{As of November 2026, the qualified statement above is established; the distinctions from the notebook wording remain essential.}
\GSfield{Assessment}{literal statement: correct with missing hypotheses; publication match: exact; verification: read the result; development: annotated; confidence: high}
\end{GSitem}

\GStalk{132}{1967-09-18}{D. B. Fuchs}{P000208-P000213}{Classifying noncompact groups over Hausdorff bases}
The Fuchs heading begins midway through P000209 and ends on P000212 before the date 25 September. The discussion changes the category of bases and bundle charts so that noncompact groups retain characteristic information invisible to ordinary homotopy classification. The nonregular ray quotient and the failure of a cone-based definition are central. The exact research continuation is \citet{gelfandfuchs1967topology}, received 15 July 1967; parts of the programme were therefore already circulating before this lecture.

\begin{GSitem}{1967-132.1}{E}{Ordinary universal principal bundles and characteristic classes}
\GSfield{Kind; source}{definition; \GSpage{P000209}}
\GSfield{The notes}{\(E_G\to B_G\) is called universal; \(\pi_k(E_G)=0\) for \(k>0\). Characteristic classes are \(H^q(B_G;A)\), and more generally \(h(B_G)\) for a contravariant homotopy functor.}
\GSfield{Statement}{For numerable principal bundles over paracompact Hausdorff bases, a contractible free universal \(G\)-space classifies bundles by homotopy classes \([X,BG]\). A class \(a\in H^q(BG;A)\) gives \(c_a(P)=f_P^*a\). The analogous statement for a contravariant homotopy functor uses the universal pullback class and its naturality. Vanishing higher homotopy groups alone must be supplemented by connectedness and suitable CW-type hypotheses to infer contractibility.}
\GSfield{Reconstruction}{Numerability, paracompactness, and the distinction between weak and actual contractibility are supplied.}
\GSfield{Fate}{This is the conventional background contrasted with the new category in \citet[Introduction]{gelfandfuchs1967topology}. The notebook uses aspherical for weak contractibility, rather than the modern condition of vanishing \(\pi_k\) only for \(k>1\).}
\GSfield{Present status}{As of November 2026, the qualified statement above is established; the distinctions from the notebook wording remain essential.}
\GSfield{Assessment}{literal statement: correct with missing hypotheses; publication match: partial; verification: read the result; development: annotated; confidence: high}
\end{GSitem}
\begin{GSitem}{1967-132.2}{E}{Loss of noncompact information under homotopy equivalence}
\GSfield{Kind; source}{theorem; \GSpage{P000209}}
\GSfield{The notes}{A maximal compact subgroup is said to be homotopy equivalent to the group, and their classifying spaces homotopy equivalent.}
\GSfield{Statement}{If \(G\) is a Lie group with finitely many connected components and \(K\) a maximal compact subgroup, \(G/K\) is contractible and \(BK\to BG\) is a homotopy equivalence in the usual classifying-space theory. This assertion is not a theorem for an arbitrary topological group merely possessing some compact subgroup.}
\GSfield{Reconstruction}{The Lie-group and finite-component hypotheses are restored.}
\GSfield{Fate}{\citet[Introduction]{gelfandfuchs1967topology} explicitly starts from this loss of information and motivates the new characteristic theory.}
\GSfield{Present status}{As of November 2026, the qualified statement above is established; the distinctions from the notebook wording remain essential.}
\GSfield{Assessment}{literal statement: correct with missing hypotheses; publication match: exact; verification: read the result; development: annotated; confidence: high}
\end{GSitem}
\begin{GSitem}{1967-132.3}{A}{A classifying theory retaining the noncompact group}
\GSfield{Kind; source}{problem; \GSpage{P000209}; I. M. Gelfand}
\GSfield{The notes}{Gelfand asks for the actual classifying space of a noncompact group, without identifying it with its maximal compact subgroup; bases should be Hausdorff, with a new local-triviality requirement agreeing with the usual one on completely regular bases.}
\GSfield{Statement}{For a closed linear Lie group \(G\), construct a category of principal \(G\)-bundles over Hausdorff spaces, together with sufficiently large universal objects, such that ordinary locally trivial bundles over completely regular bases are included and natural characteristic operations retain information about \(G\), not merely the homotopy type of \(BG\).}
\GSfield{Reconstruction}{The informal request is made precise by the definitions and theorems in the three consecutive Fuchs talks.}
\GSfield{Fate}{\citet[\S\S1--2 and 4]{gelfandfuchs1967topology} supplies the intended category, universal frame spaces, cohomological calculation, and classification of natural operations. The scope is a specified category and an arbitrarily large index set; it is not one ordinary CW space classifying all such bases.}
\GSfield{Present status}{As of November 2026, this precise linear-group programme is answered by the 1967 paper, with a more general elementary-object formulation in 1968. The broad phrase actual classifying space alone does not define an additional open problem.}
\GSfield{Assessment}{literal statement: correct with missing hypotheses; publication match: exact; verification: read the result; development: annotated; confidence: high; problem \hyperref[prob:GP-1967-01]{GP-1967-01}}
\end{GSitem}
\begin{GSitem}{1967-132.4}{A}{The rejected cone test and Kazhdan's obstruction}
\GSfield{Kind; source}{method; \GSpage{P000210}; D. A. Kazhdan}
\GSfield{The notes}{For each \(x\), an equivariant map \(E\to CG\) is required to identify the fiber with the cone base. Kazhdan notes that cone height descends to a continuous real function on \(X\); the definition is rejected.}
\GSfield{Statement}{An equivariant map to \(CG\), nonzero on a selected fiber, yields a scalar height function on the base and a local equivariant chart on its nonzero set. Over a completely regular base a local chart can be multiplied by a continuous cutoff to produce such a map. Requiring this construction over arbitrary Hausdorff bases introduces scalar separation unavailable on the ray quotient considered next.}
\GSfield{Reconstruction}{The obstruction explains why the criterion fails for the intended enlarged category. It does not assert that the cone criterion is logically inconsistent or fails for ordinary numerable bundles.}
\GSfield{Fate}{The successful frame-space replacement is given by \citet[\S1]{gelfandfuchs1967topology}; the exact cone objection is recorded in the seminar rather than promoted there to a separately numbered theorem.}
\GSfield{Present status}{As of November 2026, the qualified statement above is established; the distinctions from the notebook wording remain essential.}
\GSfield{Assessment}{literal statement: correct with missing hypotheses; publication match: partial; verification: read the result; development: annotated; confidence: high}
\end{GSitem}
\begin{GSitem}{1967-132.5}{A}{Global fiberwise linear charts over Hausdorff bases}
\GSfield{Kind; source}{definition; \GSpage{P000210-P000211}}
\GSfield{The notes}{For \(G\subset GL(n,\mathbb R)\), form \(E^\#=E\times_G\mathbb R^n\). For every prescribed base point there must be a continuous map \(E^\#\to\mathbb R^n\), linear on all fibers and invertible on the selected fiber.}
\GSfield{Statement}{Use the quotient topology on the associated vector bundle and require global continuous fiberwise linear maps that are isomorphisms on each individually chosen fiber. No support restriction is imposed. This condition implies the usual local trivialization around that fiber; it agrees with ordinary local triviality over completely regular bases. Over merely Hausdorff bases it is an additional specified condition, not a consequence of local triviality alone.}
\GSfield{Reconstruction}{The quantifier for every chosen fiber and the distinction from arbitrary ordinary Hausdorff bundles are made explicit.}
\GSfield{Fate}{This is the principal-bundle category of \citet[\S1, definition preceding Proposition 1.1]{gelfandfuchs1967topology}. Its deliberate use of global linear maps replaces the unsuccessful scalar cutoff construction.}
\GSfield{Present status}{As of November 2026, the qualified statement above is established; the distinctions from the notebook wording remain essential.}
\GSfield{Assessment}{literal statement: correct with missing hypotheses; publication match: exact; verification: read the result; development: annotated; confidence: high}
\end{GSitem}
\begin{GSitem}{1967-132.6}{E}{Recovering the principal bundle requires its reduction}
\GSfield{Kind; source}{claim; \GSpage{P000211}}
\GSfield{The notes}{The first line says that \(E\) is recovered from \(E^\#\).}
\GSfield{Statement}{The full frame bundle of a rank-\(n\) vector bundle recovers a principal \(GL(n,\mathbb R)\)-bundle. For a proper closed subgroup \(G\), recovering the original principal \(G\)-bundle requires retaining its \(G\)-reduction inside that frame bundle. The underlying associated vector bundle alone generally forgets that reduction.}
\GSfield{Reconstruction}{The missing reduction is an elementary categorical correction, not a claim that the published construction fails.}
\GSfield{Fate}{The frame construction in \citet[\S1]{gelfandfuchs1967topology} retains the group action. It does not justify an inverse from bare vector bundles for every subgroup \(G\).}
\GSfield{Present status}{As of November 2026, the qualified statement above is established; the distinctions from the notebook wording remain essential.}
\GSfield{Assessment}{literal statement: false as written; publication match: partial; verification: read the result; development: annotated; confidence: high}
\end{GSitem}
\begin{GSitem}{1967-132.7}{A}{The Hausdorff ray quotient with no nonconstant scalar functions}
\GSfield{Kind; source}{construction; \GSpage{P000211-P000212}}
\GSfield{The notes}{For the additive group \(\mathbb R\) acting by positive homotheties, the total space is written \(\mathbb R^\infty\); the quotient \(S\) is Hausdorff but every continuous real function on it is constant. A cylinder-neighborhood argument is sketched.}
\GSfield{Statement}{Let \(E=\mathbb R^{\mathbb N}\setminus\{0\}\) have the product topology and let \(t\in\mathbb R\) act by \(x\mapsto e^t x\). The quotient of positive rays is Hausdorff and every continuous map \(S=E/\mathbb R\to\mathbb R\) is constant. The quotient bundle meets the specified global linear-chart condition. It has no continuous section: the coordinates of such a section would all be constant, making its image a single point.}
\GSfield{Reconstruction}{The zero vector is removed and the additive-to-multiplicative identification is specified.}
\GSfield{Fate}{The example and the cylinder proof occur in \citet[Introduction, pp.33--35 in the Russian original]{gelfandfuchs1967topology}. Deleting zero is indispensable to the free action and is explicit in the publication.}
\GSfield{Present status}{As of November 2026, the qualified statement above is established; the distinctions from the notebook wording remain essential.}
\GSfield{Assessment}{literal statement: correct with missing hypotheses; publication match: exact; verification: read the result; development: annotated; confidence: high}
\end{GSitem}
\begin{GSitem}{1967-132.8}{E}{The proposed weighted cross-section is not a section of all rays}
\GSfield{Kind; source}{claim; \GSpage{P000211}}
\GSfield{The notes}{The notebook suggests examining the purported section \(\sum_{n\ge1}2^{-n}|x_n|/(1+|x_n|)=1/2\), to see that the resulting map is discontinuous.}
\GSfield{Statement}{This level set does not meet every nonzero positive ray. On a ray supported only at coordinate \(m\ge2\), the left side is always less than \(2^{-m}\le1/4\); for a ray supported only at coordinate 1 it tends to \(1/2\) but never reaches it. Consequently it cannot first define a global set-theoretic selector whose continuity is then tested.}
\GSfield{Reconstruction}{The one-coordinate test is an editorial check. The literal denominator \(1+|x_n|\) and level \(1/2\) were checked against the scan.}
\GSfield{Fate}{The no-section theorem is independently proved in \citet[Introduction]{gelfandfuchs1967topology}; the specific displayed selector is not needed. The failure is an immediate calculation from the handwritten formula.}
\GSfield{Present status}{As of November 2026, the qualified statement above is established; the distinctions from the notebook wording remain essential.}
\GSfield{Assessment}{literal statement: false as written; publication match: none found; verification: indirect; development: annotated; confidence: high}
\end{GSitem}
\begin{GSitem}{1967-132.9}{A}{Universal targets indexed by sufficiently large products}
\GSfield{Kind; source}{theorem; \GSpage{P000212}}
\GSfield{The notes}{Every bundle maps to the ray-bundle model built from a product of copies of \(\mathbb R\) of sufficiently large cardinality; the map need not be unique.}
\GSfield{Statement}{In the specified category, collecting enough global fiberwise linear charts yields an equivariant map into a frame-space universal bundle indexed by a set \(I\) depending on the given bundle. In rank one this gives the positive-ray construction. Universality means existence of such a bundle morphism and pullback representation; it is not uniqueness of an actual map.}
\GSfield{Fate}{\citet[Proposition 1.1]{gelfandfuchs1967topology} proves the construction. The variable cardinality is a substantive part of the statement, developed explicitly at the next lecture.}
\GSfield{Present status}{As of November 2026, the qualified statement above is established; the distinctions from the notebook wording remain essential.}
\GSfield{Assessment}{literal statement: correct; publication match: exact; verification: read the result; development: none; confidence: high; problem \hyperref[prob:GP-1967-01]{GP-1967-01}}
\end{GSitem}

\GStalk{133}{1967-09-25}{I. M. Gelfand}{P000211-P000214}{Feynman parameters and commuting interval maps}
The date 25 September appears on P000212, after the first Fuchs lecture. Two Feynman-parameter identities are followed by news about commuting continuous interval maps and a differentiable periodic-point result. The three-denominator identity lacks a factorial; the interval remark refers to the counterexample work of Boyce and Huneke and to the earlier periodic-point theorem of Schwartz. Fuchs resumes under his own heading on P000213.

\begin{GSitem}{1967-133.1}{E}{Two-denominator Feynman parametrization}
\GSfield{Kind; source}{method; \GSpage{P000212}}
\GSfield{The notes}{\(1/(ab)=\int_0^1(at+b(1-t))^{-2}\,dt\), described as a good formula of Feynman and a partial-fraction decomposition.}
\GSfield{Statement}{For positive real \(a,b\), \[\frac1{ab}=\int_0^1\frac{dt}{(ta+(1-t)b)^2}.\] The same equality holds for nonzero complex \(a,b\) when the segment \(ta+(1-t)b\) avoids zero. At \(a=b\) the integrand is constant; otherwise integrate its elementary antiderivative.}
\GSfield{Reconstruction}{The denominator-nonvanishing hypothesis and equal-parameter case are supplied.}
\GSfield{Fate}{This is the exponent-one instance of the beta-integral denominator identity; see \citet[\S5.12, Euler beta integral]{nist2026beta}. The seminar attribution names Feynman's useful parametrization, not a new result proved in 1967.}
\GSfield{Present status}{As of November 2026, the qualified statement above is established; the distinctions from the notebook wording remain essential.}
\GSfield{Assessment}{literal statement: correct with missing hypotheses; publication match: exact; verification: read the result; development: annotated; confidence: high}
\end{GSitem}
\begin{GSitem}{1967-133.2}{E}{Three denominators and the missing factorial}
\GSfield{Kind; source}{claim; \GSpage{P000212}}
\GSfield{The notes}{\(1/(abc)=\int_\Delta(at_1+bt_2+ct_3)^{-3}\,dt_1dt_2\), where \(t_i\ge0\) and \(t_1+t_2+t_3=1\).}
\GSfield{Statement}{With ordinary coordinate measure \(dt_1dt_2\) on \(\Delta=\{t_1,t_2\ge0:t_1+t_2\le1\}\), and \(a,b,c>0\), the correct identity is \[\frac1{abc}=2\int_\Delta\frac{dt_1dt_2}{(at_1+bt_2+c(1-t_1-t_2))^3}.\] More generally the \(m\)-denominator numerator factor is \((m-1)!\) for unit powers.}
\GSfield{Reconstruction}{The notebook explicitly uses \(dt_1dt_2\), so normalized probability measure on the simplex is not silently substituted.}
\GSfield{Fate}{The factor follows from the beta-integral normalization \citep[\S5.12]{nist2026beta}. Taking \(a=b=c=1\) already detects the error: the written right side equals the triangle's area \(1/2\).}
\GSfield{Present status}{As of November 2026, the qualified statement above is established; the distinctions from the notebook wording remain essential.}
\GSfield{Assessment}{literal statement: false as written; publication match: exact; verification: read the result; development: annotated; confidence: high}
\end{GSitem}
\begin{GSitem}{1967-133.3}{E}{Commuting continuous interval maps need not share a fixed point}
\GSfield{Kind; source}{theorem; \GSpage{P000212}}
\GSfield{The notes}{The common-fixed-point hypothesis for two commuting continuous maps of an interval has a counterexample.}
\GSfield{Statement}{There exist continuous maps \(f,g:[0,1]\to[0,1]\) with \(f\circ g=g\circ f\) but \(\operatorname{Fix}(f)\cap\operatorname{Fix}(g)=\varnothing\). Thus the fixed-point property of an interval for individual maps does not extend to arbitrary commuting pairs.}
\GSfield{Fate}{\citet[main construction, pp.77--92]{boyce1969commuting} gives a counterexample. His paper was received 9 January 1967, revised 1 December 1967, and published in 1969; it cites Huneke's independent 1967 announcement. The notebook therefore records circulating results, not a problem still unresolved at the lecture.}
\GSfield{Present status}{As of November 2026, the continuous common-fixed-point conjecture is false. The 1969 publication date must not conceal the 1967 circulation documented in the primary paper.}
\GSfield{Assessment}{literal statement: correct; publication match: exact; verification: read the result; development: none; confidence: high; problem \hyperref[prob:GP-1967-02]{GP-1967-02}}
\end{GSitem}
\begin{GSitem}{1967-133.4}{E}{A fixed point of one map that is periodic for the other}
\GSfield{Kind; source}{theorem; \GSpage{P000212}}
\GSfield{The notes}{If the commuting maps are differentiable, there is a point fixed for one and \ru{конечной кратности} for the other; the English translation says of finite multiplicity.}
\GSfield{Statement}{A precise established version is: if \(f:[0,1]\to[0,1]\) is continuously differentiable and \(g:[0,1]\to[0,1]\) is continuous with \(fg=gf\), then for some \(x\) and integer \(n\ge1\), \(f(x)=x\) and \(g^n(x)=x\). The Russian phrase here concerns finite period under iteration, not multiplicity of a zero.}
\GSfield{Reconstruction}{The statement uses the verified \(C^1\) formulation and identifies finite multiplicity as finite period; it does not infer a common fixed point.}
\GSfield{Fate}{The identifiable result is \citet[pp.353--355]{schwartz1965periodic}. The publisher metadata, the reference in \citet[references]{boyce1969commuting}, and the explicit statement in \citet[Introduction, pp.269--270]{alikhanikoopaei1998commuting} were checked; the Schwartz proof itself was not available in readable full text in this audit. Accordingly the exact weakening from \(C^1\) to merely differentiable maps is not asserted.}
\GSfield{Present status}{As of November 2026, the qualified statement above is established; the distinctions from the notebook wording remain essential.}
\GSfield{Assessment}{literal statement: correct with missing hypotheses; publication match: partial; verification: indirect; development: annotated; confidence: medium; problem \hyperref[prob:GP-1967-02]{GP-1967-02}}
\end{GSitem}

\GStalk{134}{1967-09-25}{D. B. Fuchs}{P000212-P000217}{Characteristic operations and universal frame spaces}
The continuation headed Fuchs starts on P000213 and runs through P000215. It computes the new characteristic cohomology from invariant forms, identifies natural operations taking principal bundles to other principal bundles, and constructs universal objects from frames in arbitrary products. These are precise counterparts of \citet[\S\S1, 2, 4]{gelfandfuchs1967topology}; the ordinary constant sheaf must be distinguished from the sheaf of continuous real functions.

\begin{GSitem}{1967-134.1}{A}{The bundle category and its characteristic operations}
\GSfield{Kind; source}{definition; \GSpage{P000213}}
\GSfield{The notes}{The free action, the global fiberwise linear-chart condition, and equivariant bundle morphisms are restated. \(h_{\mathrm{alg}}(G)\) denotes all characteristic classes formed from a contravariant functor \(h\).}
\GSfield{Statement}{Fix a closed subgroup \(G\subset GL(n,\mathbb R)\). Objects are the principal bundles over Hausdorff bases satisfying the global chart condition of 1967-132.5; morphisms are equivariant maps covering continuous base maps. A characteristic operation assigns \(c(\xi)\in h(X)\) to each object naturally under pullback. No vanishing of the chart map outside a prescribed neighborhood is required.}
\GSfield{Fate}{\citet[\S1 and the beginning of \S2]{gelfandfuchs1967topology} formalizes this category and its natural operations.}
\GSfield{Present status}{As of November 2026, the qualified statement above is established; the distinctions from the notebook wording remain essential.}
\GSfield{Assessment}{literal statement: correct; publication match: exact; verification: read the result; development: none; confidence: high}
\end{GSitem}
\begin{GSitem}{1967-134.2}{A}{Cohomology with the sheaf of continuous real functions}
\GSfield{Kind; source}{definition; \GSpage{P000213}}
\GSfield{The notes}{The coefficient sheaf called \(\mathbb R\) consists of germs of continuous real functions. The notes say \(H^q(X,\mathbb R)=0\) for paracompact \(X\).}
\GSfield{Statement}{Write \(\mathcal C_X\) for the sheaf of continuous real functions. On a paracompact Hausdorff space it is fine, so \(H^q(X;\mathcal C_X)=0\) for \(q>0\). Degree zero is \(C(X,\mathbb R)\), usually nonzero. The characteristic groups \(H^q_{\mathrm{alg}}(G;\mathbb R)\) classify natural classes with values in this sheaf cohomology over the enlarged category, not ordinary singular-cohomology characteristic classes.}
\GSfield{Reconstruction}{The notation \(\mathcal C_X\) separates the variable-function sheaf from the constant real sheaf used in the compact-dual calculation.}
\GSfield{Fate}{The notation and vanishing are the setup of \citet[\S2]{gelfandfuchs1967topology}. Positive degree is implicit in the seminar parenthesis and is restored here.}
\GSfield{Present status}{As of November 2026, the qualified statement above is established; the distinctions from the notebook wording remain essential.}
\GSfield{Assessment}{literal statement: correct with missing hypotheses; publication match: exact; verification: read the result; development: annotated; confidence: high}
\end{GSitem}
\begin{GSitem}{1967-134.3}{A}{Compact-dual calculation in the semisimple case}
\GSfield{Kind; source}{theorem; \GSpage{P000213-P000214}}
\GSfield{The notes}{For semisimple \(G\) with finitely many components, \(H^q_{\mathrm{alg}}(G)=H^q(\bar G/K;\mathbb R)\), with constant coefficients on the right; \(K\) is maximal compact and \(\bar G\) a compact form.}
\GSfield{Statement}{For connected semisimple linear \(G\), let \(\mathfrak g=\mathfrak k\oplus\mathfrak p\) be its Cartan decomposition, and choose the compact dual pair with Lie algebras \(\mathfrak g_u=\mathfrak k\oplus i\mathfrak p\) and \(\mathfrak k\). Then \(H^*_{\mathrm{alg}}(G;\mathbb R)\cong H^*(G_u/K_u;\mathbb R)\), with the compatible global compact dual and component actions included when \(G\) is disconnected.}
\GSfield{Reconstruction}{The Lie-algebra construction avoids an unjustified assertion that every chosen maximal compact subgroup embeds unchanged in an arbitrarily chosen compact global form.}
\GSfield{Fate}{This is \citet[Theorem 2.9]{gelfandfuchs1967topology}. It concerns real cohomology and the compact symmetric dual, not replacement of the original ordinary classifying space \(BG\) by \(G_u/K_u\).}
\GSfield{Present status}{As of November 2026, the qualified statement above is established; the distinctions from the notebook wording remain essential.}
\GSfield{Assessment}{literal statement: correct with missing hypotheses; publication match: exact; verification: read the result; development: annotated; confidence: high}
\end{GSitem}
\begin{GSitem}{1967-134.4}{A}{The linear-group example and the sphere ambiguity}
\GSfield{Kind; source}{claim; \GSpage{P000214}}
\GSfield{The notes}{For \(GL(n,\mathbb R)\), the notes name \(U(n)\) and \(O(n)\), then write: second order, as for \(S^2\).}
\GSfield{Statement}{The reductive compact dual for \(GL(n,\mathbb R)/O(n)\) is \(U(n)/O(n)\). The rank-two semisimple example with spherical dual is instead \(SL(2,\mathbb R)/SO(2)\), whose compact dual is \(SU(2)/SO(2)\simeq S^2\). These spaces cannot be conflated: \(\dim U(2)/O(2)=3\), whereas \(\dim S^2=2\).}
\GSfield{Reconstruction}{The missing transition from the reductive \(GL(n)\) example to \(SL(2)\) is flagged, not silently supplied as notebook wording.}
\GSfield{Fate}{\citet[Theorem 2.9 and its examples]{gelfandfuchs1967topology} supports the semisimple example. The scan does not insert \(SL(2)\) before the sphere remark, so the intended change of group is an emendation, not a resolved transcription.}
\GSfield{Present status}{As of November 2026, the qualified statement above is established; the distinctions from the notebook wording remain essential.}
\GSfield{Assessment}{literal statement: not determinable; publication match: partial; verification: read the result; development: annotated; confidence: medium}
\end{GSitem}
\begin{GSitem}{1967-134.5}{A}{Invariant forms compute the characteristic cohomology}
\GSfield{Kind; source}{theorem; \GSpage{P000214}}
\GSfield{The notes}{For a nonsemisimple group one takes \(G/K\), its \(G\)-invariant differential forms, and the cohomology of their differential complex.}
\GSfield{Statement}{For a closed linear Lie group in the category under discussion, \[H^q_{\mathrm{alg}}(G;\mathbb R)\cong H^q\bigl(\Omega^*(G/K)^G,d\bigr).\] Exterior differentiation preserves invariance. This is invariant-form, or relative Lie-algebra, cohomology; it need not equal the de Rham cohomology of the contractible space \(G/K\).}
\GSfield{Fate}{The exact result is \citet[Theorem 2.1]{gelfandfuchs1967topology}. The qualification about invariant forms explains how a nonzero answer survives when the underlying homogeneous space is contractible.}
\GSfield{Present status}{As of November 2026, the qualified statement above is established; the distinctions from the notebook wording remain essential.}
\GSfield{Assessment}{literal statement: correct; publication match: exact; verification: read the result; development: none; confidence: high}
\end{GSitem}
\begin{GSitem}{1967-134.6}{A}{Compact groups and the additive real group}
\GSfield{Kind; source}{theorem; \GSpage{P000214}}
\GSfield{The notes}{For compact groups everything is trivial. For \(\mathbb R\), the answer is \(\mathbb R\) in degrees 0 and 1 and zero above, with representatives \(1\) and \(dx\).}
\GSfield{Statement}{If \(G\) is compact, the invariant-form complex on \(G/G\) has \(H^0=\mathbb R\) and \(H^q=0\) for \(q>0\). For \(G=(\mathbb R,+)\), the complex of translation-invariant forms is \(\mathbb R\to\mathbb R\,dx\to0\) with zero differential, giving exactly the stated groups.}
\GSfield{Reconstruction}{The degree-zero exception and the zero differential are supplied explicitly.}
\GSfield{Fate}{These are direct instances of \citet[Theorem 2.1]{gelfandfuchs1967topology}. The word trivial refers to positive degrees, not degree zero.}
\GSfield{Present status}{As of November 2026, the qualified statement above is established; the distinctions from the notebook wording remain essential.}
\GSfield{Assessment}{literal statement: correct with missing hypotheses; publication match: exact; verification: read the result; development: annotated; confidence: high}
\end{GSitem}
\begin{GSitem}{1967-134.7}{A}{Natural principal H-bundles are induced by homomorphisms}
\GSfield{Kind; source}{theorem; \GSpage{P000214}}
\GSfield{The notes}{\(EH_{\mathrm{alg}}(G)=\operatorname{Hom}(G,H)\), more precisely continuous homomorphisms modulo conjugacy in \(H\).}
\GSfield{Statement}{For the specified category of principal \(G\)-bundles and a topological group \(H\), natural assignments of isomorphism classes of ordinary locally trivial principal \(H\)-bundles correspond to \(\operatorname{Hom}_{\mathrm{cont}}(G,H)/H\). A homomorphism \(\phi\) acts by extension of structure group \(E\mapsto E\times_\phi H\); conjugate homomorphisms give naturally isomorphic assignments.}
\GSfield{Fate}{This is \citet[Theorem 4.1]{gelfandfuchs1967topology}. The paper allows ordinary local triviality for \(H\)-bundles; when \(H\) is a Lie group, using the enlarged condition gives the same answer.}
\GSfield{Present status}{As of November 2026, the qualified statement above is established; the distinctions from the notebook wording remain essential.}
\GSfield{Assessment}{literal statement: correct; publication match: exact; verification: read the result; development: none; confidence: high}
\end{GSitem}
\begin{GSitem}{1967-134.8}{A}{Frames in a product and universal bundle morphisms}
\GSfield{Kind; source}{construction; \GSpage{P000215}}
\GSfield{The notes}{\(\mathcal E_G^I\) consists of \(n\)-frames in \(\mathbb R^I\) with the Tychonoff topology. \(GL(n,\mathbb R)\) and \(G\) act, and \(S_G^I=\mathcal E_G^I/G\); every bundle maps into one for sufficiently large \(I\).}
\GSfield{Statement}{Let \(F_n(\mathbb R^I)\) be ordered linearly independent \(n\)-tuples, equivalently injective linear maps \(\mathbb R^n\to\mathbb R^I\), with the subspace topology inherited from \((\mathbb R^I)^n\). Precomposition gives a free action of \(G\subset GL(n,\mathbb R)\). Taking sufficiently many global charts gives a bundle morphism into \(F_n(\mathbb R^I)\to F_n(\mathbb R^I)/G\). The total space depends on \(n,I\), while the quotient and action depend on \(G\).}
\GSfield{Fate}{\citet[\S1, Proposition 1.1]{gelfandfuchs1967topology} proves this exact universal construction.}
\GSfield{Present status}{As of November 2026, the qualified statement above is established; the distinctions from the notebook wording remain essential.}
\GSfield{Assessment}{literal statement: correct; publication match: exact; verification: read the result; development: none; confidence: high; problem \hyperref[prob:GP-1967-01]{GP-1967-01}}
\end{GSitem}
\begin{GSitem}{1967-134.9}{A}{A countable-coordinate lemma for frame spaces}
\GSfield{Kind; source}{theorem; \GSpage{P000215}}
\GSfield{The notes}{For every continuous real function on \(\mathcal E_G^I\), there is \(J\subset I\) with \(I\setminus J\) at most countable such that its restriction to \(\mathcal E_G^J\) is constant.}
\GSfield{Statement}{Embed \(F_n(\mathbb R^J)\) into \(F_n(\mathbb R^I)\) by setting coordinates outside \(J\) equal to zero. Every continuous real-valued function on the latter has a constant restriction to such a subspace for some co-countable \(J\). The statement does not guarantee a nonempty subspace when \(I\) is itself countable and too many coordinates have been removed.}
\GSfield{Reconstruction}{The zero-extension inclusion and the possibly vacuous small-index case are stated explicitly.}
\GSfield{Fate}{The exact lemma is \citet[Proposition 1.3]{gelfandfuchs1967topology}; it is one of the technical reasons large product-indexed universal spaces behave differently from familiar CW models.}
\GSfield{Present status}{As of November 2026, the qualified statement above is established; the distinctions from the notebook wording remain essential.}
\GSfield{Assessment}{literal statement: correct; publication match: exact; verification: read the result; development: annotated; confidence: high}
\end{GSitem}

\GStalk{135}{1967-10-02}{D. B. Fuchs}{P000215-P000219}{Elementary G-objects and independence of the model}
The dated Fuchs talk starts on P000216 and ends on P000217 before the programme heading. It abstracts the frame and cone constructions by choosing a pointed contractible \(G\)-space and a semigroup of admissible maps. The 1968 Gelfand--Fuchs note gives a recognizable continuation, but strengthens the local semigroup axioms and uses a more restrictive universal-space construction. Those differences prevent an unqualified exact-match verdict for the handwritten theorem.

\begin{GSitem}{1967-135.1}{A}{A pointed G-space with an admissible semigroup}
\GSfield{Kind; source}{definition; \GSpage{P000216}}
\GSfield{The notes}{\(S\subset\operatorname{Map}(T,T)\) has the compact-open topology, contains the \(G\)-action, and is stable under left and right multiplication by it. The axioms specify a fixed point, an \(S\)-contraction, faithfulness, closedness of \(G\), and nearby maps that are embeddings.}
\GSfield{Statement}{The published elementary \(G\)-object consists of a pointed space \(T\), a topological semigroup \(A\) acting faithfully and fixing the point, an embedded subgroup \(G\), and a contraction through \(A\). Crucially, a neighborhood \(W\supset G\) consists of invertible elements and admits a right-equivariant retraction \(\lambda:W\to G\), with \(\lambda(ag)=\lambda(a)g\) and \(\lambda(g)=g\). This last condition is stronger than merely saying nearby maps are topological embeddings.}
\GSfield{Reconstruction}{The published definition is stated separately from the recorded six-axiom list. No equivalence of the two axiom systems is claimed.}
\GSfield{Fate}{\citet[\S1, axioms 1--4]{gelfandfuchs1968classifying} is the subsequent formulation, received 17 April 1968 and published that year. Its stronger axiom is necessary to identify what was actually proved.}
\GSfield{Present status}{As of November 2026, the qualified statement above is established; the distinctions from the notebook wording remain essential.}
\GSfield{Assessment}{literal statement: correct with missing hypotheses; publication match: partial; verification: read the result; development: annotated; confidence: high}
\end{GSitem}
\begin{GSitem}{1967-135.2}{A}{Linear and cone elementary objects}
\GSfield{Kind; source}{construction; \GSpage{P000216-P000217}}
\GSfield{The notes}{The examples are \(T=\mathbb R^n\), with all linear maps, and \(T=CG\), with the group action and contractions to the apex. The latter is described as a product of cones with the all-apex point removed.}
\GSfield{Statement}{For a closed linear Lie subgroup \(G\), take \(T=\mathbb R^n\) and \(A=\operatorname{End}(\mathbb R^n)\). For a topological group, take the cone of pairs \((g,t)\), \(t\ge0\), with all \((g,0)\) identified, and multiplication \((g,t)(h,s)=(gh,ts)\). The cone example has invertible nonapex elements. The positive-ray example uses multiplication by \(e^u\), with the zero sequence deleted.}
\GSfield{Fate}{These are \citet[\S1, Examples 1--2]{gelfandfuchs1968classifying}. The last deleted point is essential to freeness, as in the September ray construction.}
\GSfield{Present status}{As of November 2026, the qualified statement above is established; the distinctions from the notebook wording remain essential.}
\GSfield{Assessment}{literal statement: correct with missing hypotheses; publication match: exact; verification: read the result; development: none; confidence: high}
\end{GSitem}
\begin{GSitem}{1967-135.3}{A}{Universal embeddings versus the published coordinate-invertible model}
\GSfield{Kind; source}{construction; \GSpage{P000217}}
\GSfield{The notes}{The universal total space is said to be every embedding \(T\to T^I\) whose coordinate maps lie in \(S\). In the linear example this is the frame space.}
\GSfield{Statement}{The 1968 construction instead takes tuples \((a_i)\in A^I\) with at least one coordinate invertible in \(A\), equipped with the diagonal right \(G\)-action. Such a tuple gives an embedding \(T\to T^I\), but the converse need not hold: a linear injection into a product can be assembled from singular coordinate maps. Thus the earlier all-frame model and the later coordinate-invertible model must not be identified pointwise.}
\GSfield{Reconstruction}{The distinction is explicit: the 1968 universal object is a related construction, not verbatim the set of all embeddings in the notebook.}
\GSfield{Fate}{\citet[\S2, Propositions 3--4 and Theorem 1]{gelfandfuchs1968classifying} proves universality for this published model. The all-frame linear model had already been treated in \citet[Proposition 1.1]{gelfandfuchs1967topology}.}
\GSfield{Present status}{As of November 2026, the qualified statement above is established; the distinctions from the notebook wording remain essential.}
\GSfield{Assessment}{literal statement: correct with missing hypotheses; publication match: partial; verification: read the result; development: annotated; confidence: high}
\end{GSitem}
\begin{GSitem}{1967-135.4}{A}{Finite products are finite approximations, not universal EG}
\GSfield{Kind; source}{claim; \GSpage{P000217}}
\GSfield{The notes}{For finite \(I\), the notes call the linear case the usual classifying space and the cone case Milnor's construction with a minor modification.}
\GSfield{Statement}{Finite-dimensional Stiefel and Grassmann spaces classify bundles only subject to dimension or embedding restrictions; their total spaces are not generally contractible universal \(EG\) spaces. Likewise, a finite cone-product or finite join is an approximation. Universal models require an appropriate infinite construction and its stated topology, or the sufficiently large index set tailored to the bundle.}
\GSfield{Reconstruction}{The sentence is interpreted as referring to familiar finite stages, rather than silently changing finite to infinite in the transcription.}
\GSfield{Fate}{The large-index quantifier in \citet[Proposition 1.1]{gelfandfuchs1967topology} and \citet[\S2]{gelfandfuchs1968classifying} supplies the precise interpretation. The literal finite-index universal claim is too strong.}
\GSfield{Present status}{As of November 2026, the qualified statement above is established; the distinctions from the notebook wording remain essential.}
\GSfield{Assessment}{literal statement: false as written; publication match: partial; verification: read the result; development: annotated; confidence: high}
\end{GSitem}
\begin{GSitem}{1967-135.5}{A}{Local T-triviality and universal pullback representation}
\GSfield{Kind; source}{theorem; \GSpage{P000217}}
\GSfield{The notes}{For each \(x\), the associated \(T\)-bundle has a global map to \(T\), an allowed map on every fiber and a homeomorphism at \(x\). Every bundle over a Hausdorff base maps to the universal one.}
\GSfield{Statement}{Under the published elementary-object axioms, require that for every \(x\) there is a continuous associated-bundle map to \(T\) whose restriction, written in any principal frame, belongs to \(A\), and is invertible in \(A\) on the selected fiber. Such bundles are ordinarily locally trivial; ordinary locally trivial bundles over completely regular bases satisfy this condition. Collecting these maps gives the universal morphism for sufficiently large \(I\).}
\GSfield{Reconstruction}{The stronger semigroup assumptions and invertibility requirement are retained from the publication.}
\GSfield{Fate}{\citet[Propositions 1--2 and Theorem 1]{gelfandfuchs1968classifying} proves these statements. Mere homeomorphism as a map of \(T\), without invertibility in the specified semigroup, is not the complete published hypothesis.}
\GSfield{Present status}{As of November 2026, the qualified statement above is established; the distinctions from the notebook wording remain essential.}
\GSfield{Assessment}{literal statement: correct with missing hypotheses; publication match: partial; verification: read the result; development: annotated; confidence: high; problem \hyperref[prob:GP-1967-01]{GP-1967-01}}
\end{GSitem}
\begin{GSitem}{1967-135.6}{A}{Independence of the elementary object}
\GSfield{Kind; source}{theorem; \GSpage{P000217}}
\GSfield{The notes}{\(H^q_{\mathrm{alg}}(G;\mathbb R)\) does not depend on \(T,S\); the notes add that a group \(A\) may replace \(\mathbb R\).}
\GSfield{Statement}{For a paracompact topological group \(G\) and elementary \(G\)-objects satisfying the published axioms, the characteristic groups with coefficients in a fixed topological \(G\)-module are independent of the chosen object. The coefficient sheaf is the sheaf of continuous sections of the associated module bundle. One must hold the coefficient sheaf and the naturality category fixed; the notebook does not specify which arbitrary coefficient groups or sheaves its parenthetical extension includes.}
\GSfield{Reconstruction}{The real-coefficient construction is the definite statement; the unrestricted coefficient parenthesis is not upgraded into an unproved generalization.}
\GSfield{Fate}{\citet[\S3, Theorem 2]{gelfandfuchs1968classifying} is the identifiable continuation. The stronger elementary-object axioms and unspecified coefficient extension make the match partial.}
\GSfield{Present status}{As of November 2026, the qualified statement above is established; the distinctions from the notebook wording remain essential.}
\GSfield{Assessment}{literal statement: correct with missing hypotheses; publication match: partial; verification: read the result; development: annotated; confidence: medium; problem \hyperref[prob:GP-1967-01]{GP-1967-01}}
\end{GSitem}

\GStalk{136}{1967-10-02}{I. M. Gelfand}{P000217-P000220}{Programme for future seminar talks}
The heading asking what will come next occupies the end of P000217 and P000218. It lists Uspenskii on logic; Lidskii on eigenfunction asymptotics after H\"ormander; Sinai on invariant measures; Kazhdan on discrete groups and Selberg's problem; Kushnirenko on analytic linearization; Tyurina and S. Gelfand on Pham; Anosov and Maslov; Bernstein on \(P^\lambda\); Gelfand and Graev; S. G. Krein on interpolation; and Berezin and Kirillov on Grassmann analysis. Only the analytic-linearization entry states a sufficiently definite theorem. The programme supplies no dates or additional independent MCCME records for these proposed talks.

\begin{GSitem}{1967-136.1}{E}{Analytic linearization of a semisimple action at a fixed point}
\GSfield{Kind; source}{theorem; \GSpage{P000217}}
\GSfield{The notes}{The planned Kushnirenko talk is summarized as: an analytic action of a semisimple group near a fixed point is linear.}
\GSfield{Statement}{Let a finite-dimensional real semisimple Lie algebra act by germs of real-analytic vector fields vanishing at a point of a real-analytic manifold. There are local real-analytic coordinates in which the action is its linear isotropy representation. Equivalently, an analytic local action of the corresponding connected group is locally analytically linearizable at the fixed point.}
\GSfield{Reconstruction}{The local germ formulation and connected-group qualification state the precise theorem behind the one-line announcement.}
\GSfield{Fate}{\citet[Theorem 1]{kushnirenko1967linearization} proves the result; the note was received 16 November 1966 and published in the first 1967 issue, before this programme. The smooth category cannot simply replace the analytic one.}
\GSfield{Present status}{As of November 2026, the qualified statement above is established; the distinctions from the notebook wording remain essential.}
\GSfield{Assessment}{literal statement: correct with missing hypotheses; publication match: exact; verification: read the result; development: annotated; confidence: high}
\end{GSitem}
\begin{GSitem}{1967-136.2}{F}{Programme labels that do not specify mathematical assertions}
\GSfield{Kind; source}{remark; \GSpage{P000217-P000218}}
\GSfield{The notes}{The programme mentions conditions for invariant measures, a solution of Selberg's problem, Pham's work, \(P^\lambda\), Grassmann analysis, and several names without a theorem statement.}
\GSfield{Statement}{These entries identify proposed topics and speakers only. In particular, neither the invariant-measure conditions nor which Selberg problem is meant is stated here. No mathematical quantifiers can be recovered solely from these programme labels.}
\GSfield{Fate}{The adjacent later records, and the earlier Kazhdan material in 1966-101, provide context but do not turn this plan into a uniquely stated new theorem. No exact publication match is assigned to the topic list.}
\GSfield{Present status}{As of November 2026, this remains a programme record, not an independently identifiable open question. Its named topics are covered in the talk summary and are not counted as invented extra talks.}
\GSfield{Assessment}{literal statement: not determinable; publication match: none found; verification: unresolved lead; development: none; confidence: high}
\end{GSitem}

\GStalk{137}{1967-10-16}{S. G. Krein}{P000218-P000222}{Interpolation, symmetric spaces, and Fourier coefficients}
The Krein lecture begins on P000219 and ends in the middle of P000221 before the Sinai heading. It surveys classical interpolation, the Mityagin theorem, Lorentz refinements due to Dikarev and Matsaev, and Semenov's fundamental-function criteria and Fourier applications. The historical criteria are not all valid as stated: \citet[\S3]{sharpley1975interpolation} explicitly exhibits the defect in the Semenov theory. The corrected comparison uses Boyd dilation indices, not only values of the norm on characteristic functions.

\begin{GSitem}{1967-137.1}{E}{Interpolation pairs of compatible Banach couples}
\GSfield{Kind; source}{definition; \GSpage{P000219}}
\GSfield{The notes}{All four endpoint Banach spaces are embedded in a common topological linear space on each side; every endpoint-bounded linear operator must map the intermediate pair boundedly.}
\GSfield{Statement}{Let \((E_0,E_1)\) and \((F_0,F_1)\) be compatible Banach couples. Intermediate spaces satisfy \(E_0\cap E_1\hookrightarrow E\hookrightarrow E_0+E_1\) and similarly for \(F\). The pair \((E,F)\) is an interpolation pair if each linear \(A:E_0+E_1\to F_0+F_1\), bounded from \(E_i\) to \(F_i\) for \(i=0,1\), restricts boundedly from \(E\) to \(F\). Exactness is the stronger assertion that the intermediate operator norm is at most the maximum endpoint norm.}
\GSfield{Reconstruction}{The common-domain sum space and the distinction between boundedness and exactness are supplied.}
\GSfield{Fate}{This is the standard operator-interpolation framework used by \citet[opening definitions]{mityagin1965interpolation} and \citet[\S2]{dikarevmatsaev1966exact}. The lecture later asks an exact-norm question, which must be distinguished from bounded interpolation.}
\GSfield{Present status}{As of November 2026, the qualified statement above is established; the distinctions from the notebook wording remain essential.}
\GSfield{Assessment}{literal statement: correct with missing hypotheses; publication match: exact; verification: read the result; development: annotated; confidence: high}
\end{GSitem}
\begin{GSitem}{1967-137.2}{E}{Symmetric function spaces and basic inclusions}
\GSfield{Kind; source}{definition; \GSpage{P000219}}
\GSfield{The notes}{A Banach space of summable functions on \([0,1]\) is symmetric if it has the ideal property and its norm is invariant under measure-preserving transformations; then \(L_\infty\subset E\subset L_1\).}
\GSfield{Statement}{On the Lebesgue probability interval, a rearrangement-invariant Banach function space has the ideal property \(|f|\le|g|\Rightarrow\|f\|_E\le\|g\|_E\), and equimeasurable absolute values have equal norms. In the usual Banach-function-space setting, \(L^\infty\hookrightarrow E\hookrightarrow L^1\) continuously. A Fatou property, order continuity, or full symmetry is an additional property and must not be hidden inside the two displayed axioms.}
\GSfield{Reconstruction}{Equimeasurability and continuous inclusions are the modern formulation; the extra regularity assumptions are explicitly kept separate.}
\GSfield{Fate}{The distinction matters in \citet[Theorems 1 and 3]{mityagin1965interpolation} and in \citet[\S1]{boyd1969indices}. It is precisely why the unqualified interpolation claims later in the notes require care.}
\GSfield{Present status}{As of November 2026, the qualified statement above is established; the distinctions from the notebook wording remain essential.}
\GSfield{Assessment}{literal statement: correct with missing hypotheses; publication match: exact; verification: read the result; development: annotated; confidence: high}
\end{GSitem}
\begin{GSitem}{1967-137.3}{E}{Lorentz, associate Marcinkiewicz, and Orlicz examples}
\GSfield{Kind; source}{construction; \GSpage{P000219}}
\GSfield{The notes}{The examples are \(\int_0^1f^*\,d\psi\) for monotone \(\psi\). The displayed dual norm is \[\sup_e\int_e|f|/\psi(m(e)).\] Orlicz spaces, including \(L^p\), are also listed.}
\GSfield{Statement}{For an increasing concave function \(\psi\) with \(\psi(0)=0\) and \(\psi(t)>0\) for \(t>0\), the Lorentz norm is \(\|f\|_{\Lambda_\psi}=\int_0^1f^*(t)\,d\psi(t)\). Its associate norm is \(\sup_{0<t\le1}\psi(t)^{-1}\int_0^t f^*(s)\,ds\), equivalently the supremum over measurable sets of measure \(t\). For a Young function \(\Phi\), the Luxemburg Orlicz norm is \(\inf\{c>0:\int\Phi(|f|/c)\le1\}\); \(\Phi(u)=u^p\) gives \(L^p\).}
\GSfield{Reconstruction}{Concavity, positivity, the endpoint normalization, and the meaning of dual as associate norm are supplied.}
\GSfield{Fate}{Lorentz and associate spaces are used explicitly in \citet[\S1]{dikarevmatsaev1966exact} and \citet[\S1]{boyd1969indices}. Monotonicity alone is insufficient to guarantee that the displayed Lorentz functional is a Banach norm.}
\GSfield{Present status}{As of November 2026, the qualified statement above is established; the distinctions from the notebook wording remain essential.}
\GSfield{Assessment}{literal statement: correct with missing hypotheses; publication match: exact; verification: read the result; development: annotated; confidence: high}
\end{GSitem}
\begin{GSitem}{1967-137.4}{E}{Riesz--Thorin interpolation}
\GSfield{Kind; source}{theorem; \GSpage{P000219}}
\GSfield{The notes}{The first step is named as the Riesz--Thorin theorem.}
\GSfield{Statement}{For a linear operator bounded \(L^{p_i}(\mu)\to L^{q_i}(\nu)\) with norms \(M_i\), \(1\le p_i,q_i\le\infty\), the standard complex interpolation theorem gives \(L^{p_\theta}\to L^{q_\theta}\) with norm at most \(M_0^{1-\theta}M_1^\theta\), where \(p_\theta^{-1}=(1-\theta)p_0^{-1}+\theta p_1^{-1}\) and similarly for \(q_\theta\). The operator must be consistently defined on the endpoint sum.}
\GSfield{Reconstruction}{The full familiar statement fills in the named theorem; endpoint measurability conventions are the standard sigma-finite setting.}
\GSfield{Fate}{The theorem is background for the improvement discussed in \citet[Introduction]{dikarevmatsaev1966exact}. The notebook gives no separate proof or novelty claim.}
\GSfield{Present status}{As of November 2026, the qualified statement above is established; the distinctions from the notebook wording remain essential.}
\GSfield{Assessment}{literal statement: correct with missing hypotheses; publication match: partial; verification: read the result; development: annotated; confidence: high}
\end{GSitem}
\begin{GSitem}{1967-137.5}{E}{Weak endpoint bounds and Marcinkiewicz interpolation}
\GSfield{Kind; source}{theorem; \GSpage{P000220}}
\GSfield{The notes}{The distribution-function bound is \(n_{Af}(\tau)^{1/q_i}\le C\|f\|_{p_i}/\tau\), described as weaker endpoint information.}
\GSfield{Statement}{This means \(\|Af\|_{L^{q_i,\infty}}\le C_i\|f\|_{L^{p_i}}\). For example, if \(1\le p_0<p_1<\infty\) and a sublinear operator has weak types \((p_0,p_0)\) and \((p_1,p_1)\), it is strongly bounded on \(L^p\) for \(p_0<p<p_1\). More general off-diagonal versions require their own exponent hypotheses; the two endpoint inequalities alone are not a statement about every possible intermediate exponent.}
\GSfield{Reconstruction}{A safe diagonal theorem and the precise weak-space interpretation are separated from the later improving theorem.}
\GSfield{Fate}{\citet[Introduction and \S2]{boyd1969indices} uses this precise weak-type framework, and \citet[Theorem 2]{dikarevmatsaev1966exact} gives the off-diagonal Lorentz version relevant here.}
\GSfield{Present status}{As of November 2026, the qualified statement above is established; the distinctions from the notebook wording remain essential.}
\GSfield{Assessment}{literal statement: correct with missing hypotheses; publication match: exact; verification: read the result; development: annotated; confidence: high}
\end{GSitem}
\begin{GSitem}{1967-137.6}{E}{The Dikarev--Matsaev Lorentz improvement and optimality}
\GSfield{Kind; source}{theorem; \GSpage{P000220}}
\GSfield{The notes}{For improving operators \(q_i>p_i\), the intermediate image lies in a Lorentz space narrower than \(L^q\), and this image space is minimal for all operators with the endpoint bounds.}
\GSfield{Statement}{Take \(1<p_0<p_1<\infty\), distinct \(1<q_0,q_1<\infty\), and \(0<\theta<1\). Interpolation of weak endpoint maps gives \(L^{p,p}\to L^{q,p}\), where the reciprocal exponents are interpolated linearly. Since \(L^{p,p}=L^p\) and \(L^{q,p}\subset L^q\) when \(p<q\), this improves the target. The published sharpness theorem quantifies over all endpoint-admissible operators on nonatomic measure spaces, with an infinite source measure in its stated general form.}
\GSfield{Reconstruction}{The Lorentz second index \(p\) and the measure-space restriction in the sharpness theorem are supplied; minimality is not asserted for every one fixed operator.}
\GSfield{Fate}{\citet[Theorems 2--3]{dikarevmatsaev1966exact}, received 25 September 1965 and published in 1966, gives both the target and the quantified optimality. This is earlier work being reported by Krein.}
\GSfield{Present status}{As of November 2026, the qualified statement above is established; the distinctions from the notebook wording remain essential.}
\GSfield{Assessment}{literal statement: correct with missing hypotheses; publication match: partial; verification: read the result; development: annotated; confidence: high}
\end{GSitem}
\begin{GSitem}{1967-137.7}{E}{Mityagin exact interpolation and symmetric norms}
\GSfield{Kind; source}{theorem; \GSpage{P000220}}
\GSfield{The notes}{Every contraction on \(L^1\) and \(L^\infty\) is a contraction on exactly the symmetric spaces; a two-dimensional square-and-diamond diagram motivates the statement.}
\GSfield{Statement}{In finite dimensions, operators contractive for both \(\ell^1\) and \(\ell^\infty\) contract every symmetric norm; signed coordinate permutations and their convex hull give the elementary model. In function spaces, Mityagin proves the corresponding exact interpolation theorem for minimal symmetric ideals, and separately for maximal ideals. A fully symmetric Banach function space is an exact interpolation space for \((L^1,L^\infty)\). The two symmetry axioms alone, without the relevant regularity or majorization property, must not be substituted for these hypotheses.}
\GSfield{Reconstruction}{The finite-dimensional illustration and the infinite-dimensional theorem are distinguished. Full symmetry means monotonicity under Hardy--Littlewood submajorization.}
\GSfield{Fate}{\citet[Theorems 1--3]{mityagin1965interpolation} is the exact named source, received 3 December 1963 and published in 1965. Its minimal/maximal qualifications are absent from the notebook.}
\GSfield{Present status}{As of November 2026, the qualified statement above is established; the distinctions from the notebook wording remain essential.}
\GSfield{Assessment}{literal statement: correct with missing hypotheses; publication match: exact; verification: read the result; development: annotated; confidence: high}
\end{GSitem}
\begin{GSitem}{1967-137.8}{E}{Fundamental function and dyadic ratios}
\GSfield{Kind; source}{definition; \GSpage{P000220-P000221}}
\GSfield{The notes}{\(\|\chi_e\|_E=\varphi(m(e))\); \(\ell_E=\liminf_{t\downarrow0}\varphi(2t)/\varphi(t)\), \(u_E=\limsup\), with \(1\le\ell_E\le u_E\le2\), and both equal \(2^{1/p}\) for \(L^p\).}
\GSfield{Statement}{For a rearrangement-invariant space on \([0,1]\), the fundamental function is \(\varphi_E(t)=\|\chi_{[0,t]}\|_E\). Monotonicity and the triangle inequality give \(1\le\varphi_E(2t)/\varphi_E(t)\le2\) when \(0<t\le1/2\). For \(L^p\), \(\varphi(t)=t^{1/p}\), with the standard interpretation for \(p=\infty\). These ratios measure only indicator functions, not the norm of the dilation operator on all of \(E\).}
\GSfield{Fate}{These are the invariants used in \citet[Theorem 1]{semenov1965interpolation}. Their distinction from dilation indices is made explicit in \citet[\S3]{sharpley1975interpolation}.}
\GSfield{Present status}{As of November 2026, the qualified statement above is established; the distinctions from the notebook wording remain essential.}
\GSfield{Assessment}{literal statement: correct; publication match: exact; verification: read the result; development: none; confidence: high}
\end{GSitem}
\begin{GSitem}{1967-137.9}{E}{The fundamental-function criterion fails in full generality}
\GSfield{Kind; source}{claim; \GSpage{P000221}}
\GSfield{The notes}{Semenov's first theorem is recorded as: every operator bounded on all \(L^p\), \(1<p<\infty\), is bounded on \(E\) if and only if \(1<\ell_E\le u_E<2\).}
\GSfield{Statement}{The stated criterion is false for arbitrary symmetric spaces, even with separability added. The norm of every dilation, rather than only its values on indicators, controls the general interpolation property. For a rearrangement-invariant Banach function space with the Fatou property, writing the Boyd exponents in the convention \(0\le\beta_E\le\alpha_E\le1\), the correct nonendpoint condition is \(0<\beta_E\le\alpha_E<1\). More generally all weak-type \((p,p),(q,q)\) operators act on \(E\) exactly when \(1/q<\beta_E\le\alpha_E<1/p\).}
\GSfield{Reconstruction}{The modern Boyd-index theorem is a corrected replacement, not an assertion that Semenov used those indices. The all-\(L^p\) consequence follows by choosing finite endpoints around the Boyd exponents; the Hardy operators also show necessity.}
\GSfield{Fate}{The historical claim is \citet[Theorem 1]{semenov1965interpolation}; the later general argument appears in \citet{semenov1968interpolation}. \citet[\S3, pp.122--125]{sharpley1975interpolation} explicitly disproves the underlying theorem using a space with \(\varphi(t)=\sqrt t\) but an endpoint Boyd exponent. The replacement is \citet[Theorem 1]{boyd1969indices}.}
\GSfield{Present status}{As of November 2026, the historical general criterion is refuted; its exact appearance in the literature is not evidence of correctness. Fundamental-function criteria remain valid in appropriately restricted classes where they control the dilation norms.}
\GSfield{Assessment}{literal statement: false as written; publication match: exact; verification: read the result; development: annotated; confidence: high}
\end{GSitem}
\begin{GSitem}{1967-137.10}{B}{The improving symmetric-space criterion and its limitation}
\GSfield{Kind; source}{claim; \GSpage{P000221}}
\GSfield{The notes}{For operators \(L^{1/\alpha}\to L^{1/(\alpha-\Delta)}\), \(\alpha_0<\alpha<\alpha_1\), the notes assert \(E\to E_\Delta\) when \(2^{\alpha_0}<\ell_E\le u_E<2^{\alpha_1}\), with \(\|f\|_{E_\Delta}=\|t^{-\Delta}f^*(t)\|_E\).}
\GSfield{Statement}{The displayed target functional specializes, for \(E=L^p\) and \(0<\Delta<1/p\), to \(\bigl(\int_0^1[f^*(t)]^p t^{-p\Delta}dt\bigr)^{1/p}\), the Lorentz \(L^{q,p}\) functional with \(1/q=1/p-\Delta\). This correctly explains the classical improving case. The claimed extension to every symmetric \(E\), controlled only by its fundamental-function ratios, is not a valid general theorem; additional dilation or majorization hypotheses are needed.}
\GSfield{Reconstruction}{The exponent is read as \(1/(\alpha-\Delta)\), consistent with the target functional and improving-operator discussion. The general false claim is separated from its valid \(L^p\) specialization.}
\GSfield{Fate}{\citet[Theorem 3]{semenov1967fourier} and \citet[Theorem 2]{semenov1968interpolation} are the exact historical continuation; the latter manuscript was received 7 June 1967. The inaccuracies and failed underlying representation theorem are documented in \citet[\S3]{sharpley1975interpolation}.}
\GSfield{Present status}{As of November 2026, the \(L^p\)-to-Lorentz specialization is established, while the unqualified fundamental-function extension is invalid. The printed 1968 theorem is a publication match, not a valid modern resolution.}
\GSfield{Assessment}{literal statement: false as written; publication match: exact; verification: read the result; development: annotated; confidence: high}
\end{GSitem}
\begin{GSitem}{1967-137.11}{E}{Hausdorff--Young with the conjugate exponent}
\GSfield{Kind; source}{theorem; \GSpage{P000221}}
\GSfield{The notes}{The written expression is \((|c_k|^p)^{1/p}\le\|f\|_{L^p}\), with no visible summation and with \(p\) on both sides.}
\GSfield{Statement}{On the probability circle, for \(1\le p\le2\) and \(p^{-1}+(p^{\prime})^{-1}=1\), the Fourier coefficients satisfy \(\|\widehat f\|_{\ell^{p^{\prime}}}\le\|f\|_{L^p}\). At \(p=1\) the target is \(\ell^\infty\). A bound on each individual coefficient is much weaker; an \(\ell^p\) bound for general \(p<2\) is not Hausdorff--Young.}
\GSfield{Reconstruction}{The conjugate exponent and summation or supremum are restored. The lack of a summation sign was confirmed in the scan.}
\GSfield{Fate}{This classical endpoint interpolation application is the baseline for \citet[Theorem 4]{dikarevmatsaev1966exact} and the Fourier estimates in \citet[\S2]{semenov1967fourier}.}
\GSfield{Present status}{As of November 2026, the qualified statement above is established; the distinctions from the notebook wording remain essential.}
\GSfield{Assessment}{literal statement: false as written; publication match: partial; verification: read the result; development: annotated; confidence: high}
\end{GSitem}
\begin{GSitem}{1967-137.12}{E}{Paley's rearranged-coefficient inequality}
\GSfield{Kind; source}{theorem; \GSpage{P000221}}
\GSfield{The notes}{The heading is written \ru{Пойя}; the inequality is \((\sum_{k\ge1}k^{p-2}(c_k^*)^p)^{1/p}\le C\|f\|_p\).}
\GSfield{Statement}{For \(1<p\le2\), Fourier coefficients relative to a uniformly bounded orthonormal system satisfy \[\left(\sum_{k\ge1}k^{p-2}(c_k^*)^p\right)^{1/p}\le C_p\|f\|_{L^p},\] where \(c_k^*\) is the decreasing rearrangement of absolute coefficients. This is a Lorentz \(\ell^{p^{\prime},p}\) bound; for the trigonometric system the usual normalization only changes the constant.}
\GSfield{Reconstruction}{The exponent range, orthonormal-system hypotheses, and corrected attribution are supplied.}
\GSfield{Fate}{\citet[Theorem 4]{dikarevmatsaev1966exact} gives the Lorentz estimate. \citet[Theorem 6 and the following paragraph]{semenov1967fourier} explicitly names the \(L^p\) specialization as Paley's theorem, supporting a mathematical-name emendation from the source's Poia/P'olya heading.}
\GSfield{Present status}{As of November 2026, the qualified statement above is established; the distinctions from the notebook wording remain essential.}
\GSfield{Assessment}{literal statement: correct with missing hypotheses; publication match: exact; verification: read the result; development: annotated; confidence: high}
\end{GSitem}
\begin{GSitem}{1967-137.13}{E}{Hardy--Littlewood for monotone Fourier coefficients}
\GSfield{Kind; source}{theorem; \GSpage{P000221}}
\GSfield{The notes}{Monotonically decreasing Fourier coefficients yield two-sided inequalities, attributed to Hardy and Littlewood.}
\GSfield{Statement}{For \(1<p<\infty\) and a nonnegative decreasing sequence \(a_n\to0\), the classical monotone sine-series theorem characterizes \(\sum_{n\ge1}a_n\sin nt\in L^p\) by \(\sum_{n\ge1}n^{p-2}a_n^p<\infty\), with equivalence of the corresponding norms. A parallel cosine-series formulation needs its constant coefficient included. Arbitrary rearrangements of coefficients do not give such a two-sided theorem.}
\GSfield{Reconstruction}{A definite monotone nonnegative sine-series version replaces the unspecified phrase Fourier coefficients.}
\GSfield{Fate}{\citet[Theorem 7 and its specialization paragraph]{semenov1967fourier} explicitly identifies the Hardy--Littlewood result. Its broader symmetric-space extension shares the fundamental-function difficulties discussed above; the classical \(L^p\) assertion is retained separately.}
\GSfield{Present status}{As of November 2026, the qualified statement above is established; the distinctions from the notebook wording remain essential.}
\GSfield{Assessment}{literal statement: correct with missing hypotheses; publication match: partial; verification: read the result; development: annotated; confidence: high}
\end{GSitem}
\begin{GSitem}{1967-137.14}{E}{No Hilbert interpolation space between L1 and L3/2}
\GSfield{Kind; source}{claim; \GSpage{P000221}}
\GSfield{The notes}{The notes attribute to Semenov the absence of a Hilbert interpolation space between \(L_1\) and \(L_{1.5}\).}
\GSfield{Statement}{In the nonatomic function-space setting there is no rearrangement-invariant space isomorphic to a Hilbert space with continuous inclusions \(L^{3/2}\hookrightarrow E\hookrightarrow L^1\). Indeed Hilbert type and cotype, applied to disjoint equimeasurable indicators, force \(\varphi_E(1/n)\asymp n^{-1/2}\), while \(L^{3/2}\hookrightarrow E\) forces \(\varphi_E(1/n)\le Cn^{-2/3}\), a contradiction. Consequently no Hilbert interpolation space in this symmetric category exists.}
\GSfield{Reconstruction}{The disjoint-indicator argument is editorial. It uses nonatomicity and Banach-space isomorphism to Hilbert space, not merely equality of norms. The broader abstract-category formulation is not inferred.}
\GSfield{Fate}{The symmetry framework is that of \citet{mityagin1965interpolation}. No exact early Semenov publication for this isolated assertion was identified in the 1965, 1967, and 1968 papers inspected. The elementary argument establishes the qualified claim without assigning it to an unrelated theorem.}
\GSfield{Present status}{As of November 2026, the qualified statement above is established; the distinctions from the notebook wording remain essential.}
\GSfield{Assessment}{literal statement: correct with missing hypotheses; publication match: none found; verification: indirect; development: annotated; confidence: medium}
\end{GSitem}
\begin{GSitem}{1967-137.15}{F}{Gelfand's request for an abstract description}
\GSfield{Kind; source}{question; \GSpage{P000221}; I. M. Gelfand}
\GSfield{The notes}{Gelfand asks for an abstract description of interpolation spaces; the parenthesis says it is unknown whether they are function spaces.}
\GSfield{Statement}{The notes do not specify a Banach couple, the operator category, an isometric versus isomorphic representation, or a required lattice structure. Thus several inequivalent representation and classification questions fit the wording, and no unique quantified conjecture can be extracted.}
\GSfield{Fate}{The inspected Mityagin, Semenov, Boyd, and Sharpley papers address specified function-space interpolation problems. They do not provide an exact match to this unrestricted question. It is retained as an identifiable programme remark, without an invented open-problem identifier.}
\GSfield{Present status}{As of November 2026, the historical question remains underdetermined by this record. It should not be reported as a single still-open theorem problem.}
\GSfield{Assessment}{literal statement: not determinable; publication match: thematic only; verification: unresolved lead; development: none; confidence: high}
\end{GSitem}

\GStalk{138}{1967-10-16}{Ya. G. Sinai}{P000220-P000224}{Toral automorphisms, entropy, and symbolic Markov models}
The Sinai heading begins at the bottom of P000221, after the interpolation discussion, and ends on P000223 before Leray. The exact model and recoding are those of Adler and Weiss, already published in June 1967. The general finite Markov-shift parenthesis goes beyond that special proof and leads to the later Bernoulli isomorphism theory. The continuation on 30 October explains geometric Markov partitions.

\begin{GSitem}{1967-138.1}{E}{Hyperbolic toral automorphisms and Lebesgue spectrum}
\GSfield{Kind; source}{theorem; \GSpage{P000221}}
\GSfield{The notes}{The integer matrix defines a two-torus map. Irrational eigenvalues are said to imply real eigenvalues, ergodicity, and countably multiple Lebesgue spectrum.}
\GSfield{Statement}{Let \(A\in SL(2,\mathbb Z)\) satisfy \(|\operatorname{tr}A|>2\). Its toral automorphism preserves Haar probability and is ergodic. Its Koopman operator on the mean-zero subspace of \(L^2(\mathbb T^2)\) is a countable direct sum of bilateral shifts; the constant functions separately contribute the eigenvalue \(1\). Indeed, the Fourier characters indexed by nonzero lattice points split into infinite \(A^{\mathsf T}\)-orbits. Nonrational eigenvalues alone do not imply real eigenvalues: a finite-order rotation matrix is a counterexample.}
\GSfield{Reconstruction}{The statement uses the intended real hyperbolic case; it does not redefine irrational to mean real without saying so.}
\GSfield{Fate}{This is the hyperbolic toral setting of \citet[\S1]{adlerweiss1967entropy}. The Fourier-orbit argument supplies the precise spectral meaning of the notebook sentence; the hyperbolicity hypothesis and removal of constants are indispensable.}
\GSfield{Present status}{As of November 2026, the qualified statement above is established; the distinctions from the notebook wording remain essential.}
\GSfield{Assessment}{literal statement: false as written; publication match: partial; verification: read the result; development: annotated; confidence: high}
\end{GSitem}
\begin{GSitem}{1967-138.2}{E}{Entropy classifies hyperbolic two-torus automorphisms metrically}
\GSfield{Kind; source}{problem; \GSpage{P000221-P000222}}
\GSfield{The notes}{When does a measurable \(U\) satisfy \(UT_1=T_2U\)? The entropy is \(\log\lambda\), where \(\lambda>1\), and equal entropy is asserted to imply metric similarity, although integral linear conjugacy can fail.}
\GSfield{Statement}{For hyperbolic automorphisms \(T_A,T_B\) of \(\mathbb T^2\), with Haar probability, there is a measure-preserving invertible conjugacy modulo null sets if and only if their entropies agree. Here \(h(T_A)=\log\rho(A)\). Equality of entropy does not generally imply conjugacy of the matrices in \(GL(2,\mathbb Z)\). Merely requiring a measurable map, without invertibility and preservation of measure, is a different and much weaker question.}
\GSfield{Reconstruction}{Metric similarity is interpreted as isomorphism of probability-preserving transformations modulo zero; \(\lambda\) is the modulus of the expanding eigenvalue.}
\GSfield{Fate}{This is exactly the Main Theorem of \citet[\S1]{adlerweiss1967entropy}, communicated on 24 April 1967 and published in June 1967. It was therefore already known at the October lecture, although the notes present its proof programme.}
\GSfield{Present status}{As of November 2026, the stated two-dimensional classification is settled. This occurrence is not an open entropy-classification problem for arbitrary measure-preserving systems.}
\GSfield{Assessment}{literal statement: correct with missing hypotheses; publication match: exact; verification: read the result; development: annotated; confidence: high; problem \hyperref[prob:GP-1967-03]{GP-1967-03}}
\end{GSitem}
\begin{GSitem}{1967-138.3}{E}{Topological conjugacies are affine under ergodicity}
\GSfield{Kind; source}{claim; \GSpage{P000222}}
\GSfield{The notes}{The notes say that if \(U\) is a homeomorphism then it is linear, adding in parentheses that this holds for any compact abelian groups.}
\GSfield{Statement}{If \(A,B\in GL(n,\mathbb Z)\), \(T_A\) is ergodic, and a homeomorphism \(U\) satisfies \(UT_A=T_BU\), then \(U(x)=Cx+c\), with \(C\in GL(n,\mathbb Z)\), \(CA=BC\), and \(T_Bc=c\). Thus \(U\) is linear after imposing \(U(0)=0\). Without ergodicity even the identity transformation admits many nonlinear conjugacies. The unrestricted compact-abelian parenthesis lacks its necessary hypotheses.}
\GSfield{Reconstruction}{The verified theorem is restricted to tori and states the translation term explicitly; the general compact-group extension is not asserted.}
\GSfield{Fate}{\citet[Theorem, pp.1223--1225]{adlerpalais1965homeomorphic} proves precisely the toral affine-rigidity assertion. The paper was received in November 1964 and published in 1965. Its proof does not justify deleting the translation term.}
\GSfield{Present status}{As of November 2026, the qualified statement above is established; the distinctions from the notebook wording remain essential.}
\GSfield{Assessment}{literal statement: false as written; publication match: partial; verification: read the result; development: annotated; confidence: high}
\end{GSitem}
\begin{GSitem}{1967-138.4}{E}{A primitive two-sided shift of finite type}
\GSfield{Kind; source}{definition; \GSpage{P000222}}
\GSfield{The notes}{\(\Omega_\Pi\) consists of bi-infinite sequences \((i_n)\) with \(\pi_{i_ni_{n+1}}=1\), for a finite zero-one matrix. A power of \(\Pi\) is assumed strictly positive.}
\GSfield{Statement}{For a finite zero-one matrix \(A\), let \(\Sigma_A=\{x\in\{1,\ldots,s\}^{\mathbb Z}:A_{x_nx_{n+1}}=1\}\), with the product topology and the coordinate shift \(\sigma\). Primitivity means \(A^k>0\) for some \(k\); it implies topological mixing of this shift. The Borel sigma algebra is generated by finite-coordinate cylinders.}
\GSfield{Fate}{\citet[\S\S2--3]{adlerweiss1967entropy} and \citet[\S\S2--4]{parry1964intrinsic} use this exact finite-state setup. The cylinder sigma algebra mentioned later in the notebook is part of this definition.}
\GSfield{Present status}{As of November 2026, the qualified statement above is established; the distinctions from the notebook wording remain essential.}
\GSfield{Assessment}{literal statement: correct; publication match: exact; verification: read the result; development: none; confidence: high}
\end{GSitem}
\begin{GSitem}{1967-138.5}{E}{The identity, full-shift, and golden-mean examples}
\GSfield{Kind; source}{construction; \GSpage{P000222}}
\GSfield{The notes}{The three matrices are \(I_2\), the two-by-two all-ones matrix, and \(\left(\begin{smallmatrix}1&1\\1&0\end{smallmatrix}\right)\). The last is left for examination.}
\GSfield{Statement}{For \(I_2\), the shift space consists of two fixed sequences and the matrix is not primitive. The all-ones matrix gives the full two-shift, with entropy \(\log2\). The third matrix forbids consecutive occurrences of symbol \(2\); it is primitive, has Perron eigenvalue \(\varphi=(1+\sqrt5)/2\), and has entropy \(\log\varphi\).}
\GSfield{Reconstruction}{The forbidden word and entropy of the third matrix are an editorial calculation.}
\GSfield{Fate}{These elementary examples follow directly from the finite-state framework and Perron construction in \citet[Theorems 5 and 7]{parry1964intrinsic}. The last example is worked out here as an annotation to the instruction to examine it.}
\GSfield{Present status}{As of November 2026, the qualified statement above is established; the distinctions from the notebook wording remain essential.}
\GSfield{Assessment}{literal statement: correct; publication match: exact; verification: read the result; development: annotated; confidence: high}
\end{GSitem}
\begin{GSitem}{1967-138.6}{E}{The Parry transition probabilities need both Perron vectors}
\GSfield{Kind; source}{construction; \GSpage{P000222}}
\GSfield{The notes}{The notebook writes \(\lambda e=\Pi^*e\), then \(p_{ij}=e_j\pi_{ij}/(\lambda e_i)\), and uses \(e_i\) as the stationary weights in the entropy formula.}
\GSfield{Statement}{Let \(A\) be primitive, \(Ar=\lambda r\), \(l^{\mathsf T}A=\lambda l^{\mathsf T}\), with positive vectors normalized by \(l^{\mathsf T}r=1\). Then \(P_{ij}=A_{ij}r_j/(\lambda r_i)\) is stochastic and \(w_i=l_ir_i\) is stationary. The cylinder probability is \(w_{i_0}P_{i_0i_1}\cdots P_{i_{m-1}i_m}\). A left Perron vector cannot replace \(r\) in the displayed row-transition formula for a general nonsymmetric matrix.}
\GSfield{Reconstruction}{The transpose is removed from the right-eigenvector equation, and the separate stationary vector \(w\) replaces the reused letter \(e\).}
\GSfield{Fate}{The correct construction is \citet[Theorem 5]{parry1964intrinsic}, also presented in \citet[\S3]{adlerweiss1967entropy}. Both earlier publications distinguish the two eigenvectors. A nonsymmetric numerical example is recorded in this fragment's computation file.}
\GSfield{Present status}{As of November 2026, the qualified statement above is established; the distinctions from the notebook wording remain essential.}
\GSfield{Assessment}{literal statement: false as written; publication match: exact; verification: read the result; development: annotated; confidence: high}
\end{GSitem}
\begin{GSitem}{1967-138.7}{E}{Entropy and uniqueness of the maximal stationary measure}
\GSfield{Kind; source}{theorem; \GSpage{P000222}}
\GSfield{The notes}{The entropy of a stationary Markov chain is displayed as a weighted sum of row entropies. The preceding chain is said to have entropy \(\log\lambda\) and to be the unique measure of maximal entropy.}
\GSfield{Statement}{For a finite stationary chain with distribution \(w\) and transition matrix \(P\), \(h_\mu(\sigma)=-\sum_iw_i\sum_jP_{ij}\log P_{ij}\), with \(0\log0=0\). For the Parry chain above the coboundary terms \(\log r_j-\log r_i\) cancel by stationarity, leaving \(\log\lambda\). An irreducible finite shift of finite type has this unique invariant probability of maximal entropy; primitivity additionally gives mixing.}
\GSfield{Reconstruction}{The correct stationary weights, invariance restriction, and distinction between irreducibility and primitivity are supplied.}
\GSfield{Fate}{\citet[Theorems 7 and 9]{parry1964intrinsic} proves the entropy and uniqueness results. The notebook's maximization must range over shift-invariant probabilities, rather than all measures on the cylinder sigma algebra.}
\GSfield{Present status}{As of November 2026, the qualified statement above is established; the distinctions from the notebook wording remain essential.}
\GSfield{Assessment}{literal statement: correct with missing hypotheses; publication match: exact; verification: read the result; development: annotated; confidence: high}
\end{GSitem}
\begin{GSitem}{1967-138.8}{E}{Gelfand on second-order minors and one sign change}
\GSfield{Kind; source}{theorem; \GSpage{P000222}; I. M. Gelfand}
\GSfield{The notes}{Gelfand remarks that if all entries and all second-order minors are positive, the first eigenvector is positive and the second has one sign change. He points to the action on the second exterior power.}
\GSfield{Statement}{For a real \(n\times n\) matrix, \(n\ge2\), with all entries and all ordered two-by-two minors strictly positive, the two eigenvalues of largest modulus are simple and satisfy \(\lambda_1>\lambda_2>0\), with \(\lambda_2>|\lambda_j|\) for \(j\ge3\). A first eigenvector is strictly positive; a second has exactly one sign variation after zero coordinates are omitted. Perron's theorem applied to the matrix and its second compound gives the spectral assertion. Positivity of \(v_1\wedge v_2\) makes the ratios \(v_{2i}/v_{1i}\) strictly monotone, and orthogonality to a positive left Perron vector forces their signs to cross.}
\GSfield{Reconstruction}{The statement separates the precise second-order consequence from the stronger full oscillation theorems, and makes the handling of zero coordinates explicit.}
\GSfield{Fate}{This is an elementary second-compound specialization of the oscillation method of \citet[\S\S4--5, pp.460--466]{gantmacherkrein1937matrices}. The exterior-power proof specified in the notebook is recognizable; no assumption of positivity of all higher minors is needed for this two-eigenvector conclusion.}
\GSfield{Present status}{As of November 2026, the qualified statement above is established; the distinctions from the notebook wording remain essential.}
\GSfield{Assessment}{literal statement: correct; publication match: partial; verification: read the result; development: annotated; confidence: high}
\end{GSitem}
\begin{GSitem}{1967-138.9}{E}{The four-block symbolic matrix}
\GSfield{Kind; source}{construction; \GSpage{P000223}}
\GSfield{The notes}{The alphabet has groups of sizes \(p,q,r,s\). From either \(a\)-group one may move to \(a^1\) or \(b^1\); from either \(b\)-group to \(a^2\) or \(b^2\). The equation is \(\lambda^2-(p+s)\lambda+ps-qr=0\), with \(ps-qr=1\).}
\GSfield{Statement}{Order the symbols as \(a^1,a^2,b^1,b^2\). Give all rows in the first two groups the allowed columns \(a^1,b^1\), and all rows in the last two groups the allowed columns \(a^2,b^2\). For positive block sizes the resulting matrix is primitive, has rank two, and has characteristic polynomial \(t^{p+q+r+s-2}(t^2-(p+s)t+ps-qr)\). This reduction follows by restricting to vectors constant on the two row types.}
\GSfield{Reconstruction}{Positivity of the four sizes is imposed for the simple primitive case. Zero-sized or sign-modified geometric cases need their separate conventions.}
\GSfield{Fate}{The construction and characteristic polynomial occur exactly in \citet[\S4]{adlerweiss1967entropy}. The Greek/Roman ambiguity in the notebook's transition labels does not affect the four named blocks or their destinations.}
\GSfield{Present status}{As of November 2026, the qualified statement above is established; the distinctions from the notebook wording remain essential.}
\GSfield{Assessment}{literal statement: correct; publication match: exact; verification: read the result; development: annotated; confidence: high}
\end{GSitem}
\begin{GSitem}{1967-138.10}{E}{Recoding to a full matrix with one forbidden transition}
\GSfield{Kind; source}{theorem; \GSpage{P000223}}
\GSfield{The notes}{The second model has \(N\) symbols and every transition except \(N\to1\). The two shifts are called isomorphic, and the coding is said to carry one maximal measure to the other.}
\GSfield{Statement}{Let the positive four-block matrix satisfy \(ps-qr=1\), and set \(N=p+s\). With their respective Parry measures, its shift and the shift defined by \(J_N-E_{N1}\) are measurably isomorphic modulo zero. The second matrix has characteristic polynomial \(t^{N-2}(t^2-Nt+1)\), so the Perron roots and entropies agree. The isomorphism is a variable-block measurable coding, not a global topological conjugacy.}
\GSfield{Reconstruction}{The category of isomorphism and the relation \(N=p+s\) are made explicit. Transposing the single forbidden transition merely changes the symbol/reversal convention.}
\GSfield{Fate}{This is precisely the determinant-one theorem in \citet[\S4]{adlerweiss1967entropy}. The paper proves the almost-everywhere invertible block coding and preservation of the maximal measures; equality of the eigenvalues alone would not prove it.}
\GSfield{Present status}{As of November 2026, the qualified statement above is established; the distinctions from the notebook wording remain essential.}
\GSfield{Assessment}{literal statement: correct with missing hypotheses; publication match: exact; verification: read the result; development: annotated; confidence: high; problem \hyperref[prob:GP-1967-03]{GP-1967-03}}
\end{GSitem}
\begin{GSitem}{1967-138.11}{B}{The parenthetical extension to all equal-entropy Markov shifts}
\GSfield{Kind; source}{claim; \GSpage{P000223}}
\GSfield{The notes}{A parenthesis refers to isomorphism of any two \(\Omega_\Pi\) with equal entropy, without another proof.}
\GSfield{Statement}{For primitive finite matrices, equip each two-sided shift with its Parry probability. These mixing finite-state Markov systems are Bernoulli, and two are measure-theoretically isomorphic exactly when their entropies agree. This is a statement about probability-preserving systems; equality of entropy does not classify shifts of finite type up to topological conjugacy.}
\GSfield{Reconstruction}{The notebook parenthesis is interpreted in its surrounding maximal-measure and primitive-matrix setting.}
\GSfield{Fate}{The general claim is supplied by the later Bernoulli theory: the 1970 Friedman--Ornstein Markov result and Ornstein's equal-entropy isomorphism theorem are explicitly presented together in \citet[\S5, Theorems 4--5 and references 3, 5]{ornstein1975applications}. That primary author exposition was read; the 1970 original proof was not accessible in this audit. Adler--Weiss's June 1967 proof verifies the preceding special family only.}
\GSfield{Present status}{As of November 2026, the qualified measurable statement is established. No unrestricted entropy-only topological classification is asserted.}
\GSfield{Assessment}{literal statement: correct with missing hypotheses; publication match: partial; verification: read the result; development: annotated; confidence: high}
\end{GSitem}

\GStalk{139}{1967-10-23}{Jean Leray}{P000222-P000225}{Algebraic singular supports of parameter-dependent periods}
The Leray heading begins near the bottom of P000223 and ends midway through P000224, before the date 30 October. Its relative cycles, moving hyperplanes, higher residues, and algebraic singular-support theorem match Leray's 1967 complement to Nilsson. The singular support is algebraic; the period function itself need not be algebraic. The notebook does not state the separate slow-growth conjecture appearing in Leray's paper.

\begin{GSitem}{1967-139.1}{E}{Projective relative periods}
\GSfield{Kind; source}{definition; \GSpage{P000223}}
\GSfield{The notes}{\(J(t)=\int_{\gamma^q(t)}\omega(t,\xi)\), with \(t\in\mathbb C^n\), \(\xi\in\mathbb C^{\ell+1}\), a cycle relative to \(W(t)\), and a degree-zero homogeneous form vanishing on \(W(t)\). Homogeneity is said to put the form on \(\mathbb P^\ell\).}
\GSfield{Statement}{A differential form on \(\mathbb C^{\ell+1}\setminus\{0\}\) descends to projective space when it is invariant under scalar multiplication and horizontal for the Euler vector field. For a closed relative \(q\)-form \(\omega_t\) with zero pullback to \(W_t\), integrate over a relative real \(q\)-cycle with boundary in \(W_t\), avoiding the polar locus. The period depends on its relative homology class, and a local family of cycles must be transported continuously away from singular parameters.}
\GSfield{Reconstruction}{Invariance alone is insufficient for a general differential form; horizontality and closedness are explicit editorial hypotheses.}
\GSfield{Fate}{This is the setup of \citet[\S1, pp.313--315]{leray1967complement}. The projective-basic condition, closedness, and relative-cycle interpretation make precise the notebook's compressed notation.}
\GSfield{Present status}{As of November 2026, the qualified statement above is established; the distinctions from the notebook wording remain essential.}
\GSfield{Assessment}{literal statement: correct with missing hypotheses; publication match: partial; verification: read the result; development: annotated; confidence: high}
\end{GSitem}
\begin{GSitem}{1967-139.2}{A}{The degeneracy locus of a moving hyperplane arrangement}
\GSfield{Kind; source}{construction; \GSpage{P000224}}
\GSfield{The notes}{\(W(t)=\bigcup_kW_k(t)\) is a union of hyperplanes depending algebraically on \(t\). The set \(S_0[W]\) is where the dimension of an intersection changes.}
\GSfield{Statement}{For a finite algebraic family of projective hyperplanes, the ranks of the matrices defining every subfamily are generically constant. Their rank-drop conditions, given by minors, produce an algebraic discriminant after coefficients and branch choices are placed on the appropriate algebraic parameter space. Outside it the incidence type is locally constant. Singularities involving the polar variety require additional tangency or stratified nontransversality conditions, not only hyperplane intersection ranks.}
\GSfield{Reconstruction}{Generic rank and algebraic-branch qualifications are supplied; the arrangement discriminant is not equated with all possible period singularities.}
\GSfield{Fate}{\citet[\S1 and \S\S3--4]{leray1967complement} describes the arrangement singular set and the additional appui conditions. The notebook's dimension-change test names only the arrangement part of the full exceptional locus.}
\GSfield{Present status}{As of November 2026, the qualified statement above is established; the distinctions from the notebook wording remain essential.}
\GSfield{Assessment}{literal statement: correct with missing hypotheses; publication match: partial; verification: read the result; development: annotated; confidence: high}
\end{GSitem}
\begin{GSitem}{1967-139.3}{A}{Iterated residues and the cycle on the complete intersection}
\GSfield{Kind; source}{method; \GSpage{P000224}}
\GSfield{The notes}{The notes identify \(\gamma^{q-m}(t)\) with \(\bigcap_jV_j(t)\), write poles \(v_j^{-1-\beta_j}\), and use a quotient of derivatives by powers of \(dv_j\) as a residue notation.}
\GSfield{Statement}{For transverse smooth hypersurfaces \(V_j=\{v_j=0\}\) and a meromorphic \(q\)-form \(\eta\) with simple logarithmic poles, its iterated residue is a \((q-m)\)-form on their complete intersection. A real \((q-m)\)-cycle \(\gamma\) there has an iterated tubular \(q\)-cycle \(\delta^m\gamma\), and, with compatible orientations, \(\int_{\delta^m\gamma}\eta=(2\pi i)^m\int_\gamma\operatorname{Res}_{V_1,\ldots,V_m}\eta\). Higher pole orders are handled by Leray residue classes and pole-order differentiation with the corresponding factorials. Powers of a one-form in the notebook denominator are not ordinary exterior powers.}
\GSfield{Reconstruction}{The simple-pole residue identity states a precise special case. The higher-pole notebook expression is retained as symbolic residue notation, not expanded into an unjustified differential quotient.}
\GSfield{Fate}{The exact higher-pole notation and its reduction to a tubular integral occur in \citet[\S1, formula (3), and \S20, p.352]{leray1967complement}. The published cycle lies on the complete intersection; it is not identical to that algebraic variety.}
\GSfield{Present status}{As of November 2026, the qualified statement above is established; the distinctions from the notebook wording remain essential.}
\GSfield{Assessment}{literal statement: false as written; publication match: partial; verification: read the result; development: annotated; confidence: high}
\end{GSitem}
\begin{GSitem}{1967-139.4}{A}{Algebraic singular support of the period}
\GSfield{Kind; source}{theorem; \GSpage{P000224}}
\GSfield{The notes}{The theorem says that \(J(t)\) has algebraic singular support: it is holomorphic on a covering of the complement of an algebraic set. Leray describes that set in the hyperplane case.}
\GSfield{Statement}{In Leray's projective relative-period setting, assume the integrand has algebraic singular support and impose his non-appui hypothesis for the moving arrangement against the polar and auxiliary algebraic varieties. Then there is an algebraic hypersurface \(S\) in parameter space such that the analytically continued period is holomorphic on the universal covering of its complement. Its singular support lies in the union of the arrangement discriminant and the specified incidence/tangency locus. This permits nontrivial monodromy and transcendental periods.}
\GSfield{Reconstruction}{The omitted non-appui hypothesis is retained by name and its role is explained; the algebraicity assertion concerns the exceptional set, not the value of the integral.}
\GSfield{Fate}{This is \citet[Theorem 1, pp.316--317]{leray1967complement}; its Corollary 1.1 identifies the absolute-cycle case with Nilsson's earlier theorem. The publisher dates \citet{nilsson1964growth} to October 1964, despite later bibliographies sometimes using the volume year 1965. Leray's general relative statement appeared in 1967. Only the absolute-case attribution, not Nilsson's full proof, was independently checked in the latter source.}
\GSfield{Present status}{As of November 2026, the stated singular-support theorem is established. The notes do not pose a slow-growth question, and none is manufactured from the separate hypothetical Theorem 3 in Leray's paper.}
\GSfield{Assessment}{literal statement: correct with missing hypotheses; publication match: partial; verification: read the result; development: annotated; confidence: high}
\end{GSitem}

\GStalk{140}{1967-10-30}{Ya. G. Sinai}{P000223-P000230}{Anosov diffeomorphisms and geometric Markov partitions}
The actual dated heading starts midway through P000224 and the talk ends on P000228. P000229 begins a new calendar-year block and does not belong to this talk despite the catalogue overlap. Stable and unstable foliations lead to rectangles, Markov coding, leaf measures, and the two-torus entropy proof. The first Sinai paper was received on 20 October 1967, ten days before this continuation; the construction paper followed on 3 November. The algebraic-classification question repeats GP-1966-01.

\begin{GSitem}{1967-140.1}{E}{Uniform hyperbolicity and the compact-distance caveat}
\GSfield{Kind; source}{definition; \GSpage{P000224-P000225}}
\GSfield{The notes}{The \(U\)-condition is illustrated by exponential separation of nearby points in one time direction, followed by contracting and expanding submanifolds of dimensions \(k,\ell\), with \(k+\ell=n\) and uniform transversality.}
\GSfield{Statement}{A \(C^1\) diffeomorphism \(T\) of a compact Riemannian manifold is Anosov if \(TM=E^s\oplus E^u\) is a continuous invariant splitting and there are \(C>0\), \(0<\theta<1\) such that \(\|DT^nv_s\|\le C\theta^n\|v_s\|\) and \(\|DT^{-n}v_u\|\le C\theta^n\|v_u\|\) for all \(n\ge0\). Uniform exponential expansion is a derivative or leafwise local statement. Distances between two fixed distinct points in a compact space cannot grow exponentially for all time.}
\GSfield{Reconstruction}{The derivative definition replaces the globally impossible distance inequality. One may also use intrinsic leaf metrics as in Sinai's publication.}
\GSfield{Fate}{The standard splitting and its stable/unstable interpretation are recalled by \citet[\S1]{sinai1968markov} and, in modern notation, \citet[Definition 2.1]{spatzier2016rigidity}. These were the established Anosov framework by the lecture, not a new consequence of the displayed metric inequality.}
\GSfield{Present status}{As of November 2026, the qualified statement above is established; the distinctions from the notebook wording remain essential.}
\GSfield{Assessment}{literal statement: false as written; publication match: partial; verification: read the result; development: annotated; confidence: high}
\end{GSitem}
\begin{GSitem}{1967-140.2}{E}{Stable manifolds and the choice of contraction constants}
\GSfield{Kind; source}{theorem; \GSpage{P000224-P000225}}
\GSfield{The notes}{Points with exponentially shrinking forward distance form \(\Gamma^{(c)}(x)\); backward shrinking gives \(\Gamma^{(p)}(x)\). The sets are said not to depend on \(a,c\).}
\GSfield{Statement}{For an Anosov diffeomorphism, \(W^s(x)=\{y:d(T^ny,T^nx)\to0\text{ as }n\to+\infty\}\) and the analogous \(W^u(x)\) are immersed leaves tangent to \(E^s,E^u\), of complementary dimensions. Their local plaques vary continuously and are uniformly transverse. For \(C^r\) dynamics the individual leaves have the corresponding smoothness, although transverse regularity is weaker. Changing to admissible weaker exponential estimates does not change the global asymptotic leaf; a fixed local plaque or a prescribed overly fast rate can change the selected set.}
\GSfield{Reconstruction}{Global leaves, local plaques, and the range of admissible exponential rates are distinguished.}
\GSfield{Fate}{\citet[\S1, pp.64--65]{sinai1968markov} explicitly distinguishes leafwise metrics and allowable contraction constants. The notebook's independence of constants is valid only in that standard asymptotic sense.}
\GSfield{Present status}{As of November 2026, the qualified statement above is established; the distinctions from the notebook wording remain essential.}
\GSfield{Assessment}{literal statement: correct with missing hypotheses; publication match: partial; verification: read the result; development: annotated; confidence: high}
\end{GSitem}
\begin{GSitem}{1967-140.3}{E}{Linear examples and structural stability}
\GSfield{Kind; source}{construction; \GSpage{P000225}}
\GSfield{The notes}{An integral two-by-two matrix of determinant one is given as an example. Smooth periodic nonlinear perturbations remain \(U\)-diffeomorphisms and have the same rotation number; a later sentence attributes topological conjugacy under small perturbations to Anosov.}
\GSfield{Statement}{A toral map \(T_A\), \(A\in GL(n,\mathbb Z)\), is Anosov precisely when no eigenvalue has modulus one. Every sufficiently small \(C^1\) perturbation is Anosov and is topologically conjugate to \(T_A\); more generally every Anosov diffeomorphism is structurally stable in this sense. In dimension two, determinant one alone is insufficient. The stable and unstable asymptotic directions supply the intended rotation/slope information, but there is no single circle rotation number for an arbitrary toral map.}
\GSfield{Reconstruction}{Smallness is in \(C^1\), hyperbolicity is explicit, and the informal rotation-number phrase is not promoted to a new invariant.}
\GSfield{Fate}{The structural-stability theorem and algebraic examples are part of \citet{smale1967dynamical}; the precise modern theorem is stated in \citet[Theorem 2.2]{spatzier2016rigidity}. The notebook itself correctly attributes the perturbation result to Anosov.}
\GSfield{Present status}{As of November 2026, the qualified statement above is established; the distinctions from the notebook wording remain essential.}
\GSfield{Assessment}{literal statement: correct with missing hypotheses; publication match: partial; verification: read the result; development: annotated; confidence: high}
\end{GSitem}
\begin{GSitem}{1967-140.4}{E}{Lattice-preserving Lie-group automorphisms}
\GSfield{Kind; source}{construction; \GSpage{P000225}}
\GSfield{The notes}{An automorphism of a Lie group preserves a discrete cocompact subgroup and descends to the quotient. It is \(U\) if its Lie-algebra map has no eigenvalue of modulus one.}
\GSfield{Statement}{Let \(G\) be a simply connected Lie group with lattice \(\Gamma\), and let \(\Phi\in\operatorname{Aut}(G)\) satisfy \(\Phi(\Gamma)=\Gamma\). Its induced diffeomorphism of \(G/\Gamma\) is Anosov when \(d\Phi_e\) has no eigenvalue on the unit circle; the invariant tangent framing makes this criterion necessary as well. Nilpotent \(G\) gives the standard nilmanifold algebraic examples.}
\GSfield{Fate}{The algebraic construction is described in \citet[algebraic examples]{smale1967dynamical} and is the model class in \citet[\S2]{spatzier2016rigidity}. The group, lattice-preservation, and compactness conditions are integral to the statement.}
\GSfield{Present status}{As of November 2026, the qualified statement above is established; the distinctions from the notebook wording remain essential.}
\GSfield{Assessment}{literal statement: correct; publication match: partial; verification: read the result; development: none; confidence: high}
\end{GSitem}
\begin{GSitem}{1967-140.5}{E}{Nilpotence does not guarantee an Anosov lattice automorphism}
\GSfield{Kind; source}{claim; \GSpage{P000225}}
\GSfield{The notes}{The parenthesis states that such automorphisms exist if and only if the Lie algebra is nilpotent.}
\GSfield{Statement}{Nilpotence alone is not sufficient. For the three-dimensional Heisenberg group with its standard lattice, a lattice-preserving automorphism induces \(B\in GL(2,\mathbb Z)\) on the abelianization and acts on the one-dimensional center by \(\det B=\pm1\). Hence its derivative always has an eigenvalue of modulus one and cannot define an Anosov automorphism. The existence of an appropriate lattice-preserving hyperbolic automorphism is an additional arithmetic and algebraic condition.}
\GSfield{Reconstruction}{This is an editorial counterexample, with its symbolic bracket calculation included in the fragment's checks.}
\GSfield{Fate}{The restricted nature of algebraic examples is consistent with \citet[\S2]{spatzier2016rigidity}. The explicit Heisenberg bracket calculation is supplied here as a direct counterexample to the notebook's converse; no general classification of Anosov Lie algebras is claimed.}
\GSfield{Present status}{As of November 2026, the qualified statement above is established; the distinctions from the notebook wording remain essential.}
\GSfield{Assessment}{literal statement: false as written; publication match: partial; verification: read the result; development: annotated; confidence: high}
\end{GSitem}
\begin{GSitem}{1967-140.6}{C}{Are all Anosov diffeomorphisms algebraic up to conjugacy?}
\GSfield{Kind; source}{problem; \GSpage{P000225}}
\GSfield{The notes}{The main problem is \ru{топ. подобность алгебраич.}; small perturbations are already handled by Anosov's theorem.}
\GSfield{Statement}{Is every Anosov diffeomorphism of a compact connected manifold topologically conjugate to an affine Anosov automorphism of an infranilmanifold? The distinction between determining the underlying manifold and conjugating its dynamics is essential. Smooth conjugacy is a stronger assertion and is not the question recorded here.}
\GSfield{Fate}{This repeats 1966-94.4 and reuses GP-1966-01. \citet[Conjecture 2.3 and Theorem 2.4]{spatzier2016rigidity} states the precise conjecture and the Franks--Manning result on infranilmanifolds. \citet[\S1, Background]{damjanovic2026symmetries} still describes algebraicity of general Anosov dynamics as conjectural. Sinai's two-torus case appears as Theorem 5.2 of \citet{sinai1968markov}.}
\GSfield{Present status}{As of November 2026, the general classification remains open. The two-torus case and the Franks--Manning classification when the manifold is already an infranilmanifold are established; structural stability solves only a neighborhood of a known example.}
\GSfield{Assessment}{literal statement: correct; publication match: exact; verification: read the result; development: none; confidence: high; problem \hyperref[prob:GP-1966-01]{GP-1966-01}}
\GSfield{Worth developing}{Further work is worthwhile on explicit manifold classes or additional dynamical hypotheses, since these address the remaining global classification rather than reproving local structural stability.}
\end{GSitem}
\begin{GSitem}{1967-140.7}{F}{Hyperbolicity on the nonwandering set: an incomplete quantifier}
\GSfield{Kind; source}{conjecture; \GSpage{P000225}}
\GSfield{The notes}{The manifold is split into wandering and nonwandering points, and the \(U\)-condition is asserted on the latter, with a margin note saying this too is a hypothesis.}
\GSfield{Statement}{For a homeomorphism, a point is nonwandering when every neighborhood meets one of its positive iterates; its complement is the wandering set. It is false that the nonwandering set of every smooth diffeomorphism is hyperbolic: the identity is an immediate counterexample. A density or genericity conjecture would require a topology, a class of systems, and an approximation quantifier, none of which is stated here.}
\GSfield{Fate}{The terminology belongs to the hyperbolic-dynamics programme surveyed by \citet{smale1967dynamical}, but this fragment does not determine which conjecture was intended. It must not be automatically identified with the earlier Axiom A or periodic-density questions merely because the words nonwandering and hyperbolic occur.}
\GSfield{Present status}{As of November 2026, the literal universal assertion is false and the intended qualified conjecture cannot be uniquely recovered. No new problem-index entry is assigned.}
\GSfield{Assessment}{literal statement: not determinable; publication match: thematic only; verification: indirect; development: none; confidence: medium}
\end{GSitem}
\begin{GSitem}{1967-140.8}{E}{Absolute continuity of holonomy}
\GSfield{Kind; source}{theorem; \GSpage{P000226}}
\GSfield{The notes}{The canonical map between nearby stable plaques, obtained along unstable plaques, need not be smooth but is absolutely continuous; natural measures are said to be carried into one another.}
\GSfield{Statement}{For a \(C^{1+\alpha}\) Anosov diffeomorphism, \(\alpha>0\), stable and unstable holonomies between appropriate smooth transverse plaques are absolutely continuous with respect to the induced Riemannian volume classes. Thus a null set is carried to a null set in both directions. This does not assert equality of the two Riemannian volume measures or differentiability of the holonomy.}
\GSfield{Reconstruction}{The modern \(C^{1+\alpha}\) regularity condition and preservation of measure classes, rather than measure values, are explicit.}
\GSfield{Fate}{\citet[\S1, p.65]{sinai1968markov} states exactly the equivalence of the induced volume classes, attributing the established argument to Anosov and noting that a smooth invariant global measure is not required. The notebook's warning about lack of smoothness is substantive.}
\GSfield{Present status}{As of November 2026, the qualified statement above is established; the distinctions from the notebook wording remain essential.}
\GSfield{Assessment}{literal statement: correct with missing hypotheses; publication match: exact; verification: read the result; development: annotated; confidence: high}
\end{GSitem}
\begin{GSitem}{1967-140.9}{A}{Rectangles from the local product structure}
\GSfield{Kind; source}{definition; \GSpage{P000226}}
\GSfield{The notes}{A parallelogram consists of stable and unstable leaf pieces, canonically identified within each family. Its boundary splits into expanding and contracting parts; finitely many such pieces cover the manifold with disjoint interiors.}
\GSfield{Statement}{A sufficiently small rectangle is a set closed under the local bracket \([x,y]=W^s_{\mathrm{loc}}(x)\cap W^u_{\mathrm{loc}}(y)\), with a product description by stable and unstable plaques. A rectangular partition has finitely many such closed pieces, covers the manifold, and has pairwise disjoint interiors. Its two boundary types come from the boundaries in the two local-product factors. In higher dimensions these rectangles need not be Euclidean polyhedra.}
\GSfield{Reconstruction}{The local bracket and non-Euclidean meaning of rectangle are made explicit.}
\GSfield{Fate}{This is \citet[\S1, Definitions 1.1--1.2]{sinai1968markov}; the published admissibility conditions also control plaque-boundary measure. The schematic figures in the notebook do not assert piecewise linear boundaries.}
\GSfield{Present status}{As of November 2026, the qualified statement above is established; the distinctions from the notebook wording remain essential.}
\GSfield{Assessment}{literal statement: correct with missing hypotheses; publication match: exact; verification: read the result; development: annotated; confidence: high}
\end{GSitem}
\begin{GSitem}{1967-140.10}{A}{Markov boundary inclusions}
\GSfield{Kind; source}{definition; \GSpage{P000226}}
\GSfield{The notes}{The partition is Markov when, for some \(m>0\), \(\Gamma_p(T^{-m}\alpha)\subset\Gamma_p(\alpha)\) and \(\Gamma_c(T^m\alpha)\subset\Gamma_c(\alpha)\).}
\GSfield{Statement}{For a sufficiently small rectangular partition, use the notebook's stable/unstable boundary conventions and require those two inclusions. In the usual one-step formulation, if \(x\in\operatorname{int}R_i\) and \(Tx\in\operatorname{int}R_j\), then \(T(W^s(x,R_i))\subset W^s(Tx,R_j)\) and \(T(W^u(x,R_i))\supset W^u(Tx,R_j)\). The direction of inclusion is essential.}
\GSfield{Reconstruction}{The plaque formulation supplies a coordinate-free interpretation while retaining the recorded boundary convention.}
\GSfield{Fate}{\citet[Definition 1.3]{sinai1968markov} gives the same iterate-based boundary definition. \citet[Definitions 1.1 and 2.2]{sinai1968construction} uses it in the existence proof.}
\GSfield{Present status}{As of November 2026, the qualified statement above is established; the distinctions from the notebook wording remain essential.}
\GSfield{Assessment}{literal statement: correct; publication match: exact; verification: read the result; development: annotated; confidence: high}
\end{GSitem}
\begin{GSitem}{1967-140.11}{A}{Existence of arbitrarily fine Markov partitions}
\GSfield{Kind; source}{theorem; \GSpage{P000226}}
\GSfield{The notes}{The technical theorem asserts many Markov partitions: each given rectangular partition can be approximated within any \(\varepsilon\), allowing an iterate \(m\). The figures constrain how endpoints lie on opposite boundaries.}
\GSfield{Statement}{Every Anosov diffeomorphism of a compact manifold has Markov partitions with arbitrarily small diameters. The construction starts with controlled local-product coverings, adjusts stable and unstable boundary pieces by convergent displacements, and forms a compatible partition. An unrestricted distance on every possible preassigned partition is not specified by the notebook and should not be inferred from the small-mesh theorem alone.}
\GSfield{Reconstruction}{The definite proved small-mesh theorem is separated from the informal universal approximation wording.}
\GSfield{Fate}{The theorem is announced in \citet[\S1, p.66]{sinai1968markov}, received 20 October 1967. The detailed proof is \citet[\S1 in dimension two; \S2 in general]{sinai1968construction}, received 3 November 1967 and published in issue 3 of 1968. Its controlled covering-distance construction explains the stronger approximation language in the talk, without defining the notebook's unspecified metric on all partitions.}
\GSfield{Present status}{As of November 2026, the qualified statement above is established; the distinctions from the notebook wording remain essential.}
\GSfield{Assessment}{literal statement: correct with missing hypotheses; publication match: partial; verification: read the result; development: annotated; confidence: high}
\end{GSitem}
\begin{GSitem}{1967-140.12}{A}{Reducing an iterate-Markov partition to one step}
\GSfield{Kind; source}{construction; \GSpage{P000226-P000227}}
\GSfield{The notes}{Take \(\alpha_1=\alpha\vee T^{-1}\alpha\vee\cdots\vee T^{-m+1}\alpha\); the union of transformed boundary pieces shows that \(m=1\) now suffices.}
\GSfield{Statement}{If the rectangular partition is Markov for the boundary inclusions at iterate \(m\), its indicated common refinement, interpreted by the local-product components of nonempty intersections, satisfies the one-step inclusions. The boundary is contained in the union of the finitely many transformed boundaries, so the iterate condition closes the last term under one further iterate.}
\GSfield{Reconstruction}{Taking the relevant nonempty local-product components avoids treating an arbitrary disconnected intersection as a single small chart.}
\GSfield{Fate}{This construction is stated explicitly in \citet[\S1, immediately after the existence theorem]{sinai1968markov}. It is the same refinement written across the notebook page break.}
\GSfield{Present status}{As of November 2026, the qualified statement above is established; the distinctions from the notebook wording remain essential.}
\GSfield{Assessment}{literal statement: correct; publication match: exact; verification: read the result; development: annotated; confidence: high}
\end{GSitem}
\begin{GSitem}{1967-140.13}{A}{Coding by admissible itineraries}
\GSfield{Kind; source}{method; \GSpage{P000227}}
\GSfield{The notes}{\(\pi_{ij}=1\) when \(V_i\cap TV_j\ne\varnothing\). An admissible sequence is mapped to \(\bigcap_nT^n\overline V_{j_n}\), a single point for a sufficiently fine partition. Boundary points can have several codes; uniqueness holds on a set of second category.}
\GSfield{Statement}{Using open rectangle interiors in the transition rule, a sufficiently fine Markov partition gives a continuous surjection from a two-sided shift of finite type onto the manifold. The intersection of the closed itinerary rectangles is one point, and the coding intertwines the shift with \(T\) or \(T^{-1}\) according to the time-index convention. Off the union of all iterates of partition boundaries the itinerary is unique; this is a residual set when the boundaries are nowhere dense. Global injectivity is not asserted.}
\GSfield{Reconstruction}{The interiors and time orientation are specified; nonempty intersections caused only by shared closed boundaries are not automatically allowed transitions.}
\GSfield{Fate}{This is the construction and conclusion of \citet[\S3, Theorem 3.1]{sinai1968markov}. The notebook already recognizes the boundary obstruction, so replacing its coding by a homeomorphism would lose a genuine qualification.}
\GSfield{Present status}{As of November 2026, the qualified statement above is established; the distinctions from the notebook wording remain essential.}
\GSfield{Assessment}{literal statement: correct with missing hypotheses; publication match: exact; verification: read the result; development: annotated; confidence: high}
\end{GSitem}
\begin{GSitem}{1967-140.14}{A}{Primitivity and density of the invariant leaves}
\GSfield{Kind; source}{theorem; \GSpage{P000227}}
\GSfield{The notes}{A strictly positive power of \(\Pi\) is required, and this is said to be equivalent to density of every stable and unstable leaf.}
\GSfield{Statement}{Within the Markov coding of a connected Anosov system, the mixing case has a primitive transition matrix. Sinai's theorem identifies this condition with density of every stable and every unstable leaf. His word transitive in this part of the paper is defined by the positive-power matrix condition, which is stronger than mere irreducibility of a general finite matrix. It is an additional dynamical hypothesis, not part of the definition of every Anosov diffeomorphism.}
\GSfield{Reconstruction}{The original terminology is translated into primitivity/mixing and kept separate from the unresolved general transitivity questions.}
\GSfield{Fate}{\citet[\S3, Definition 3.2 and Theorem 3.2]{sinai1968markov} gives the recorded equivalence. The precise published convention prevents accidental conflation of periodic irreducible chains with primitive ones.}
\GSfield{Present status}{As of November 2026, the qualified statement above is established; the distinctions from the notebook wording remain essential.}
\GSfield{Assessment}{literal statement: correct with missing hypotheses; publication match: exact; verification: read the result; development: annotated; confidence: high}
\end{GSitem}
\begin{GSitem}{1967-140.15}{A}{Holonomy-invariant measures on stable and unstable leaves}
\GSfield{Kind; source}{construction; \GSpage{P000227}}
\GSfield{The notes}{Each stable leaf has a sigma-finite measure invariant under canonical holonomy, with \(\lambda_c(TA)=h\lambda_c(A)\), \(h<1\), an algebraic number. The finite-dimensional construction is \(h^{-1}\vec\lambda=\Pi\vec\lambda\).}
\GSfield{Statement}{For the mixing Anosov system under discussion, there are nonzero locally finite, hence sigma-finite, stable-leaf measures invariant under unstable holonomy, with \(\mu^s(TA)=\lambda^{-1}\mu^s(A)\), where \(\lambda>1\) is the Perron eigenvalue of a one-step Markov matrix. The analogous unstable-leaf measures scale by \(\lambda\) under \(T\). Initial plaque masses satisfy the appropriate Perron eigenvector equation, and subdivision extends them to the leaves. Thus the scaling factors are algebraic. They need not be equivalent to Riemannian leaf volume.}
\GSfield{Reconstruction}{The notebook's \(h\) is \(\lambda^{-1}\), a scaling factor rather than entropy. The separate unstable family and singularity possibility are made explicit.}
\GSfield{Fate}{This is exactly \citet[\S5, Theorem 5.1, pp.84--86]{sinai1968markov}; it is a general statement in that mixing setting, preceding the two-torus specialization in Theorem 5.2. For linear toral examples the leaf measures arise from Haar measure, as the notebook remarks.}
\GSfield{Present status}{As of November 2026, the qualified statement above is established; the distinctions from the notebook wording remain essential.}
\GSfield{Assessment}{literal statement: correct with missing hypotheses; publication match: exact; verification: read the result; development: annotated; confidence: high}
\end{GSitem}
\begin{GSitem}{1967-140.16}{A}{An invariant measure with a finite Markov model}
\GSfield{Kind; source}{theorem; \GSpage{P000227}}
\GSfield{The notes}{As a corollary there is an invariant measure for which \(T\) is a finite Markov chain. The illustrated construction then specializes to dimension two.}
\GSfield{Statement}{For a mixing finite Markov coding, the Parry probability pushes forward to an invariant probability. When the coding boundaries have zero measure, the resulting system is isomorphic modulo zero to the finite-state stationary Markov shift. For a hyperbolic toral automorphism, Haar probability has exactly this Markov property and is the maximal-entropy measure. For nonlinear dynamics the constructed maximal-entropy measure need not be smooth volume or the measure with absolutely continuous unstable conditionals.}
\GSfield{Reconstruction}{The precise verified toral theorem is separated from the general coding implication. Maximal-entropy measure is not silently identified with the different invariant measure constructed from smooth unstable conditionals.}
\GSfield{Fate}{\citet[Theorem 5.4 and its proof]{sinai1968markov} proves the Haar/Markov assertion for every hyperbolic toral automorphism; Theorem 5.1 supplies the leaf-measure construction used in the surrounding talk. The boundary-null and measure-class distinctions explain what is required when the notebook states its broader corollary without qualifications.}
\GSfield{Present status}{As of November 2026, the qualified statement above is established; the distinctions from the notebook wording remain essential.}
\GSfield{Assessment}{literal statement: correct with missing hypotheses; publication match: partial; verification: read the result; development: annotated; confidence: high}
\end{GSitem}
\begin{GSitem}{1967-140.17}{E}{Entropy, the Perron root, and the repeated metric classification}
\GSfield{Kind; source}{theorem; \GSpage{P000228}}
\GSfield{The notes}{The page first gives \(h(T)=\log\lambda\), then writes \(h(T)=\lambda\), and returns to metric isomorphism through the standard symbolic model.}
\GSfield{Statement}{For a hyperbolic two-torus automorphism, Haar entropy is \(\log\rho(A)\). More generally, in the mixing Anosov Markov construction, \(h_{\mathrm{top}}(T)=\log\lambda\), where \(\lambda\) is the Perron root and the stable-leaf scaling factor is \(\lambda^{-1}\). The equality \(h(T)=\lambda\) on the next notebook line omits the logarithm if the same symbol \(\lambda\) is intended. Equal Haar entropy classifies the two-torus automorphisms metrically as in 1967-138.2.}
\GSfield{Reconstruction}{The repeated theorem is separated from the erroneous second entropy line and from the stable-leaf scaling symbol.}
\GSfield{Fate}{\citet[Theorem 5.3]{sinai1968markov} identifies topological entropy with the logarithm of the leaf/Perron factor. The metric classification itself is the already published Main Theorem of \citet{adlerweiss1967entropy}; its recurrence here remains expository.}
\GSfield{Present status}{As of November 2026, the qualified statement above is established; the distinctions from the notebook wording remain essential.}
\GSfield{Assessment}{literal statement: false as written; publication match: partial; verification: read the result; development: annotated; confidence: high; problem \hyperref[prob:GP-1967-03]{GP-1967-03}}
\end{GSitem}
\begin{GSitem}{1967-140.18}{E}{The two-rectangle torus partition and its connected refinement}
\GSfield{Kind; source}{construction; \GSpage{P000228}}
\GSfield{The notes}{Two parallelograms have sides along eigenlines through zero. The coarse itinerary can contain many points. Intersections \(V_i\cap TV_j\) are refined into connected components in four groups of sizes \(p,q,r,s\); one of the four displayed intersections is repeated.}
\GSfield{Statement}{For a real hyperbolic two-torus automorphism, a suitable two-rectangle partition is built from segments in its stable and unstable eigendirections. The component refinement of all four intersections \(V_i\cap TV_j\), \(i,j\in\{1,2\}\), gives a generating Markov partition with the four-block transition structure in 1967-138.9. The two coarse labels alone need not generate pointwise. All four ordered pairs are required; the duplicated off-diagonal line in the notes cannot be retained literally.}
\GSfield{Reconstruction}{The four intersections are stated as the complete indexed family, without pretending the ambiguous overwritten second line identifies the author's intended ordering of \(q,r\).}
\GSfield{Fate}{This is the geometric step in \citet[\S5]{adlerweiss1967entropy}, and it is expressly credited to Adler and Weiss in \citet[\S1, p.66]{sinai1968markov}. The source provides a proof-level exact match, not just thematic similarity to later symbolic dynamics.}
\GSfield{Present status}{As of November 2026, the qualified statement above is established; the distinctions from the notebook wording remain essential.}
\GSfield{Assessment}{literal statement: false as written; publication match: exact; verification: read the result; development: annotated; confidence: high; problem \hyperref[prob:GP-1967-03]{GP-1967-03}}
\end{GSitem}

\GStalk{141}{}{}{P000228-P000230}{Calendar heading: 1968}
The actual heading is \ru{1968 год} on P000229. It begins the calendar-year continuation of the 1967--1968 academic seminar; the catalogue range overlaps Sinai's preceding talk. There is no independent mathematical item in this heading.


\GStalk{142}{1968-02-12}{I. M. Gelfand}{P000228-P000230}{Positive functionals on countably normed algebras}
On P000229 Gelfand asks which countably normed algebras admit a character-space integral representation for positive functionals. The comparison is with the commutative Banach-algebra representation theorem and Bochner's theorem. Completeness, continuity, commutativity and the class of measures are suppressed in the notebook; those distinctions control the modern moment problem.

\begin{GSitem}{1967-142.1}{E}{Countably normed algebra and continuity of multiplication}
\GSfield{Kind; source}{definition; \GSpage{P000229}}
\GSfield{The notes}{The notes write \(\|xy\|_i\leq\|x\|_k\|y\|_\ell\) for a countably normed ring \(R\).}
\GSfield{Statement}{A countably normed complex algebra is an algebra with a Hausdorff locally convex topology defined by a countable family of norms. The displayed condition, interpreted as: for each \(i\) there exist \(k,\ell\) such that it holds for all \(x,y\), implies joint continuity of multiplication. Completeness makes this a Fr\'echet algebra. For the subsequent moment question one additionally needs a commutative algebra, a continuous involution and a specified continuous dual.}
\GSfield{Reconstruction}{The quantifiers and the distinction between complete and incomplete algebras are supplied. No assumption that every norm is submultiplicative is inserted.}
\GSfield{Fate}{The definition is the setting of the question, not a new theorem. Modern comparisons distinguish the topology from the positivity cone; see \citet[\S\S2.2--2.3]{infusino2024moment}.}
\GSfield{Present status}{As of November 2026, this remains standard terminology, with completeness and involution hypotheses to be stated separately.}
\GSfield{Assessment}{literal statement: correct with missing hypotheses; publication match: exact; verification: read the result; development: annotated; confidence: high}
\end{GSitem}
\begin{GSitem}{1967-142.2}{E}{Positive functionals as integrals over Hermitian characters}
\GSfield{Kind; source}{theorem; \GSpage{P000229}}
\GSfield{The notes}{A theorem labelled Gelfand--Naimark says \(f(x^*x)\geq0\Rightarrow f(x)=\int x(M)\,d\sigma(M)\), with \(M\) symmetric maximal ideals.}
\GSfield{Statement}{For a unital commutative complex \(C^*\)-algebra \(A\), every positive linear functional \(L\) is uniquely \(L(a)=\int_{X(A)}\chi(a)\,d\mu(\chi)\), where \(X(A)\) is the compact space of unital characters and \(\mu\) is a finite positive regular Borel measure. Every character is Hermitian: \(\chi(a^*)=\overline{\chi(a)}\).}
\GSfield{Reconstruction}{The precise \(C^*\)-algebra version is used to avoid attributing the assertion to arbitrary normed noncommutative rings. The notebook's symmetric maximal ideals correspond to kernels of Hermitian characters in the commutative setting.}
\GSfield{Fate}{This is the classical representation paradigm explicitly invoked by the notes. Its infinite-dimensional moment analogues require additional positivity or continuity assumptions; \citet[Abstract]{schmudgen2018moment}.}
\GSfield{Present status}{As of November 2026, the stated classical version is established; the literal phrase ``normed ring does not specify adequate hypotheses.}
\GSfield{Assessment}{literal statement: correct with missing hypotheses; publication match: partial; verification: publisher metadata checked; development: annotated; confidence: medium}
\end{GSitem}
\begin{GSitem}{1967-142.3}{F}{Axioms ensuring moment representations in countably normed algebras}
\GSfield{Kind; source}{question; \GSpage{P000229}}
\GSfield{The notes}{\ru{Перенос --- ? (Не всегда: надо охар. аксиомат. класс).} The question asks when the preceding representation extends.}
\GSfield{Statement}{Fix the convention that \(A\) is a commutative unital complex Fr\'echet \(*\)-algebra: its topology is Hausdorff, complete and defined by countably many norms, multiplication is jointly continuous, and the involution is continuous. Let \(X_h(A)\) be its continuous unital Hermitian characters with the weak topology. Characterize those \(A\) for which every continuous complex-linear functional \(L\) satisfying \(L(a^*a)\ge0\) for every \(a\in A\) has a finite positive Radon measure \(\mu_L\) on \(X_h(A)\) such that \(L(a)=\int\chi(a)\,d\mu_L(\chi)\) for all \(a\in A\). All character evaluations in this identity must be integrable. This is a universal property of the algebra, not a question about one chosen \(L\).}
\GSfield{Reconstruction}{Commutativity, the unit, completeness, continuity of the involution and restriction to continuous characters fix one explicit version of the notebook's proposed axiomatization. These choices are supplied by the editor; the notes do not uniquely determine all of them. Cylinder-measure representation is weaker than the stated Radon representation.}
\GSfield{Fate}{\citet[Abstract]{schmudgen2018moment}, circulated in December 2017, treats cylinder-measure representations under stronger positivity or Carleman assumptions. \citet[\S\S2.2--2.3]{infusino2024moment}, circulated in January 2023 and published in 2024, gives Radon criteria for algebras generated by a nuclear space. These results address specified hypotheses, not every algebra in the notebook.}
\GSfield{Present status}{As of November 2026, the notebook does not determine a unique algebraic/topological category for this axiomatization request. The stated universal Radon-representation property is an explicit editorial comparison problem; the checked nuclear and cylinder-measure results do not certify that it is a recognized unresolved conjecture. No precise open problem is inferred from this occurrence.}
\GSfield{Assessment}{literal statement: not determinable; publication match: thematic only; verification: read the result; development: annotated; confidence: medium}
\end{GSitem}
\begin{GSitem}{1967-142.4}{E}{Convolution algebra and Bochner representation}
\GSfield{Kind; source}{remark; \GSpage{P000229}}
\GSfield{The notes}{The example is \(L^1\) with convolution and \(x^*(t)=\overline{x(-t)}\), followed by the Fourier integral of a positive measure.}
\GSfield{Statement}{For the additive group \(\mathbb R\), a bounded continuous positive-definite function \(\varphi\) is the Fourier transform \(\varphi(t)=\int_{\mathbb R}e^{it\lambda}\,d\mu(\lambda)\) of a unique finite positive measure. It defines a bounded positive functional on \(L^1(\mathbb R)\) by \(L(f)=\int f(t)\varphi(t)\,dt\), up to the chosen Fourier-sign convention.}
\GSfield{Reconstruction}{The notebook does not distinguish the function being transformed from the argument of the functional; the two roles are separated here.}
\GSfield{Fate}{The notes identify this as Bochner's already known theorem. It supplies a concrete spectral model for the preceding moment question, rather than an additional generalization.}
\GSfield{Present status}{As of November 2026, the example is standard and does not settle the countably normed case.}
\GSfield{Assessment}{literal statement: correct with missing hypotheses; publication match: none found; verification: indirect; development: annotated; confidence: high}
\end{GSitem}

\GStalk{143}{1968-02-12}{I. N. Bernstein}{P000228-P000233}{Grothendieck topologies and sheaf cohomology}
The substantive text is P000229--P000231. It introduces fibre products, covering axioms, presheaves, sheaves, the augmented \v{C}ech complex and derived-functor cohomology. The group example is only announced. This anticipates Geronimus's September exposition of Mumford, whose 1965 text already presents this machinery.

\begin{GSitem}{1967-143.1}{E}{Fibre product by its universal property}
\GSfield{Kind; source}{definition; \GSpage{P000229}}
\GSfield{The notes}{An object \(V=V_1\times_U V_2\) is drawn with commuting projections and a unique factorization from every test object.}
\GSfield{Statement}{Given \(V_1\to U\leftarrow V_2\), the fibre product represents the functor \(W\mapsto\operatorname{Hom}(W,V_1)\times_{\operatorname{Hom}(W,U)}\operatorname{Hom}(W,V_2)\). It is unique up to unique isomorphism, when it exists, and depends on the two morphisms to \(U\), not only on the three objects.}
\GSfield{Fate}{Expository category theory; compare \citet[p.~38, footnote]{mumford1965picard}.}
\GSfield{Present status}{As of November 2026, this is the standard definition.}
\GSfield{Assessment}{literal statement: correct; publication match: exact; verification: read the result; development: annotated; confidence: high}
\end{GSitem}
\begin{GSitem}{1967-143.2}{E}{Grothendieck pretopology}
\GSfield{Kind; source}{definition; \GSpage{P000229-P000230}}
\GSfield{The notes}{The notebook specifies a category, distinguished covering families, fibre products, identity covers, composition and base change.}
\GSfield{Statement}{A pretopology is a class of families \(\{U_i\to U\}\) containing isomorphisms as covers, stable under composition of covering families and under pullback along every morphism, whenever the required pullbacks exist. It generates a Grothendieck topology of covering sieves.}
\GSfield{Reconstruction}{``Pretopology distinguishes these covering-family axioms from the equivalent sieve-based formulation; it does not change the notes' content.}
\GSfield{Fate}{Already standard; \citet[p.~38]{mumford1965picard}.}
\GSfield{Present status}{As of November 2026, the definition remains in use.}
\GSfield{Assessment}{literal statement: correct; publication match: exact; verification: read the result; development: annotated; confidence: high}
\end{GSitem}
\begin{GSitem}{1967-143.3}{E}{Local-homeomorphism site of a space}
\GSfield{Kind; source}{construction; \GSpage{P000230}}
\GSfield{The notes}{For an ordinary topological space, objects are spaces locally homeomorphic to it; coverings are surjective families.}
\GSfield{Statement}{For a space \(X\), take local homeomorphisms \(U\to X\), maps over \(X\), and jointly surjective families as covers. Fibre products over \(X\) exist and remain local homeomorphisms. Ordinary open subsets form a basis of this site, yielding the same category of sheaves as the usual topology of \(X\).}
\GSfield{Reconstruction}{The structural map to the fixed base \(X\) and joint surjectivity are made explicit.}
\GSfield{Fate}{Expository enlargement of ordinary topology; \citet[\S1 and pp.~44--45]{mumford1965picard}.}
\GSfield{Present status}{As of November 2026, the comparison is standard.}
\GSfield{Assessment}{literal statement: correct; publication match: exact; verification: read the result; development: annotated; confidence: high}
\end{GSitem}
\begin{GSitem}{1967-143.4}{E}{Fibre products of circular coverings}
\GSfield{Kind; source}{remark; \GSpage{P000230}}
\GSfield{The notes}{The circular example says \(U\times_U U=U\) and compares coverings of degrees two and three.}
\GSfield{Statement}{If both maps \(S^1\to S^1\) have degree two, their pullback is the disjoint union of two circles, each mapping by degree two to either factor. For degrees \(m,n>0\), it has \(\gcd(m,n)\) components and each component covers the common base with degree \(\operatorname{lcm}(m,n)\). In particular degrees two and three give one degree-six covering.}
\GSfield{Reconstruction}{The equation \(U\times_U U\cong U\) is valid for identity structure maps; it is false for two copies of the degree-two cover suggested by the drawing. The unclear arrow labels are not silently repaired.}
\GSfield{Fate}{The corrected calculation follows directly from the pullback definition: pairs satisfy \(z^m=w^n\). It is an elementary emendation, not later research.}
\GSfield{Present status}{As of November 2026, the mathematical alternatives are settled, but the exact unlabelled arrows in the drawing remain ambiguous.}
\GSfield{Assessment}{literal statement: not determinable; publication match: none found; verification: indirect; development: annotated; confidence: high}
\end{GSitem}
\begin{GSitem}{1967-143.5}{E}{Presheaves and the sheaf equalizer condition}
\GSfield{Kind; source}{definition; \GSpage{P000230-P000231}}
\GSfield{The notes}{A presheaf assigns abelian groups contravariantly; compatible local sections must glue uniquely for a sheaf.}
\GSfield{Statement}{An abelian presheaf on a site \(\mathcal C\) is a functor \(F:\mathcal C^{\mathrm{op}}\to\mathrm{Ab}\). It is a sheaf when, for every cover, \(F(U)\to\prod_iF(U_i)\rightrightarrows\prod_{i,j}F(U_i\times_UU_j)\) is an equalizer. This combines uniqueness and existence of descent.}
\GSfield{Fate}{Expository; \citet[p.~39]{mumford1965picard}.}
\GSfield{Present status}{As of November 2026, this is the standard sheaf axiom.}
\GSfield{Assessment}{literal statement: correct; publication match: exact; verification: read the result; development: annotated; confidence: high}
\end{GSitem}
\begin{GSitem}{1967-143.6}{E}{Augmented Cech complex}
\GSfield{Kind; source}{construction; \GSpage{P000231}}
\GSfield{The notes}{The groups are indexed \(C^0=P(U), C^1=\prod P(U_i), C^2=\prod P(U_i\times_UU_j)\), with alternating restrictions.}
\GSfield{Statement}{The augmented complex is \(F(U)\to\check C^0(\mathcal U,F)\to\check C^1(\mathcal U,F)\to\cdots\), where \(\check C^r=\prod_{i_0,\ldots,i_r}F(U_{i_0}\times_U\cdots\times_UU_{i_r})\) and the differential is the alternating sum of omitted-factor pullbacks. A sheaf makes the first two terms exact; higher exactness needs an acyclic cover or another additional hypothesis.}
\GSfield{Reconstruction}{The notebook shifts the augmented complex by one compared with common modern \v{C}ech notation. No assertion that every cover computes derived cohomology is added.}
\GSfield{Fate}{The construction belongs to the pre-existing sheaf formalism; \citet[\S1]{mumford1965picard}.}
\GSfield{Present status}{As of November 2026, the indexing convention is harmless when stated.}
\GSfield{Assessment}{literal statement: correct; publication match: exact; verification: read the result; development: annotated; confidence: high}
\end{GSitem}
\begin{GSitem}{1967-143.7}{E}{Derived-functor cohomology of a site}
\GSfield{Kind; source}{method; \GSpage{P000231}}
\GSfield{The notes}{Global sections are left exact and their derived functors define cohomology.}
\GSfield{Statement}{For a site with terminal object \(e\), \(\Gamma(F)=F(e)\) on abelian sheaves is left exact. The category has enough injectives, and \(H^q(\mathcal C,F)=R^q\Gamma(F)\). Without a representable terminal object one uses the terminal sheaf to define global sections.}
\GSfield{Reconstruction}{The existence of enough injectives and the terminal-object convention are supplied.}
\GSfield{Fate}{Expository; \citet[pp.~37,39]{mumford1965picard}; see also \citet[\S\S2--3]{grothendieck1957homological}.}
\GSfield{Present status}{As of November 2026, the definition is established.}
\GSfield{Assessment}{literal statement: correct; publication match: exact; verification: read the result; development: annotated; confidence: high}
\end{GSitem}
\begin{GSitem}{1967-143.8}{F}{An announced finite-group example}
\GSfield{Kind; source}{remark; \GSpage{P000231}}
\GSfield{The notes}{The last line announces a topology associated with representations of a finite group, without giving its objects or covers.}
\GSfield{Statement}{No precise theorem is recoverable from this line alone. A natural later comparison is the site of \(\Gamma\)-sets with jointly surjective covers, whose abelian sheaves are \(\Gamma\)-modules; this is explicitly constructed in record 151.}
\GSfield{Reconstruction}{The \(\Gamma\)-set interpretation is a cross-reference, not a replacement for the missing February construction.}
\GSfield{Fate}{Compare \citet[pp.~39--41]{mumford1965picard} and occurrences 1968-151.2--1968-151.4; the isolated phrase does not establish that this was the exact category intended.}
\GSfield{Present status}{As of November 2026, the later example is identified, while this announcement remains incomplete.}
\GSfield{Assessment}{literal statement: not determinable; publication match: partial; verification: read the result; development: annotated; confidence: medium}
\end{GSitem}

\GStalk{144}{1968-03-11}{I. M. Gelfand; A. A. Kirillov}{P000231-P000234}{Lie fields and virtual Laplace operators}
P000232--P000233 develop the birational study of enveloping algebras via differential operators on the basic affine space. A central extension by virtual Laplace operators is asserted to have a Weyl division ring. This is the close precursor of the 1969 Gelfand--Kirillov paper, whereas the unextended general conjecture must be separated from it. The talk ends with questions about division-ring extensions and global holomorphic differential operators.

\begin{GSitem}{1967-144.1}{E}{Enveloping algebra, centre and differential-operator realization}
\GSfield{Kind; source}{definition; \GSpage{P000232}}
\GSfield{The notes}{\(U(\mathfrak g)\) is generated by \(e_i\) with \([e_i,e_j]=\sum c_{ij}^ke_k\); central elements are called Laplace or Casimir operators. A group action yields differential operators, faithfully on the group itself.}
\GSfield{Statement}{For a finite-dimensional complex Lie algebra \(\mathfrak g\), \(U(\mathfrak g)=T(\mathfrak g)/(xy-yx-[x,y])\). If a Lie group with this Lie algebra acts smoothly or algebraically, differentiation gives a representation by differential operators, which may have a kernel. The representation by invariant differential operators on the group is faithful. For connected \(G\), the image of \(Z(U(\mathfrak g))\) commutes with the differentiated and group actions.}
\GSfield{Reconstruction}{Connectedness is relevant to passage from infinitesimal to group invariance. PBW supplies the faithfulness assertion and the filtration underlying the later fractions.}
\GSfield{Fate}{Background recalled in \citet[Introduction]{gelfandkirillov1966corps} and \citet[\S\S2--3]{gelfandkirillov1969structure}.}
\GSfield{Present status}{As of November 2026, these are standard algebraic constructions.}
\GSfield{Assessment}{literal statement: correct; publication match: exact; verification: read the result; development: annotated; confidence: high}
\end{GSitem}
\begin{GSitem}{1967-144.2}{E}{Faithfulness on the basic affine space}
\GSfield{Kind; source}{theorem; \GSpage{P000232}}
\GSfield{The notes}{The action of \(U(\mathfrak g)\) on \(N_-\backslash G\) is said to be faithful.}
\GSfield{Statement}{Let \(G\) be a simply connected complex semisimple algebraic group and \(N_-\) a maximal unipotent subgroup. Right translation on \(\mathcal O(N_-\backslash G)\) differentiates to an injective homomorphism \(U(\mathfrak g)\to\mathcal D(N_-\backslash G)\). The coordinate ring contains every finite-dimensional irreducible representation, and these representations separate elements of \(U(\mathfrak g)\).}
\GSfield{Reconstruction}{The algebraic category, regular functions and simply connected group are specified. ``Nilpotent group is replaced by the intended maximal unipotent subgroup.}
\GSfield{Fate}{It is part of the representation-theoretic setup in \citet[\S\S3--6]{gelfandkirillov1969structure}, not the new fraction-field theorem.}
\GSfield{Present status}{As of November 2026, the statement is established.}
\GSfield{Assessment}{literal statement: correct with missing hypotheses; publication match: exact; verification: read the result; development: annotated; confidence: high}
\end{GSitem}
\begin{GSitem}{1967-144.3}{E}{Finite-translation functions and the compact-group completeness theorem}
\GSfield{Kind; source}{remark; \GSpage{P000232}}
\GSfield{The notes}{A function is called regular when its translates span a finite-dimensional space; on a compact group such functions approximate all functions. The sine function on a line is an example.}
\GSfield{Statement}{On a compact group, finite sums of matrix coefficients of continuous finite-dimensional representations are uniformly dense in continuous functions. For a complex semisimple algebraic group, regular functions are finite sums of algebraic matrix coefficients. The finite-translation characterization must be understood with the relevant continuity, holomorphicity or algebraicity conditions; on a real additive line, sine is representative but not a polynomial regular function.}
\GSfield{Reconstruction}{The notebook uses ``regular across analytic and algebraic categories. These two meanings are distinguished.}
\GSfield{Fate}{Classical representation-theoretic background; the algebraic formulation is recalled in \citet[\S3]{gelfandkirillov1969structure}.}
\GSfield{Present status}{As of November 2026, the density and algebraic matrix-coefficient statements are established in their specified categories.}
\GSfield{Assessment}{literal statement: correct with missing hypotheses; publication match: exact; verification: read the result; development: annotated; confidence: high}
\end{GSitem}
\begin{GSitem}{1967-144.4}{E}{Affinization and rationality of the basic affine space}
\GSfield{Kind; source}{construction; \GSpage{P000232}}
\GSfield{The notes}{Take the invariant ring of functions \(f(n_-g)=f(g)\), its spectrum \(X\), and note that \(N_-\backslash G\) is birational to projective space.}
\GSfield{Statement}{The basic affine space \(A=N_-\backslash G\) is quasi-affine; \(X=\operatorname{Spec}\mathcal O(G)^{N_-}\) is its affinization. The big Bruhat cell gives a dense open subset of \(A\) isomorphic to \(T\times N_+\), hence its function field is purely transcendental of degree \(\operatorname{rank}G+|\Phi^+|\). It is consequently rational, as is \(X\).}
\GSfield{Reconstruction}{The quasi-affine quotient is not identified pointwise with its affine completion. ``Birational to projective space means rationality, not compactness.}
\GSfield{Fate}{Explicit setup in \citet[\S\S3--4]{gelfandkirillov1969structure}.}
\GSfield{Present status}{As of November 2026, this construction and rationality are established.}
\GSfield{Assessment}{literal statement: correct; publication match: exact; verification: read the result; development: annotated; confidence: high}
\end{GSitem}
\begin{GSitem}{1967-144.5}{E}{Ore fractions of an enveloping algebra}
\GSfield{Kind; source}{construction; \GSpage{P000232}}
\GSfield{The notes}{The notes stress the nontrivial existence of a division ring of quotients and that single fractions \(QP^{-1}\) suffice.}
\GSfield{Statement}{PBW implies that \(U(\mathfrak g)\) is a left and right Noetherian domain. It satisfies both Ore conditions and embeds in its classical division ring \(D(\mathfrak g)\). Each element can be represented as a single right fraction \(QP^{-1}\), and also as a single left fraction, with nonzero denominator.}
\GSfield{Reconstruction}{The Ore condition, rather than formal commutative fraction manipulation, justifies the two-step description.}
\GSfield{Fate}{This repeats the theme of 1965-80.1--1965-80.3; see \citet[\S1]{gelfandkirillov1966corps}.}
\GSfield{Present status}{As of November 2026, the construction is standard for finite-dimensional Lie algebras in characteristic zero.}
\GSfield{Assessment}{literal statement: correct; publication match: exact; verification: read the result; development: annotated; confidence: high}
\end{GSitem}
\begin{GSitem}{1967-144.6}{E}{Weyl division rings with central parameters}
\GSfield{Kind; source}{definition; \GSpage{P000232-P000233}}
\GSfield{The notes}{\(D_{n,k}\) uses \(x_i,\partial/\partial x_i\) and commuting parameters \(y_1,\ldots,y_k\); the two integers are invariants.}
\GSfield{Statement}{Over \(\mathbb C\), define \(D_{n,k}=\operatorname{Frac}(A_n(\mathbb C)\otimes\mathbb C[y_1,\ldots,y_k])\), with \([\partial_i,x_j]=\delta_{ij}\) and all other generator commutators zero. Its centre is \(\mathbb C(y_1,\ldots,y_k)\). For a semisimple Lie algebra the proposed parameters are \(k=\operatorname{rank}\mathfrak g\), \(n=|\Phi^+|\), so \(2n+k=\dim\mathfrak g\).}
\GSfield{Reconstruction}{The field and algebraic independence of central parameters are stated explicitly.}
\GSfield{Fate}{Defined and its numerical invariants studied in \citet[\S\S1,4]{gelfandkirillov1966corps}; compare 1965-80.5--1965-80.6.}
\GSfield{Present status}{As of November 2026, the standard division rings and the distinction of their parameters are established.}
\GSfield{Assessment}{literal statement: correct; publication match: exact; verification: read the result; development: annotated; confidence: high}
\end{GSitem}
\begin{GSitem}{1967-144.7}{B}{The unextended Gelfand--Kirillov conjecture for semisimple algebras}
\GSfield{Kind; source}{conjecture; \GSpage{P000233}}
\GSfield{The notes}{The conjectural isomorphism \(D(\mathfrak g)\simeq D_{n,k}\) is described as almost proved; a later parenthesis says that it is unknown for \(U\), except the unimodular case.}
\GSfield{Statement}{For a complex semisimple Lie algebra \(\mathfrak g\), ask whether \(\operatorname{Frac}U(\mathfrak g)\cong D_{|\Phi^+|,\operatorname{rank}\mathfrak g}\). This assertion concerns the original division ring, with no finite extension of its centre.}
\GSfield{Reconstruction}{The historical term unimodular is read in its matrix-group sense \(SL_m\), not as the modern condition \(\operatorname{tr}(\operatorname{ad}x)=0\), which every semisimple Lie algebra has.}
\GSfield{Fate}{Type \(A\) was proved in \citet[\S3]{gelfandkirillov1966corps}. \citet[Theorem~1.1]{premet2010modular} disproved the conjecture for simple types \(B_\ell\;(\ell\ge3),D_\ell\;(\ell\ge4),E_6,E_7,E_8,F_4\). Premet states that the results were announced at the INI Algebraic Lie Theory final conference in June 2009; the preprint appeared on 15 July 2009 and the article in 2010.}
\GSfield{Present status}{As of November 2026, the blanket semisimple assertion is false. The cited 2013 discussion still lists types \(C\) and \(G_2\) as open; no later resolution of those cases was verified in this search. The commutative rationality theorem for \(G_2\) in \citet[Theorem~2]{anderson2013g2} is not a resolution of the division-ring conjecture.}
\GSfield{Assessment}{literal statement: false as written; publication match: exact; verification: read the result; development: annotated; confidence: high; problem \hyperref[prob:GP-1967-21]{GP-1967-21}}
\end{GSitem}
\begin{GSitem}{1967-144.8}{E}{Nilpotent and algebraic solvable cases}
\GSfield{Kind; source}{remark; \GSpage{P000233}}
\GSfield{The notes}{For nilpotent algebras the conjecture is proved; for algebraic solvable algebras the speaker says ``probably.}
\GSfield{Statement}{For a finite-dimensional nilpotent Lie algebra over an algebraically closed field of characteristic zero, its enveloping division ring is a Weyl division ring with central transcendental parameters, with \(k\) the index and \(2n=\dim\mathfrak g-k\). The notebook's assertion about all algebraic solvable algebras is a separate tentative extension.}
\GSfield{Reconstruction}{The nilpotent theorem is isolated from the unproved-at-this-point solvable remark. No theorem for arbitrary nonalgebraic solvable Lie algebras is inferred.}
\GSfield{Fate}{The nilpotent case is proved in \citet[\S2]{gelfandkirillov1966corps}. The precise publication intended by the tentative solvable parenthesis was not established in the sources opened for this item.}
\GSfield{Present status}{As of November 2026, the nilpotent assertion is established. This entry makes no comprehensive November 2026 classification claim for solvable Lie fields.}
\GSfield{Assessment}{literal statement: correct with missing hypotheses; publication match: partial; verification: read the result; development: annotated; confidence: medium}
\end{GSitem}
\begin{GSitem}{1967-144.9}{A}{Central extension by virtual Laplace operators}
\GSfield{Kind; source}{construction; \GSpage{P000233}}
\GSfield{The notes}{The new algebra is \(\widetilde U=U\otimes_Z\widetilde Z\); \(\widetilde Z/Z\) is called algebraic with Weyl Galois group.}
\GSfield{Statement}{For a split semisimple Lie algebra in characteristic zero, the Harish-Chandra description identifies the centre with Weyl-invariant polynomials on a Cartan dual, using the shifted Weyl action. Adjoining the Cartan coordinate functions gives a finite central algebra extension \(\widetilde U\). On fraction fields the induced extension \(\mathbb C(\mathfrak h^*)/\mathbb C(\mathfrak h^*)^W\) is Galois with group \(W\). These added commuting coordinates are the virtual Laplace operators.}
\GSfield{Reconstruction}{Galois is asserted for fraction fields; the extension of rings can ramify. The shift convention is part of the Harish-Chandra identification.}
\GSfield{Fate}{This construction is developed in \citet[\S\S7--8]{gelfandkirillov1969structure}; the paper was received on 30 September 1968 and appeared in 1969.}
\GSfield{Present status}{As of November 2026, the construction is established and remains distinct from the original enveloping algebra.}
\GSfield{Assessment}{literal statement: correct with missing hypotheses; publication match: exact; verification: read the result; development: annotated; confidence: high}
\end{GSitem}
\begin{GSitem}{1967-144.10}{A}{The Weyl fraction field after central extension}
\GSfield{Kind; source}{theorem; \GSpage{P000233}}
\GSfield{The notes}{The theorem explicitly states that the fraction division ring of \(\widetilde U(\mathfrak g)\) is \(D_{n,k}\).}
\GSfield{Statement}{With \(\widetilde U\) the virtual-Laplace extension above, \(\operatorname{Frac}\widetilde U(\mathfrak g)\cong D_{n,k}\), where \(n=|\Phi^+|\) and \(k=\operatorname{rank}\mathfrak g\). The isomorphism is a statement after the prescribed finite extension of the centre.}
\GSfield{Reconstruction}{The central extension is retained in both the hypothesis and conclusion.}
\GSfield{Fate}{Exact publication match: \citet[\S13, Theorem~3]{gelfandkirillov1969structure}, received September 1968 and published in 1969.}
\GSfield{Present status}{As of November 2026, this theorem is established; its truth does not imply the unextended conjecture, which fails in several types.}
\GSfield{Assessment}{literal statement: correct; publication match: exact; verification: read the result; development: annotated; confidence: high}
\end{GSitem}
\begin{GSitem}{1967-144.11}{F}{Algebraic extensions of Weyl division rings}
\GSfield{Kind; source}{question; \GSpage{P000233}}
\GSfield{The notes}{\ru{Алг. расшир.} of \(D_{n,k}\) are queried, with algebraic varieties mentioned as the commutative analogy.}
\GSfield{Statement}{The note asks for a geometric or structural understanding of extensions of \(D_{n,k}\), but specifies neither finite left/right degree, centrality, embeddings, nor the equivalence relation on extensions. These choices lead to different mathematical problems, so no unique classification assertion can be reconstructed.}
\GSfield{Reconstruction}{The preceding central Galois extension is an example, not a definition of every extension meant by the question.}
\GSfield{Fate}{\citet[Introduction and \S13]{gelfandkirillov1969structure} explains the motivating finite central extension. A precise general extension problem matching this one-line question was not identified.}
\GSfield{Present status}{As of November 2026, the remark remains under-specified; no purported solution of noncommutative birational geometry is assigned to it.}
\GSfield{Assessment}{literal statement: not determinable; publication match: thematic only; verification: read the result; development: annotated; confidence: low}
\end{GSitem}
\begin{GSitem}{1967-144.12}{F}{Global holomorphic differential operators on compact complex manifolds}
\GSfield{Kind; source}{question; \GSpage{P000233}}
\GSfield{The notes}{The question asks to describe regular differential operators on compact complex manifolds, followed by three vector fields on the Riemann sphere.}
\GSfield{Statement}{For a compact complex manifold \(X\), study the filtered algebra \(H^0(X,\mathcal D_X)\), where \(\mathcal D_X\) consists of finite-order holomorphic differential operators acting on \(\mathcal O_X\). Its order-one part contains \(H^0(X,T_X)\) and constants; higher symbols lie in \(H^0(X,\operatorname{Sym}^m T_X)\), but not every symbol necessarily lifts.}
\GSfield{Reconstruction}{This fixes ``regular as holomorphic and the coefficient bundle as trivial. The notebook gives no manifold class or requested invariant, so a universal classification target cannot be extracted.}
\GSfield{Fate}{The differential-operator viewpoint is explained in \citet[Introduction]{gelfandkirillov1966corps}; its basic affine-space results do not answer this compact-manifold question.}
\GSfield{Present status}{As of November 2026, the object is meaningful, but the requested general classification is not precise enough to assign a solved/open theorem.}
\GSfield{Assessment}{literal statement: not determinable; publication match: thematic only; verification: read the result; development: annotated; confidence: medium}
\end{GSitem}
\begin{GSitem}{1967-144.13}{E}{Global vector fields on the projective line}
\GSfield{Kind; source}{remark; \GSpage{P000233}}
\GSfield{The notes}{The sphere has \(\partial_z,z\partial_z,z^2\partial_z\).}
\GSfield{Statement}{Every global holomorphic vector field on \(\mathbb P^1\) has the form \((a+bz+cz^2)\partial_z\). In the coordinate \(w=1/z\), this becomes \(-(aw^2+bw+c)\partial_w\), which is regular at infinity. The three displayed fields form a basis of \(H^0(\mathbb P^1,T_{\mathbb P^1})\cong\mathfrak{sl}_2\).}
\GSfield{Reconstruction}{The observation describes order-one operators; it does not itself classify all higher-order operators on all compact manifolds.}
\GSfield{Fate}{This is the elementary sphere example in the notes. The coordinate calculation directly verifies its global regularity.}
\GSfield{Present status}{As of November 2026, the example is correct and established.}
\GSfield{Assessment}{literal statement: correct; publication match: none found; verification: indirect; development: annotated; confidence: high}
\end{GSitem}

\GStalk{145}{1968-03-11}{L. G. Makar-Limanov}{P000232-P000235}{Dixmier on the first Weyl algebra}
The actual discussion starts near the bottom of P000233 and ends before Palamodov's heading on P000234. It reports Dixmier's automorphism theorem, strictly semisimple elements and nonpolynomial centralizers. Arnold inserts a Poisson-bracket comparison. The open division-ring automorphism question repeats GP-1965-08 and is kept separate from the algebra automorphism theorem.

\begin{GSitem}{1967-145.1}{E}{The first Weyl algebra}
\GSfield{Kind; source}{definition; \GSpage{P000233}}
\GSfield{The notes}{The ring of quantum mechanics is generated by \(p,q\) with \([p,q]=1\).}
\GSfield{Statement}{\(A_1(\mathbb C)=\mathbb C\langle p,q\rangle/(pq-qp-1)\), represented faithfully by \(p=d/dx\) and \(q=x\) on \(\mathbb C[x]\). Its monomials \(p^a q^b\), \(a,b\ge0\), form a vector-space basis. Its Ore fraction division ring is \(D_1\).}
\GSfield{Reconstruction}{The commutator convention is fixed and maintained in the later eigenvalue calculation.}
\GSfield{Fate}{The notation and algebra are those of \citet[\S1]{dixmier1968weyl}; compare 1965-80.5.}
\GSfield{Present status}{As of November 2026, the construction is standard.}
\GSfield{Assessment}{literal statement: correct; publication match: exact; verification: read the result; development: annotated; confidence: high}
\end{GSitem}
\begin{GSitem}{1967-145.2}{E}{Polynomial shears generate the automorphism group of A1}
\GSfield{Kind; source}{theorem; \GSpage{P000233}}
\GSfield{The notes}{The automorphisms are generated by \(p\mapsto p, q\mapsto q+f(p)\) and \(q\mapsto q,p\mapsto p+g(q)\), or by one shear family and Fourier exchange.}
\GSfield{Statement}{Over \(\mathbb C\), \(\operatorname{Aut}_{\mathbb C}A_1\) is generated by the two polynomial shear families \((p,q)\mapsto(p,q+f(p))\), \((p,q)\mapsto(p+g(q),q)\). Equivalently, one may use \((p,q)\mapsto(p+g(q),q)\) together with \((p,q)\mapsto(-q,p)\); linear shears generate the required determinant-one linear transformations.}
\GSfield{Reconstruction}{The statement is for automorphisms, not arbitrary endomorphisms. It does not assert the Dixmier endomorphism conjecture.}
\GSfield{Fate}{This is explicitly an exposition of Dixmier's work, published in 1968: \citet[Theorem~8.10]{dixmier1968weyl}. The notebook supplies evidence that it circulated by 11 March 1968.}
\GSfield{Present status}{As of November 2026, the generation theorem is established.}
\GSfield{Assessment}{literal statement: correct; publication match: exact; verification: read the result; development: annotated; confidence: high}
\end{GSitem}
\begin{GSitem}{1967-145.3}{C}{Automorphisms of the Weyl division ring}
\GSfield{Kind; source}{question; \GSpage{P000233}}
\GSfield{The notes}{For \(D_{1,0}\), the analogous automorphism description is unknown.}
\GSfield{Statement}{Describe \(\operatorname{Aut}_{\mathbb C}D_1\), where \(D_1=\operatorname{Frac}A_1(\mathbb C)\). The subgroup stabilizing the rational function subfield \(\mathbb C(q)\) consists exactly of maps \(q\mapsto\theta(q)=(aq+b)/(cq+d)\), \(p\mapsto\theta^{\prime}(q)^{-1}p+r(q)\), where \(ad-bc\ne0\) and \(r\in\mathbb C(q)\). This subgroup does not exhaust the full problem.}
\GSfield{Reconstruction}{The stabilizer formula is a precisely bounded modern result, not part of the notebook's assertion of ignorance.}
\GSfield{Fate}{This repeats 1965-80.10 and retains GP-1965-08. \citet[\S5.3.3, final Remark, p.~66]{dumas2011invariants} gives the exact rational-subfield stabilizer and explicitly documents the unresolved full automorphism-group structure. The pseudodifferential-series embedding used in the elementary centralizer proof is described in \S5.2.2, pp.~59--60. These are known special results, not new research claims.}
\GSfield{Present status}{As of November 2026, no complete description of \(\operatorname{Aut}_{\mathbb C}D_1\) was found in this audit; its open status is explicit in \citet[p.~66]{dumas2011invariants}. The checked 2026 result of \citet[Corollary~8.12]{lauritzenthomsen2026duality}, first posted on 11 October and revised on 2 November 2026, establishes equality of finite left and right indices for noncommutative division subrings. That theorem does not classify automorphisms and does not supersede the full-group question. The rational-subfield stabilizer remains a bounded, established reduction.}
\GSfield{Assessment}{literal statement: correct; publication match: partial; verification: read the result; development: developed in this chapter; confidence: high; problem \hyperref[prob:GP-1965-08]{GP-1965-08}}
\GSfield{Worth developing}{Worth developing the rational-subfield stabilizer and its centralizer proof, because this gives a checkable reduction without claiming to classify all division-ring automorphisms.}
\end{GSitem}
\begin{GSitem}{1967-145.4}{E}{Arnold: polynomial canonical maps and the Poisson bracket}
\GSfield{Kind; source}{remark; \GSpage{P000234}; V. I. Arnold}
\GSfield{The notes}{Arnold writes the usual two-variable Poisson bracket and considers polynomials \(P,Q\) for which it equals one.}
\GSfield{Statement}{On \(\mathbb C^2\) with coordinates \(p,q\), \(\{P,Q\}=P_pQ_q-P_qQ_p\). If \(\{P,Q\}=1\), the polynomial map \((p,q)\mapsto(P,Q)\) has Jacobian determinant one and is locally biholomorphic and symplectic. If it is additionally a polynomial automorphism, it is a global polynomial canonical transformation.}
\GSfield{Reconstruction}{The distinction between local invertibility and a global polynomial inverse is essential. The notes do not state or solve the Jacobian conjecture.}
\GSfield{Fate}{The bracket calculation is the classical analogue offered in the seminar; \citet[Introduction and \S8]{dixmier1968weyl} provides the Weyl-algebra comparison.}
\GSfield{Present status}{As of November 2026, the local conclusion follows directly from the inverse-function theorem; no global-inverse claim is warranted.}
\GSfield{Assessment}{literal statement: correct with missing hypotheses; publication match: thematic only; verification: read the result; development: annotated; confidence: high}
\end{GSitem}
\begin{GSitem}{1967-145.5}{E}{The qualified relationship with associated graded Poisson algebras}
\GSfield{Kind; source}{remark; \GSpage{P000234}}
\GSfield{The notes}{The note says there is no connection, even with grading.}
\GSfield{Statement}{For the Bernstein filtration \(\deg p=\deg q=1\), \(\operatorname{gr}A_1\cong\mathbb C[p,q]\), and the commutator induces its canonical Poisson bracket of degree \(-2\). A general polynomial shear need not preserve this filtration, so taking an associated graded map does not directly compare all automorphisms.}
\GSfield{Reconstruction}{This distinguishes a valid warning about arbitrary automorphisms from the overly broad denial of any graded relationship.}
\GSfield{Fate}{The filtered construction follows from the relation \([p,q]=1\); the Newton/degree analysis in \citet[\S\S2--3]{dixmier1968weyl} is relevant background.}
\GSfield{Present status}{As of November 2026, the blanket literal statement is false; the more restricted warning is correct.}
\GSfield{Assessment}{literal statement: false as written; publication match: partial; verification: read the result; development: annotated; confidence: medium}
\end{GSitem}
\begin{GSitem}{1967-145.6}{E}{Strictly semisimple elements and the missing affine parameters}
\GSfield{Kind; source}{theorem; \GSpage{P000234}}
\GSfield{The notes}{An element is completely semisimple when its adjoint has a complete eigenbasis; it is then said to be carried to \(pq\).}
\GSfield{Statement}{For a nonscalar \(x\in A_1(\mathbb C)\) such that \(\operatorname{ad}x\) is diagonalizable on the whole algebra, there are an automorphism \(\alpha\), a nonzero \(a\in\mathbb C\), and \(b\in\mathbb C\) with \(\alpha(x)=a pq+b\). The coefficient and scalar term cannot generally be removed by an automorphism.}
\GSfield{Reconstruction}{``Complete system of eigenvalues is read as an eigenvector decomposition. Scalars are excluded from the strictly semisimple case.}
\GSfield{Fate}{Exact corrected form: \citet[Theorem~9.2 and Corollary~9.3]{dixmier1968weyl}.}
\GSfield{Present status}{As of November 2026, the corrected theorem is established; the notebook's single normal form \(pq\) is too restrictive.}
\GSfield{Assessment}{literal statement: false as written; publication match: exact; verification: read the result; development: annotated; confidence: high}
\end{GSitem}
\begin{GSitem}{1967-145.7}{E}{Eigenvalues of the Euler element under the stated commutator convention}
\GSfield{Kind; source}{claim; \GSpage{P000234}}
\GSfield{The notes}{The formula gives \(p^a q^b\mapsto(a-b)p^a q^b\).}
\GSfield{Statement}{With \([p,q]=1\), \([pq,p]=-p\), \([pq,q]=q\), hence \([pq,p^a q^b]=(b-a)p^a q^b\). Thus the Weyl monomials are eigenvectors of \(\operatorname{ad}(pq)\), with eigenvalues \(b-a\).}
\GSfield{Reconstruction}{The sign is fixed by the preceding definition, not by a different convention silently substituted for it.}
\GSfield{Fate}{Direct commutator computation confirms the correction; compare \citet[\S5.3, equation~(4)]{dixmier1968weyl}.}
\GSfield{Present status}{As of November 2026, the printed \(a-b\) has the opposite sign under the notebook's convention.}
\GSfield{Assessment}{literal statement: false as written; publication match: exact; verification: read the result; development: annotated; confidence: high}
\end{GSitem}
\begin{GSitem}{1967-145.8}{E}{Commutative centralizers of nonscalar Weyl elements}
\GSfield{Kind; source}{theorem; \GSpage{P000234}}
\GSfield{The notes}{A result attributed to Amitsur says the centralizer \(C(x)\) is commutative.}
\GSfield{Statement}{For every nonscalar \(x\in A_1(\mathbb C)\), \(C_{A_1}(x)\) is commutative. It is a finitely generated commutative algebra of transcendence degree one. For a scalar \(x\), its centralizer is all of \(A_1\), so that exception is necessary.}
\GSfield{Reconstruction}{The nonscalar hypothesis and the precise ambient algebra are supplied; the note does not claim this for every arbitrary noncommutative algebra.}
\GSfield{Fate}{\citet[\S4.2 and Theorem~5.1]{dixmier1968weyl} states the centralizer results and credits Amitsur for commutativity.}
\GSfield{Present status}{As of November 2026, the nonscalar theorem is established.}
\GSfield{Assessment}{literal statement: correct with missing hypotheses; publication match: exact; verification: read the result; development: annotated; confidence: high}
\end{GSitem}
\begin{GSitem}{1967-145.9}{E}{Elliptic and cuspidal centralizers}
\GSfield{Kind; source}{construction; \GSpage{P000234}}
\GSfield{The notes}{The notes say centralizers need not be polynomial rings and display \(X^3=Y^2+\alpha\) for arbitrary \(\alpha\in\mathbb C\).}
\GSfield{Statement}{There are commuting elements \(X,Y\in A_1(\mathbb C)\) satisfying \(Y^2=X^3-\alpha\), whose generated commutative algebra has that sole defining relation and occurs as a centralizer in Dixmier's construction. For \(\alpha\ne0\) this is the affine coordinate ring of a nonsingular elliptic curve with one point removed; for \(\alpha=0\) it is a cuspidal cubic ring, not \(\mathbb C[t]\).}
\GSfield{Reconstruction}{Commutativity and absence of additional algebraic relations are needed for the geometric conclusion; the displayed equation by itself would not prove it.}
\GSfield{Fate}{Exact family in \citet[\S\S5.5--5.6]{dixmier1968weyl}.}
\GSfield{Present status}{As of November 2026, the examples are established and refute the assertion that every nonscalar centralizer is a polynomial ring.}
\GSfield{Assessment}{literal statement: correct with missing hypotheses; publication match: exact; verification: read the result; development: annotated; confidence: high}
\end{GSitem}

\GStalk{146}{1968-03-11}{V. P. Palamodov}{P000233-P000236}{Hormander: weighted d-bar existence and approximation}
The discussion on P000234--P000235 before the heading 18/III reports H\"ormander's \(L^2\) existence and approximation theory. The estimate has a recognizable 1965 source. The growth-class approximation and Martineau indicator remarks are more compressed; their weight conventions and limiting operations require separate scrutiny.

\begin{GSitem}{1967-146.1}{E}{Weighted Dolbeault spaces and plurisubharmonic exhaustion}
\GSfield{Kind; source}{definition; \GSpage{P000234}}
\GSfield{The notes}{The Levi form \(L(\phi)\) is nonnegative, and \(\|f\|_\phi^2=\int_\Omega |f|^2e^{-\phi}dV\). Compact sublevel sets are said to imply pseudoconvexity.}
\GSfield{Statement}{For an open \(\Omega\subset\mathbb C^n\) and a real \(C^2\) function \(\phi\), its Levi form is \(\sum_{j,k}\phi_{j\bar k}\xi_j\overline{\xi_k}\). Nonnegativity is equivalent to plurisubharmonicity. The weighted space \(L^2_{p,q}(\Omega,e^{-\phi})\) consists of measurable forms with finite displayed norm. A plurisubharmonic exhaustion, whose sublevels are relatively compact, makes \(\Omega\) pseudoconvex.}
\GSfield{Reconstruction}{Exhaustion is a global hypothesis; it does not follow from nonnegativity of the Levi form alone.}
\GSfield{Fate}{These are the definitions used in \citet[\S\S2.1--2.2]{hormander1965l2}.}
\GSfield{Present status}{As of November 2026, the formulation is standard.}
\GSfield{Assessment}{literal statement: correct; publication match: exact; verification: read the result; development: annotated; confidence: high}
\end{GSitem}
\begin{GSitem}{1967-146.2}{E}{A weighted L2 solution with logarithmic loss}
\GSfield{Kind; source}{theorem; \GSpage{P000234}}
\GSfield{The notes}{For \(q\ge1\), \(\bar\partial f=0\) and \(f\in L^2(\phi)\), the notes claim \(\bar\partial g=f\), \(g\in L^2(\phi+2\log(1+|z|^2))\), and \(\sqrt q\|g\|\le\|f\|\).}
\GSfield{Statement}{Let \(\Omega\subset\mathbb C^n\) be pseudoconvex, \(\phi\in C^2(\Omega)\) plurisubharmonic, \(1\le q\le n\), and let \(f\) be a distributionally \(\bar\partial\)-closed \((p,q)\)-form with \(\int |f|^2e^{-\phi}<\infty\). Then a solution \(g\) satisfies \(\int |g|^2e^{-\phi}(1+|z|^2)^{-2}\le q^{-1}\int|f|^2e^{-\phi}\).}
\GSfield{Reconstruction}{The norm on the left uses the modified weight, while that on the right uses the original weight. With the usual Euclidean coefficient norm, the lower Levi eigenvalue of \(2\log(1+|z|^2)\) even gives the stronger factor \(1/(2q)\); the displayed weaker inequality is therefore consistent.}
\GSfield{Fate}{A direct specialization of \citet[Theorems~2.2.1 and 2.2.1\(^{\prime}\)]{hormander1965l2}, published three years before the seminar.}
\GSfield{Present status}{As of November 2026, the existence theorem is established with the specified weights.}
\GSfield{Assessment}{literal statement: correct with missing hypotheses; publication match: exact; verification: read the result; development: annotated; confidence: high}
\end{GSitem}
\begin{GSitem}{1967-146.3}{F}{Density of entire functions with weighted growth}
\GSfield{Kind; source}{claim; \GSpage{P000235}}
\GSfield{The notes}{The union of entire functions in \(L^2(\phi+c\log(1+|z|^2))\) is called nontrivial and dense in all entire functions for compact convergence; the note identifies this with a pointwise exponential-polynomial bound.}
\GSfield{Statement}{Fix \(\phi\) and define \(\mathcal A_\phi=\bigcup_{c>0}\{f\in\mathcal O(\mathbb C^n):\int|f|^2e^{-\phi}(1+|z|^2)^{-c}<\infty\}\). The recorded assertion is density of \(\mathcal A_\phi\) in \(\mathcal O(\mathbb C^n)\). Under the explicitly sufficient extra hypothesis that \(\phi\) is bounded below, every polynomial belongs to this union, so density follows from Taylor approximation. For arbitrary weights, integral and pointwise bounds require their precise regularization and factor-of-two conventions.}
\GSfield{Reconstruction}{The bounded-below case is a verified elementary reconstruction, not asserted as the entire intended theorem. The quoted pointwise bound \(e^\phi\) cannot simply be substituted for the weight \(e^{-\phi}\) in an \(L^2\) norm.}
\GSfield{Fate}{\citet[Theorem~2.3.5 and Corollary~2.3.6]{hormander1965l2} gives a precise Runge approximation theorem, but its statement is not the identical fixed-growth-class assertion. The exact version intended here was not established from the opened source.}
\GSfield{Present status}{As of November 2026, the elementary bounded-below case is settled; the full intended growth-weight statement remains unidentified in this audit.}
\GSfield{Assessment}{literal statement: not determinable; publication match: partial; verification: read the result; development: annotated; confidence: medium}
\end{GSitem}
\begin{GSitem}{1967-146.4}{E}{Martineau: realization of regularized growth indicators}
\GSfield{Kind; source}{theorem; \GSpage{P000235}}
\GSfield{The notes}{The notebook attributes to ``Martino'' the assertion that indicators are all radially homogeneous plurisubharmonic functions, with \(i(z)=\lim_{\tau\to\infty}\tau^{-k}\log|f(\tau z)|\).}
\GSfield{Statement}{Fix \(\rho>0\). For an entire function \(f\) on \(\mathbb C^n\) satisfying \(|f(z)|\le C\exp(A|z|^\rho)\), form the radial upper indicator \(h_f(z)=\limsup_{r\to\infty}r^{-\rho}\log|f(rz)|\) and its upper-semicontinuous regularization \(h_f^*\). The regularized indicators, allowing the identically minus-infinity case for the zero function, are exactly the plurisubharmonic functions positively homogeneous of degree \(\rho\): \(h(tz)=t^\rho h(z)\) for \(t>0\). Thus any such nontrivial \(h\) is realized by an entire function of the indicated finite growth order.}
\GSfield{Reconstruction}{The limsup, upper-semicontinuous regularization and finite-type growth convention are essential. The scan's stroke above lim is too unclear to certify a transcription correction. The precise theorem uses one plurisubharmonic condition; the English phrase ``plurisubharmonic harmonic'' must not be read as imposing harmonicity as well.}
\GSfield{Fate}{Exact identification: \citet[Theorem~1, p.~8-02]{martineau1966croissance}. The title page dates the lecture 6 June 1966 and the redaction June 1967, both before the seminar. The introduction acknowledges errors in the earlier presentation and subsequent corrections; the opened corrected exposition is the source used here. This is not the different P\'olya--Ehrenpreis--Martineau theorem on analytic functionals.}
\GSfield{Present status}{As of November 2026, the regularized-indicator characterization is established. The unregularized ordinary-limit reading is not the theorem.}
\GSfield{Assessment}{literal statement: correct with missing hypotheses; publication match: exact; verification: read the result; development: annotated; confidence: high}
\end{GSitem}

\GStalk{147}{1968-03-18}{V. P. Palamodov}{P000234-P000236}{Fixed-degree Dolbeault cohomology and Levi conditions}
The continuation starts at 18/III on P000235 and occupies P000235--P000236. It treats convex weights, a one-variable growth example, fixed-degree Levi signatures, finite-dimensional cohomology, vanishing, Runge restriction, and the weighted integration-by-parts identity. A bounded-disc claim is false literally; Gelfand's short ``domains of regularity'' question is too incomplete to equate with a modern regularity conjecture.

\begin{GSitem}{1967-147.1}{E}{Convexity of weights depending only on real parts}
\GSfield{Kind; source}{remark; \GSpage{P000235}}
\GSfield{The notes}{If \(\phi=\phi(x)\), it is plurisubharmonic exactly when it is convex; the notes also recall subharmonicity on complex lines.}
\GSfield{Statement}{Let \(\phi(z)=u(\operatorname{Re}z)\), with \(u\in C^2(\mathbb R^n,\mathbb R)\). Then \(\phi_{z_j\bar z_k}=\frac14u_{x_jx_k}\); thus \(\phi\) is plurisubharmonic exactly when \(u\) is convex. More generally a continuous function is plurisubharmonic exactly when its restriction to every complex affine line is subharmonic wherever defined.}
\GSfield{Reconstruction}{The regularity is stated for the Hessian proof; the distributional formulation extends the equivalence.}
\GSfield{Fate}{This is standard background for the weights in \citet[\S2.2]{hormander1965l2}.}
\GSfield{Present status}{As of November 2026, the statement is established.}
\GSfield{Assessment}{literal statement: correct; publication match: exact; verification: read the result; development: annotated; confidence: high}
\end{GSitem}
\begin{GSitem}{1967-147.2}{E}{Solving d-bar with exponential type and polynomial growth}
\GSfield{Kind; source}{remark; \GSpage{P000235}}
\GSfield{The notes}{For \(\bar\partial u=f\) on the plane, \(|f(z)|\le C e^{a|\operatorname{Im}z|}(1+|z|^2)^k\); a solution in this growth class is called nontrivial.}
\GSfield{Statement}{For measurable \(f\) satisfying this bound, with \(a\ge0\) and finite \(k\), the weighted \(L^2\) method gives a distributional solution with the same exponential type and a possibly larger polynomial exponent. One chooses a plurisubharmonic weight comparable to \(2a|\operatorname{Im}z|+c\log(1+|z|^2)\); the logarithmic loss is absorbed by increasing \(c\).}
\GSfield{Reconstruction}{The phrase ``this class'' is interpreted as allowing some polynomial exponent, not preserving the exact given \(k\). Interior estimates or a local Cauchy decomposition are needed to convert the weighted integral estimate to pointwise growth.}
\GSfield{Fate}{This illustrates \citet[\S2.4]{hormander1965l2}, rather than a new independent theorem of Palamodov.}
\GSfield{Present status}{As of November 2026, the growth-class solvability is established with a polynomial loss; no fixed-exponent estimate is claimed.}
\GSfield{Assessment}{literal statement: correct with missing hypotheses; publication match: partial; verification: read the result; development: annotated; confidence: high}
\end{GSitem}
\begin{GSitem}{1967-147.3}{E}{The bounded right-hand-side assertion on the unit disc}
\GSfield{Kind; source}{claim; \GSpage{P000235}}
\GSfield{The notes}{The notes say that for bounded \(f\) on the unit disc, \(\partial_z u=f\) (or \(\partial_{\bar z}u=f\)) need not have a bounded solution.}
\GSfield{Statement}{For every \(f\in L^\infty(\mathbb D)\), \(u(z)=\pi^{-1}\int_{\mathbb D}f(\zeta)/(z-\zeta)\,dA(\zeta)\) solves \(\bar\partial u=f\) distributionally and satisfies \(\|u\|_\infty\le4\|f\|_\infty\). Conjugating the kernel gives the corresponding result for \(\partial_z\). Thus the recorded nonexistence claim is false without another omitted requirement, such as a boundary condition or a different data norm.}
\GSfield{Reconstruction}{The scan confirms the negative assertion. The denominator's bar ambiguity does not rescue it, since both operators have the same bounded-data conclusion.}
\GSfield{Fate}{The correction follows from \(\bar\partial(1/(\pi z))=\delta_0\) and the bound \(\mathbb D\subset B(z,2)\), \(\pi^{-1}\int_{B(z,2)}|z-\zeta|^{-1}dA=4\). This is an elementary verification, not a later solution of a research problem.}
\GSfield{Present status}{As of November 2026, the literal assertion is refuted by the displayed solution operator.}
\GSfield{Assessment}{literal statement: false as written; publication match: none found; verification: indirect; development: annotated; confidence: high}
\end{GSitem}
\begin{GSitem}{1967-147.4}{F}{A measure-growth condition of Carleson type}
\GSfield{Kind; source}{remark; \GSpage{P000235}}
\GSfield{The notes}{A measure \(f\) is drawn near the disc boundary with \(|f(S\cap\Omega)|\le C R_S\), said to suffice for solvability.}
\GSfield{Statement}{A standard precise growth condition would be \(|\mu|(B(\zeta,r)\cap\mathbb D)\le Cr\) for boundary-centred discs, where \(|\mu|\) is total variation. The notebook does not specify disc centres, total variation, boundary versus interior boundedness, or the solution's trace. Those choices are essential to a \(\bar\partial\) theorem with measure data.}
\GSfield{Reconstruction}{The name ``Carleson type'' describes the scaling; it is not an exact identification with a corona theorem or a bounded-interior solution theorem.}
\GSfield{Fate}{Searches for the displayed measure condition and bounded \(\bar\partial\) solvability did not determine a unique theorem matching the drawing. The nearby \(L^2\) theorem in \citet[\S2.2]{hormander1965l2} does not by itself establish this measure-data statement.}
\GSfield{Present status}{As of November 2026, the intended norm and measure theorem remain unidentified.}
\GSfield{Assessment}{literal statement: not determinable; publication match: none found; verification: unresolved lead; development: annotated; confidence: low}
\end{GSitem}
\begin{GSitem}{1967-147.5}{E}{The boundary signature condition aq}
\GSfield{Kind; source}{definition; \GSpage{P000235}}
\GSfield{The notes}{On the complex tangent hyperplane, the Levi form has at least \(q+1\) negative or \(n-q\) positive eigenvalues.}
\GSfield{Statement}{For a defining function \(\rho\) with \(d\rho\ne0\) on \(\partial\Omega\), condition \(a_q\) requires that the Hermitian form \(\partial\bar\partial\rho\) restricted to \(\ker\partial\rho\) have at least \(q+1\) strictly negative eigenvalues or at least \(n-q\) strictly positive eigenvalues at each boundary point, with \(\Omega=\{\rho<0\}\).}
\GSfield{Reconstruction}{The normal orientation and strictness are fixed. The complex tangent space has dimension \(n-1\).}
\GSfield{Fate}{Exact definition in \citet[Definition~3.3.4]{hormander1965l2}.}
\GSfield{Present status}{As of November 2026, the condition is established terminology for the ensuing estimates.}
\GSfield{Assessment}{literal statement: correct; publication match: exact; verification: read the result; development: annotated; confidence: high}
\end{GSitem}
\begin{GSitem}{1967-147.6}{E}{Pseudoconvexity and the incorrect index a0}
\GSfield{Kind; source}{claim; \GSpage{P000235}}
\GSfield{The notes}{The boundary is called pseudoconvex if \(a_0\) holds, glossed as \(n-1\) positive eigenvalues.}
\GSfield{Statement}{A \(C^2\) boundary is Levi pseudoconvex if the tangential Levi form is nonnegative. Strict pseudoconvexity means all \(n-1\) tangential eigenvalues are positive; this is the positive branch of \(a_1\), and it implies \(a_q\) for \(q\ge1\). It is not the positive branch of \(a_0\), which would demand \(n\) positive eigenvalues on an \((n-1)\)-dimensional space.}
\GSfield{Reconstruction}{Both the index and the distinction between weak and strict pseudoconvexity must be corrected.}
\GSfield{Fate}{Compare \citet[Definition~3.3.4 and Theorem~3.5.1]{hormander1965l2}; the notebook's index is inconsistent with its immediately preceding definition.}
\GSfield{Present status}{As of November 2026, the literal identification is false; the corrected strict-pseudoconvex statement is standard.}
\GSfield{Assessment}{literal statement: false as written; publication match: exact; verification: read the result; development: annotated; confidence: high}
\end{GSitem}
\begin{GSitem}{1967-147.7}{E}{Finite-dimensional weighted cohomology and closed range}
\GSfield{Kind; source}{theorem; \GSpage{P000235}}
\GSfield{The notes}{For a bounded domain with \(C^3\) boundary satisfying \(a_q\), the cohomology is finite dimensional and the image is closed.}
\GSfield{Statement}{Let \(\Omega\) be relatively compact in a Hermitian complex \(n\)-manifold, with \(C^3\) boundary satisfying \(a_q\), and let the real weight be continuous on \(\overline\Omega\). For the maximal weak \(\bar\partial\) operators in the \((p,q-1),(p,q),(p,q+1)\) weighted \(L^2\) spaces, \(\operatorname{im}\bar\partial\) at degree \(q\) is closed, and \(\ker\bar\partial/\operatorname{im}\bar\partial\) at degree \(q\) is finite dimensional.}
\GSfield{Reconstruction}{The Hilbert-space domains and the weight's boundary control are supplied; arbitrary singular weights are not covered by this statement.}
\GSfield{Fate}{Exact match with \citet[Theorem~3.4.1]{hormander1965l2}.}
\GSfield{Present status}{As of November 2026, the theorem is established under these boundary and weight hypotheses.}
\GSfield{Assessment}{literal statement: correct with missing hypotheses; publication match: exact; verification: read the result; development: annotated; confidence: high}
\end{GSitem}
\begin{GSitem}{1967-147.8}{F}{Gelfand asks about domains of regularity}
\GSfield{Kind; source}{question; \GSpage{P000235}; I. M. Gelfand}
\GSfield{The notes}{\ru{Гельфанд: области регулярности --- ?} Palamodov reportedly says that nobody is working on this.}
\GSfield{Statement}{The phrase does not identify a function space, operator domain, boundary regularity scale or necessary-and-sufficient condition. It could concern solvability domains or regularity of solutions, but no unique mathematical question can be recovered.}
\GSfield{Reconstruction}{No identification with later global regularity of the \(\bar\partial\)-Neumann problem, or with a particular counterexample, is made.}
\GSfield{Fate}{The surrounding lecture has the definite scope of \citet[\S3.4]{hormander1965l2}; its literature does not determine which additional question Gelfand meant. The historical report about other researchers is not accepted as a verified priority claim.}
\GSfield{Present status}{As of November 2026, the intended question remains unrecoverable from this sentence.}
\GSfield{Assessment}{literal statement: not determinable; publication match: none found; verification: unresolved lead; development: annotated; confidence: low}
\end{GSitem}
\begin{GSitem}{1967-147.9}{E}{Global exhaustion and vanishing in a fixed degree}
\GSfield{Kind; source}{theorem; \GSpage{P000235-P000236}}
\GSfield{The notes}{A real function \(\Psi\) on the whole domain, with one more positive Levi direction, is said to force zero cohomology.}
\GSfield{Statement}{Let \(X\) be a complex \(n\)-manifold with a \(C^2\) exhaustion \(\Psi\). Assume that outside a compact sublevel its regular levels satisfy \(a_q\), and everywhere its full Levi form has at least \(q+1\) negative or at least \(n-q+1\) positive eigenvalues. Then \(H^q(X,\Omega_X^p)=0\). In the positive branch this is the familiar \(q\)-complete vanishing setting.}
\GSfield{Reconstruction}{Proper exhaustion and the level-set hypothesis are restored. Existence of an arbitrary local function with a suitable Hessian is not enough.}
\GSfield{Fate}{Exact fixed-degree form in \citet[Theorem~3.4.10]{hormander1965l2}.}
\GSfield{Present status}{As of November 2026, the theorem is established with these global hypotheses.}
\GSfield{Assessment}{literal statement: correct with missing hypotheses; publication match: exact; verification: read the result; development: annotated; confidence: high}
\end{GSitem}
\begin{GSitem}{1967-147.10}{E}{Andreotti--Grauert extension to complex spaces}
\GSfield{Kind; source}{remark; \GSpage{P000236}}
\GSfield{The notes}{The notes attribute all these theorems for complex spaces to Andreotti and Grauert.}
\GSfield{Statement}{For a strongly \(q\)-pseudoconvex complex space \(X\) and a coherent analytic sheaf \(\mathcal F\), the relevant higher cohomology groups are finite dimensional in the prescribed degree range; for \(q\)-complete spaces, \(H^j(X,\mathcal F)=0\) for \(j\ge q\). Pseudoconcave variants have ranges involving the depth of \(\mathcal F\), so they cannot be obtained merely by copying every smooth-domain formula.}
\GSfield{Reconstruction}{This states the positive branch explicitly and flags the additional depth requirements for the opposite branch.}
\GSfield{Fate}{Already published in \citet[\S5, finiteness and vanishing theorems]{andreottigrauert1962finitude}.}
\GSfield{Present status}{As of November 2026, the coherent-sheaf results are established; the blanket phrase requires degree and depth qualifications.}
\GSfield{Assessment}{literal statement: correct with missing hypotheses; publication match: exact; verification: read the result; development: annotated; confidence: high}
\end{GSitem}
\begin{GSitem}{1967-147.11}{E}{Runge restriction and extension of solvability}
\GSfield{Kind; source}{theorem; \GSpage{P000236}}
\GSfield{The notes}{For \(\omega=\{\Psi<0\}\subset\Omega\), with a Levi condition outside \(\omega\) and \(\Psi\to\infty\) at the boundary, cohomology restricts injectively. Solvability on \(\omega\) then gives solvability on \(\Omega\).}
\GSfield{Statement}{A precise sufficient version is: let \(X\) have a \(C^3\) exhaustion \(\Psi\) whose regular levels outside \(X_c=\{\Psi<c\}\) satisfy \(a_q\). Then the restriction \(H^q(X,\Omega_X^p)\to H^q(X_c,\Omega_X^p)\) is injective. For relatively compact smooth sublevels the analogous statement holds for the weighted \(L^2\) cohomology with weights continuous on the closure.}
\GSfield{Reconstruction}{The condition is imposed on level-set tangent hyperplanes, not replaced without justification by a signature of the full Hessian. The strict positive full-Hessian branch in the notes implies this level condition; the negative branch requires care.}
\GSfield{Fate}{\citet[Theorems~3.4.6 and 3.4.9]{hormander1965l2}.}
\GSfield{Present status}{As of November 2026, the restriction theorem is established; it is an injectivity statement, not unconditional vanishing.}
\GSfield{Assessment}{literal statement: correct with missing hypotheses; publication match: exact; verification: read the result; development: annotated; confidence: high}
\end{GSitem}
\begin{GSitem}{1967-147.12}{E}{Weighted integration by parts and boundary Levi form}
\GSfield{Kind; source}{method; \GSpage{P000236}}
\GSfield{The notes}{The final displayed identity contains the weighted \(\bar\partial\) and adjoint norms, interior Hessian terms, derivative norms and the boundary Levi integral.}
\GSfield{Statement}{For a smooth \((0,q)\)-form \(u\) on a smooth bounded domain, satisfying the \(\bar\partial\)-Neumann boundary condition \(\sum_j\rho_{z_j}u_{jK}=0\), the Morrey--Kohn--H\"ormander identity expresses \(\|\bar\partial u\|_\phi^2+\|\bar\partial_\phi^*u\|_\phi^2\) as the sum of coefficient \(\bar z\)-derivative norms, the weighted interior Levi form \(\sum\phi_{j\bar k}u_{jK}\overline{u_{kK}}\), and the boundary Levi form \(\sum\rho_{j\bar k}u_{jK}\overline{u_{kK}}\), with boundary measure divided by \(|\nabla\rho|\) unless \(\rho\) is normalized.}
\GSfield{Reconstruction}{The boundary condition, normalization and the weighted formal adjoint are necessary for this precise identity; the general \((p,q)\) version is coefficientwise.}
\GSfield{Fate}{The proof uses \citet[Prop.~2.1.2, eq.~(2.1.14)]{hormander1965l2}.}
\GSfield{Present status}{As of November 2026, the identity is established and explains the role of the Levi signatures.}
\GSfield{Assessment}{literal statement: correct with missing hypotheses; publication match: exact; verification: read the result; development: annotated; confidence: high}
\end{GSitem}
\begin{GSitem}{1967-147.13}{E}{Approximation of closed forms across a Runge exhaustion}
\GSfield{Kind; source}{theorem; \GSpage{P000236}}
\GSfield{The notes}{Closed forms in \(\omega\) are approximated by closed forms in \(\Omega\).}
\GSfield{Statement}{Under the hypotheses of the preceding Runge theorem, \(\bar\partial\)-closed \((p,q-1)\)-forms on the inner sublevel are limits of restrictions of \(\bar\partial\)-closed forms on the larger domain. In the relatively compact setting the approximation is in the specified weighted \(L^2\) norm; along an exhaustion one obtains the corresponding local form of approximation.}
\GSfield{Reconstruction}{The degree is \(q-1\), while the cohomological injection is in degree \(q\). No uniform-up-to-boundary approximation of arbitrary forms is inferred.}
\GSfield{Fate}{Exact source: \citet[Theorems~3.4.7 and 3.4.9]{hormander1965l2}.}
\GSfield{Present status}{As of November 2026, the approximation theorem is established with the stated degree and norm.}
\GSfield{Assessment}{literal statement: correct with missing hypotheses; publication match: exact; verification: read the result; development: annotated; confidence: high}
\end{GSitem}
\begin{GSitem}{1967-147.14}{E}{Existence on a Levi-pseudoconvex bounded domain}
\GSfield{Kind; source}{theorem; \GSpage{P000236}}
\GSfield{The notes}{A \(C^2\) defining function with nonnegative tangential Levi form is followed by ``solvability follows'' from the identity.}
\GSfield{Statement}{If \(\Omega\Subset\mathbb C^n\) has \(C^2\) pseudoconvex boundary and \(1\le q\le n\), every distributionally closed \(f\in L^2_{p,q}(\Omega)\) has a solution \(u\in L^2_{p,q-1}(\Omega)\), with \(\|u\|_2\le C_\Omega\|f\|_2\). One may use the strictly plurisubharmonic bounded weight \(|z|^2\), obtain a weighted estimate, and compare weighted and unweighted norms on the bounded domain.}
\GSfield{Reconstruction}{This is solvability in \(L^2\), not smoothness of the solution up to the boundary. The notebook's comparison with grad and curl concerns the integration-by-parts mechanism.}
\GSfield{Fate}{\citet[Theorem~2.2.3]{hormander1965l2}; the weighted identity is the analytic tool, not an additional regularity theorem.}
\GSfield{Present status}{As of November 2026, the stated existence result is established.}
\GSfield{Assessment}{literal statement: correct with missing hypotheses; publication match: exact; verification: read the result; development: annotated; confidence: high}
\end{GSitem}

\GStalk{148}{}{}{P000237-P000238}{Programme for the 1968--1969 seminar}
P000237 lists proposed subjects: Mumford's moduli spaces (Geronimus); Kirillov on characters; Kazhdan on Dwork; Griffiths on geometry, topology and differential equations; the Spencer resolution (Palamodov); Gromov; Buslaev on eigenvalues including Maslov's work; Grauert--Remmert on \(p\)-adic analysis with Shnirelman's earlier Cauchy-integral ideas; Nilsson; Vershik on factors of types II and III, or Sinai and \(K^0\); and H\"ormander on Lie groups and hypoellipticity. This is a programme, not eleven mathematical theorems. Gelfand's substantive remarks below it are assigned to record 149.


\GStalk{149}{1968-09-02}{I. M. Gelfand}{P000237-P000239}{Programme remarks: Dwork and Gromov}
The mathematical remarks on P000237 concern arithmetic rationality, \(p\)-adic operators, foliations and the passage from germs to jets. The rationality slogan omits decisive global hypotheses; the foliation slogan likewise cannot be treated as an unrestricted theorem. The final psychology anecdote contains no mathematical assertion and is retained only as evidence of the programme discussion. The date is supplied by the catalogue.

\begin{GSitem}{1968-149.1}{E}{Dwork rationality and the missing global hypothesis}
\GSfield{Kind; source}{claim; \GSpage{P000237}}
\GSfield{The notes}{A power series convergent near zero over the complex numbers and for all \(p\)-adic metrics is said to be rational.}
\GSfield{Statement}{Local convergence alone is insufficient: \(e^t=\sum t^n/n!\) has positive radius in \(\mathbb C\) and in every \(\mathbb Q_p\), but is not rational. A precise sufficient Dwork criterion is that \(F\in\mathbb Z[[t]]\) have positive complex radius and admit meromorphic continuation to the entire \(p\)-adic affine line for one prime \(p\); then \(F\in\mathbb Q(t)\).}
\GSfield{Reconstruction}{Integer coefficients and global \(p\)-adic meromorphic continuation are restored. The rationality theorem is not a statement about arbitrary convergent germs at all places.}
\GSfield{Fate}{\citet[Introduction and \S4]{dwork1960rationality} supplies the arithmetic-analytic argument. It was published in 1960 and available in Russian in 1961, before the seminar.}
\GSfield{Present status}{As of November 2026, the corrected criterion and rationality of finite-field zeta functions are established; the local-convergence slogan is false.}
\GSfield{Assessment}{literal statement: false as written; publication match: partial; verification: read the result; development: annotated; confidence: high}
\end{GSitem}
\begin{GSitem}{1968-149.2}{E}{Zeta functions through p-adic completely continuous operators}
\GSfield{Kind; source}{method; \GSpage{P000237}}
\GSfield{The notes}{The idea is to represent the zeta function using completely continuous operators over a \(p\)-adic field.}
\GSfield{Statement}{For a variety \(X/\mathbb F_q\), \(Z(X,t)=\exp(\sum_{m\ge1}\#X(\mathbb F_{q^m})t^m/m)\). Dwork's method expresses point counts through traces of suitable \(p\)-adic operators, and the resulting Fredholm determinants give \(p\)-adic meromorphic continuation. Combined with arithmetic coefficient control and the rationality criterion, this proves \(Z(X,t)\in\mathbb Q(t)\).}
\GSfield{Reconstruction}{A trace or Fredholm formula alone does not automatically prove rationality of an infinite-dimensional determinant.}
\GSfield{Fate}{The analytic continuation and arithmetic steps are separated in \citet[\S\S3--4]{dwork1960rationality}.}
\GSfield{Present status}{As of November 2026, the method and theorem are established; no new claim is attributed to this brief programme remark.}
\GSfield{Assessment}{literal statement: correct; publication match: exact; verification: read the result; development: annotated; confidence: high}
\end{GSitem}
\begin{GSitem}{1968-149.3}{F}{Plane distributions and the foliation existence slogan}
\GSfield{Kind; source}{remark; \GSpage{P000237}}
\GSfield{The notes}{For one-dimensional foliations the question is said to reduce to vector fields; for higher codimension, existence is asserted whenever a family of planes exists.}
\GSfield{Statement}{A smooth nowhere-zero line distribution is involutive and determines a one-dimensional foliation. A higher-rank plane distribution need not be involutive. The meaningful deformation question asks whether a given distribution is homotopic through plane distributions to the tangent distribution of a foliation, under specified manifold and normal-bundle hypotheses.}
\GSfield{Reconstruction}{The rank-one observation is exact. The higher-rank assertion lacks the hypotheses of a general theorem and cannot be read as integrability of the given distribution.}
\GSfield{Fate}{\citet[\S4.4, Theorem~4.4.2 and Corollary~4.4.3]{gromov1969stable}, received July 1968 and published in 1969, treats qualified foliation existence on open manifolds. It does not justify the unqualified notebook sentence.}
\GSfield{Present status}{As of November 2026, the broad announcement remains under-specified; Frobenius integrability and homotopy-to-foliation are distinct assertions.}
\GSfield{Assessment}{literal statement: correct with missing hypotheses; publication match: partial; verification: read the result; development: annotated; confidence: medium}
\end{GSitem}
\begin{GSitem}{1968-149.4}{B}{Germs, jets and a weak homotopy equivalence}
\GSfield{Kind; source}{claim; \GSpage{P000237}}
\GSfield{The notes}{The sheaf of germs and the sheaf of jets are called weakly homotopy equivalent; a global set of coefficients is said to permit construction of a function.}
\GSfield{Statement}{A precise related theorem is the open invariant relation theorem: on an open smooth manifold, for the natural bundles and open diffeomorphism-invariant differential relations of Gromov's construction, the holonomic-section map to formal solutions is a weak homotopy equivalence. This is not true for an arbitrary differential relation or an arbitrary prescribed jet section.}
\GSfield{Reconstruction}{Open manifold, openness, invariance and the specified natural-bundle setting are necessary qualifications. A formal solution is deformed to a holonomic one; its coefficients need not be realized without deformation.}
\GSfield{Fate}{\citet[Theorem~4.1.1]{gromov1969stable}, received on 8 July 1968 and printed in 1969. Record 162 gives a fuller seminar exposition of this subject.}
\GSfield{Present status}{As of November 2026, the qualified theorem is established; the unrestricted slogan is not valid.}
\GSfield{Assessment}{literal statement: correct with missing hypotheses; publication match: partial; verification: read the result; development: annotated; confidence: high}
\end{GSitem}

\GStalk{150}{1968-09-02}{A. Yu. Geronimus}{P000237-P000242}{Mumford: moduli, universality and automorphisms}
P000238--P000240 introduce compact Riemann surfaces and the elliptic modular quotient, then compare a set of isomorphism classes with a universal family and a representable functor. Automorphisms produce nontrivial isotrivial families, obstructing a fine parameter space. The proposed remedy is Mumford's category and topology of families, not merely the coarse quotient. Kirillov's intervention supplies the functorial formulation.

\begin{GSitem}{1968-150.1}{E}{Genus zero and genus one uniformization}
\GSfield{Kind; source}{remark; \GSpage{P000238}}
\GSfield{The notes}{All genus-zero compact Riemann surfaces are spheres; a genus-one surface has universal cover the complex plane and a lattice of deck translations.}
\GSfield{Statement}{A connected compact Riemann surface of genus zero is biholomorphic to \(\mathbb P^1\). A connected compact genus-one surface, after choosing a point as origin, is \(E=\mathbb C/\Lambda\), where \(\Lambda=\mathbb Z\omega_1+\mathbb Z\omega_2\) is a rank-two lattice with \(\operatorname{Im}(\omega_2/\omega_1)>0\). The marked point equips it with an elliptic-curve group law.}
\GSfield{Reconstruction}{The notebook initially discusses curves without a marked point; pointing is necessary for the elliptic group structure used later.}
\GSfield{Fate}{Classical background to \citet[\S\S4,7]{mumford1965picard}.}
\GSfield{Present status}{As of November 2026, the classification is established.}
\GSfield{Assessment}{literal statement: correct with missing hypotheses; publication match: exact; verification: read the result; development: annotated; confidence: high}
\end{GSitem}
\begin{GSitem}{1968-150.2}{E}{Lattice homothety and the modular group action}
\GSfield{Kind; source}{construction; \GSpage{P000238}}
\GSfield{The notes}{Lattices differing by nonzero complex scaling give isomorphic curves; a change of integral basis yields a fractional-linear transformation.}
\GSfield{Statement}{Pointed complex elliptic curves \(\mathbb C/\Lambda\) and \(\mathbb C/\Lambda'\) are isomorphic exactly when \(\Lambda'=a\Lambda\) for some \(a\in\mathbb C^*\). After normalization \(\Lambda_\tau=\mathbb Z+\mathbb Z\tau\), \(\operatorname{Im}\tau>0\), the basis-change group \(\Gamma=SL_2(\mathbb Z)\) acts by \(\tau\mapsto(a\tau+b)/(c\tau+d)\).}
\GSfield{Reconstruction}{The denominator is \(c\tau+d\), as confirmed by the later formula on P000250--P000251. The P000238 denominator appears to repeat \(b\).}
\GSfield{Fate}{The same presentation is used in \citet[\S7, pp.~76--78]{mumford1965picard}.}
\GSfield{Present status}{As of November 2026, the lattice criterion and action are established.}
\GSfield{Assessment}{literal statement: false as written; publication match: exact; verification: read the result; development: annotated; confidence: high}
\end{GSitem}
\begin{GSitem}{1968-150.3}{E}{Fundamental region and the coarse j-line}
\GSfield{Kind; source}{construction; \GSpage{P000238}}
\GSfield{The notes}{The upper half-plane modulo \(SL_2(\mathbb Z)\) is depicted by the usual modular fundamental region and called the plane of elliptic curves.}
\GSfield{Statement}{A standard closed fundamental region is \(\{\tau\in\mathbb H:|\operatorname{Re}\tau|\le1/2,\ |\tau|\ge1\}\), with the prescribed boundary identifications. The analytic coarse quotient \(\mathbb H/SL_2(\mathbb Z)\) is the complex affine line via the \(j\)-invariant. It classifies isomorphism classes of individual elliptic curves, but does not represent their family functor.}
\GSfield{Reconstruction}{The boundary identifications and the word coarse prevent the drawing from being interpreted as a universal family.}
\GSfield{Fate}{\citet[\S4 and \S7]{mumford1965picard} makes this coarse/family distinction.}
\GSfield{Present status}{As of November 2026, the coarse classification is established.}
\GSfield{Assessment}{literal statement: correct; publication match: exact; verification: read the result; development: annotated; confidence: high}
\end{GSitem}
\begin{GSitem}{1968-150.4}{E}{A family listing isomorphism classes is not yet universal}
\GSfield{Kind; source}{definition; \GSpage{P000238-P000239}}
\GSfield{The notes}{A first definition asks for a family \(X\to Y\) whose fibres are the pairwise nonisomorphic objects of a class \(K\), one of each; uniqueness is then discussed.}
\GSfield{Statement}{A family with one fibre in each isomorphism class is a set-theoretic parametrization. It need not be universal for families, need not exist with the proposed regularity, and has no categorical uniqueness merely from this condition. A moduli problem must also specify which families over each test base are admissible and how they pull back.}
\GSfield{Reconstruction}{This preserves the progression of the lecture from a provisional definition to a functorial one.}
\GSfield{Fate}{The distinction motivates \citet[Introduction and \S3]{mumford1965picard}.}
\GSfield{Present status}{As of November 2026, the provisional definition is insufficient for a fine moduli problem.}
\GSfield{Assessment}{literal statement: correct with missing hypotheses; publication match: partial; verification: read the result; development: annotated; confidence: high}
\end{GSitem}
\begin{GSitem}{1968-150.5}{E}{Pullback of an analytic family}
\GSfield{Kind; source}{construction; \GSpage{P000239}}
\GSfield{The notes}{The fibre product construction induces a family over a new base and has a universal property.}
\GSfield{Statement}{For a holomorphic family \(\pi:X\to S\) and a holomorphic map \(f:T\to S\), the induced family is \(X\times_ST\to T\), with fibre \(X_{f(t)}\) over \(t\). In the moduli category, a morphism over \(f\) includes an isomorphism from the source family to this pullback; these isomorphisms carry the automorphism data.}
\GSfield{Reconstruction}{Analytic spaces, rather than only manifolds, are allowed so that the required fibre products exist.}
\GSfield{Fate}{Expository construction in \citet[\S3]{mumford1965picard}.}
\GSfield{Present status}{As of November 2026, the construction is standard.}
\GSfield{Assessment}{literal statement: correct; publication match: exact; verification: read the result; development: annotated; confidence: high}
\end{GSitem}
\begin{GSitem}{1968-150.6}{E}{Kirillov: representability of the moduli functor}
\GSfield{Kind; source}{definition; \GSpage{P000239}; A. A. Kirillov}
\GSfield{The notes}{Kirillov reformulates universality by assigning to a base the families over it and comparing this with maps into a moduli object.}
\GSfield{Statement}{Let \(\mathcal F(S)\) be the set of isomorphism classes of admissible families over \(S\). Pullback makes \(\mathcal F\) contravariant. A fine moduli space is an object \(M\) with a natural isomorphism \(\operatorname{Hom}(-,M)\cong\mathcal F\); the universal family corresponds to \(\mathrm{id}_M\). A representing object is unique up to unique isomorphism by the Yoneda argument.}
\GSfield{Reconstruction}{Isomorphism classes and naturality are essential. A family category with nontrivial automorphisms contains more information than the set-valued functor.}
\GSfield{Fate}{This is the standard representability reformulation underlying \citet[\S3]{mumford1965picard}, not a new conjecture of Kirillov.}
\GSfield{Present status}{As of November 2026, the formulation and uniqueness statement are established.}
\GSfield{Assessment}{literal statement: correct with missing hypotheses; publication match: exact; verification: read the result; development: annotated; confidence: high}
\end{GSitem}
\begin{GSitem}{1968-150.7}{E}{Fine universality versus weaker parameter families}
\GSfield{Kind; source}{remark; \GSpage{P000239-P000240}}
\GSfield{The notes}{The definitions oscillate between an inducing map existing for every family and being unique; a weaker moduli space is retained when universality fails.}
\GSfield{Statement}{Fine representability requires a unique classifying map for every family and compatibility with pullback. Existence without uniqueness is a weaker parameter property; local existence near each base point is a versality property and is weaker again. A coarse moduli space has its own universal property for maps out of the moduli problem, and is not defined merely by dropping uniqueness from a universal-family assertion.}
\GSfield{Reconstruction}{The three notions are separated rather than identifying all weakened formulations with coarse moduli.}
\GSfield{Fate}{\citet[\S\S3--4]{mumford1965picard} supplies the relevant categories of modular families.}
\GSfield{Present status}{As of November 2026, these distinctions are standard and essential for the elliptic example.}
\GSfield{Assessment}{literal statement: correct with missing hypotheses; publication match: partial; verification: read the result; development: annotated; confidence: high}
\end{GSitem}
\begin{GSitem}{1968-150.8}{E}{Nonexistence of a fine moduli space for all pointed elliptic curves}
\GSfield{Kind; source}{claim; \GSpage{P000240}}
\GSfield{The notes}{Automorphisms yield a nontrivial family with all fibres isomorphic. Its map to a one-class-per-point space would be constant, forcing a trivial family.}
\GSfield{Statement}{The functor of isomorphism classes of all pointed complex elliptic families is not represented by the coarse \(j\)-line. For a fixed elliptic curve \(E\), the family over \(\mathbb C^*\) obtained by descending \(\mathbb C\times E\) under \((t,x)\mapsto(t+2\pi i,-x)\) has constant \(j\) but monodromy \(-I\) on \(H_1(E,\mathbb Z)\). It is not the product family, while pullback along a constant map from a fine one-point-per-class parameter space would be a product.}
\GSfield{Reconstruction}{This explicit isotrivial twist makes the notebook's obstruction precise; the zero section descends.}
\GSfield{Fate}{This is the obstruction discussed in \citet[\S\S3--4, especially pp.~58--60]{mumford1965picard}.}
\GSfield{Present status}{As of November 2026, the nonrepresentability obstruction is established; the moduli stack retains the missing automorphisms.}
\GSfield{Assessment}{literal statement: correct; publication match: exact; verification: read the result; development: annotated; confidence: high; problem \hyperref[prob:GP-1968-01]{GP-1968-01}}
\end{GSitem}
\begin{GSitem}{1968-150.9}{E}{Rigidification by a homology basis, not merely polarization}
\GSfield{Kind; source}{remark; \GSpage{P000240}}
\GSfield{The notes}{The notes say a universal family exists for ``polarized'' curves, explained parenthetically as a chosen homology basis.}
\GSfield{Statement}{For pointed complex elliptic curves with an oriented integral basis of \(H_1\), the upper half-plane is a fine analytic parameter space and the family \((\mathbb H\times\mathbb C)/\mathbb Z^2\) is universal. This is a marking or rigidification. A principal polarization alone does not eliminate the automorphism \([-1]\); for an unpointed genus-one curve, a homology marking alone also leaves translations.}
\GSfield{Reconstruction}{The parenthetical basis is taken as the intended data, while the terminology ``polarized'' is corrected. Both the point and the basis are required for the stated rigidification.}
\GSfield{Fate}{The analytic marked family is used in \citet[\S7]{mumford1965picard}.}
\GSfield{Present status}{As of November 2026, the marked fine-moduli construction is established; the polarization-only claim would be false.}
\GSfield{Assessment}{literal statement: false as written; publication match: exact; verification: read the result; development: annotated; confidence: high}
\end{GSitem}
\begin{GSitem}{1968-150.10}{E}{The moduli object represented by a topology of families}
\GSfield{Kind; source}{construction; \GSpage{P000240}}
\GSfield{The notes}{A terminal object is formally adjoined when none exists; its sheaves and cohomology are then studied.}
\GSfield{Statement}{The family category must retain morphisms together with fibrewise isomorphisms. Given representable isomorphism functors and an appropriate covering rule, one can study its associated site or stack and define sheaves and cohomology there. This replaces the absent universal family over an ordinary coarse parameter space, rather than asserting that formal adjunction alone creates all necessary products and descent.}
\GSfield{Reconstruction}{The modern stack terminology is explanatory reconstruction; the notebook and Mumford use a topology with a formal final object.}
\GSfield{Fate}{The precise construction is \citet[\S3]{mumford1965picard}, expanded in records 151--154.}
\GSfield{Present status}{As of November 2026, the method is established.}
\GSfield{Assessment}{literal statement: correct with missing hypotheses; publication match: exact; verification: read the result; development: annotated; confidence: high}
\end{GSitem}
\begin{GSitem}{1968-150.11}{E}{Automorphism groups of pointed elliptic curves}
\GSfield{Kind; source}{remark; \GSpage{P000240}}
\GSfield{The notes}{Most curves have ``one automorphism'', \(x\mapsto-x\), and some have extra rotations.}
\GSfield{Statement}{For a pointed complex elliptic curve, the generic origin-preserving automorphism group is cyclic of order two, consisting of the identity and \([-1]\). The square lattice has a group of order four, and the hexagonal lattice a group of order six. Thus ``one'' must mean one nonidentity automorphism in the generic case.}
\GSfield{Reconstruction}{The origin-preserving condition is retained; without a marked point, all translations are automorphisms.}
\GSfield{Fate}{\citet[\S4, pp.~57--59]{mumford1965picard}.}
\GSfield{Present status}{As of November 2026, the automorphism calculation is established.}
\GSfield{Assessment}{literal statement: correct with missing hypotheses; publication match: exact; verification: read the result; development: annotated; confidence: high}
\end{GSitem}

\GStalk{151}{1968-09-09}{A. Yu. Geronimus}{P000240-P000246}{Sites of group actions and families}
P000241--P000244 repeat the site axioms and construct the \(\Gamma\)-set example, identify its sheaves with modules, compare an analytic local-isomorphism topology with ordinary topology, and construct products in a moduli category using isomorphism spaces. The principal corrections are the target of a base-changed cover, products versus direct sums over infinitely many orbits, and the finite-group hypothesis in the character formula.

\begin{GSitem}{1968-151.1}{E}{Pretopology axioms and base change revisited}
\GSfield{Kind; source}{definition; \GSpage{P000241}}
\GSfield{The notes}{The fibre product is a categorical product over a base; a covering \(X_i\to X\) pulled back along \(Y\to X\) is written as a covering of \(X\).}
\GSfield{Statement}{The base-changed covering is \(\{X_i\times_XY\to Y\}\), not a family with target \(X\). Identity, composite-cover and pullback-cover axioms define the pretopology. A functor between the displayed site categories must preserve the specified finite limits and covers for the intended induced sheaf construction; arrow conventions for geometric morphisms must then be stated separately.}
\GSfield{Reconstruction}{The displayed target is emended using the universal property and the same axiom as correctly repeated on P000245.}
\GSfield{Fate}{\citet[pp.~38,41]{mumford1965picard}; compare 1967-143.1--1967-143.2.}
\GSfield{Present status}{As of November 2026, the axioms are standard, with the corrected base-change target.}
\GSfield{Assessment}{literal statement: false as written; publication match: exact; verification: read the result; development: annotated; confidence: high}
\end{GSitem}
\begin{GSitem}{1968-151.2}{E}{The site of Gamma-sets}
\GSfield{Kind; source}{construction; \GSpage{P000241}}
\GSfield{The notes}{Objects are \(\Gamma\)-sets, pullbacks are pairs with equal images, covers are jointly surjective, and the terminal object is a point.}
\GSfield{Statement}{For a discrete group \(\Gamma\), let \(\mathcal T_\Gamma\) be the category of left \(\Gamma\)-sets and equivariant maps. A family covers if its images jointly equal the target. The fibre product has the diagonal action, and the one-point set with trivial action is terminal. This is a site even when \(\Gamma\) is infinite.}
\GSfield{Fate}{Exact example in \citet[pp.~39--41]{mumford1965picard}.}
\GSfield{Present status}{As of November 2026, the construction is established.}
\GSfield{Assessment}{literal statement: correct; publication match: exact; verification: read the result; development: annotated; confidence: high}
\end{GSitem}
\begin{GSitem}{1968-151.3}{E}{Sheaves on the group site are modules}
\GSfield{Kind; source}{theorem; \GSpage{P000241}}
\GSfield{The notes}{A sheaf is determined by \(A=F(\Gamma)\), with \(F(\Gamma/H)=A^H\).}
\GSfield{Statement}{The category of abelian sheaves on \(\mathcal T_\Gamma\) is equivalent to the category of \(\mathbb Z[\Gamma]\)-modules. A module \(A\) gives \(F_A(S)=\operatorname{Map}_\Gamma(S,A)\); in particular \(F_A(\Gamma/H)=A^H\). The action on \(F(\Gamma)\) comes from equivariant right translations and contravariance, with a consistent left-action convention.}
\GSfield{Reconstruction}{The acting subgroup in the invariant formula is \(H\); the adjacent unclear letter does not change the displayed equation.}
\GSfield{Fate}{\citet[pp.~40--41]{mumford1965picard}.}
\GSfield{Present status}{As of November 2026, the equivalence and orbit formula are established.}
\GSfield{Assessment}{literal statement: correct; publication match: exact; verification: read the result; development: annotated; confidence: high}
\end{GSitem}
\begin{GSitem}{1968-151.4}{E}{The value on a disjoint union is a product}
\GSfield{Kind; source}{claim; \GSpage{P000242}}
\GSfield{The notes}{For a disjoint union of orbits the notes say to take a direct sum.}
\GSfield{Statement}{For any disjoint union \(S=\coprod_{\alpha}S_\alpha\), the sheaf condition gives \(F_A(S)=\prod_\alpha F_A(S_\alpha)\). A direct sum agrees only for a finite family of abelian groups. In particular the constant coefficient module can yield arbitrarily supported sections on an infinite disjoint union.}
\GSfield{Reconstruction}{The scan confirms the Russian words for direct sum; the mathematical correction is product.}
\GSfield{Fate}{The product, explicitly printed in \citet[p.~40]{mumford1965picard}, also follows immediately from \(\operatorname{Map}_\Gamma\).}
\GSfield{Present status}{As of November 2026, the infinite-union direct-sum statement is false.}
\GSfield{Assessment}{literal statement: false as written; publication match: exact; verification: read the result; development: annotated; confidence: high}
\end{GSitem}
\begin{GSitem}{1968-151.5}{E}{Cohomology of the terminal group-site object}
\GSfield{Kind; source}{theorem; \GSpage{P000242}}
\GSfield{The notes}{\(H^q(\cdot,F)=H^q(\Gamma,A)\).}
\GSfield{Statement}{Under the sheaf-module equivalence, global sections are \(A^\Gamma\). Their right derived functors are group cohomology: \(H^q(\mathcal T_\Gamma,F_A)\cong H^q(\Gamma,A)\) for every \(q\ge0\). The terminal object's underlying set being a point does not make these groups vanish.}
\GSfield{Fate}{\citet[p.~41]{mumford1965picard}.}
\GSfield{Present status}{As of November 2026, the comparison is established.}
\GSfield{Assessment}{literal statement: correct; publication match: exact; verification: read the result; development: annotated; confidence: high}
\end{GSitem}
\begin{GSitem}{1968-151.6}{E}{H2 with integer coefficients and one-dimensional characters}
\GSfield{Kind; source}{claim; \GSpage{P000242}}
\GSfield{The notes}{\(H^2(\Gamma,\mathbb Z)\) is identified with all one-dimensional complex representations of \(\Gamma\), without a group restriction.}
\GSfield{Statement}{If \(\Gamma\) is finite, the exact sequence \(0\to\mathbb Z\to\mathbb C\xrightarrow{\exp(2\pi i\,\cdot)}\mathbb C^*\to1\), with trivial actions, gives \(H^2(\Gamma,\mathbb Z)\cong H^1(\Gamma,\mathbb C^*)=\operatorname{Hom}(\Gamma,\mathbb C^*)\), since positive-degree cohomology with coefficients in \(\mathbb C\) vanishes by averaging. More generally the same conclusion needs the relevant vanishing of \(H^1\) and \(H^2\) with \(\mathbb C\) coefficients.}
\GSfield{Reconstruction}{The claim fails for \(\Gamma=\mathbb Z\): \(H^2=0\), whereas its complex characters form \(\mathbb C^*\).}
\GSfield{Fate}{This is a direct application of the derived-functor formalism recalled in \citet[p.~41]{mumford1965picard}; the additional finite-group hypothesis is not visible in this September note.}
\GSfield{Present status}{As of November 2026, the corrected finite-group formula is established; the unrestricted one is false.}
\GSfield{Assessment}{literal statement: false as written; publication match: partial; verification: read the result; development: annotated; confidence: high}
\end{GSitem}
\begin{GSitem}{1968-151.7}{E}{The analytic local-isomorphism topology}
\GSfield{Kind; source}{construction; \GSpage{P000242}}
\GSfield{The notes}{Analytic manifolds and analytic maps are used, with jointly surjective unramified local coverings; fibre products are declared to exist.}
\GSfield{Statement}{For a fixed complex analytic space \(X\), the analytic \(\acute{e}\)tale site has local biholomorphisms \(U\to X\) as objects and jointly surjective families as covers. Its fibre products exist. If one instead uses all holomorphic maps among all analytic manifolds, arbitrary fibre products need not remain manifolds; the larger category of analytic spaces repairs this, as the next lecture explicitly acknowledges.}
\GSfield{Reconstruction}{Discreteness of fibres alone does not mean unramified: \(z\mapsto z^2\) has discrete fibres but ramifies at zero.}
\GSfield{Fate}{\citet[\S2, pp.~44--45]{mumford1965picard}; the internal correction is occurrence 1968-152.1.}
\GSfield{Present status}{As of November 2026, the analytic-space site is established; the blanket manifold pullback assertion is false.}
\GSfield{Assessment}{literal statement: correct with missing hypotheses; publication match: partial; verification: read the result; development: annotated; confidence: high}
\end{GSitem}
\begin{GSitem}{1968-151.8}{E}{Comparison with ordinary analytic sheaf cohomology}
\GSfield{Kind; source}{theorem; \GSpage{P000242}}
\GSfield{The notes}{The topology's cohomology on \(X\) equals the cohomology of the restricted sheaf on the usual topology because coverings split locally.}
\GSfield{Statement}{The inclusion of ordinary open subsets of a complex analytic space \(X\) into its local-biholomorphism site induces an equivalence of sheaf categories. For an abelian sheaf \(F\), \(H^q(X_{\mathrm{an},\acute{e}t},F)\cong H^q(X,F|_{\mathrm{opens}})\). Locally each local biholomorphism is an isomorphism onto an ordinary open subset, providing the comparison basis.}
\GSfield{Reconstruction}{This is an analytic local-isomorphism comparison; it is not an assertion that algebraic \(\acute{e}\)tale cohomology with arbitrary coefficients equals ordinary topology over every field.}
\GSfield{Fate}{\citet[pp.~44--45]{mumford1965picard}.}
\GSfield{Present status}{As of November 2026, the comparison is established.}
\GSfield{Assessment}{literal statement: correct; publication match: exact; verification: read the result; development: annotated; confidence: high}
\end{GSitem}
\begin{GSitem}{1968-151.9}{E}{The site of families and cohomology of an individual family}
\GSfield{Kind; source}{construction; \GSpage{P000243}}
\GSfield{The notes}{Objects are analytic families \(\pi:X\to S\), morphisms induce pullbacks, and covers are covers of bases. The cohomology of \(\pi\) is that of \(S\).}
\GSfield{Statement}{For a fixed family \(\pi\), the local-isomorphism site of its pullbacks \(X\times_SU\to U\) is equivalent to the local-isomorphism site of \(S\). Thus a sheaf restricted to this slice has the corresponding cohomology \(H^q(S,F_\pi)\). Morphisms of families include a Cartesian square and its chosen fibrewise isomorphism.}
\GSfield{Reconstruction}{The identification is for the slice over the particular family. It does not identify the cohomology of the entire moduli problem with the cohomology of one arbitrary base.}
\GSfield{Fate}{The family-site construction is \citet[\S3]{mumford1965picard}; compare its induced-topology definition on p.~42.}
\GSfield{Present status}{As of November 2026, the slice comparison is established.}
\GSfield{Assessment}{literal statement: correct; publication match: exact; verification: read the result; development: annotated; confidence: high}
\end{GSitem}
\begin{GSitem}{1968-151.10}{E}{Products over the formal moduli object are isomorphism spaces}
\GSfield{Kind; source}{construction; \GSpage{P000243-P000244}}
\GSfield{The notes}{The direct product of two families cannot simply have base \(S_1\times S_2\); the notes later count automorphisms as multiplicities. An earlier sentence says a curve's product with itself is itself.}
\GSfield{Statement}{For proper smooth pointed genus-one families, or smooth proper curves of genus at least two, the product over the moduli object has base \(I=\operatorname{Isom}_{S_1\times S_2}(\operatorname{pr}_1^*X_1,\operatorname{pr}_2^*X_2)\). A point consists of \((s_1,s_2,\alpha)\) with \(\alpha:X_{1,s_1}\cong X_{2,s_2}\). Its fibre over \((s_1,s_2)\) is empty if the curves are nonisomorphic and otherwise a torsor under their automorphism group. A self-product therefore has one point for each automorphism, not necessarily one point.}
\GSfield{Reconstruction}{The Isom space is an analytic space; smoothness or local biholomorphy requires additional conditions on the families. The lower-page multiplicity remark corrects the preceding set-level simplification.}
\GSfield{Fate}{Exact construction in \citet[\S3, pp.~49--53]{mumford1965picard}.}
\GSfield{Present status}{As of November 2026, the Isom construction is established and retains the automorphism information lost by a coarse quotient.}
\GSfield{Assessment}{literal statement: correct with missing hypotheses; publication match: exact; verification: read the result; development: annotated; confidence: high}
\end{GSitem}
\begin{GSitem}{1968-151.11}{E}{Coverings of the formal terminal object}
\GSfield{Kind; source}{definition; \GSpage{P000244}}
\GSfield{The notes}{A family \(\{\pi_i\to M\}\) covers if for every \(\pi\), the products \(\pi_i\times_M\pi\to\pi\) form a genuine covering. All families together need not be a covering.}
\GSfield{Statement}{An atlas family covers the moduli object when every base change by an admissible family is represented by a jointly surjective family of local isomorphisms of bases. This requires both local deformation completeness and the appropriate unramified condition. Merely listing a representative of every curve, or including every possible family, does not ensure that each component map is eligible as an \(\acute{e}\)tale cover.}
\GSfield{Reconstruction}{The covering class is fixed before applying joint surjectivity; adding non-\(\acute{e}\)tale families is not automatically permitted.}
\GSfield{Fate}{\citet[\S3, pp.~50--54]{mumford1965picard}.}
\GSfield{Present status}{As of November 2026, the base-change test is established.}
\GSfield{Assessment}{literal statement: correct; publication match: exact; verification: read the result; development: annotated; confidence: high}
\end{GSitem}

\GStalk{152}{1968-09-16}{A. Yu. Geronimus}{P000244-P000249}{Analytic spaces, cohomology comparison and moduli atlases}
On P000245--P000247 Geronimus explicitly repairs the previous lecture by replacing analytic manifolds by analytic spaces. He repeats the site and sheaf axioms, recalls sheafification and derived cohomology, gives a local-acyclicity comparison criterion, and formulates covering and local-extension conditions for a moduli atlas. These are expository parts of the Mumford programme, not new general existence theorems for arbitrary moduli problems.

\begin{GSitem}{1968-152.1}{E}{Analytic spaces are required for general fibre products}
\GSfield{Kind; source}{remark; \GSpage{P000245}}
\GSfield{The notes}{The opening apology says that everywhere one should read analytic spaces, not analytic manifolds, because fibre products must exist.}
\GSfield{Statement}{The category of complex analytic spaces admits fibre products of holomorphic maps. The subcategory of manifolds does not: the pullback of \(z\mapsto z^2\) along the map of a point to zero is the nonreduced analytic space with local ring \(\mathbb C\{z\}/(z^2)\). Even among reduced spaces, the equation \(x^2=y^2\) gives a singular pullback.}
\GSfield{Reconstruction}{Nonreduced structure is retained in the first example; passing only to point sets would erase the obstruction. This is the lecturer's own correction, not a modern conjectural emendation.}
\GSfield{Fate}{The distinction is already built into \citet[\S\S2--3]{mumford1965picard}. It corrects occurrence 1968-151.7.}
\GSfield{Present status}{As of November 2026, the category correction is established and necessary.}
\GSfield{Assessment}{literal statement: correct; publication match: exact; verification: read the result; development: annotated; confidence: high}
\end{GSitem}
\begin{GSitem}{1968-152.2}{E}{Fibre products and the three covering axioms repeated}
\GSfield{Kind; source}{definition; \GSpage{P000245}}
\GSfield{The notes}{The universal square and its unique factorization are redrawn, followed by composition, identity and pullback of covers.}
\GSfield{Statement}{For maps \(X\to Z\leftarrow Y\), \(X\times_ZY\) represents compatible pairs of maps into \(X,Y\). A pretopology has identity covers, composites of covering families as covers, and base-changed covers \(X_i\times_XY\to Y\). These statements include the existence of the indicated pullbacks.}
\GSfield{Reconstruction}{The undeciphered letter above the self-map \(X\to X\) is understood as the identity map, as required by the axiom; no further mathematical content is assigned to the cancelled word.}
\GSfield{Fate}{Already given in \citet[p.~38]{mumford1965picard}; this is a recurrence of 1967-143.1--1967-143.2 and 1968-151.1.}
\GSfield{Present status}{As of November 2026, the formulation is standard.}
\GSfield{Assessment}{literal statement: correct; publication match: exact; verification: read the result; development: annotated; confidence: high}
\end{GSitem}
\begin{GSitem}{1968-152.3}{E}{Abelian presheaves and effective gluing}
\GSfield{Kind; source}{definition; \GSpage{P000245}}
\GSfield{The notes}{A presheaf is contravariant, and a sheaf has a unique global section inducing every compatible family of local sections.}
\GSfield{Statement}{An abelian presheaf is a functor \(F:\mathcal C^{\mathrm{op}}\to\mathrm{Ab}\). For each cover \(\{U_i\to U\}\), the map \(F(U)\to\prod_iF(U_i)\) identifies \(F(U)\) with the equalizer of the two restrictions to \(\prod_{i,j}F(U_i\times_UU_j)\). Thus compatibility is equality after the two projections, and both existence and uniqueness are required.}
\GSfield{Fate}{\citet[p.~39]{mumford1965picard}; a repeated definition rather than an additional theorem.}
\GSfield{Present status}{As of November 2026, the definition is established.}
\GSfield{Assessment}{literal statement: correct; publication match: exact; verification: read the result; development: annotated; confidence: high}
\end{GSitem}
\begin{GSitem}{1968-152.4}{E}{Canonical sheafification}
\GSfield{Kind; source}{construction; \GSpage{P000246}}
\GSfield{The notes}{A parenthesis asserts a canonical procedure turning every presheaf into a sheaf for any topology.}
\GSfield{Statement}{For a site, the inclusion of abelian sheaves into abelian presheaves has a left adjoint \(F\mapsto aF\), characterized by \(\operatorname{Hom}(aF,G)\cong\operatorname{Hom}(F,G)\) for sheaves \(G\). For abelian groups sheafification is exact. In particular a presheaf may be nonzero while its associated sheaf vanishes because its sections vanish locally.}
\GSfield{Reconstruction}{The universal property and exactness specify the construction used in the following comparison criterion; no claim that sections are unchanged is made.}
\GSfield{Fate}{This is standard machinery in \citet[\S1]{mumford1965picard} and the sheaf-theoretic framework of \citet[\S3]{grothendieck1957homological}.}
\GSfield{Present status}{As of November 2026, the construction is established.}
\GSfield{Assessment}{literal statement: correct; publication match: exact; verification: read the result; development: annotated; confidence: high}
\end{GSitem}
\begin{GSitem}{1968-152.5}{E}{Derived cohomology, exact sequences and injective acyclicity}
\GSfield{Kind; source}{theorem; \GSpage{P000246}}
\GSfield{The notes}{The notes prescribe \(H^0(U,F)=F(U)\), a long exact sequence for short exact sequences, enough injectives, and zero cohomology for injective sheaves.}
\GSfield{Statement}{The category of abelian sheaves on a site has enough injectives. Deriving the left-exact sections functor gives \(H^q(U,F)\), with \(H^0(U,F)=F(U)\), functorial connecting morphisms, and the long exact cohomology sequence. An injective sheaf \(I\) has \(H^q(U,I)=0\) for \(q>0\); its degree-zero sections need not vanish.}
\GSfield{Reconstruction}{The omitted positive-degree qualification is restored. Cohomology on \(U\) means cohomology in the slice site.}
\GSfield{Fate}{\citet[pp.~37--39]{mumford1965picard}; \citet[\S\S2--3]{grothendieck1957homological}.}
\GSfield{Present status}{As of November 2026, the corrected derived-functor assertions are established.}
\GSfield{Assessment}{literal statement: correct with missing hypotheses; publication match: exact; verification: read the result; development: annotated; confidence: high}
\end{GSitem}
\begin{GSitem}{1968-152.6}{E}{Local-acyclicity criterion for comparison of sites}
\GSfield{Kind; source}{method; \GSpage{P000246}}
\GSfield{The notes}{For a structure-compatible functor \(f:T_1\to T_2\), the notation \(f^*F(U)=F(f(U))\) is used. If the sheaf associated with \(U\mapsto H^i(f(U),F)\) is zero, the displayed cohomology map is an isomorphism.}
\GSfield{Statement}{Let a functor \(u:\mathcal C\to\mathcal D\) induce a geometric morphism \(f:\operatorname{Sh}(\mathcal D)\to\operatorname{Sh}(\mathcal C)\) whose direct image is \(f_*F(U)=F(uU)\), and suppose the global-sections functors compose in the usual way. If \(R^if_*F=0\) for every \(i>0\), the Leray spectral sequence gives \(H^q(\mathcal C,f_*F)\cong H^q(\mathcal D,F)\). The higher direct image is computed by sheafifying the presheaf of cohomology over the inverse-image slices; when these are represented by \(uU\), this is \(U\mapsto H^i(\mathcal D/uU,F)\). The same criterion applies to the corresponding slices over an object.}
\GSfield{Reconstruction}{The notebook's star labels precomposition, which is the direct-image functor for the geometric morphism in this convention. The positive-degree condition and the precise slice interpretation are supplied; a condition at \(i=0\) would generally force the coefficient sheaf itself to vanish.}
\GSfield{Fate}{The exact application appears in \citet[\S7, p.~80, vanishing of \(R^iy_*\) for \(i>0\)]{mumford1965picard}; the homological argument is the Leray formalism of \citet[\S3]{grothendieck1957homological}.}
\GSfield{Present status}{As of November 2026, the criterion is established under the indicated functorial and acyclicity hypotheses.}
\GSfield{Assessment}{literal statement: correct with missing hypotheses; publication match: exact; verification: read the result; development: annotated; confidence: high}
\end{GSitem}
\begin{GSitem}{1968-152.7}{E}{A terminal moduli object and its candidate atlas}
\GSfield{Kind; source}{construction; \GSpage{P000246}}
\GSfield{The notes}{A new category \(T_M\) adjoins a terminal object. Candidate covering families must contain every isomorphism type, with each type occurring only discretely.}
\GSfield{Statement}{For a moduli category with representable isomorphism functors, an atlas is a family of objects whose maps to the moduli object are locally isomorphisms in the chosen topology and are jointly surjective after every base change. Meeting every isomorphism class is the pointwise-surjectivity part; discrete occurrence is a necessary unramified-type condition in this setting. Those two conditions alone do not prove local universality or the existence of the required pullbacks.}
\GSfield{Reconstruction}{The local extension requirement on the next page supplies information beyond the two enumerated conditions. Formal adjunction of a terminal object is not treated as an existence proof for a fine moduli space.}
\GSfield{Fate}{The precise modular-family construction is \citet[\S3, pp.~49--54]{mumford1965picard}; compare 1968-151.10--1968-151.11.}
\GSfield{Present status}{As of November 2026, the atlas construction is established for the stated moduli problem, with the local hypotheses retained.}
\GSfield{Assessment}{literal statement: correct with missing hypotheses; publication match: exact; verification: read the result; development: annotated; confidence: high}
\end{GSitem}
\begin{GSitem}{1968-152.8}{E}{Local extension property of a modular family}
\GSfield{Kind; source}{method; \GSpage{P000247}}
\GSfield{The notes}{A morphism from a subfamily over a submanifold is required to extend locally to a neighbourhood in the base of an arbitrary family; the picture calls it a tubular neighbourhood.}
\GSfield{Statement}{A sufficient precise version is the following. Let \(X\to S\) be a modular analytic family such that, for every test family \(Y\to T\), its isomorphism space \(I=\operatorname{Isom}(Y,X)\to T\) is a local biholomorphism. A family morphism defined on an analytic subspace \(T_0\subset T\) gives a section of \(I\to T\) over \(T_0\). Near each point of \(T_0\), the inverse of the appropriate local biholomorphic branch extends that section, hence extends the family morphism.}
\GSfield{Reconstruction}{This is a local assertion, not a global extension theorem or existence of a holomorphic tubular neighbourhood for every analytic subspace. The displayed crossed arrow does not uniquely encode a stronger formal deformation criterion.}
\GSfield{Fate}{The isomorphism-space criterion and modular families are treated in \citet[\S3, pp.~50--54]{mumford1965picard}. The notebook gives an informal extension formulation of this mechanism.}
\GSfield{Present status}{As of November 2026, the precise local lifting statement is established; the broad tubular-neighbourhood wording needs the reconstruction above.}
\GSfield{Assessment}{literal statement: correct with missing hypotheses; publication match: partial; verification: read the result; development: annotated; confidence: medium}
\end{GSitem}

\GStalk{153}{1968-09-16}{V. P. Palamodov}{P000247-P000251}{Jet bundles and the Spencer resolution}
The mathematical text is P000248--P000249. Palamodov introduces the jet prolongation of smooth functions, the holonomicity equations and the finite Spencer resolution. The notes expressly direct readers to Sweeney's Acta Mathematica paper, which was printed in June 1968 after a September 1966 submission. The listed researchers include Sweeney, Goldschmidt, Guillemin, Quillen, Sternberg and Singer; the unclear name is not treated as a separate attribution. The construction for unrestricted jets is separated from exactness and solvability for a general differential equation.

\begin{GSitem}{1968-153.1}{E}{Resolving compatibility conditions for a differential system}
\GSfield{Kind; source}{method; \GSpage{P000248}}
\GSfield{The notes}{The Spencer resolution is presented as a method for solving \(p(x,D)u=w\), with a matrix \(p\) and vector-valued unknown and data.}
\GSfield{Statement}{A differential operator \(P:\mathcal E\to\mathcal F\) of order \(\mu\) factors as \(\mathcal E\xrightarrow{j_\mu}J_\mu E\xrightarrow{p_\mu}F\). Prolonged equations and their compatibility operators yield Spencer complexes under constant-rank and regularity hypotheses. Their formal exactness is distinct from exactness on smooth sections. A locally exact compatibility complex gives local solvability for data satisfying its next compatibility equation; the unrestricted-jet resolution alone does not prove this for every \(P\).}
\GSfield{Reconstruction}{The hypotheses are restored at the level of the method. Neither formal integrability nor the use of jets is presented as an unconditional smooth existence theorem.}
\GSfield{Fate}{The expressly named source is \citet[Introduction; Chapter~I, \S3 and Proposition~8.3]{sweeney1968neumann}, received 4 September 1966, revised 30 November 1967, and printed June 1968. It gives exactness results for regular constant-coefficient systems and real-analytic germs under its stated hypotheses. Compare \citet[\S\S1.2,1.6]{spencer1969overdetermined}.}
\GSfield{Present status}{As of November 2026, the method and its specified applications are established; no general smooth-solvability claim follows from the note.}
\GSfield{Assessment}{literal statement: correct with missing hypotheses; publication match: partial; verification: read the result; development: annotated; confidence: high}
\end{GSitem}
\begin{GSitem}{1968-153.2}{E}{Finite jet bundles and an injective map of sheaves}
\GSfield{Kind; source}{construction; \GSpage{P000248}}
\GSfield{The notes}{For the sheaf \(\mathcal O\) of smooth-function germs, \(j_\mu:\mathcal O\to J_\mu\) is said to be an embedding. The zeroth-order case is followed by an ambiguous zero symbol.}
\GSfield{Statement}{On a smooth \(n\)-manifold, \(J_\mu\) is the vector bundle of \(\mu\)-jets of scalar functions. A local section of this bundle has coordinates \(\phi_I(x)\), \(0\le |I|\le\mu\), using unnormalized derivatives. Prolongation sends \(f\) to the section \(j_\mu f=(\partial^If)_I\). As a map of sheaves of real vector spaces it is injective because its zeroth component is \(f\). It is not the evaluation map from a function germ to its finite jet at one point, which is not injective. For \(\mu=0\), \(J_0\cong\mathcal O\) and \(j_0\) is the identity.}
\GSfield{Reconstruction}{Bundles, sheaves of their sections, and individual jet fibres are distinguished. The subsequent map at order zero is to zero, yielding \(0\to\mathcal O\xrightarrow{\mathrm{id}}\mathcal O\to0\).}
\GSfield{Fate}{\citet[Chapter~I, \S1 and Proposition~6.2]{sweeney1968neumann}; \citet[\S\S1.1,1.3]{spencer1969overdetermined}.}
\GSfield{Present status}{As of November 2026, the construction and sheaf-level injectivity are established.}
\GSfield{Assessment}{literal statement: correct with missing hypotheses; publication match: exact; verification: read the result; development: annotated; confidence: high}
\end{GSitem}
\begin{GSitem}{1968-153.3}{E}{Holonomic jet sections are characterized by derivative compatibility}
\GSfield{Kind; source}{theorem; \GSpage{P000248-P000249}}
\GSfield{The notes}{The notes first compare \(\partial\phi_0/\partial x_j\) with \(\phi_j\), then require \(\partial_j\phi_I=\phi_{I+1_j}\) for every \(|I|\le\mu-1\).}
\GSfield{Statement}{A smooth local jet section \(\phi=(\phi_I)_{|I|\le\mu}\) is holonomic, meaning \(\phi=j_\mu f\) for a smooth function \(f\), if and only if \(\partial_j\phi_I=\phi_{I+e_j}\) for all \(j\) and \(|I|\le\mu-1\). Necessity follows by differentiation; sufficiency follows inductively from \(f=\phi_0\). For normalized Taylor coefficients the formula would instead contain the corresponding multi-index factors.}
\GSfield{Reconstruction}{The derivative rather than Taylor-coefficient normalization is fixed by the first-order example. Gelfand's intervention is the coordinate presentation, not a separate conjecture.}
\GSfield{Fate}{This is the kernel statement in \citet[Chapter~I, Proposition~5.1(ii)]{sweeney1968neumann} and \citet[equation~(1.3.2)]{spencer1969overdetermined}.}
\GSfield{Present status}{As of November 2026, the criterion is established and directly verified by induction.}
\GSfield{Assessment}{literal statement: correct; publication match: exact; verification: read the result; development: annotated; confidence: high}
\end{GSitem}
\begin{GSitem}{1968-153.4}{E}{The Spencer operator as actual minus formal differentiation}
\GSfield{Kind; source}{construction; \GSpage{P000249}}
\GSfield{The notes}{The first operator is \(\sum_j(\partial_j\phi_I-\phi_{I+e_j})dx_j\), and its continuation is called \(D=d-\delta\).}
\GSfield{Statement}{For a jet-valued differential form \(\phi\in\Omega^r\otimes J_k\), the Spencer operator is locally \((D\phi)_I=\sum_j dx_j\wedge(\partial_j\phi_I-\phi_{I+e_j})\), \(|I|\le k-1\). Thus \(D:\Omega^r\otimes J_k\to\Omega^{r+1}\otimes J_{k-1}\). It is a globally defined first-order differential operator; in these coordinates it is the difference of actual and formal exterior differentiation. The two summands separately depend on coordinates, whereas their difference does not.}
\GSfield{Reconstruction}{The formula includes exterior degree, jet order, and the order of wedge factors. It is not linear over smooth functions for the usual module structures.}
\GSfield{Fate}{\citet[Chapter~I, \S5, equation~(5.4)]{sweeney1968neumann}; \citet[equation~(1.3.1)\(^{\prime}\)]{spencer1969overdetermined}.}
\GSfield{Present status}{As of November 2026, the operator is established; its invariant character is essential to the resolution.}
\GSfield{Assessment}{literal statement: correct; publication match: exact; verification: read the result; development: annotated; confidence: high}
\end{GSitem}
\begin{GSitem}{1968-153.5}{E}{The one-form formula requires a wedge-order correction}
\GSfield{Kind; source}{claim; \GSpage{P000249}}
\GSfield{The notes}{The displayed continuation is \(d(\phi_I^jdx_j)-\phi_{I+e_s}^jdx_j\wedge dx_s\), with the derivative-index correction also unclear.}
\GSfield{Statement}{With \(D=d-\delta\) as in the preceding degree-zero formula, the correct expression is 
\[
D(\phi_I^jdx_j)=\sum_{s,j}(\partial_s\phi_I^j-\phi_{I+e_s}^j)\,dx_s\wedge dx_j.
\]
Equivalently the formal term is \(+\sum_{s,j}\phi_{I+e_s}^jdx_j\wedge dx_s\). Retaining the printed minus sign together with the reversed wedge order is inconsistent with the previous definition and does not generally give \(D^2=0\).}
\GSfield{Reconstruction}{The scan confirms the displayed order \(dx_j\wedge dx_s\); the choice \(e_s\) is dictated by the differentiation direction. The repair is mathematical, not a claim that the transcription omitted a legible sign.}
\GSfield{Fate}{The corrected order is explicit in \citet[equation~(1.3.1)\(^{\prime}\)]{spencer1969overdetermined}; compare \citet[Chapter~I, equation~(5.4)]{sweeney1968neumann}.}
\GSfield{Present status}{As of November 2026, the literal sign-and-order combination is false; the corrected formula gives the standard differential.}
\GSfield{Assessment}{literal statement: false as written; publication match: exact; verification: read the result; development: annotated; confidence: high}
\end{GSitem}
\begin{GSitem}{1968-153.6}{E}{The finite exact Spencer sequence for unrestricted jets}
\GSfield{Kind; source}{theorem; \GSpage{P000249}}
\GSfield{The notes}{The sequence \(\mathcal O\to J_\mu\to T^*\otimes J_{\mu-1}\to\Lambda^2T^*\otimes J_{\mu-2}\to\cdots\) is said to terminate and be exact.}
\GSfield{Statement}{For every \(\mu\ge0\), the sequence of sheaves of real vector spaces 
\[
0\longrightarrow\mathcal C^\infty_M\xrightarrow{j_\mu}J_\mu\xrightarrow D\Omega^1\otimes J_{\mu-1}\xrightarrow D\cdots\xrightarrow D\Omega^n\otimes J_{\mu-n}\longrightarrow0
\]
is exact, with \(J_k=0\) for \(k<0\). The identities \(d^2=\delta^2=0\) and \(d\delta+\delta d=0\) give \(D^2=0\); exactness follows inductively in jet order using the exact polynomial Koszul complex for the highest-order symbols. The same result holds for the jet bundle of any smooth vector bundle.}
\GSfield{Reconstruction}{This resolves the sheaf of unrestricted smooth sections. For a subbundle of jets defined by a differential equation, exactness is a separate analytic issue, as explained in 1968-153.1. The individual maps are not morphisms of the usual smooth-function module structures.}
\GSfield{Fate}{Exact publication match: \citet[Chapter~I, Proposition~6.2, p.~238]{sweeney1968neumann}; reproduced as equation~(1.3.3) with proof in \citet[\S1.3]{spencer1969overdetermined}.}
\GSfield{Present status}{As of November 2026, the finite jet resolution is established and was available before the seminar.}
\GSfield{Assessment}{literal statement: correct; publication match: exact; verification: read the result; development: annotated; confidence: high}
\end{GSitem}
\begin{GSitem}{1968-153.7}{F}{The unresolved comparison with a higher-order pseudo-object}
\GSfield{Kind; source}{remark; \GSpage{P000248}}
\GSfield{The notes}{A marginal comparison reads either \ru{псевдогр.} or \ru{псевдооп.}, followed by ``one order higher''.}
\GSfield{Statement}{No mathematical assertion beyond the legible jet construction can be uniquely recovered from this phrase. The first reading might refer to pseudogroups and the second to an operator construction, but the note does not supply a definition or a formula deciding between them.}
\GSfield{Reconstruction}{Neither reading is silently promoted to an additional theorem; the names listed above it belong to the surrounding Spencer-theory attribution.}
\GSfield{Fate}{\citet[Introduction and bibliography]{sweeney1968neumann} explicitly connects Spencer theory with pseudogroups, making that reading plausible but not conclusive. The alternative remains a transcription ambiguity.}
\GSfield{Present status}{As of November 2026, the marginal comparison remains unidentified; it is not an open mathematical problem.}
\GSfield{Assessment}{literal statement: not determinable; publication match: thematic only; verification: unresolved lead; development: annotated; confidence: low}
\end{GSitem}

\GStalk{154}{1968-09-23}{A. Yu. Geronimus; I. R. Shafarevich}{P000249-P000254}{The elliptic moduli site and a foliation example}
P000250--P000251 applies Mumford's machinery to pointed elliptic curves: the upper-half-plane family, the modular action, level-three comparison and the class of order twelve. The integral cohomology formula needs its degree-zero and periodic-even-degree corrections, and the assertion over general fields concerns the Picard invariant rather than ordinary integral Betti cohomology. On P000252 Shafarevich proposes a site of transverse patches for a differential equation and makes a compressed Euler-characteristic remark. The following Kazhdan and Palamodov headings on P000253--P000254 belong to records 155--156 and are excluded here.

\begin{GSitem}{1968-154.1}{E}{Pointed elliptic families and Cartesian morphisms}
\GSfield{Kind; source}{definition; \GSpage{P000250}}
\GSfield{The notes}{The objects are analytic families of complex tori with a distinguished section \(\sigma:S\to X\).}
\GSfield{Statement}{An analytic family of pointed elliptic curves is a proper smooth holomorphic map \(\pi:E\to S\) with connected genus-one fibres, together with a holomorphic section \(e:S\to E\). Morphisms consist of a map of bases and a section-preserving isomorphism with the pulled-back family. The section makes each fibre an elliptic group curve and removes translation automorphisms, but retains the origin-preserving finite automorphisms.}
\GSfield{Reconstruction}{Properness, smoothness and connected fibres specify the intended family category; arbitrary maps with pointwise torus fibres are not silently admitted.}
\GSfield{Fate}{The category is the elliptic-curve case of \citet[\S4]{mumford1965picard}; see also \citet[\S2]{hain2008moduli}.}
\GSfield{Present status}{As of November 2026, the definition is established.}
\GSfield{Assessment}{literal statement: correct with missing hypotheses; publication match: exact; verification: read the result; development: annotated; confidence: high}
\end{GSitem}
\begin{GSitem}{1968-154.2}{E}{The coarse j-line has no smooth family realizing its identity parameter}
\GSfield{Kind; source}{claim; \GSpage{P000250}}
\GSfield{The notes}{The quotient \(S'=\mathbb H/\Gamma\) is called coarse; the notes say there is a family with base \(S'\), then explain that the family over \(\mathbb H\) is not induced from it and that isotrivial twists are missed.}
\GSfield{Statement}{For \(\Gamma=SL_2(\mathbb Z)\), the coarse quotient is the \(j\)-line. There is no smooth pointed elliptic family over the whole coarse line whose \(j\)-invariant is the identity. Locally near the square or hexagonal elliptic point, a family over a disc has a lifted period map to \(\mathbb H\); the coarse map is locally a square or cube of a local coordinate, so it cannot compose to a simple parameter on that disc. Separately, nontrivial isotrivial twists show that a pointwise classifying coarse map cannot classify all families.}
\GSfield{Reconstruction}{The statement ``there is a family'' is false if it means one realizing each coarse parameter, as the surrounding discussion suggests. Trivial families over the line certainly exist but do not have the required classifying map. The two obstructions are distinguished.}
\GSfield{Fate}{The exceptional local parameters are discussed in \citet[\S4, p.~59]{mumford1965picard}. The local lifting obstruction and isotrivial construction are explicit in \citet[Exercises~25--26, \S3.6]{hain2008moduli}. This repeats GP-1968-01.}
\GSfield{Present status}{As of November 2026, the coarse-versus-fine distinction and nonexistence are established.}
\GSfield{Assessment}{literal statement: false as written; publication match: exact; verification: read the result; development: annotated; confidence: high; problem \hyperref[prob:GP-1968-01]{GP-1968-01}}
\end{GSitem}
\begin{GSitem}{1968-154.3}{E}{The marked universal family and its modular action}
\GSfield{Kind; source}{construction; \GSpage{P000250-P000251}}
\GSfield{The notes}{The family is \((\mathbb H\times\mathbb C)/((\tau,z)\sim(\tau,z+m\tau+n))\). The group acts by \(\tau'=(a\tau+b)/(c\tau+d)\), \(z'=z/(c\tau+d)\).}
\GSfield{Statement}{The quotient \(\mathcal E_{\mathbb H}=(\mathbb H\times\mathbb C)/\mathbb Z^2\), with translations \(z\mapsto z+m\tau+n\), is a holomorphic family with zero section and a distinguished oriented homology basis. Every \(\gamma=\left(\begin{smallmatrix}a&b\\c&d\end{smallmatrix}\right)\in SL_2(\mathbb Z)\) lifts by the displayed formula because \((c\tau+d)^{-1}(\mathbb Z\tau+\mathbb Z)=\mathbb Z\gamma\tau+\mathbb Z\). The quotient family over the analytic stack \([\mathbb H/SL_2(\mathbb Z)]\) is universal in the family-category sense.}
\GSfield{Reconstruction}{The notebook's underlined ``problem'' is the task of specifying this lifted action, answered on the next page; it is not a newly open problem. Stack notation is modern reconstruction.}
\GSfield{Fate}{\citet[\S7]{mumford1965picard}; an explicit modern formulation is \citet[\S3.6, Theorem~3.18]{hain2008moduli}.}
\GSfield{Present status}{As of November 2026, the construction and action are established.}
\GSfield{Assessment}{literal statement: correct; publication match: exact; verification: read the result; development: annotated; confidence: high}
\end{GSitem}
\begin{GSitem}{1968-154.4}{E}{Three marked boundary points represent two exceptional orbits}
\GSfield{Kind; source}{remark; \GSpage{P000250}}
\GSfield{The notes}{The fundamental-domain drawing marks three fixed points and says that branching occurs there.}
\GSfield{Statement}{For the effective \(PSL_2(\mathbb Z)\)-action on \(\mathbb H\), there are two orbits with nontrivial stabilizers, represented by \(i\) and \(e^{2\pi i/3}\), with orders two and three. The two boundary representatives of the hexagonal orbit account for the third marked point in a closed fundamental domain. For \(SL_2(\mathbb Z)\), \(-I\) fixes every base point, so the generic stabilizer is \(C_2\), and the exceptional stabilizers are \(C_4,C_6\).}
\GSfield{Reconstruction}{Counting points drawn in one closed domain is separated from counting orbits and from the non-effective central action.}
\GSfield{Fate}{\citet[\S4]{mumford1965picard}; \citet[\S3.5, Example~3.14]{hain2008moduli}.}
\GSfield{Present status}{As of November 2026, the stabilizer and branching descriptions are established after this qualification.}
\GSfield{Assessment}{literal statement: correct with missing hypotheses; publication match: exact; verification: read the result; development: annotated; confidence: high}
\end{GSitem}
\begin{GSitem}{1968-154.5}{E}{Constant coefficients and the cohomology of one family}
\GSfield{Kind; source}{construction; \GSpage{P000251}}
\GSfield{The notes}{The coefficient sheaf has one copy of \(\mathbb Z\) for each connected component, and \(H^q(\pi,\mathbb Z)=H^q(S,\mathbb Z)\).}
\GSfield{Statement}{On the local-isomorphism slice over a family \(\pi:E\to S\), the constant coefficient sheaf assigns locally constant integer-valued functions to each base. If a base \(U\) is locally connected, this group is \(\prod_{\pi_0(U)}\mathbb Z\). The slice-site comparison identifies \(H^q(\pi,\mathbb Z)\) with ordinary sheaf cohomology \(H^q(S,\mathbb Z)\). These are not the cohomology groups of the total space \(E\), nor yet those of the terminal moduli object.}
\GSfield{Reconstruction}{The finite three-summand drawing is harmless; for an infinite set of components one must take a product. Connected components are open for the analytic spaces under discussion.}
\GSfield{Fate}{The comparison uses \citet[\S\S2--3]{mumford1965picard} and repeats 1968-151.8--1968-151.9 with specified coefficients.}
\GSfield{Present status}{As of November 2026, the slice identification is established.}
\GSfield{Assessment}{literal statement: correct; publication match: exact; verification: read the result; development: annotated; confidence: high}
\end{GSitem}
\begin{GSitem}{1968-154.6}{E}{The level-three subsite and comparison with group cohomology}
\GSfield{Kind; source}{method; \GSpage{P000251}}
\GSfield{The notes}{A subcategory of \(\Gamma\)-sets has point stabilizers inside \(\Gamma_3\). The assignment \(S\mapsto(\mathbb H\times S)/\Gamma\) induces a cohomology isomorphism after adjoining terminal objects.}
\GSfield{Statement}{Let \(\Gamma=SL_2(\mathbb Z)\) and \(\Gamma(3)=\ker(\Gamma\to SL_2(\mathbb Z/3))\). This subgroup is torsion-free. For a discrete \(\Gamma\)-set \(S\) with all stabilizers contained in \(\Gamma(3)\), the diagonal action on \(\mathbb H\times S\) is free and its quotient is a complex analytic base carrying the descended elliptic family. These objects form the comparison subsite. Contractibility of \(\mathbb H\), local acyclicity and the covering criterion give \(H^q(\mathcal M_{1,1}^{\mathrm{an}},\mathbb Z)\cong H^q(\Gamma,\mathbb Z)\).}
\GSfield{Reconstruction}{The freeness is for the diagonal action on these admissible objects; the full modular action on the upper half-plane is not free. The formal terminal object is adjoined because the one-point \(\Gamma\)-set does not satisfy the stabilizer restriction.}
\GSfield{Fate}{Exact construction and comparison: \citet[\S7, pp.~79--80]{mumford1965picard}. The classifying-space interpretation is \citet[Proposition~3.13]{hain2008moduli}.}
\GSfield{Present status}{As of November 2026, the comparison is established.}
\GSfield{Assessment}{literal statement: correct; publication match: exact; verification: read the result; development: annotated; confidence: high}
\end{GSitem}
\begin{GSitem}{1968-154.7}{E}{Integral cohomology of the elliptic moduli site is periodic}
\GSfield{Kind; source}{claim; \GSpage{P000251}}
\GSfield{The notes}{The formula is transcribed as \(H^q(M,\mathbb Z)=0\) for \(q\ne2\), and \(\mathbb Z_{12}\) for \(q=2\); the degree sign is heavily retraced.}
\GSfield{Statement}{For the uncompactified analytic moduli stack of pointed elliptic curves,
\[
H^q(\mathcal M_{1,1}^{\mathrm{an}},\mathbb Z)=
\begin{cases}\mathbb Z,&q=0,\\\mathbb Z/12,&q>0\text{ even},\\0,&q\text{ odd}.\end{cases}
\]
More precisely its ring is \(\mathbb Z[c]/(12c)\), with \(|c|=2\). This differs both from the coarse \(j\)-line, whose positive-degree integral cohomology vanishes, and from the compactified moduli stack.}
\GSfield{Reconstruction}{The higher-degree correction follows from \(SL_2(\mathbb Z)\cong C_4*_{C_2}C_6\). In every positive even degree the Mayer--Vietoris map \(\mathbb Z/4\oplus\mathbb Z/6\to\mathbb Z/2\) is \((a,b)\mapsto a-b\) modulo two, is surjective, and has cyclic kernel of order twelve. The odd groups vanish; degree zero is \(\mathbb Z\). Choosing the degree-two class restricting to generators of both cyclic groups gives the displayed ring. This is an explicit verification of a correction, not a novel cohomology theorem.}
\GSfield{Fate}{The moduli-to-group comparison is \citet[\S7]{mumford1965picard}; the degree-one and degree-two computation and group presentation are in \citet[Proposition~3.13 and presentation~(3)]{hain2008moduli}. The periodic extension is verified by the amalgam calculation above.}
\GSfield{Present status}{As of November 2026, the corrected integral cohomology is established; the literal concentration in degree two is false.}
\GSfield{Assessment}{literal statement: false as written; publication match: exact; verification: read the result; development: annotated; confidence: high}
\end{GSitem}
\begin{GSitem}{1968-154.8}{B}{The assertion over other fields concerns the Picard group}
\GSfield{Kind; source}{claim; \GSpage{P000251}}
\GSfield{The notes}{The notes say that the answer remains true over any field of characteristic different from two.}
\GSfield{Statement}{For every field \(k\), the algebraic Picard group of the stack of elliptic curves is \(\operatorname{Pic}(\mathcal M_{1,1,k})\cong\mathbb Z/12\), generated by the Hodge line bundle. This is a statement about invertible sheaves. Ordinary integral Betti cohomology is the complex-analytic invariant of the preceding item and is not defined by that formula for an arbitrary field.}
\GSfield{Reconstruction}{The text switches invariants without explaining the comparison. Mumford's theorem as published assumes characteristic different from both two and three. Its Picard conclusion over fields of characteristic three, and indeed two, has a later exact extension; this does not extend the entire displayed Betti-cohomology table to all fields.}
\GSfield{Fate}{\citet[\S6, Main Theorem, p.~73]{mumford1965picard} proves the characteristic-not-two-or-three case. \citet[Theorem~1.1]{fultonolsson2010picard}, circulated on 17 April 2007 and published in 2010, proves \(\operatorname{Pic}(\mathcal M_{1,1,S})\cong\mathbb Z/12\oplus\operatorname{Pic}(\mathbb A^1_S)\) for reduced \(S\), or when two is invertible on \(S\). For \(S=\operatorname{Spec}k\), the second summand is zero.}
\GSfield{Present status}{As of November 2026, the algebraic Picard assertion is established for every field. The literal general-field cohomology statement is under-specified until its invariant and topology are fixed.}
\GSfield{Assessment}{literal statement: not determinable; publication match: partial; verification: read the result; development: annotated; confidence: high}
\end{GSitem}
\begin{GSitem}{1968-154.9}{E}{What the invariant modulo twelve means for a family}
\GSfield{Kind; source}{remark; \GSpage{P000251}}
\GSfield{The notes}{A boxed conclusion says that families have an invariant modulo twelve.}
\GSfield{Statement}{The universal Hodge line bundle \(\lambda\) generates the cyclic Picard group of order twelve. For a smooth pointed elliptic family \(\pi:E\to S\), its pullback is \(\lambda_S=\pi_*\Omega^1_{E/S}\), and the nonvanishing discriminant gives a canonical trivialization of \(\lambda_S^{\otimes12}\). In the complex analytic setting, \(c_1(\lambda_S)\in H^2(S,\mathbb Z)\) is the pullback of the universal order-twelve class and is killed by twelve.}
\GSfield{Reconstruction}{A family over an arbitrary base does not automatically yield a scalar in \(\mathbb Z/12\). It yields a line-bundle class, or its cohomology class on that base, whose universal source has order twelve.}
\GSfield{Fate}{\citet[\S6]{mumford1965picard}; \citet[Introduction, equations~(1-3)--(1-7)]{fultonolsson2010picard}; the analytic generator is \citet[Theorem~6.6]{hain2008moduli}.}
\GSfield{Present status}{As of November 2026, the Hodge-line interpretation and its order-twelve relation are established.}
\GSfield{Assessment}{literal statement: correct with missing hypotheses; publication match: exact; verification: read the result; development: annotated; confidence: high}
\end{GSitem}
\begin{GSitem}{1968-154.10}{B}{Shafarevich: a site of transverse patches and holonomy}
\GSfield{Kind; source}{construction; \GSpage{P000252}; I. R. Shafarevich}
\GSfield{The notes}{For an ordinary differential equation, objects are transverse patches and maps follow trajectories to another patch. A fibre product is drawn.}
\GSfield{Statement}{For a nonsingular smooth foliation, choose a complete transversal \(T\). Germs of holonomy maps between open subsets of \(T\) form an \(\acute{e}\)tale groupoid \(G\). Equivariant sheaves on \(T\), or the corresponding local site with ordinary open covers and holonomy identifications, retain the transverse topology and isotropy of the foliation. Fibre products of local transverse maps exist locally after restricting patches. Singular points require an additional definition and are not covered by this nonsingular construction.}
\GSfield{Reconstruction}{The holonomy-groupoid interpretation makes the drawing precise but is a modern reconstruction. The note does not say whether path data, holonomy germs or only pointwise trajectory incidence define its morphisms; those choices can produce different groupoids.}
\GSfield{Fate}{A precise later development is \citet[\S\S1.1,2.1--2.2 and Theorem~3.1]{moerdijk1991foliations}: the classifying topos of an \(\acute{e}\)tale groupoid agrees in homotopy and cohomology with its classifying space. This matches the transverse-site method, not the unqualified singular-flow or Euler formula in the next remarks. The paper appeared in 1991; no earlier circulation date was verified.}
\GSfield{Present status}{As of November 2026, the nonsingular holonomy-site construction is established; the exact morphism convention of this notebook example remains only partially identified.}
\GSfield{Assessment}{literal statement: correct with missing hypotheses; publication match: partial; verification: read the result; development: annotated; confidence: medium}
\end{GSitem}
\begin{GSitem}{1968-154.11}{E}{For a genuine leaf fibration the transverse cohomology is that of the base}
\GSfield{Kind; source}{claim; \GSpage{P000252}; I. R. Shafarevich}
\GSfield{The notes}{The abbreviated sentence says that, for a fibration, the cohomology of the base and the system are \ru{трив.}, and mentions a final object.}
\GSfield{Statement}{If a foliation is the foliation by connected fibres of a surjective smooth submersion \(p:M\to B\) with manifold leaf space \(B\), its holonomy groupoid on a complete transversal is equivalent to the \(\acute{e}\)tale groupoid presenting \(B\). Its classifying topos is therefore equivalent to ordinary sheaves on \(B\), and the corresponding cohomology groups agree. They need not vanish: a nonempty connected base has \(H^0(B,\mathbb Z)=\mathbb Z\).}
\GSfield{Reconstruction}{The abbreviation may mean that the identification is trivial, rather than that both cohomologies vanish. The Russian and scan do not decide the omitted predicate. The existence of a terminal object by itself never forces positive-degree cohomology to vanish.}
\GSfield{Fate}{The groupoid invariance is \citet[\S2.2]{moerdijk1991foliations}; for a groupoid presenting a space the classifying topos is its usual sheaf topos, as explained in \S2.4 of that paper. The fibration case follows directly.}
\GSfield{Present status}{As of November 2026, the base comparison is established; a reading asserting vanishing is false and the notebook wording remains ambiguous.}
\GSfield{Assessment}{literal statement: not determinable; publication match: partial; verification: read the result; development: annotated; confidence: medium}
\end{GSitem}
\begin{GSitem}{1968-154.12}{F}{Euler characteristic as a difference of cycle and singular-point counts}
\GSfield{Kind; source}{claim; \GSpage{P000252}; I. R. Shafarevich}
\GSfield{The notes}{For an equation on the plane and the constant sheaf of its Grothendieck topology, the Euler characteristic is said to be the difference between the number of limit cycles and the number of singular points.}
\GSfield{Statement}{No precise theorem is determined without specifying the vector-field class, treatment of singular leaves, cohomology theory and support condition, finiteness of the relevant groups, and which signed difference is intended. For example, for the constant nonzero horizontal vector field on \(\mathbb R^2\), the ordinary transverse leaf space is \(\mathbb R\), with ordinary constant-sheaf Euler characteristic one, while both stated counts are zero. Thus the unrestricted ordinary-cohomology interpretation is false.}
\GSfield{Reconstruction}{The counterexample does not identify the author's intended compact-support, reduced or relative invariant. It demonstrates why the omitted convention is substantive. The phrase does not supply a definite open problem, and no problem-index entry is invented.}
\GSfield{Fate}{Searches combining the named speaker, transverse sites, Euler characteristic, singular points and limit cycles did not identify a publication with these exact hypotheses and conclusion. \citet[Theorem~3.1]{moerdijk1991foliations} provides the general cohomological framework but no such cycle-count formula.}
\GSfield{Present status}{As of November 2026, the intended invariant and theorem remain unidentified; the broad ordinary-cohomology reading is refuted by the stated elementary example.}
\GSfield{Assessment}{literal statement: not determinable; publication match: none found; verification: unresolved lead; development: annotated; confidence: low}
\end{GSitem}

\GStalk{155}{1968-09-30}{D. A. Kazhdan}{P000252-P000255}{Ihara: curves, coverings, and congruence monodromy}
The actual heading is on P000253. Kazhdan identifies Ihara by journal and pages 107--121, then begins the analogy between hyperbolic uniformization and curves over finite fields. The surviving mathematical content is introductory: coordinate rings, points, morphisms, and unramified fibres. None of Ihara\textquoteright s three precise congruence-monodromy conjectures is written out, so they are not imported into the problem index. The paper was received in August 1967 and published in 1968 \citep[Introduction]{ihara1968monodromy}.

\begin{GSitem}{1968-155.1}{E}{Hyperbolic uniformization and unramified coverings}
\GSfield{Kind; source}{construction; \GSpage{P000253}}
\GSfield{The notes}{\(X=\Gamma\backslash M\), with \(M\) the Lobachevsky plane; \(\Gamma\) describes both \(X\) and its unramified coverings.}
\GSfield{Statement}{For a torsion-free discrete subgroup \(\Gamma<\operatorname{PSL}_2(\mathbb R)\), the projection \(\mathbb H\to\Gamma\backslash\mathbb H\) is a universal covering. Connected coverings correspond, up to isomorphism, to conjugacy classes of subgroups of \(\Gamma\); finite covers correspond to finite-index subgroups. Compact Riemann surfaces of genus at least two admit such uniformization.}
\GSfield{Reconstruction}{Torsion-freeness, connectedness, genus restriction, and conjugacy of subgroups make the notebook shorthand precise.}
\GSfield{Fate}{This is Ihara\textquoteright s classical starting analogy, not his new finite-field correspondence \citep[Introduction]{ihara1968monodromy}.}
\GSfield{Present status}{As of November 2026: established mathematics predating this talk; the qualification in the statement is essential.}
\GSfield{Assessment}{literal statement: correct with missing hypotheses; publication match: partial; verification: read the result; development: annotated; confidence: high}
\end{GSitem}
\begin{GSitem}{1968-155.2}{E}{Coordinate rings, closed points, and contravariance}
\GSfield{Kind; source}{definition; \GSpage{P000253}}
\GSfield{The notes}{\(A=k[x,y]/(f)\); maximal ideals are called geometric points, and \(Y\to X\) induces \(A\to B\).}
\GSfield{Statement}{For an affine plane curve \(X=\operatorname{Spec}(k[x,y]/(f))\), closed points correspond to maximal ideals, while \(\bar k\)-valued geometric points are homomorphisms \(A\to\bar k\). An affine morphism \(\operatorname{Spec}B\to\operatorname{Spec}A\) is equivalent to a \(k\)-algebra homomorphism \(A\to B\). A closed point may have several geometric points above it.}
\GSfield{Fate}{Pre-existing algebraic geometry used in Ihara\textquoteright s finite-field setting \citep[\S0]{ihara1968monodromy}.}
\GSfield{Present status}{As of November 2026: established mathematics predating this talk; the qualification in the statement is essential.}
\GSfield{Assessment}{literal statement: false as written; publication match: partial; verification: read the result; development: annotated; confidence: high}
\end{GSitem}
\begin{GSitem}{1968-155.3}{E}{Reduced fibres as a criterion for unramifiedness}
\GSfield{Kind; source}{claim; \GSpage{P000253}}
\GSfield{The notes}{For a maximal ideal \(m\subset A\), if \(B/mB\) has no nilpotents, the map is unramified at \(m\).}
\GSfield{Statement}{For a finite morphism of finite-type schemes over a perfect field, at a closed point \(x\) the morphism is unramified at every point above \(x\) if and only if its fibre algebra is a finite product of finite separable extensions of \(\kappa(x)\). Over finite fields the residue field is perfect, so reducedness of that finite fibre algebra suffices. Flatness is an additional condition for an \emph{\'{e}tale} morphism.}
\GSfield{Reconstruction}{Finiteness and separability are supplied. A reduced fibre over an imperfect field need not be unramified.}
\GSfield{Fate}{The finite unramified extensions are the objects of Ihara\textquoteright s Conjecture 2; the fibre test is older background \citep[p.~107]{ihara1968monodromy}.}
\GSfield{Present status}{As of November 2026: established mathematics predating this talk; the qualification in the statement is essential.}
\GSfield{Assessment}{literal statement: correct with missing hypotheses; publication match: partial; verification: read the result; development: annotated; confidence: high}
\end{GSitem}

\GStalk{156}{1968-10-14}{V. P. Palamodov}{P000253-P000257}{Jets along the diagonal and the Spencer complex}
The talk begins on P000254 and continues on P000255. Palamodov introduces jets and the trivial Spencer resolution; Gelfand reformulates jets as functions modulo powers of the diagonal ideal and sketches a many-point version. Sweeney, Malgrange, and Quillen are bibliographical leads. The explicit Sweeney citation resolves to \citet[Chapter~I]{sweeney1968neumann}; the notebook does not specify which Malgrange or Quillen work it means.

\begin{GSitem}{1968-156.1}{E}{Jet bundles and holonomic jet sections}
\GSfield{Kind; source}{definition; \GSpage{P000254}}
\GSfield{The notes}{\(O\hookrightarrow J^k\); a jet is a family \(\{\varphi_i:|i|\leq k\}\) transforming like derivatives.}
\GSfield{Statement}{For a smooth manifold \(X\), \(J^k(X,\mathbb R)_x\) consists of germs of smooth functions at \(x\) modulo equality of all derivatives through order \(k\). The map \(j^k:C^\infty_X\to\mathcal J^k\) sends a function to its jet section. It is injective as a map of sheaves of real vector spaces, but is not \(C^\infty_X\)-linear for the usual module structure on jet sections.}
\GSfield{Reconstruction}{The sheaf consists of arbitrary smooth sections of the jet bundle; holonomic sections are the image of the jet map.}
\GSfield{Fate}{The construction is explained in \citet[Chapter~I, \S1]{sweeney1968neumann}, published before this October talk.}
\GSfield{Present status}{As of November 2026: established mathematics predating this talk; the qualification in the statement is essential.}
\GSfield{Assessment}{literal statement: correct with missing hypotheses; publication match: exact; verification: read the result; development: annotated; confidence: high}
\end{GSitem}
\begin{GSitem}{1968-156.2}{E}{Jets from a power of the diagonal ideal}
\GSfield{Kind; source}{construction; \GSpage{P000254}; I. M. Gelfand}
\GSfield{The notes}{\(\varphi_1\sim\varphi_2\) when \(\varphi_1-\varphi_2\in\mathcal J^{k+1}\), with \(\mathcal J\) the ideal of functions vanishing on the diagonal of \(X\times X\).}
\GSfield{Statement}{Let \(\mathcal I_\Delta\) be the sheaf of smooth functions vanishing on the diagonal. The sheaf \(C^\infty_{X\times X}/\mathcal I_\Delta^{k+1}\), restricted to a neighbourhood germ of \(\Delta\) and viewed over the first projection, is the sheaf of order-\(k\) jets in the second variable. Locally its elements are finite Taylor polynomials in \(y-x\).}
\GSfield{Reconstruction}{The choice of projection and the germ along the diagonal are supplied; the two projections yield two module structures.}
\GSfield{Fate}{This is the principal-parts description of the jets used in \citet[Chapter~I, \S1]{sweeney1968neumann}.}
\GSfield{Present status}{As of November 2026: established mathematics predating this talk; the qualification in the statement is essential.}
\GSfield{Assessment}{literal statement: correct with missing hypotheses; publication match: exact; verification: read the result; development: annotated; confidence: high}
\end{GSitem}
\begin{GSitem}{1968-156.3}{E}{Many-point jets and the first-order splitting}
\GSfield{Kind; source}{construction; \GSpage{P000254-P000255}; I. M. Gelfand}
\GSfield{The notes}{Gelfand replaces \(X\times X\) by \(X^{m+1}\), uses the small diagonal, and says \((1,m)\)-jet equals a scalar function plus a differential form.}
\GSfield{Statement}{With the first factor distinguished, \(C^\infty_{X^{m+1}}/\mathcal I_\Delta^{k+1}\) records total normal Taylor degree at most \(k\). Its first-order part is \(C^\infty_X\oplus (T^*X)^{\oplus m}\) as a module: one scalar and \(m\) covectors. For \(m=1,k=1\) this is a function plus one one-form. The displayed \((1,m)\) wording does not describe a general higher jet.}
\GSfield{Reconstruction}{The normal bundle of the small diagonal and its chosen first-factor splitting give the precise interpretation; no antisymmetrization is specified in the notes.}
\GSfield{Fate}{The two-point case is standard jet theory \citep[Chapter~I, \S1]{sweeney1968neumann}; the many-point notation is a notebook reformulation.}
\GSfield{Present status}{As of November 2026: established mathematics predating this talk; the qualification in the statement is essential.}
\GSfield{Assessment}{literal statement: false as written; publication match: partial; verification: read the result; development: annotated; confidence: high}
\end{GSitem}
\begin{GSitem}{1968-156.4}{E}{Two Taylor descriptions of a first jet}
\GSfield{Kind; source}{claim; \GSpage{P000254}; I. M. Gelfand}
\GSfield{The notes}{\(\varphi(x,y)=\varphi(x,x)+\partial_x\varphi|_{y=x}(y-x)\), followed by the analogous expression with \(x,y\) reversed.}
\GSfield{Statement}{Modulo \(\mathcal I_\Delta^2\), the correct formula with \(x\) fixed is \(\varphi(x,y)=\varphi(x,x)+\partial_y\varphi(x,x)(y-x)\); with \(y\) fixed it is \(\varphi(x,y)=\varphi(y,y)+\partial_x\varphi(y,y)(x-y)\). Both are identities of first-order jets, not exact equalities of arbitrary smooth functions.}
\GSfield{Reconstruction}{Equality is interpreted modulo the squared diagonal ideal, and the differentiated variables are corrected.}
\GSfield{Fate}{Ordinary Taylor expansion; the scan confirms the derivative variables were interchanged in the notebook.}
\GSfield{Present status}{As of November 2026: established mathematics predating this talk; the qualification in the statement is essential.}
\GSfield{Assessment}{literal statement: false as written; publication match: partial; verification: read the result; development: annotated; confidence: high}
\end{GSitem}
\begin{GSitem}{1968-156.5}{E}{The trivial Spencer resolution}
\GSfield{Kind; source}{theorem; \GSpage{P000255}}
\GSfield{The notes}{\(0\to O_X\xrightarrow{j}J^k\xrightarrow{d_x}T^*\otimes J^{k-1}\xrightarrow{d_x}\Lambda^2T^*\otimes J^{k-2}\to\cdots\), called the trivial Spencer resolution.}
\GSfield{Statement}{For the trivial scalar equation, the Spencer operator is \((D\sigma)_\alpha=d\sigma_\alpha-\sum_i\sigma_{\alpha+e_i}\,dx_i\), \(|\alpha|<k\). It extends to forms, satisfies \(D^2=0\), and the sheaf sequence \(0\to C^\infty_X\xrightarrow{j^k}\mathcal J^k\xrightarrow D\Lambda^1T^*X\otimes\mathcal J^{k-1}\xrightarrow D\cdots\xrightarrow D\Lambda^nT^*X\otimes\mathcal J^{k-n}\to0\) is exact, interpreting negative-order jet terms as zero.}
\GSfield{Reconstruction}{The Spencer differential is distinguished from componentwise exterior differentiation, and the negative-order terms are set to zero as in the cited proposition.}
\GSfield{Fate}{Exactly \citet[Chapter~I, Proposition~6.2, p.~238]{sweeney1968neumann}, published before the talk. The quotient construction for the second Spencer sequence is a separate construction and is unnecessary in this displayed first sequence.}
\GSfield{Present status}{As of November 2026: established mathematics predating this talk; the qualification in the statement is essential.}
\GSfield{Assessment}{literal statement: correct with missing hypotheses; publication match: exact; verification: read the result; development: annotated; confidence: high}
\end{GSitem}
\begin{GSitem}{1968-156.6}{E}{Extension from trivial jets to linear PDE systems}
\GSfield{Kind; source}{method; \GSpage{P000255}}
\GSfield{The notes}{Replace \(O_X\) by the sheaf of solutions of an arbitrary linear system. Sweeney, Malgrange, Quillen are named.}
\GSfield{Statement}{For a linear differential operator one forms prolonged equation bundles and their Spencer complexes. Formal integrability, regularity of the prolongations, and appropriate local exactness hypotheses are needed before the complex resolves the solution sheaf. An arbitrary linear system does not automatically inherit the exact trivial complex.}
\GSfield{Reconstruction}{This gives a precise conditional method, not an assertion of local solvability for every smooth system.}
\GSfield{Fate}{\citet[Chapter~I, \S\S7--8]{sweeney1968neumann} separates stability and exactness; the other two notebook references remain unspecified.}
\GSfield{Present status}{As of November 2026: established mathematics predating this talk; the qualification in the statement is essential.}
\GSfield{Assessment}{literal statement: correct with missing hypotheses; publication match: partial; verification: read the result; development: annotated; confidence: high}
\end{GSitem}

\GStalk{157}{1968-10-27}{A. A. Kirillov}{P000255-P000257}{Traces of integral operators with operator-valued kernels}
On P000256 Kirillov states a trace formula for a weakly jointly continuous operator-valued kernel and singles out positivity for the converse. This is an exact match to Lemma~8 of \citet[pp.~49--50 of the Russian original]{kirillov1968characters}. That paper had already appeared in issue~2 of 1968, so this occurrence is expository (E), even though its 1967 announcement may have class~A. The catalogue date is retained; 27 October 1968 was Sunday, and the notebook date is not silently shifted to Monday.

\begin{GSitem}{1968-157.1}{E}{Operator-valued integral kernels}
\GSfield{Kind; source}{construction; \GSpage{P000256}}
\GSfield{The notes}{\(\mathcal Kf(x)=\int_MK(x,y)f(y)\,d\mu(y)\) on \(L^2(M,H)\), with a Hilbert space or Hermitian bundle and weak joint continuity.}
\GSfield{Statement}{Let \(M\) be a compact smooth manifold with a finite smooth positive measure and \(H\) a separable Hilbert space. A jointly weakly continuous kernel \(K:M^2\to\mathcal B(H)\) is uniformly bounded and determines a bounded integral operator on \(L^2(M,H)\). The same construction works in finite-rank Hermitian bundles using fibre maps.}
\GSfield{Reconstruction}{A finite smooth measure, separability, and the distinction between weak continuity and trace-norm continuity are supplied.}
\GSfield{Fate}{The setup is exactly that of \citet[Lemma~8]{kirillov1968characters}.}
\GSfield{Present status}{As of November 2026: established mathematics predating this talk; the qualification in the statement is essential.}
\GSfield{Assessment}{literal statement: correct with missing hypotheses; publication match: exact; verification: read the result; development: annotated; confidence: high}
\end{GSitem}
\begin{GSitem}{1968-157.2}{E}{Diagonal trace formula and positive converse}
\GSfield{Kind; source}{theorem; \GSpage{P000256}}
\GSfield{The notes}{\(\operatorname{tr}\mathcal K=\int_M\operatorname{tr}K(x,x)d\mu(x)\). Existence of the right-hand side is necessary, and for positive operators sufficient.}
\GSfield{Statement}{Under the preceding kernel hypotheses, if \(\mathcal K\) is trace class then \(K(x,x)\) is trace class almost everywhere, the diagonal trace is integrable, and the displayed equality holds. If \(\mathcal K\geq0\), almost-everywhere existence of the diagonal trace together with its finite integral also implies that \(\mathcal K\) is trace class.}
\GSfield{Fate}{Already proved in \citet[Lemma~8, Russian pp.~49--50]{kirillov1968characters}.}
\GSfield{Present status}{As of November 2026: established mathematics predating this talk; the qualification in the statement is essential.}
\GSfield{Assessment}{literal statement: correct with missing hypotheses; publication match: exact; verification: read the result; development: annotated; confidence: high}
\end{GSitem}
\begin{GSitem}{1968-157.3}{E}{Why the almost-everywhere diagonal statement matters}
\GSfield{Kind; source}{remark; \GSpage{P000256}}
\GSfield{The notes}{The nontrivial point is that trace class of \(\mathcal K\) implies existence of \(\operatorname{tr}K(x,x)\) almost everywhere; scalar kernels are referred to Gohberg--Krein and Carleman.}
\GSfield{Statement}{The trace-class condition concerns an operator on \(L^2(M,H)\), not pointwise operators on \(H\). Weak continuity and the integral representation bridge those two levels. Without a regular representative of a merely measurable kernel, its values on the diagonal can be changed arbitrarily on a product-measure null set, so a naive diagonal formula has no invariant meaning.}
\GSfield{Reconstruction}{The measurable-kernel caution explains the role of the recorded continuity hypothesis; it is not a new open problem.}
\GSfield{Fate}{Kirillov\textquoteright s Lemma~8 proves this distinction rather than assuming pointwise trace class \citep{kirillov1968characters}.}
\GSfield{Present status}{As of November 2026: established mathematics predating this talk; the qualification in the statement is essential.}
\GSfield{Assessment}{literal statement: correct with missing hypotheses; publication match: exact; verification: read the result; development: annotated; confidence: high}
\end{GSitem}

\GStalk{158}{1968-10-27}{I. M. Gelfand}{P000255-P000257}{Schedule notices; no separately identifiable intervention}
The catalogue supplies a separate Gelfand record, but P000256 names Kirillov and Weinstein and contains scheduling notes (Kirillov on traces, a skipped meeting, and Kazhdan on Dwork). No independent mathematical passage is explicitly attributed to Gelfand in the surviving page. This record is retained as a plan with uncertain attribution and receives no invented occurrence item.


\GStalk{159}{1968-10-27}{Alan Weinstein [catalogue attribution; probably Alexander Weinstein]}{P000255-P000258}{Intermediate eigenvalue problems and a delay-equation notice}
The actual passage on P000256--P000257 concerns Rayleigh quotients, finite-dimensional constraints, and lower spectral bounds through intermediate problems. These are Alexander Weinstein\textquoteright s subject, with a directly matching 1963 paper, whereas the catalogue says Alan Weinstein; the scan says only Prof.~Weinstein (USA). The name is therefore not silently emended. A separate brief notice on P000257 attributes periodic delay equations and control applications to a Polish speaker, written Rolevich; its precise attribution and theorem are unresolved.

\begin{GSitem}{1968-159.1}{E}{Rayleigh quotients and constrained extrema}
\GSfield{Kind; source}{method; \GSpage{P000256}}
\GSfield{The notes}{\(R(u)=(Au,u)/(u,u)\), followed by an Euler equation and the constraint \(u\perp\mathcal P_n\).}
\GSfield{Statement}{For a densely defined semibounded self-adjoint operator \(A\) with compact resolvent, the Rayleigh quotient is minimized on the unit sphere of its form domain. The unconstrained Euler equation is \(Au=\lambda u\); with constraints \(u\perp p_j\), it becomes \((A-\lambda)u=\sum_jc_jp_j\). The min--max principle characterizes the discrete eigenvalues.}
\GSfield{Reconstruction}{Semiboundedness, form domains, compact resolvent, and Lagrange multipliers make the variational statement precise.}
\GSfield{Fate}{The Rayleigh and constrained-minimum formulation occurs in \citet[\S\S1--2; Russian translation pp.~91--93]{weinstein1963intermediate}, predating the talk.}
\GSfield{Present status}{As of November 2026: established mathematics predating this talk; the qualification in the statement is essential.}
\GSfield{Assessment}{literal statement: correct with missing hypotheses; publication match: partial; verification: read the result; development: annotated; confidence: high}
\end{GSitem}
\begin{GSitem}{1968-159.2}{E}{Finite constraint determinants and intermediate lower bounds}
\GSfield{Kind; source}{method; \GSpage{P000256-P000257}}
\GSfield{The notes}{\(\mathcal P=\bigcup_n\mathcal P_n\); minimizing on \(u\perp\mathcal P_n\) gives a determinant equation. Classical bounds come from one side, Weinstein\textquoteright s from the other.}
\GSfield{Statement}{For \(\lambda\notin\sigma(A)\), constrained eigenvalues may be detected by the finite matrix \(G_{ij}(\lambda)=\langle(A-\lambda)^{-1}p_j,p_i\rangle\), through \(\det G(\lambda)=0\). Persistent eigenvalues in \(\sigma(A)\) need separate treatment. In intermediate-problem methods one chooses ordered, explicitly tractable quadratic forms converging to the target; their min--max eigenvalues give certified bounds, with convergence requiring density in the relevant form norm.}
\GSfield{Reconstruction}{The resolvent determinant is an explicit reconstruction of the ellipsis, valid only off the original spectrum; it is not a recovered formula.}
\GSfield{Fate}{\citet[\S\S2--4]{weinstein1963intermediate} treats the generalized intermediate problems and their meromorphic spectral functions. The specific finite-matrix display fills the notebook ellipsis and is labeled reconstruction.}
\GSfield{Present status}{As of November 2026: established mathematics predating this talk; the qualification in the statement is essential.}
\GSfield{Assessment}{literal statement: correct with missing hypotheses; publication match: partial; verification: read the result; development: annotated; confidence: high}
\end{GSitem}
\begin{GSitem}{1968-159.3}{E}{Extension to Schrödinger operators}
\GSfield{Kind; source}{claim; \GSpage{P000257}}
\GSfield{The notes}{The constrained minimum exists for the Schrödinger operator, not only for completely continuous operators.}
\GSfield{Statement}{For a semibounded self-adjoint Schrödinger operator, min--max values below the essential spectrum are discrete eigenvalues and are attained when they lie strictly below its bottom. Compact resolvent is one sufficient global condition. An arbitrary Schrödinger operator on a noncompact domain need not attain the infimum of its Rayleigh quotient.}
\GSfield{Reconstruction}{Spectral isolation or compactness is essential; the notes do not specify a potential or domain.}
\GSfield{Fate}{The operator formulation belongs to the classical intermediate-problem framework \citep{weinstein1963intermediate}.}
\GSfield{Present status}{As of November 2026: established mathematics predating this talk; the qualification in the statement is essential.}
\GSfield{Assessment}{literal statement: correct with missing hypotheses; publication match: partial; verification: read the result; development: annotated; confidence: high}
\end{GSitem}
\begin{GSitem}{1968-159.4}{F}{The Rolewicz periodic delay-equation notice}
\GSfield{Kind; source}{remark; \GSpage{P000257}}
\GSfield{The notes}{An unnamed Pole, written Rolevich, discusses periodic solutions of delay equations, reduction to ordinary equations with periodic coefficients, and control applications.}
\GSfield{Statement}{No delay length, equation, phase space, regularity, periodicity relation, or reduction map is specified. The strongest faithful statement is a programme notice about such a reduction; a general equivalence between delay equations and finite-dimensional ODEs would be false.}
\GSfield{Reconstruction}{Rolevich plausibly refers to Stefan Rolewicz; the matching literature also includes Danuta Przeworska-Rolewicz. No first name is supplied by the source.}
\GSfield{Fate}{Two close primary leads are \citet{przeworskarolewicz1968periodic} and \citet{przeworskarolewiczrolewicz1969control}. Publisher issue pages verify their titles, dates, authors, and pages; the PDF downloads were inaccessible. They support the programme identification but do not establish the exact reduction or which author spoke.}
\GSfield{Present status}{As of November 2026: the exact assertion or attribution remains unresolved; no unsolved mathematical problem is inferred.}
\GSfield{Assessment}{literal statement: not determinable; publication match: thematic only; verification: publisher metadata checked; development: annotated; confidence: medium}
\end{GSitem}

\GStalk{160}{1968-11-04}{A. A. Kirillov}{P000256-P000262}{Characters, coadjoint orbits, and Plancherel measure}
The actual talk occupies P000257--P000260. It moves from distribution characters and the affine group to symplectic coadjoint orbits, orbital Fourier transforms, integrality, the compact rank-one example, and Plancherel measure. A Pfaffian exercise and brief remarks on matrix products, tensor products, and Rodrigues formulas close the page. The unrestricted all-group wording is separated from the nilpotent and compact theorems actually established in \citet{kirillov1962unitary,kirillov1968characters}.

\begin{GSitem}{1968-160.1}{E}{Distribution characters and their test-function domain}
\GSfield{Kind; source}{definition; \GSpage{P000257}}
\GSfield{The notes}{\(T(\varphi)=\int_G\varphi(g)T(g)\,dg\); if its trace exists, \(\operatorname{tr}T(\varphi)=\langle\chi,\varphi\rangle\). The examples appear no worse than measures, but no general theorem is asserted.}
\GSfield{Statement}{For a strongly continuous unitary representation \(\pi\) of a Lie group with a fixed left Haar measure, \(\pi(\varphi)\) is bounded for \(\varphi\in C_c^\infty(G)\). A distribution character on this test space is defined when these operators are trace class and \(\varphi\mapsto\operatorname{Tr}\pi(\varphi)\) is distribution-continuous. These conditions are additional assumptions; the trace is not defined for an arbitrary irreducible representation of an arbitrary Lie group.}
\GSfield{Reconstruction}{The continuity required for a distribution, the Haar choice, and the distinction between a bounded and trace-class operator are made explicit.}
\GSfield{Fate}{The trace formulation and its need for hypotheses are the starting point of \citet[Introduction and \S2]{kirillov1968characters}.}
\GSfield{Present status}{As of November 2026: established mathematics predating this talk; the qualification in the statement is essential.}
\GSfield{Assessment}{literal statement: correct with missing hypotheses; publication match: partial; verification: read the result; development: annotated; confidence: high}
\end{GSitem}
\begin{GSitem}{1968-160.2}{E}{The affine group and a restricted test-function ideal}
\GSfield{Kind; source}{remark; \GSpage{P000257}}
\GSfield{The notes}{For \(\left(\begin{smallmatrix}a&b\\0&1\end{smallmatrix}\right)\), the test functions must satisfy \(\int\varphi(a,b)\,db=0\) for every \(a\); the preceding sentence says examples have no trace for any \(\varphi\).}
\GSfield{Statement}{For the connected affine group \(a>0\), its infinite-dimensional irreducibles are realized on the positive or negative half-line in the dual of the translation subgroup. The condition that the Fourier transform in \(b\) vanish at frequency zero removes the endpoint obstruction to a trace. The literal claim that no test function has a trace is too strong: the zero test function always does, and restricted ideals can contain nonzero examples.}
\GSfield{Reconstruction}{The half-line realization explains the recorded cancellation condition; the notebook does not state a complete necessary-and-sufficient trace-class criterion.}
\GSfield{Fate}{The relevant reduction through a commutative normal subgroup is discussed in \citet[\S2]{kirillov1968characters}; this does not certify an unrestricted character on all test functions.}
\GSfield{Present status}{As of November 2026: established mathematics predating this talk; the qualification in the statement is essential.}
\GSfield{Assessment}{literal statement: false as written; publication match: partial; verification: read the result; development: annotated; confidence: high}
\end{GSitem}
\begin{GSitem}{1968-160.3}{E}{The symplectic form on a coadjoint orbit}
\GSfield{Kind; source}{construction; \GSpage{P000258}}
\GSfield{The notes}{\(0\to\mathcal H\to\mathfrak g\to T_F\Omega\to0\); \(B_F(\xi,\eta)=(2\pi)^{-1}\langle F,[X,Y]\rangle\), independent of lifts.}
\GSfield{Statement}{For a coadjoint orbit \(\Omega\subset\mathfrak g^*\), the stabilizer Lie algebra is \(\mathfrak g_F=\{X:\langle F,[X,Y]\rangle=0\text{ for all }Y\}\). The displayed bilinear form descends to \(\mathfrak g/\mathfrak g_F\), is nondegenerate, is invariant under the coadjoint action, and is closed by the Jacobi identity. A consistent sign convention for fundamental vector fields is understood.}
\GSfield{Fate}{\citet[\S1, Lemmas~1--3]{kirillov1968characters} proves these statements; the normalization agrees with the notebook.}
\GSfield{Present status}{As of November 2026: established mathematics predating this talk; the qualification in the statement is essential.}
\GSfield{Assessment}{literal statement: correct; publication match: exact; verification: read the result; development: annotated; confidence: high}
\end{GSitem}
\begin{GSitem}{1968-160.4}{E}{Liouville measure and the orbital Fourier distribution}
\GSfield{Kind; source}{construction; \GSpage{P000258}}
\GSfield{The notes}{\(\beta=e^B=1+B+\cdots+B^k/k!\), with only its top term integrated; \(\mathcal J_\Omega(X)=\int_\Omega e^{i\langle F,X\rangle}\beta(F)\).}
\GSfield{Statement}{On a \(2k\)-dimensional symplectic coadjoint orbit, \(B^k/k!\) is an invariant positive density after the symplectic orientation is chosen. Its Fourier transform is interpreted distributionally whenever integration against the relevant Fourier transforms converges; a noncompact orbital integral need not be an absolutely convergent function of \(X\).}
\GSfield{Reconstruction}{The top-degree component, orientation, and distributional meaning are specified; a growth condition must precede a tempered-distribution claim.}
\GSfield{Fate}{This is the measure used in the orbit formulas of \citet[\S1]{kirillov1968characters}.}
\GSfield{Present status}{As of November 2026: established mathematics predating this talk; the qualification in the statement is essential.}
\GSfield{Assessment}{literal statement: correct with missing hypotheses; publication match: exact; verification: read the result; development: annotated; confidence: high}
\end{GSitem}
\begin{GSitem}{1968-160.5}{F}{The unqualified universal orbital character formula}
\GSfield{Kind; source}{claim; \GSpage{P000258-P000259}}
\GSfield{The notes}{The assertion \(\chi(\exp X)=\mathcal J_\Omega(X)\) is marked false; it is replaced, for every irreducible representation, by \(\chi(\exp X)=p_\Omega(X)^{-1}\mathcal J_\Omega(X)\), \(p_\Omega(0)=1\), smooth.}
\GSfield{Statement}{A precise local character formula must specify a class of groups and representations, the orbit or cycle attached to each representation, a test-function domain, the exponential-chart Jacobian, and whether the factor is smooth and nonvanishing on that chart. Those data are not specified for the all-group assertion. The nilpotent and compact instances below are precise; they do not imply this universal statement.}
\GSfield{Reconstruction}{No globally defined quotient of distributions is asserted: such a quotient is not an operation in distribution theory.}
\GSfield{Fate}{\citet[Introduction, \S\S2--3]{kirillov1968characters} proves selected cases and reductions, not this all-group assertion. The already noted affine trace obstruction prevents treating the displayed expression as a distribution identity on an unstated common test space. This is a recurrence of the broad announcement indexed at 1967-131.1, without a precise historical problem identifier.}
\GSfield{Present status}{As of November 2026: the exact assertion or attribution remains unresolved; no unsolved mathematical problem is inferred.}
\GSfield{Assessment}{literal statement: not determinable; publication match: partial; verification: read the result; development: annotated; confidence: medium}
\end{GSitem}
\begin{GSitem}{1968-160.6}{F}{Haar-density factors in generic character formulas}
\GSfield{Kind; source}{claim; \GSpage{P000259}}
\GSfield{The notes}{For generic orbits the factor is \(\sqrt{\rho(X)}\) if an invariant measure exists, and a fourth root of a product of left and right densities otherwise.}
\GSfield{Statement}{In exponential coordinates the left and right Haar densities of a Lie group may be written, up to orientation and normalization, as \(j_L(X)=\det((1-e^{-\operatorname{ad}X})/\operatorname{ad}X)\) and \(j_R(X)=\det((e^{\operatorname{ad}X}-1)/\operatorname{ad}X)\). Their geometric mean removes the modular exponential factor. This identifies a plausible local density normalization but does not establish that its square root is the correction factor for every generic irreducible character.}
\GSfield{Reconstruction}{The Haar Jacobians are supplied as a local calculation. The notebook does not define generic, specify the representation class, or distinguish Haar measure on the group from an invariant measure in exponential coordinates.}
\GSfield{Fate}{The exponential-unimodular calculation and compact formula in \citet[\S\S2--3]{kirillov1968characters} support restricted cases; no all-group generic-orbit theorem is certified. The left/right interpretation of the handwritten subscripts is itself uncertain.}
\GSfield{Present status}{As of November 2026: the exact assertion or attribution remains unresolved; no unsolved mathematical problem is inferred.}
\GSfield{Assessment}{literal statement: not determinable; publication match: partial; verification: read the result; development: annotated; confidence: medium}
\end{GSitem}
\begin{GSitem}{1968-160.7}{E}{The simply connected nilpotent orbit formula}
\GSfield{Kind; source}{theorem; \GSpage{P000259}}
\GSfield{The notes}{\(p_\Omega=1\) for nilpotent groups; orbits do correspond to representations in that case.}
\GSfield{Statement}{For a connected simply connected nilpotent real Lie group, equivalence classes of irreducible unitary representations are parametrized by coadjoint orbits. With compatible Haar, Fourier, and symplectic normalizations, the distribution character pulled back by the global exponential diffeomorphism is the Fourier transform of the orbit Liouville measure; the exponential Jacobian is one.}
\GSfield{Fate}{This is pre-existing: \citet[Theorems~5.1--5.2 and~7.4]{kirillov1962unitary}, received 17 October 1961 and published in 1962.}
\GSfield{Present status}{As of November 2026: established mathematics predating this talk; the qualification in the statement is essential.}
\GSfield{Assessment}{literal statement: correct with missing hypotheses; publication match: exact; verification: read the result; development: annotated; confidence: high}
\end{GSitem}
\begin{GSitem}{1968-160.8}{E}{Compact characters require the rho-shifted regular orbit}
\GSfield{Kind; source}{theorem; \GSpage{P000259}}
\GSfield{The notes}{For an arbitrary compact group all orbits are said to be in general position; \(p_\Omega=\sqrt\rho\) is suggested.}
\GSfield{Statement}{For a compact connected semisimple group and an irreducible representation of dominant integral highest weight \(\lambda\), the character formula uses the regular coadjoint orbit through \(\lambda+\rho\), where \(\rho\) is half the sum of positive roots. Near zero, the correction is the analytic square root \(j(X)^{1/2}\) of \(j(X)=\det_{\mathfrak g}((1-e^{-\operatorname{ad}X})/\operatorname{ad}X)\). Singular coadjoint orbits certainly exist; the zero orbit is one.}
\GSfield{Reconstruction}{The word all is restricted to the regular shifted orbits attached by this character formula, not to every orbit in the dual.}
\GSfield{Fate}{The precise compact result is \citet[\S3, Theorem~2]{kirillov1968characters}; its coordinates \(f_k\) and highest weights \(f_k-1\) exhibit the shift.}
\GSfield{Present status}{As of November 2026: established mathematics predating this talk; the qualification in the statement is essential.}
\GSfield{Assessment}{literal statement: false as written; publication match: exact; verification: read the result; development: annotated; confidence: high}
\end{GSitem}
\begin{GSitem}{1968-160.9}{F}{The proposed product of modified hyperbolic sines}
\GSfield{Kind; source}{conjecture; \GSpage{P000259}}
\GSfield{The notes}{\(\operatorname{sth}x=\sinh(x/2)/(x/2)\); \(p_\Omega(X)=\prod_\alpha\operatorname{sth}\langle\alpha,X\rangle\) is marked conjecture, with \(\alpha\) called roots.}
\GSfield{Statement}{For a compact connected semisimple group, on a Cartan subalgebra and with complexified roots, \(j(X)^{1/2}=\prod_{\alpha>0}\sinh(\alpha(X)/2)/(\alpha(X)/2)\), normalized to one at zero. The product is over positive roots; taking both signs squares it. Outside that setting the notebook does not identify a root multiset depending on \(\Omega\), or the intended generality of the conjecture.}
\GSfield{Reconstruction}{The precise positive-root formula is supplied as a verified special case; it is not asserted to recover the full conjecture.}
\GSfield{Fate}{The compact factor is already present in \citet[\S3]{kirillov1968characters}. No exact general root-product theorem was identified from this notation; the general extension is therefore not entered as a newly solved or open problem.}
\GSfield{Present status}{As of November 2026: the exact assertion or attribution remains unresolved; no unsolved mathematical problem is inferred.}
\GSfield{Assessment}{literal statement: not determinable; publication match: partial; verification: read the result; development: annotated; confidence: medium}
\end{GSitem}
\begin{GSitem}{1968-160.10}{F}{Integrality and characters of the orbit fundamental group}
\GSfield{Kind; source}{claim; \GSpage{P000259}}
\GSfield{The notes}{Choose \(\Omega\) with \([B]\in H^2(\Omega,\mathbb Z)\); it corresponds to as many representations as characters of \(\pi_1(\Omega)\), which is said to be commutative.}
\GSfield{Statement}{For a connected symplectic manifold, an integral class \([B]\) is the condition for a Hermitian line bundle with a unitary connection of curvature \(2\pi iB\). When one such connection exists, isomorphism classes with this curvature form a torsor for flat unitary line bundles, classified by \(\operatorname{Hom}(\pi_1(\Omega),U(1))\). This geometric statement neither requires \(\pi_1\) itself to be abelian nor provides, by itself, a bijection with the unitary dual of an arbitrary Lie group.}
\GSfield{Reconstruction}{The line-bundle classification is a modern geometric interpretation of the integral-form and fundamental-group data, explicitly separated from the stronger representation claim.}
\GSfield{Fate}{The compact and nilpotent orbit constructions in \citet{kirillov1962unitary,kirillov1968characters} do not establish the proposed universal representation count. Equivariance, polarization, and representation-theoretic admissibility cannot be omitted.}
\GSfield{Present status}{As of November 2026: the exact assertion or attribution remains unresolved; no unsolved mathematical problem is inferred.}
\GSfield{Assessment}{literal statement: not determinable; publication match: partial; verification: read the result; development: annotated; confidence: medium}
\end{GSitem}
\begin{GSitem}{1968-160.11}{E}{The SU(2) sphere calculation and normalization}
\GSfield{Kind; source}{theorem; \GSpage{P000260}}
\GSfield{The notes}{On a sphere of radius \(r\), \(B=(2\pi r)^{-1}d\sigma\); integrality gives \(2r=n\). The orbital transform is written \(\sin(rx)/x\).}
\GSfield{Statement}{With the stated normalization, \(\int_{S_r^2}B=2r\), so prequantization requires \(2r\in\mathbb Z\). For \(|X|=x\), direct integration gives \(\int_{S_r^2}e^{i\langle F,X\rangle}B(F)=2\sin(rx)/x\), with value \(2r\) at zero. Choosing \(r=(n+1)/2\) for the highest-weight-\(n\) representation and dividing by \(\sin(x/2)/(x/2)\) gives \(\sin((n+1)x/2)/\sin(x/2)\), of dimension \(n+1\).}
\GSfield{Reconstruction}{The Euclidean normalization of the Lie-algebra coordinate is fixed, and the notebook integral is completed by direct spherical integration.}
\GSfield{Fate}{This is the rank-one specialization of \citet[\S3, Theorem~2]{kirillov1968characters}. The factor two and the shift are required by evaluation at the identity.}
\GSfield{Present status}{As of November 2026: established mathematics predating this talk; the qualification in the statement is essential.}
\GSfield{Assessment}{literal statement: false as written; publication match: exact; verification: read the result; development: annotated; confidence: high}
\end{GSitem}
\begin{GSitem}{1968-160.12}{E}{The scope of the displayed Plancherel formula}
\GSfield{Kind; source}{claim; \GSpage{P000260}}
\GSfield{The notes}{\(\|f\|_2^2=\int\operatorname{tr}(T(f)T(f)^*)\,d\mu(T)\) is called universal for all groups; the question is to calculate \(d\mu\).}
\GSfield{Statement}{For a second-countable unimodular type-I locally compact group, the Plancherel transform admits this Hilbert--Schmidt norm formula over the unitary dual, initially for a suitable dense subspace of \(L^1\cap L^2\). Unimodularity and a suitable measurable decomposition are essential to this form. Nonunimodular groups require operator-valued formal-degree corrections, and arbitrary non-type-I groups cannot be handled by the same ordinary parametrization of irreducibles.}
\GSfield{Reconstruction}{The modern qualified theorem is stated; the measure-calculation programme lacks a particular new group or measure and therefore creates no separate open problem.}
\GSfield{Fate}{For the nilpotent case this predates the talk \citep[\S7]{kirillov1962unitary}. The notebook\textquoteright s all-group formula is not established by that example or by the exponential-unimodular calculation in \citet[\S2]{kirillov1968characters}.}
\GSfield{Present status}{As of November 2026: established mathematics predating this talk; the qualification in the statement is essential.}
\GSfield{Assessment}{literal statement: false as written; publication match: partial; verification: read the result; development: annotated; confidence: high}
\end{GSitem}
\begin{GSitem}{1968-160.13}{E}{Extracting the polynomial square root of a skew determinant}
\GSfield{Kind; source}{problem; \GSpage{P000260}}
\GSfield{The notes}{Problem: for a skew-symmetric matrix, extract a polynomial square root of its determinant; the two-by-two example is given as \(-a^2\).}
\GSfield{Statement}{For a skew-symmetric \(2m\times2m\) matrix over a field of characteristic different from two, \(\det A=\operatorname{Pf}(A)^2\). The Pfaffian is the signed sum \[\operatorname{Pf}(A)=\sum_\pi\operatorname{sgn}(\pi)\prod_{j=1}^m a_{\pi(2j-1),\pi(2j)}.\] Here the sum runs over permutations satisfying \(\pi(2j-1)<\pi(2j)\) and \(\pi(1)<\pi(3)<\cdots<\pi(2m-1)\); these index unordered pairings without division. In characteristic zero this also equals the sum over all permutations divided by \(2^m m!\), and its exterior-form description is \(\omega^m=m!\operatorname{Pf}(A)\,dx_1\wedge\cdots\wedge dx_{2m}\). Odd skew determinants vanish.}
\GSfield{Reconstruction}{The characteristic restriction, even size, and pairing-polynomial formula are supplied. The factorial quotient and exterior-power formula are restricted to characteristic zero. This elementary reconstruction receives no historical problem identifier.}
\GSfield{Fate}{This is the classical Pfaffian identity underlying the top-degree form on P000258; the notebook poses an exercise, not an unsolved problem. The displayed example has determinant \(+a^2\).}
\GSfield{Present status}{As of November 2026: established mathematics predating this talk; the qualification in the statement is essential.}
\GSfield{Assessment}{literal statement: false as written; publication match: exact; verification: read the result; development: annotated; confidence: high}
\end{GSitem}
\begin{GSitem}{1968-160.14}{F}{Matrix spectra and tensor products: an unspecified duality}
\GSfield{Kind; source}{remark; \GSpage{P000260}}
\GSfield{The notes}{A duality is drawn between products of positive matrices (with unitary inserted) and tensor products of representations of the unitary group; Perelomov and Popov are linked to a convex region.}
\GSfield{Statement}{The spectra of two matrices alone do not determine the spectrum of their product without information on their relative eigenspaces. A probability distribution would additionally require a law for those eigenspaces. Tensor-product decomposition requires specified representations and multiplicity data. The notes do not distinguish an attainable-spectra problem, a random-matrix law, or a semiclassical limit, and do not define the convex region.}
\GSfield{Reconstruction}{No identification with a particular Horn problem or asymptotic tensor-product theorem is made.}
\GSfield{Fate}{\citet[Abstract]{perelomovpopov1968casimir} is a verified contemporary paper by the named authors, but it concerns Casimir eigenvalues and is only a thematic lead, not an exact publication match to this duality.}
\GSfield{Present status}{As of November 2026: the exact assertion or attribution remains unresolved; no unsolved mathematical problem is inferred.}
\GSfield{Assessment}{literal statement: not determinable; publication match: thematic only; verification: publisher metadata checked; development: annotated; confidence: low}
\end{GSitem}
\begin{GSitem}{1968-160.15}{E}{Rodrigues formulas for Legendre and Jacobi polynomials}
\GSfield{Kind; source}{claim; \GSpage{P000260}}
\GSfield{The notes}{\([(x^2-1)^n]^{(n)}\) is Legendre; \([(x-1)^{m_1}(x+1)^{m_2}]^{(m_3)}\) is Jacobi.}
\GSfield{Statement}{For \(n\ge0\), \(P_n(x)=(2^n n!)^{-1}(d/dx)^n(x^2-1)^n\). More generally, \(P_n^{(\alpha,\beta)}(x)=(-1)^n[2^n n!(1-x)^\alpha(1+x)^\beta]^{-1}(d/dx)^n[(1-x)^{n+\alpha}(1+x)^{n+\beta}]\); for the orthogonality problem one assumes \(\alpha,\beta>-1\). The raw derivative in the notebook is a weighted polynomial, not in general the normalized Jacobi polynomial itself.}
\GSfield{Reconstruction}{Normalization and the Jacobi weight divisor are restored.}
\GSfield{Fate}{Classical formulas, checked in the NIST \citet[\S18.5(ii), Table~18.5.1]{nist2026dlmf}; their origin predates the talk.}
\GSfield{Present status}{As of November 2026: established mathematics predating this talk; the qualification in the statement is essential.}
\GSfield{Assessment}{literal statement: correct with missing hypotheses; publication match: exact; verification: read the result; development: annotated; confidence: high}
\end{GSitem}
\begin{GSitem}{1968-160.16}{F}{Fractional Rodrigues operations and hypergeometric functions}
\GSfield{Kind; source}{remark; \GSpage{P000260}}
\GSfield{The notes}{When the derivative order \(m_3\) is fractional, a hypergeometric function is obtained; \(x(x-1)\) is also suggested.}
\GSfield{Statement}{A fractional derivative requires an integration endpoint, a branch, a kernel, and an admissible parameter domain. For one explicit choice, the Riemann--Liouville integral satisfies \(I_{0+}^{\mu}[x^{b-1}(1-x)^{-a}]=\Gamma(b)x^{b+\mu-1}\,{}_2F_1(a,b;b+\mu;x)/\Gamma(b+\mu)\) when \(\Re b,\Re\mu>0\) and \(0<x<1\). This exhibits the intended connection, but is one reconstruction, not a uniquely recoverable notebook formula.}
\GSfield{Reconstruction}{The displayed integral identity follows by the binomial series and beta integral. It does not claim that the notebook chose this fractional-calculus convention.}
\GSfield{Fate}{The ordinary Jacobi--hypergeometric relation is classical \citep[\S18.5(iii)]{nist2026dlmf}; the fractional operation intended here is unidentified.}
\GSfield{Present status}{As of November 2026: the exact assertion or attribution remains unresolved; no unsolved mathematical problem is inferred.}
\GSfield{Assessment}{literal statement: not determinable; publication match: partial; verification: read the result; development: annotated; confidence: medium}
\end{GSitem}

\GStalk{161}{1968-11-18}{I. M. Gelfand}{P000260-P000264}{A graded reading programme in representation theory}
The programme runs from P000261 to the upper part of P000263. Level zero covers finite and compact groups, rotations, and Lie algebras; level one proceeds through classical unitary representations, characters, induction, nilpotent groups, rank-one examples, and operator algebras. Level two adds non-Archimedean groups, symmetric-space Plancherel theory, horospheres, Borel--Weil, type-I theory, enveloping algebras, orbit methods, Kazhdan\textquoteright s work, and lattices. The closing promise of a list in Uspekhi is a publication plan, not the list itself. Individual bibliography entries without mathematical assertions are preserved here rather than inflated into occurrence items.

\begin{GSitem}{1968-161.1}{E}{Representation theory organized by prerequisites and examples}
\GSfield{Kind; source}{method; \GSpage{P000261-P000263}}
\GSfield{The notes}{Three levels are distinguished; compact-group fundamentals precede explicit noncompact examples, which precede structural harmonic analysis and enveloping algebras.}
\GSfield{Statement}{This is a pedagogical method: learn finite/compact representation theory and concrete rank-one models before distribution characters, induced representations, Plancherel theory, type-I decomposition, and orbit/enveloping-algebra methods. It contains no claim that a single listed text proves a complete classification for all groups.}
\GSfield{Fate}{The nilpotent component is concretely anchored by \citet{kirillov1962unitary}; the programme itself is source evidence about teaching and seminar priorities.}
\GSfield{Present status}{As of November 2026: established mathematics predating this talk; the qualification in the statement is essential.}
\GSfield{Assessment}{literal statement: correct; publication match: partial; verification: read the result; development: annotated; confidence: high}
\end{GSitem}
\begin{GSitem}{1968-161.2}{F}{Berezin and the phrase existence of all representations}
\GSfield{Kind; source}{claim; \GSpage{P000261}}
\GSfield{The notes}{Berezin, Transactions of the Moscow Mathematical Society volume 6 and a letter in volume 12: description through Laplace operators, and existence of all representations proved.}
\GSfield{Statement}{The class of semisimple groups, the precise category of representations, the role of infinitesimal characters, and the admissible spectral parameters are not given. Existence for all algebraically imaginable Laplace eigenvalues is not a justified reading. The remark therefore cannot be converted into a precise all-group unitary-dual theorem.}
\GSfield{Fate}{The volume-level bibliographical lead is retained, but no theorem in the two Berezin references was read that matches the unqualified wording. This is an unresolved identification, not an open classification problem.}
\GSfield{Present status}{As of November 2026: the exact assertion or attribution remains unresolved; no unsolved mathematical problem is inferred.}
\GSfield{Assessment}{literal statement: not determinable; publication match: none found; verification: unresolved lead; development: annotated; confidence: low}
\end{GSitem}
\begin{GSitem}{1968-161.3}{E}{Enveloping algebras and the continuing orbit programme}
\GSfield{Kind; source}{remark; \GSpage{P000262-P000263}}
\GSfield{The notes}{Gelfand--Kirillov, the French Publications Mathématiques paper and Functional Analysis volume 3 issue 1, are listed; Kirillov and Kostant are named for general representations and characters.}
\GSfield{Statement}{The reference links the seminar\textquoteright s earlier study of enveloping algebras and their division rings with the orbit-method talks. A bibliography citation alone does not repeat every conjecture of the cited work, nor assert that the general orbit programme is complete.}
\GSfield{Reconstruction}{Only the connection explicitly conveyed by the reading list is indexed; no earlier problem identifier is repeated solely because its literature is named.}
\GSfield{Fate}{The French reference identifies \citet{gelfandkirillov1966corps}; the 1969 volume reference may be prospective or a later annotation, and is not evidence that the article had already been published in November 1968.}
\GSfield{Present status}{As of November 2026: established mathematics predating this talk; the qualification in the statement is essential.}
\GSfield{Assessment}{literal statement: correct; publication match: partial; verification: publisher metadata checked; development: annotated; confidence: high}
\end{GSitem}

\GStalk{162}{1968-11-18}{M. L. Gromov}{P000262-P000268}{Differential relations, local flexibility, and global existence}
Gromov begins in the lower half of P000263 and continues through P000266. Jets turn PDEs into differential relations; deformation lifting replaces uniqueness, and local solvability plus suitable invariance is proposed as the route to global solutions. The precise open-manifold theorem appears in \citet[Theorem~4.1.1]{gromov1969stable}, received 8 July 1968 and published in 1969. The closing isometric-mapping statements have illegible dimensions and missing hypotheses and are not all identified with that theorem.

\begin{GSitem}{1968-162.1}{E}{Differential relations and holonomic solutions}
\GSfield{Kind; source}{definition; \GSpage{P000263-P000264}}
\GSfield{The notes}{For a bundle \(\mathcal L\to M\), let \(\mathcal L^r\) be its jet extension and \(A\subset X^r\) a differential equation. The solution line writes \(\gamma\in\Gamma(\mathcal L)\), \(\operatorname{Im}\gamma\subset A\).}
\GSfield{Statement}{For a smooth fibre bundle \(E\to M\), an order-\(r\) differential relation is a subset \(\mathcal R\subset J^rE\). A formal solution is a section of \(\mathcal R\to M\); a genuine solution is a smooth section \(s:M\to E\) such that \(j^rs(x)\in\mathcal R\) for every \(x\). Its jet section is holonomic. A homogeneous linear differential equation of constant rank defines a vector subbundle; arbitrary equations need not be subbundles.}
\GSfield{Reconstruction}{The missing jet prolongation in the solution line is restored; regularity is supplied for the subbundle claim.}
\GSfield{Fate}{The terminology is formalized in \citet[\S4.1]{gromov1969stable}; the underlying jet construction already appears in Palamodov\textquoteright s record156.}
\GSfield{Present status}{As of November 2026: established mathematics predating this talk; the qualification in the statement is essential.}
\GSfield{Assessment}{literal statement: false as written; publication match: partial; verification: read the result; development: annotated; confidence: high}
\end{GSitem}
\begin{GSitem}{1968-162.2}{E}{Open and closed differential relations}
\GSfield{Kind; source}{remark; \GSpage{P000263}}
\GSfield{The notes}{Classical equations usually define closed submanifolds; \(Lu\ne0\) gives an open set and \(Lu=0\) a closed one.}
\GSfield{Statement}{If a differential operator is represented by a continuous bundle map \(L:J^rE\to F\), then \(L^{-1}(0)\) is closed and \(L^{-1}(F\setminus0)\) is open, provided the zero section is closed. Neither preimage is automatically a submanifold. Openness allows small perturbations of jets to preserve a relation.}
\GSfield{Fate}{The open-relation condition is an explicit hypothesis of \citet[Theorem~4.1.1]{gromov1969stable}.}
\GSfield{Present status}{As of November 2026: established mathematics predating this talk; the qualification in the statement is essential.}
\GSfield{Assessment}{literal statement: correct with missing hypotheses; publication match: partial; verification: read the result; development: annotated; confidence: high}
\end{GSitem}
\begin{GSitem}{1968-162.3}{E}{The metric pullback equation and the immersion condition}
\GSfield{Kind; source}{construction; \GSpage{P000264}}
\GSfield{The notes}{\(f:(M_1,g_1)\to(M_2,g_2)\), \(f^*g_2=g_1\), as a section of \(M_1\times M_2\); add the open condition that it be an embedding.}
\GSfield{Statement}{The equation \(g_{2,f(x)}(df_xv,df_xw)=g_{1,x}(v,w)\) is a first-order relation. For positive-definite \(g_1\), every solution is an immersion. Full rank of \(df\) is an open first-jet condition; global injectivity and being a homeomorphism onto the image are not conditions on a single jet. Thus embedding cannot be imposed as the stated local open relation.}
\GSfield{Reconstruction}{Immersion and embedding are distinguished; the metric equation itself is a closed, not open, relation.}
\GSfield{Fate}{Gromov\textquoteright s published open-relation framework treats immersion rank precisely \citep[\S4.2]{gromov1969stable}; it does not identify embeddings with full-rank jets.}
\GSfield{Present status}{As of November 2026: established mathematics predating this talk; the qualification in the statement is essential.}
\GSfield{Assessment}{literal statement: false as written; publication match: partial; verification: read the result; development: annotated; confidence: high}
\end{GSitem}
\begin{GSitem}{1968-162.4}{A}{Parametric solvability as homotopy lifting}
\GSfield{Kind; source}{method; \GSpage{P000264}}
\GSfield{The notes}{For a family \(A_N\), well-posedness means that the family of solution spaces over \(N\) is a Serre fibration; finite-parameter deformations should lift.}
\GSfield{Statement}{For a topological family \(\mathcal R_n\), define \(\mathcal S=\{(s,n):j^rs(M)\subset\mathcal R_n\}\). The proposed property is that \(\mathcal S\to N\) has the homotopy-lifting property for disks. The full product \(\Gamma(E)\times N\to N\) always has that property and therefore cannot encode solvability; the restricted solution space is required. This homotopical condition does not assert analytic uniqueness or quantitative continuous dependence.}
\GSfield{Reconstruction}{The total space is restricted to actual solutions; the notes use well-posedness in a homotopy-theoretic sense.}
\GSfield{Fate}{The Serre-fibration and flexibility machinery is developed by \citet[\S\S1--3]{gromov1969stable}, circulated by its July 1968 submission and printed in 1969.}
\GSfield{Present status}{As of November 2026: the identified result is established, with the scope stated here.}
\GSfield{Assessment}{literal statement: correct with missing hypotheses; publication match: partial; verification: read the result; development: annotated; confidence: high}
\end{GSitem}
\begin{GSitem}{1968-162.5}{E}{Nash inversion as the local analytic input}
\GSfield{Kind; source}{method; \GSpage{P000264}}
\GSfield{The notes}{The local part of the metric example is attributed to Nash\textquoteright s implicit-function theorems; the global part is unknown.}
\GSfield{Statement}{For the metric-induction operator \(f\mapsto f^*g_{\mathrm{eucl}}\), local solvability near a suitably free embedding uses an inverse procedure that compensates for loss of derivatives by smoothing. A classical Banach-space implicit-function theorem cannot simply be applied when the linearized right inverse loses derivatives. A global existence principle needs further geometric hypotheses.}
\GSfield{Reconstruction}{The free-embedding and derivative-loss conditions explain the method; they are not written out in the notebook.}
\GSfield{Fate}{The named source is Nash\textquoteright s 1956 embedding paper \citep{nash1956embedding}; its identity and scope are confirmed in \citet[reference~3]{choquet1958nash}. The accessible Nash PDF was opened but not text-extractable, so no numbered local theorem is asserted here.}
\GSfield{Present status}{As of November 2026: established mathematics predating this talk; the qualification in the statement is essential.}
\GSfield{Assessment}{literal statement: correct with missing hypotheses; publication match: partial; verification: publisher metadata checked; development: annotated; confidence: high}
\end{GSitem}
\begin{GSitem}{1968-162.6}{E}{C1 isometric approximation and the embedding distinction}
\GSfield{Kind; source}{theorem; \GSpage{P000265}}
\GSfield{The notes}{C1-Nash: deforming the metric deforms the immersion; an immersion for one metric can be deformed for any metric, and the last line says obtain an embedding.}
\GSfield{Statement}{The classical Nash--Kuiper approximation theorem says that a short \(C^1\) immersion of an \(n\)-dimensional Riemannian manifold into Euclidean \(\mathbb R^q\), \(q\ge n+1\), admits arbitrarily close \(C^0\) approximations by \(C^1\) isometric immersions. With an initial embedding and the appropriate embedding version, embeddings can be retained. Deforming a mere immersion does not automatically eliminate self-intersections. The full parametric lifting statement requires its own formulation.}
\GSfield{Reconstruction}{Shortness, positive codimension, approximation topology, and the initial embedding condition are supplied.}
\GSfield{Fate}{Nash (1954) and Kuiper (1955) predate this talk; their work is documented in the contemporary Bourbaki exposition \citep[title and references~1--2]{choquet1958nash}. This does not certify the notebook\textquoteright s unqualified Serre-fibration assertion.}
\GSfield{Present status}{As of November 2026: established mathematics predating this talk; the qualification in the statement is essential.}
\GSfield{Assessment}{literal statement: false as written; publication match: partial; verification: publisher metadata checked; development: annotated; confidence: high}
\end{GSitem}
\begin{GSitem}{1968-162.7}{A}{Local extension and microflexibility}
\GSfield{Kind; source}{method; \GSpage{P000265}}
\GSfield{The notes}{The Whitney--Smale continuation principle extends a deformation from part of a region; microscopic realizability is an implicit-function theorem, and often micro implies macro.}
\GSfield{Statement}{For a sheaf of local solutions, microflexibility is a local short-time homotopy-extension property over compact pairs. Flexibility asks for the corresponding full homotopy lifting. Microflexibility alone is not a universal global existence theorem; the passage to flexibility uses additional invariance and geometric hypotheses.}
\GSfield{Reconstruction}{The two lifting properties are distinguished instead of identifying an informal micro-to-macro arrow with a theorem for all PDEs.}
\GSfield{Fate}{\citet[\S3, especially the flexibility propositions]{gromov1969stable} gives a precise sheaf formulation of this programme; received July 1968, published 1969.}
\GSfield{Present status}{As of November 2026: the identified result is established, with the scope stated here.}
\GSfield{Assessment}{literal statement: correct with missing hypotheses; publication match: partial; verification: read the result; development: annotated; confidence: high}
\end{GSitem}
\begin{GSitem}{1968-162.8}{A}{The open invariant relation theorem}
\GSfield{Kind; source}{theorem; \GSpage{P000265-P000266}}
\GSfield{The notes}{Invariance under all diffeomorphisms is called stability; in the stable case micro implies macro, and Theorem2 gives global existence on a noncompact parallelizable connected manifold.}
\GSfield{Statement}{Let \(M\) be a smooth manifold without boundary and with no compact connected component. Let \(E\to M\) be a natural bundle associated to the frame bundle, and let \(\mathcal R\subset J^rE\) be an open relation invariant under the induced action of local diffeomorphisms. Then \(j^r:\operatorname{Sol}(\mathcal R)\to\Gamma(\mathcal R)\) is a weak homotopy equivalence. In particular a formal solution is sufficient for existence of a genuine solution.}
\GSfield{Reconstruction}{Openness, the natural bundle, formal solutions, and the meaning of open manifold are supplied from the theorem; a general stable closed relation is not included.}
\GSfield{Fate}{Exact qualified publication match: \citet[Theorem~4.1.1]{gromov1969stable}, received 8 July 1968 and published in 1969. Record1968-149.4 contains the earlier germs-to-jets programme; this is a more explicit occurrence of the same theorem line, not a new open problem.}
\GSfield{Present status}{As of November 2026: the identified result is established, with the scope stated here.}
\GSfield{Assessment}{literal statement: correct with missing hypotheses; publication match: exact; verification: read the result; development: annotated; confidence: high}
\end{GSitem}
\begin{GSitem}{1968-162.9}{A}{Parallelizability and nonempty fibres}
\GSfield{Kind; source}{claim; \GSpage{P000265-P000266}}
\GSfield{The notes}{Theorem2 lists parallelizable noncompact connected \(M\), a submanifold or subbundle \(A\) with nonempty fibres, micro, and stability; a solution follows.}
\GSfield{Statement}{In the preceding natural-bundle setting, a chosen parallelization trivializes the relevant finite-order frame data (after the usual jet-frame splitting). If the invariant model relation is nonempty, the associated relation bundle admits a section, so Gromov\textquoteright s theorem yields a solution. Nonempty fibres of an arbitrary bundle do not by themselves give a global section; neither an arbitrary submanifold relation nor diffeomorphism invariance alone suffices.}
\GSfield{Reconstruction}{The formal-section step is made explicit. The source does not justify discarding the open-relation or naturality hypotheses.}
\GSfield{Fate}{This is a conditional corollary of \citet[Theorem~4.1.1]{gromov1969stable}; the notebook abbreviations do not define a stronger general theorem.}
\GSfield{Present status}{As of November 2026: the identified result is established, with the scope stated here.}
\GSfield{Assessment}{literal statement: correct with missing hypotheses; publication match: partial; verification: read the result; development: annotated; confidence: high}
\end{GSitem}
\begin{GSitem}{1968-162.10}{E}{Plateau surfaces and reparametrization invariance}
\GSfield{Kind; source}{remark; \GSpage{P000265}}
\GSfield{The notes}{The Plateau problem is an example of invariance under diffeomorphisms; minimal films \(\mathbb R^2\to\mathbb R^3\) are mentioned.}
\GSfield{Statement}{The geometric minimal-surface condition is invariant under changes of source coordinates when expressed intrinsically, but a particular coordinate PDE may include a gauge. Such invariance does not establish openness, microflexibility, an unrestricted h-principle, or existence for arbitrary boundary data.}
\GSfield{Fate}{The additional hypotheses in \citet[Theorem~4.1.1]{gromov1969stable} explain why the Plateau example is an invariance remark rather than an application of the open-relation theorem.}
\GSfield{Present status}{As of November 2026: established mathematics predating this talk; the qualification in the statement is essential.}
\GSfield{Assessment}{literal statement: correct with missing hypotheses; publication match: partial; verification: read the result; development: annotated; confidence: high}
\end{GSitem}
\begin{GSitem}{1968-162.11}{E}{Why an open manifold has no top-dimensional handles}
\GSfield{Kind; source}{method; \GSpage{P000266}}
\GSfield{The notes}{Noncompact implies there is a function without a maximum; this is attached to micro implies macro.}
\GSfield{Statement}{A smooth connected noncompact \(n\)-manifold admits a Morse exhaustion with no critical points of index \(n\). In the open-manifold proof, this allows the required induction along a handle decomposition of index at most \(n-1\). Merely having a function with no global maximum would be much weaker and is not the needed conclusion.}
\GSfield{Reconstruction}{The phrase without max is interpreted as no local maxima of the Morse exhaustion, not simply unboundedness.}
\GSfield{Fate}{This is the explicit topological input in the proof of \citet[Theorem~4.1.1]{gromov1969stable}; it is used as existing background.}
\GSfield{Present status}{As of November 2026: established mathematics predating this talk; the qualification in the statement is essential.}
\GSfield{Assessment}{literal statement: correct with missing hypotheses; publication match: exact; verification: read the result; development: annotated; confidence: high}
\end{GSitem}
\begin{GSitem}{1968-162.12}{E}{The analytic local embedding dimension}
\GSfield{Kind; source}{theorem; \GSpage{P000266}}
\GSfield{The notes}{An \(n\)-dimensional Riemannian manifold embeds locally in \(\mathbb R^{n(n+1)/2}\), and in general a lower dimension does not suffice.}
\GSfield{Statement}{The Janet--Cartan local theorem applies to real-analytic positive-definite Riemannian metrics: near each point there is a real-analytic isometric embedding into \(\mathbb R^{n(n+1)/2}\). This dimension is a generic local analytic bound, not a statement that every particular metric needs it. The same assertion for arbitrary smooth metrics does not follow from the analytic theorem.}
\GSfield{Reconstruction}{Analyticity and the generic meaning of the lower-bound remark are supplied.}
\GSfield{Fate}{Janet (1926) and Cartan (1927) are identified in the contemporary reference list of \citet[references~7--8]{choquet1958nash}. The original articles were not opened, so this attribution is checked through that source.}
\GSfield{Present status}{As of November 2026: established mathematics predating this talk; the qualification in the statement is essential.}
\GSfield{Assessment}{literal statement: correct with missing hypotheses; publication match: partial; verification: indirect; development: annotated; confidence: high}
\end{GSitem}
\begin{GSitem}{1968-162.13}{F}{Approximation by two maps with the same induced metric}
\GSfield{Kind; source}{claim; \GSpage{P000266}}
\GSfield{The notes}{Two maps \(f_1,f_2:M\to\mathbb R^{?s}\) are C0-approximated by smooth maps \(\varphi_1,\varphi_2\) inducing the same metric; the dimension and a formula involving \(0<\lambda<1\) are obscured.}
\GSfield{Statement}{A definite theorem would require the target dimension, compactness assumptions, ranks of the initial maps, allowed induced metric, regularity, and the meaning of the parameter condition. The equality \(\varphi_1^*g_{\mathrm{eucl}}=\varphi_2^*g_{\mathrm{eucl}}\) alone is clear. No unique precise quantified approximation theorem can be recovered from the surviving line.}
\GSfield{Fate}{The published flexibility programme \citep{gromov1969stable} is relevant, but its open invariant relation theorem does not by itself prove this closed metric equality. No exact publication match was established.}
\GSfield{Present status}{As of November 2026: the exact assertion or attribution remains unresolved; no unsolved mathematical problem is inferred.}
\GSfield{Assessment}{literal statement: not determinable; publication match: thematic only; verification: read the result; development: annotated; confidence: low}
\end{GSitem}
\begin{GSitem}{1968-162.14}{F}{Large-dimension approximation with homothetic induced metric}
\GSfield{Kind; source}{claim; \GSpage{P000266}}
\GSfield{The notes}{For large \(t\ge s\), \(t=?+An\); any map is approximated by a similar one with induced metric \(\lambda g\). The target exponent and a cancelled word are unreadable.}
\GSfield{Statement}{The surviving assertion suggests approximation by maps whose induced metric is a scalar multiple of a prescribed metric, but does not determine \(t\), the constant \(A\), whether \(\lambda\) is fixed or freely chosen, the regularity class, or the approximation topology. These choices change the mathematical statement.}
\GSfield{Fate}{Searches by Gromov, stable mappings, isometric approximation, and dimension \(s+An\) led to \citet{gromov1969stable}, but not to a theorem matching all the unreadable data. This is F, not a claim that a defined problem remains unsolved.}
\GSfield{Present status}{As of November 2026: the exact assertion or attribution remains unresolved; no unsolved mathematical problem is inferred.}
\GSfield{Assessment}{literal statement: not determinable; publication match: thematic only; verification: read the result; development: annotated; confidence: low}
\end{GSitem}

\GStalk{164}{1968-11-25}{O. A. Oleinik}{P000266-P000273}{Radkevich on hypoelliptic second-order operators}
Oleinik\textquoteright s heading is on P000267, before Kazhdan\textquoteright s heading on P000268, despite the catalogue numbering. She reviews H\"ormander\textquoteright s sum-of-squares criterion and announces a symbol-derivative/commutator approach for general nonnegative second-order forms. Radkevich\textquoteright s short paper was received 19 November 1968, six days before the talk, and printed in 1969 \citep{radkevich1969hormander}. Gelfand\textquoteright s question about locally integrable solutions at the top of P000268 belongs here.

\begin{GSitem}{1968-164.1}{E}{Hypoellipticity and the second-order class}
\GSfield{Kind; source}{definition; \GSpage{P000267}}
\GSfield{The notes}{\(Pu=a^{ij}u_{x_ix_j}+b^iu_{x_i}+cu\); hypoelliptic is followed by a question mark.}
\GSfield{Statement}{On an open set \(\Omega\subset\mathbb R^n\), a differential operator \(P\) with smooth coefficients is hypoelliptic if for every open \(U\subset\Omega\) and every distribution \(u\in\mathcal D^{\prime}(U)\), \(Pu\in C^\infty(U)\) implies \(u\in C^\infty(U)\). The announced real second-order results concern a symmetric nonnegative principal matrix \((a^{ij})\), not arbitrary complex principal forms.}
\GSfield{Reconstruction}{Distributional and local quantifiers, smoothness, and the nonnegative form are supplied.}
\GSfield{Fate}{The definition and coefficient class are explicit in \citet[pp.~139--140]{oleinikradkevich1971local}; H\"ormander\textquoteright s preceding result supplies the comparison class.}
\GSfield{Present status}{As of November 2026: established mathematics predating this talk; the qualification in the statement is essential.}
\GSfield{Assessment}{literal statement: correct with missing hypotheses; publication match: partial; verification: read the result; development: annotated; confidence: high}
\end{GSitem}
\begin{GSitem}{1968-164.2}{E}{Hörmander brackets imply hypoellipticity and a positive Sobolev gain}
\GSfield{Kind; source}{theorem; \GSpage{P000267}}
\GSfield{The notes}{For \(P=\sum X_i^2+X_0+c\), if at every point there are \(n\) independent commutators then \(P\) is hypoelliptic, using \(\|u\|_\varepsilon\le C(\|u\|_0+\|Pu\|_0)\), \(\varepsilon>0\).}
\GSfield{Statement}{Let \(X_0,\ldots,X_r\) be smooth real vector fields on \(\Omega\subset\mathbb R^n\), and \(c\) smooth. If their iterated Lie brackets span \(T_x\Omega\) at every point, then \(P=\sum_{j=1}^rX_j^2+X_0+c\) is hypoelliptic. For each compact \(K\Subset\Omega\), there are \(C,\varepsilon>0\) such that \(\|u\|_{H^\varepsilon}\le C(\|Pu\|_{L^2}+\|u\|_{L^2})\) for \(u\in C_c^\infty(K)\); localized higher-order estimates yield smoothness.}
\GSfield{Fate}{Already \citet[Theorem~1.1 and its proof]{hormander1967hypoelliptic}, in the 1967 volume (printed February 1968). Radkevich\textquoteright s revised estimates appeared as \citet[Theorems~1,3]{radkevich1969hormander}.}
\GSfield{Present status}{As of November 2026: established mathematics predating this talk; the qualification in the statement is essential.}
\GSfield{Assessment}{literal statement: correct with missing hypotheses; publication match: exact; verification: read the result; development: annotated; confidence: high}
\end{GSitem}
\begin{GSitem}{1968-164.3}{F}{The unproved necessity comment has ambiguous scope}
\GSfield{Kind; source}{remark; \GSpage{P000267}}
\GSfield{The notes}{After \(\varepsilon>0\), necessity not proved is written; a later box says that for the estimate the condition is necessary.}
\GSfield{Statement}{The source does not unambiguously distinguish necessity of bracket generation for hypoellipticity, necessity for a fixed positive Sobolev-gain estimate, or necessity of the subsequent symbol criterion. These are different assertions. Hypoellipticity with smooth coefficients cannot in general be characterized solely by the rank of the vector-field family.}
\GSfield{Reconstruction}{No problem-index entry is created by choosing one of the competing scopes of necessity.}
\GSfield{Fate}{\citet[Theorem~2, example~(6), and concluding paragraph of \S2]{oleinikradkevich1971local} explicitly separates smooth-coefficient phenomena from the more restrictive analytic criterion. This supplies a limitation, not an exact solution of an identifiable notebook necessity question.}
\GSfield{Present status}{As of November 2026: the exact assertion or attribution remains unresolved; no unsolved mathematical problem is inferred.}
\GSfield{Assessment}{literal statement: not determinable; publication match: partial; verification: read the result; development: annotated; confidence: medium}
\end{GSitem}
\begin{GSitem}{1968-164.4}{A}{An elliptic commutator system built from symbol derivatives}
\GSfield{Kind; source}{theorem; \GSpage{P000267}}
\GSfield{The notes}{Use \(P^{(j)}=a^{ij}D_i\) and \(P_{(j)}=\sum\partial_{x_j}a_{ik}\,\partial_{x_i}\partial_{x_k}\,\mathcal E_{-1}\); retain principal parts and their commutators until an elliptic system is obtained.}
\GSfield{Statement}{Let \(L=\sum a^{ij}\partial_i\partial_j+\sum b^i\partial_i+c\) have real smooth coefficients and symmetric nonnegative matrix \(a\). Put \(p(x,\xi)=\sum a^{ij}(x)\xi_i\xi_j\) and let \(Q\) be the first-order drift in the divergence-form rewriting. Form first-order pseudodifferential operators with principal symbols \(\partial_{\xi_j}p\) and \(|\xi|^{-1}\partial_{x_j}p\), together with \(Q\). Suppose on every compact \(K\Subset\Omega\) finitely many of their iterated commutators have principal symbols \(q_J\) satisfying \(\sum_J|q_J(x,\xi)|^2\ge c_K|\xi|^2\) for \(x\in K, |\xi|\ge1\). Then \(L\) is hypoelliptic and has localized positive-gain estimates.}
\GSfield{Reconstruction}{The drift term, nonnegative characteristic form, finite uniform ellipticity on compacts, and the normalization of symbol derivatives are taken from the published formulation. The notes omit these details.}
\GSfield{Fate}{The completed criterion is \citet[Theorem~6, pp.~145--147]{oleinikradkevich1971local}; the 1969 announcements are listed there as references~14--16. The original two-page \citet[Theorems~2--3]{radkevich1969hormander} proves a related general commutator criterion.}
\GSfield{Present status}{As of November 2026: the identified result is established, with the scope stated here.}
\GSfield{Assessment}{literal statement: correct with missing hypotheses; publication match: partial; verification: read the result; development: annotated; confidence: high}
\end{GSitem}
\begin{GSitem}{1968-164.5}{A}{Toward a formulation containing Hörmander’s theorem}
\GSfield{Kind; source}{remark; \GSpage{P000267}}
\GSfield{The notes}{There is as yet no formulation including H\"ormander\textquoteright s theorem; the first-order case should work without pseudodifferential operators.}
\GSfield{Statement}{The programme is to place vector-field sum-of-squares estimates and general second-order nonnegative-form estimates in a common commutator framework. For differential first-order generators, brackets remain differential; derivative-of-symbol operators normalized by order-minus-one factors generally require pseudodifferential calculus.}
\GSfield{Reconstruction}{The chronology reflects the paper submission separately from its 1969 publication.}
\GSfield{Fate}{\citet[Theorems~1--3]{radkevich1969hormander}, received 19 November 1968, treats both the H\"ormander case and a higher-order extension; the systematic synthesis is \citet[Theorems~1,5--6]{oleinikradkevich1971local}. This is a near-contemporary development, not evidence that every general criterion was already stated at the talk.}
\GSfield{Present status}{As of November 2026: the identified result is established, with the scope stated here.}
\GSfield{Assessment}{literal statement: correct with missing hypotheses; publication match: partial; verification: read the result; development: annotated; confidence: high}
\end{GSitem}
\begin{GSitem}{1968-164.6}{F}{When are all generalized solutions locally integrable?}
\GSfield{Kind; source}{question; \GSpage{P000268, opening line before Kazhdan}; I. M. Gelfand}
\GSfield{The notes}{Gelfand: for applications it is important to know when every solution is a summable function.}
\GSfield{Statement}{One possible precise property is \(\ker(P:\mathcal D^{\prime}(U)\to\mathcal D^{\prime}(U))\subset L^1_{\mathrm{loc}}(U)\) for every open \(U\). Hypoellipticity implies it, but the notebook does not specify homogeneous versus inhomogeneous equations, the allowed right-hand sides, or which coefficient class is to be characterized. Global integrability would be a different condition.}
\GSfield{Reconstruction}{The displayed kernel inclusion is explicitly one local homogeneous interpretation, not a claimed transcription.}
\GSfield{Fate}{The regularity criteria above give sufficient conditions \citep{oleinikradkevich1971local}, but they do not identify an exact necessary-and-sufficient theorem for this abbreviated question. Searches of Oleinik--Radkevich regularity work did not resolve its missing scope.}
\GSfield{Present status}{As of November 2026: the exact assertion or attribution remains unresolved; no unsolved mathematical problem is inferred.}
\GSfield{Assessment}{literal statement: not determinable; publication match: partial; verification: read the result; development: annotated; confidence: medium}
\end{GSitem}

\GStalk{163}{1968-11-25}{D. A. Kazhdan}{P000266-P000268}{Dwork rationality and non-Archimedean Fredholm theory}
The heading begins after Gelfand\textquoteright s intervention on P000268 and the talk continues through P000271. Closed-point Euler products and extension-field point counts lead to Dwork\textquoteright s rationality theorem and the Borel--Dwork criterion. The final page supplies the non-Archimedean tools: compact operators, Fredholm determinants, Teichm\"uller representatives, analytic zeros, and integrality of reciprocal roots. These are an exposition of work already published in 1960--1962, not new conjectures of November 1968.

\begin{GSitem}{1968-163.1}{E}{Scheme points, closed points, and geometric points}
\GSfield{Kind; source}{definition; \GSpage{P000268}}
\GSfield{The notes}{\(A=k[x_1,\ldots,x_n]/\mathfrak J\); maximal ideals give closed points, prime ideals give points; \(A/I\) is a finite extension of \(k\). A map \(A\to K\) is called a geometric point.}
\GSfield{Statement}{For a finitely generated \(k\)-algebra \(A\), points of \(\operatorname{Spec}A\) are prime ideals and closed points are maximal ideals. Zariski\textquoteright s lemma says that the residue field at a maximal ideal is finite over \(k\). A geometric point is a morphism from the spectrum of an algebraically closed field; a homomorphism to an arbitrary field is a field-valued point, whose kernel need not be maximal if the field extension is transcendental. For \(\mathbb R[x]/(x^2+1)\), one closed point yields two \(\mathbb C\)-valued points.}
\GSfield{Reconstruction}{The residue-field finiteness assertion is restricted to closed points, and geometric versus arbitrary field-valued points is distinguished.}
\GSfield{Fate}{This is pre-existing background for the closed-point definition used by \citet[\S1]{dwork1960rationality}. It also corrects the shorthand at record155.}
\GSfield{Present status}{As of November 2026: established mathematics predating this talk; the qualification in the statement is essential.}
\GSfield{Assessment}{literal statement: correct with missing hypotheses; publication match: partial; verification: read the result; development: annotated; confidence: high}
\end{GSitem}
\begin{GSitem}{1968-163.2}{E}{Reduced schemes and the word variety}
\GSfield{Kind; source}{remark; \GSpage{P000268}}
\GSfield{The notes}{An algebraic variety is distinguished from a scheme by absence of nilpotent elements.}
\GSfield{Statement}{An affine scheme \(\operatorname{Spec}A\) is reduced precisely when \(A\) has no nonzero nilpotents. A variety has additional convention-dependent hypotheses, commonly finite type, separatedness, and sometimes irreducibility. Reducedness alone is not a definition of an arbitrary variety. Finite-field point counts and the associated zeta function are unchanged by passage to the reduction.}
\GSfield{Reconstruction}{The source\textquoteright s local terminology is translated into the precise property reduced, without imposing one universal convention for variety.}
\GSfield{Fate}{Dwork\textquoteright s zeta function counts field-valued points \citep[\S1]{dwork1960rationality}; maps to fields kill nilpotents, which explains this last assertion directly.}
\GSfield{Present status}{As of November 2026: established mathematics predating this talk; the qualification in the statement is essential.}
\GSfield{Assessment}{literal statement: false as written; publication match: partial; verification: read the result; development: annotated; confidence: high}
\end{GSitem}
\begin{GSitem}{1968-163.3}{E}{Finite fields and Frobenius}
\GSfield{Kind; source}{theorem; \GSpage{P000268}}
\GSfield{The notes}{\(\mathbb F_{p^n}\) is obtained by a displayed polynomial with overwritten exponents; its automorphism group is cyclic, generated by \(x\mapsto x^p\).}
\GSfield{Statement}{Inside an algebraic closure of \(\mathbb F_p\), the roots of \(X^{p^n}-X\) form the unique field of cardinality \(p^n\). Its Galois group over \(\mathbb F_p\) is cyclic of order \(n\), generated by Frobenius. To present this field as a quotient of \(\mathbb F_p[X]\), one chooses an irreducible polynomial of degree \(n\); the entire polynomial \(X^{p^n}-X\) does not define a field quotient when \(n>1\).}
\GSfield{Reconstruction}{A correct root-set construction is supplied independently of the unreadable fraction; it is not offered as a recovered reading.}
\GSfield{Fate}{This is the elementary finite-field structure used throughout \citet[\S1]{dwork1960rationality}. The corrupted displayed quotient cannot be certified as a defining polynomial.}
\GSfield{Present status}{As of November 2026: established mathematics predating this talk; the qualification in the statement is essential.}
\GSfield{Assessment}{literal statement: not determinable; publication match: partial; verification: read the result; development: annotated; confidence: high}
\end{GSitem}
\begin{GSitem}{1968-163.4}{E}{The closed-point Euler product}
\GSfield{Kind; source}{definition; \GSpage{P000268-P000269}}
\GSfield{The notes}{Degree of a closed point is its residue-field degree; \(\zeta(t,A)=\prod_x(1-t^{\deg x})^{-1}\), with \(t\) formal.}
\GSfield{Statement}{For a finite-type scheme \(X\) over \(\mathbb F_q\), let \(|X|\) be its set of closed points and \(\deg x=[\kappa(x):\mathbb F_q]\). Then \(Z(X,t)=\prod_{x\in|X|}(1-t^{\deg x})^{-1}\in1+t\mathbb Z[[t]]\). Finitely many closed points have each bounded degree, so each coefficient is a finite sum. The product is initially formal, without a claim of convergence on the complex unit circle.}
\GSfield{Fate}{This is Dwork\textquoteright s input zeta function \citep[\S1]{dwork1960rationality}.}
\GSfield{Present status}{As of November 2026: established mathematics predating this talk; the qualification in the statement is essential.}
\GSfield{Assessment}{literal statement: correct with missing hypotheses; publication match: exact; verification: read the result; development: annotated; confidence: high}
\end{GSitem}
\begin{GSitem}{1968-163.5}{E}{Dwork rationality}
\GSfield{Kind; source}{theorem; \GSpage{P000269}}
\GSfield{The notes}{For every affine algebraic variety over a finite field the Euler product is rational; this was a Weil conjecture and was known earlier for curves.}
\GSfield{Statement}{For every finite-type scheme \(X/\mathbb F_q\), \(Z(X,t)\) is the Taylor expansion at zero of an element of \(\mathbb Q(t)\). The affine statement extends by finite stratification and multiplicativity for disjoint locally closed strata; nilpotents do not affect point counts. This is the rationality part, not the entire collection of Weil conjectures.}
\GSfield{Reconstruction}{The finite-type scheme extension is explained; no purity, functional-equation, or Riemann-hypothesis conclusion is attributed to Dwork\textquoteright s rationality theorem.}
\GSfield{Fate}{Established before the talk by \citet[main theorem]{dwork1960rationality}, published in 1960; the Russian translation appeared in 1961.}
\GSfield{Present status}{As of November 2026: established mathematics predating this talk; the qualification in the statement is essential.}
\GSfield{Assessment}{literal statement: correct; publication match: exact; verification: read the result; development: annotated; confidence: high}
\end{GSitem}
\begin{GSitem}{1968-163.6}{E}{Affine-line closed points and irreducible-polynomial counts}
\GSfield{Kind; source}{construction; \GSpage{P000269}}
\GSfield{The notes}{For \(\mathbb F_p[x]\), \(Z=1/(1-pt)\), and the degree-one, two, three, four factors have exponents \(p,(p^2-p)/2,(p^3-p)/3,(p^4-p^2)/4\).}
\GSfield{Statement}{The number of degree-\(d\) closed points of \(\mathbb A^1_{\mathbb F_q}\) is \(b_d=d^{-1}\sum_{e\mid d}\mu(e)q^{d/e}\). Indeed \(q^r=\sum_{d\mid r}db_d\), because a degree-\(d\) closed point contributes \(d\) points over \(\mathbb F_{q^r}\) exactly when \(d\mid r\). Thus \(Z(\mathbb A^1,t)=(1-qt)^{-1}\). Every displayed Euler factor should have negative exponent in the product.}
\GSfield{Reconstruction}{The general degree formula is derived from Frobenius orbits; the fourth factor is interpreted as a denominator like the preceding factors.}
\GSfield{Fate}{The example is pre-existing elementary finite-field counting, consistent with \citet[\S1]{dwork1960rationality}; M\"obius inversion fills the notebook\textquoteright s omitted general expression.}
\GSfield{Present status}{As of November 2026: established mathematics predating this talk; the qualification in the statement is essential.}
\GSfield{Assessment}{literal statement: correct with missing hypotheses; publication match: partial; verification: read the result; development: annotated; confidence: high}
\end{GSitem}
\begin{GSitem}{1968-163.7}{E}{Poles of formal Euler factors are not poles of the continued zeta function}
\GSfield{Kind; source}{remark; \GSpage{P000269}}
\GSfield{The notes}{After \(Z(\mathbb A^1,t)=1/(1-pt)\), the note says the infinite expansion has poles at one.}
\GSfield{Statement}{Each Euler factor \((1-t^d)^{-1}\) has poles at roots of unity, including one. The infinite product for the affine line is initially convergent only for \(|t|<1/q\), and its rational continuation has its sole pole at \(t=1/q\). Poles of individual factors outside the convergence region cannot be read as poles of that continuation.}
\GSfield{Reconstruction}{The unclear scope of poles at one is interpreted as a comment on the individual factors, not silently attributed to the rational function.}
\GSfield{Fate}{The affine-line identity in the notebook itself resolves the distinction; the formal/analytic separation is consistent with \citet[\S1]{dwork1960rationality}.}
\GSfield{Present status}{As of November 2026: established mathematics predating this talk; the qualification in the statement is essential.}
\GSfield{Assessment}{literal statement: not determinable; publication match: partial; verification: read the result; development: annotated; confidence: high}
\end{GSitem}
\begin{GSitem}{1968-163.8}{E}{Uniqueness of formal Euler factorization}
\GSfield{Kind; source}{question; \GSpage{P000269}}
\GSfield{The notes}{Uniqueness? The degree is chosen at the last possible step to match a coefficient; hence every series can be expanded.}
\GSfield{Statement}{Every series \(F(t)\in1+t\mathbb Z[[t]]\) has a unique formal expansion \(F(t)=\prod_{n\ge1}(1-t^n)^{-b_n}\) with \(b_n\in\mathbb Z\). Choose \(b_n\) recursively to cancel the coefficient of \(t^n\) after the earlier factors are removed; uniqueness follows from the first degree at which two sequences differ. The exponents need not be nonnegative, so an arbitrary such series need not be the zeta function of a variety.}
\GSfield{Reconstruction}{The coefficient ring, constant term, allowance of negative exponents, and formal-product meaning are supplied.}
\GSfield{Fate}{The notebook itself gives the recursive method. This is a solved formal-algebra exercise in the exposition and receives no new historical problem identifier.}
\GSfield{Present status}{As of November 2026: established mathematics predating this talk; the qualification in the statement is essential.}
\GSfield{Assessment}{literal statement: correct with missing hypotheses; publication match: exact; verification: read the result; development: annotated; confidence: high}
\end{GSitem}
\begin{GSitem}{1968-163.9}{E}{The exponential point-counting formula}
\GSfield{Kind; source}{theorem; \GSpage{P000270}}
\GSfield{The notes}{\(Z=\exp(\sum_eN_et^e/e)\), with \(N_e\) the number of homomorphisms \(A\to\mathbb F_{p^e}\); integrality is less obvious in this form.}
\GSfield{Statement}{For finite-type \(X/\mathbb F_q\), \(N_r=\#X(\mathbb F_{q^r})\) and \(Z(X,t)=\exp(\sum_{r\ge1}N_rt^r/r)\). Taking formal logarithms of the Euler product proves the identity through \(N_r=\sum_{d\mid r}db_d\). Integral coefficients follow directly from the Euler product; arbitrary integer sequences \(N_r\) do not automatically give integral exponential coefficients.}
\GSfield{Reconstruction}{The divisibility relation with closed-point counts explains both the logarithm and the integrality caveat.}
\GSfield{Fate}{This identity is the starting formula in \citet[\S1]{dwork1960rationality}.}
\GSfield{Present status}{As of November 2026: established mathematics predating this talk; the qualification in the statement is essential.}
\GSfield{Assessment}{literal statement: correct; publication match: exact; verification: read the result; development: annotated; confidence: high}
\end{GSitem}
\begin{GSitem}{1968-163.10}{E}{Spec Z and the usual Riemann zeta function}
\GSfield{Kind; source}{remark; \GSpage{P000270}}
\GSfield{The notes}{Applying the construction to \(A=\mathbb Z\) gives the usual zeta function.}
\GSfield{Statement}{For \(\operatorname{Spec}\mathbb Z\), the norm of the closed point \((p)\) is \(p\), and its arithmetic zeta function is \(\prod_p(1-p^{-s})^{-1}=\zeta(s)\) for \(\Re s>1\). There is no single finite base field and hence no common integer degree producing this product in one formal variable \(t\). The analogy uses residue-field norms.}
\GSfield{Reconstruction}{Norms and the complex convergence half-plane replace the finite-field degree variable.}
\GSfield{Fate}{The assertion is a classical Euler-product analogy, not an application of Dwork rationality to the Riemann zeta function \citep[Introduction]{dwork1960rationality}.}
\GSfield{Present status}{As of November 2026: established mathematics predating this talk; the qualification in the statement is essential.}
\GSfield{Assessment}{literal statement: correct with missing hypotheses; publication match: partial; verification: read the result; development: annotated; confidence: high}
\end{GSitem}
\begin{GSitem}{1968-163.11}{E}{The Borel–Dwork rationality criterion needs compatible radii}
\GSfield{Kind; source}{theorem; \GSpage{P000270}}
\GSfield{The notes}{The integral-coefficient series converges in a complex disc and is meromorphic p-adically; therefore it is automatically rational.}
\GSfield{Statement}{Let \(F(t)\in\mathbb Z[[t]]\). Suppose it is holomorphic in the complex disc \(|t|<R\) and represents a quotient of convergent power series on the \(\mathbb C_p\)-disc \(|t|_p<r\), with \(|p|_p=p^{-1}\) and \(Rr>1\). Then \(F\in\mathbb Q(t)\). Small positive convergence radii in both norms alone do not suffice; p-adic meromorphy on arbitrarily large discs does suffice once the complex radius is positive.}
\GSfield{Reconstruction}{The crucial inequality between radii, the integral coefficient condition, and the meaning of meromorphic are restored.}
\GSfield{Fate}{Exact pre-talk source: \citet[\S4, Theorem~2; Russian translation p.~68]{dwork1960rationality}. A modern exposition states a slightly broader meromorphic-complex version \citep[Theorem~6.8]{thorne2010padic}.}
\GSfield{Present status}{As of November 2026: established mathematics predating this talk; the qualification in the statement is essential.}
\GSfield{Assessment}{literal statement: correct with missing hypotheses; publication match: exact; verification: read the result; development: annotated; confidence: high}
\end{GSitem}
\begin{GSitem}{1968-163.12}{E}{The p-adic valuation ring, residue field, and convergence criterion}
\GSfield{Kind; source}{definition; \GSpage{P000270}}
\GSfield{The notes}{\(\mathbb Z_p=\{|x|_p\le1\}\), \(M_p=\{|x|_p<1\}\), residue field \(\mathbb F_p\); \(\sum x_n\) converges iff \(x_n\to0\). Banach spaces carry over.}
\GSfield{Statement}{The valuation ring of \(\mathbb Q_p\) is \(\mathbb Z_p\), its maximal ideal is \(p\mathbb Z_p\), and its residue field is \(\mathbb F_p\). In any complete ultrametric normed vector space a series converges if and only if its terms tend to zero: the ultrametric triangle inequality bounds each finite tail by its largest term. A non-Archimedean Banach space is a complete normed vector space whose norm satisfies that inequality.}
\GSfield{Reconstruction}{Completeness is essential for the implication from vanishing terms to convergence.}
\GSfield{Fate}{The coefficient-field and Banach conventions appear in \citet[\S1]{serre1962endomorphismes}; these are older basic facts.}
\GSfield{Present status}{As of November 2026: established mathematics predating this talk; the qualification in the statement is essential.}
\GSfield{Assessment}{literal statement: correct; publication match: exact; verification: read the result; development: annotated; confidence: high}
\end{GSitem}
\begin{GSitem}{1968-163.13}{E}{Completely continuous operators and p-adic Fredholm determinants}
\GSfield{Kind; source}{theorem; \GSpage{P000271}}
\GSfield{The notes}{In a Banach space a completely continuous operator is nuclear, hence \(\det(1-tA)\) is an entire p-adic function.}
\GSfield{Statement}{Let \(K\) be a complete nontrivially valued non-Archimedean field and let \(E\) be an orthonormalizable Banach space, isomorphic to \(c_0(I,K)\). If \(A\) is a norm limit of finite-rank operators, its Fredholm determinant is well defined by finite-rank approximation and is an entire power series in \(t\). For \(K\) discretely valued, equivalent norms reduce general Banach spaces to this setting. Nuclear is used in Serre\textquoteright s non-Archimedean sense, not the real/complex trace-class condition.}
\GSfield{Reconstruction}{The coefficient field, complete-continuity convention, and orthonormalizability are explicit.}
\GSfield{Fate}{Precisely \citet[\S\S4--5, Proposition~7]{serre1962endomorphismes}; received 18 December 1961 and published in 1962.}
\GSfield{Present status}{As of November 2026: established mathematics predating this talk; the qualification in the statement is essential.}
\GSfield{Assessment}{literal statement: correct with missing hypotheses; publication match: exact; verification: read the result; development: annotated; confidence: high}
\end{GSitem}
\begin{GSitem}{1968-163.14}{E}{C\_p and Teichmüller representatives}
\GSfield{Kind; source}{construction; \GSpage{P000271}}
\GSfield{The notes}{\(\Omega\) is the completion of the algebraic closure of \(\mathbb Q_p\); its valuation ring modulo the maximal ideal is \(F\), and \(F^*\to Z_\Omega^*\) has canonical representatives.}
\GSfield{Statement}{Write \(\mathbb C_p=\widehat{\overline{\mathbb Q}_p}\). Its valuation ring has residue field \(\overline{\mathbb F}_p\). If \(a\in\mathbb F_{p^m}^*\), its Teichm\"uller lift is the unique \((p^m-1)\)-st root of unity reducing to \(a\). These lifts are compatible in finite subfields and give a multiplicative section \(\overline{\mathbb F}_p^*\to\mathcal O_{\mathbb C_p}^*\), extended by \([0]=0\). The section is generally not additive.}
\GSfield{Reconstruction}{The completion is metric, the residue field is identified, and multiplicative is distinguished from a ring section.}
\GSfield{Fate}{The field and root-of-unity lifts are part of the p-adic construction in \citet[\S\S1--2]{dwork1960rationality}.}
\GSfield{Present status}{As of November 2026: established mathematics predating this talk; the qualification in the statement is essential.}
\GSfield{Assessment}{literal statement: correct with missing hypotheses; publication match: exact; verification: read the result; development: annotated; confidence: high}
\end{GSitem}
\begin{GSitem}{1968-163.15}{E}{Nonzero entire p-adic functions have finitely many zeros on bounded discs}
\GSfield{Kind; source}{theorem; \GSpage{P000271}}
\GSfield{The notes}{An entire series \(\sum c_kt^k\), with coefficients tending sufficiently fast to zero, has finitely many zeros in every finite-radius disc; Shnirelman, maximum principle, Schwarz lemma.}
\GSfield{Statement}{Over \(\mathbb C_p\), an entire power series satisfies \(|c_k|_pR^k\to0\) for every real \(R>0\). If it is not identically zero, it has only finitely many zeros, counted with multiplicity, in every closed disc of finite radius. This follows by the non-Archimedean Weierstrass/Newton-polygon argument, not from topological compactness of the disc.}
\GSfield{Reconstruction}{The coefficient growth, nonzero condition, and closed-disc interpretation are made explicit.}
\GSfield{Fate}{\citet[\S4, Propositions~1--3 and the Shnirelman discussion]{dwork1960rationality} records the analytic facts used in rationality. The notebook omits the exclusion of the zero function.}
\GSfield{Present status}{As of November 2026: established mathematics predating this talk; the qualification in the statement is essential.}
\GSfield{Assessment}{literal statement: correct with missing hypotheses; publication match: exact; verification: read the result; development: annotated; confidence: high}
\end{GSitem}
\begin{GSitem}{1968-163.16}{E}{The unit ball in Q\_p is compact}
\GSfield{Kind; source}{claim; \GSpage{P000271}}
\GSfield{The notes}{The unit disc in \(\mathbb Q_p\) is not compact.}
\GSfield{Statement}{The closed unit ball \(\mathbb Z_p\subset\mathbb Q_p\) is compact: it is complete and has a finite covering by residue classes modulo \(p^n\) at every scale. The closed unit ball of \(\mathbb C_p\) is not compact, because its infinite residue field yields infinitely many pairwise distance-one points. Thus \(\mathbb C_p\), the preceding field \(\Omega\), may have been intended, but the scan writes \(\mathbb Q_p\).}
\GSfield{Reconstruction}{The possible intended field is flagged rather than silently substituted.}
\GSfield{Fate}{The valuation-field distinction is standard background \citep[\S1]{serre1962endomorphismes}; the elementary covering argument directly settles the literal statement.}
\GSfield{Present status}{As of November 2026: established mathematics predating this talk; the qualification in the statement is essential.}
\GSfield{Assessment}{literal statement: false as written; publication match: partial; verification: read the result; development: annotated; confidence: high}
\end{GSitem}
\begin{GSitem}{1968-163.17}{E}{Integral power-series coefficients force integral reciprocal roots in reduced form}
\GSfield{Kind; source}{theorem; \GSpage{P000271}}
\GSfield{The notes}{\(F(z)=\prod(1-\alpha_i z)/\prod(1-\beta_jz)=1+\sum c_kz^k\), \(c_k\in\mathbb Z\); algebraic \(\alpha_i,\beta_j\) must be algebraic integers.}
\GSfield{Statement}{If \(F\in\mathbb Q(z)\cap(1+z\mathbb Z[[z]])\) is expressed as a reduced quotient, normalized with numerator and denominator constant terms one, then every reciprocal zero and reciprocal pole is an algebraic integer. Indeed both \(F\) and its formal inverse have p-adically integral coefficients; neither can have a zero or pole at \(|z|_p<1\). Apply this at every algebraic embedding and every prime. Cancellation must be removed: \((1-z/2)/(1-z/2)=1\) is an immediate counterexample without that condition.}
\GSfield{Reconstruction}{Reducedness, normalization, and the inverse-series argument supply the missing hypotheses and handle numerator zeros as well as poles.}
\GSfield{Fate}{The notebook\textquoteright s valuation argument is made complete here; it is classical arithmetic of rational integral series, consistent with the integrality framework of \citet[\S4]{dwork1960rationality}.}
\GSfield{Present status}{As of November 2026: established mathematics predating this talk; the qualification in the statement is essential.}
\GSfield{Assessment}{literal statement: correct with missing hypotheses; publication match: partial; verification: read the result; development: annotated; confidence: high}
\end{GSitem}

\GStalk{165}{}{}{P000271-P000273}{1969}
The year heading marks the calendar transition to the meeting dated 10 February 1969 on P000272. It contains no independent mathematical assertion and receives no occurrence item.


\GStalk{166}{1969-02-10}{G. A. Margulis}{P000271-P000273}{Counting orbit points in a hyperbolic disc}
The Margulis passage is only the opening part of P000272; Gelfand\textquoteright s 3-adic arithmetic begins on the same page and belongs to record167. The notice describes lifts of a point to the universal cover, a leading count expressed through area and boundary length, a claimed square-root remainder, and the use of horocycle ergodicity. The leading asymptotic has an exact 1969 publication match; the claimed error requires separate treatment. Linnik\textquoteright s sphere-point work is invoked as an analogy, not a theorem transferring one count to the other.

\begin{GSitem}{1968-166.1}{A}{Hyperbolic orbit counting and the leading asymptotic}
\GSfield{Kind; source}{theorem; \GSpage{P000272, first paragraph}}
\GSfield{The notes}{Integer points are preimages on the universal cover; the count is area plus a term proportional to boundary length, both of the same order.}
\GSfield{Statement}{Let \(M=\Gamma\backslash\mathbb H^2\) be a compact hyperbolic surface of curvature \(-1\), with \(\Gamma\) torsion free, and fix \(z,w\in\mathbb H^2\). Then \(N_\Gamma(R;z,w)=\#\{\gamma\in\Gamma:d(z,\gamma w)\le R\}\sim\pi e^R/\operatorname{area}(M)\), equivalently \(N_\Gamma(R;z,w)\sim\operatorname{area}(B_R)/\operatorname{area}(M)\). Here \(\operatorname{area}(B_R)=2\pi(\cosh R-1)\) and \(\operatorname{length}(\partial B_R)=2\pi\sinh R\); both grow like \(\pi e^R\).}
\GSfield{Reconstruction}{Compactness, curvature, lattice normalization, and the constant are supplied. No independent lower-order boundary term is inferred from two quantities of the same exponential order.}
\GSfield{Fate}{\citet[Theorems~1--2]{margulis1969applications}, received 6 June 1969 and published later that year, establishes the leading asymptotic and its constant. It does not provide the notebook\textquoteright s unspecified two-term area-plus-boundary formula.}
\GSfield{Present status}{As of November 2026: the identified result is established, with the scope stated here.}
\GSfield{Assessment}{literal statement: correct with missing hypotheses; publication match: partial; verification: read the result; development: annotated; confidence: high}
\end{GSitem}
\begin{GSitem}{1968-166.2}{C}{The square-root remainder and the hyperbolic circle problem}
\GSfield{Kind; source}{claim; \GSpage{P000272, remainder line}}
\GSfield{The notes}{The remainder is of the order of the square root of the first terms. Neither coefficients nor a precise error expression are written.}
\GSfield{Statement}{The precise modern question reconstructed here is the pointwise hyperbolic circle problem. For a fixed cofinite Fuchsian group and fixed points, subtract the full spectral main term, including contributions from Laplace eigenvalues below \(1/4\). Is the remaining error \(O_{\varepsilon,\Gamma,z,w}(X^{1/2+\varepsilon})\) for every \(\varepsilon>0\), where \(X\asymp e^R\)? This is weaker than a literal uniform \(O(e^{R/2})\) claim and more specific than the notebook\textquoteright s unspecified first terms.}
\GSfield{Reconstruction}{The group class, pointwise quantifiers, spectral subtraction, and epsilon loss are explicit modern additions. The notebook itself does not state this fully qualified conjecture.}
\GSfield{Fate}{\citet[Theorem~6]{margulis1969applications} distinguishes the count from complementary-series contributions. The later standard conjecture and the classical \(O(X^{2/3})\) barrier are stated in \citet[Introduction, equations~(1.3)--(1.4)]{chatzakosetal2025hyperbolic}, first posted 16 June 2025. Their Theorem~1.1 improves the error only for specified Heegner points; it does not solve the general pointwise problem. \citet[Theorems~1.1--1.2]{biro2026localsquare}, first posted 13 April 2026, studies a local \(L^2\) average and a conditional improvement of its exponent; it does not settle the pointwise conjecture.}
\GSfield{Present status}{As of November 2026: the reconstructed general pointwise \(X^{1/2+\varepsilon}\) bound remains open in the searched primary literature. The 2025 special-point logarithmic saving and the 2026 local-mean conditional improvement are different partial results \citep[Introduction and Theorem~1.1]{chatzakosetal2025hyperbolic}\citep[\S1.1]{biro2026localsquare}. The original unspecified \(O(e^{R/2})\) announcement is not verified.}
\GSfield{Assessment}{literal statement: not determinable; publication match: partial; verification: read the result; development: annotated; confidence: high; problem \hyperref[prob:GP-1968-20]{GP-1968-20}}
\GSfield{Worth developing}{Worth developing now: separating the main spectral terms from the remainder exposes the exact obstruction to reading the notebook as a proved square-root bound.}
\end{GSitem}
\begin{GSitem}{1968-166.3}{A}{Horocycle ergodicity in counting}
\GSfield{Kind; source}{method; \GSpage{P000272, method line}}
\GSfield{The notes}{Ergodicity of the horospherical flow is used in an important way.}
\GSfield{Statement}{For compact negatively curved manifolds, the geodesic flow has stable and unstable horospherical foliations. Suitable transverse or leafwise measures and their asymptotic invariance allow one to relate dynamical equidistribution to counts in the universal cover. This explains the counting method without asserting that bare ergodicity alone supplies a quantitative error rate.}
\GSfield{Reconstruction}{The published measure construction is distinguished from an effective mixing estimate, which would be needed for a specific remainder.}
\GSfield{Fate}{\citet[Theorem~5 and Remark~6]{margulis1969applications} explicitly says that the horospherical measure construction is used to prove the counting theorems; received June 1969, published 1969.}
\GSfield{Present status}{As of November 2026: the identified result is established, with the scope stated here.}
\GSfield{Assessment}{literal statement: correct with missing hypotheses; publication match: partial; verification: read the result; development: annotated; confidence: high}
\end{GSitem}
\begin{GSitem}{1968-166.4}{F}{The analogy with Linnik’s integer points on spheres}
\GSfield{Kind; source}{remark; \GSpage{P000272, Linnik notice}}
\GSfield{The notes}{Linnik has work on how many points lie on a sphere.}
\GSfield{Statement}{The notice connects homogeneous dynamics with arithmetic lattice-point problems. It does not specify a radius sequence, congruence restriction, counting theorem, or equidistribution claim, and does not state that the hyperbolic and Euclidean spherical problems have identical errors.}
\GSfield{Reconstruction}{Only the explicit analogy is retained; no theorem from Linnik\textquoteright s broader work is imported.}
\GSfield{Fate}{This is a thematic return to the earlier Linnik discussion indexed as 1964-14.2. The current line is too short to identify an additional precise result or recurring problem.}
\GSfield{Present status}{As of November 2026: the exact assertion or attribution remains unresolved; no unsolved mathematical problem is inferred.}
\GSfield{Assessment}{literal statement: not determinable; publication match: thematic only; verification: unresolved lead; development: annotated; confidence: medium}
\end{GSitem}

\GStalk{167}{1969-02-10}{I. M. Gelfand}{P000271-P000274}{3-adic arithmetic for schoolchildren}
The mathematical passage is on P000272, after Margulis and a list of prospective speakers. Gelfand introduces left-infinite ternary expansions, carries, the expansions of \(-1\) and \(1/2\), and the obstruction to division in base ten. The comparison with real decimal expansions is pedagogical, not a claim that the real numbers are mathematically defective. The surrounding catalogue pages contain other records.

\begin{GSitem}{1968-167.1}{E}{Left-infinite digits and the 3-adic topology}
\GSfield{Kind; source}{construction; \GSpage{P000272}}
\GSfield{The notes}{\ru{3-адические числа (для школьников)}; infinitely many digits extend to the left, and the rules and topology are introduced by analogy.}
\GSfield{Statement}{\(\mathbb Z_3=\varprojlim_k\mathbb Z/3^k\mathbb Z\). Each element has a unique expansion \(\sum_{j\geq0}a_j3^j\), \(a_j\in\{0,1,2\}\), convergent for \(|x|_3=3^{-v_3(x)}\). Addition and multiplication are compatible finite carry operations modulo \(3^k\); \(\mathbb Q_3=\mathbb Z_3[1/3]\). The digit construction first describes the ring of integers, not every element of the field without a finite right-hand tail.}
\GSfield{Reconstruction}{The inverse limit, valuation and distinction between the integer ring and field make the informal construction precise.}
\GSfield{Fate}{Classical arithmetic, already known before the seminar; \citet[Chs.~3--4]{gouvea2020padic} is a modern exposition, not the first publication.}
\GSfield{Present status}{As of November 2026: established in the stated scope; this occurrence does not pose an unresolved problem.}
\GSfield{Assessment}{literal statement: correct with missing hypotheses; publication match: exact; verification: publisher metadata checked; development: annotated; confidence: high}
\end{GSitem}
\begin{GSitem}{1968-167.2}{E}{The ternary expansions of minus one and one half}
\GSfield{Kind; source}{claim; \GSpage{P000272}}
\GSfield{The notes}{The sums display \(\ldots2222=-1\) and \(\ldots11112=1/2\); the latter is obtained by multiplying by two and carrying.}
\GSfield{Statement}{In \(\mathbb Z_3\), \(\sum_{j\geq0}2\cdot3^j=-1\) and \(2+\sum_{j\geq1}3^j=1/2\). Indeed the finite truncations are \(3^{k+1}-1\) and \((3^{k+1}+1)/2\); their errors tend to zero in the 3-adic norm. The displays are expansions in base three.}
\GSfield{Reconstruction}{Finite-truncation verification is supplied here; no research novelty is claimed.}
\GSfield{Fate}{These are elementary instances of the geometric-series calculation in the 3-adic ring; the modern reference is \citet[Ch.~4]{gouvea2020padic}. The finite-truncation identities also verify the notebook directly.}
\GSfield{Present status}{As of November 2026: established in the stated scope; this occurrence does not pose an unresolved problem.}
\GSfield{Assessment}{literal statement: correct; publication match: exact; verification: publisher metadata checked; development: annotated; confidence: high}
\end{GSitem}
\begin{GSitem}{1968-167.3}{E}{Decimal identifications and the nonfield of 10-adic integers}
\GSfield{Kind; source}{remark; \GSpage{P000272}}
\GSfield{The notes}{The real expansion \(0.999\ldots=1\) is contrasted with a disconnected digit space. The last line says \ru{10-адические --- деления нет} and notes the absence of \(1/5\).}
\GSfield{Statement}{The real value map from digit sequences to \([0,1]\) is continuous and identifies the familiar terminating/nonterminating pairs; its image is connected although its domain is totally disconnected. In contrast, \(\mathbb Z_{10}=\varprojlim_k\mathbb Z/10^k\mathbb Z\simeq\mathbb Z_2\times\mathbb Z_5\) by the Chinese remainder theorem. Its units are precisely the elements invertible modulo ten; \(5\) is not a unit, so \(1/5\notin\mathbb Z_{10}\). Division by units does exist.}
\GSfield{Reconstruction}{The product-ring formulation and restriction of the phrase no division to nonunits are editorial precisifications.}
\GSfield{Fate}{This is classical positional arithmetic. The prime-base setting is treated by \citet[Ch.~3]{gouvea2020padic}; the product description is the direct inverse-limit form of the Chinese remainder theorem.}
\GSfield{Present status}{As of November 2026: established in the stated scope; this occurrence does not pose an unresolved problem.}
\GSfield{Assessment}{literal statement: correct with missing hypotheses; publication match: exact; verification: publisher metadata checked; development: annotated; confidence: high}
\end{GSitem}

\GStalk{168}{1969-03-31}{P. M. Bleher}{P000272-P000277}{Hörmander spectral asymptotics and their sharpness}
The heading on P000273 reads 31/III; the surrounding year is 1969, so the date can be resolved despite the missing normalized date in the supplied list. P000273--P000275 explain \citet{hormander1968spectral}: the local Weyl law, the sphere obstruction to an improved unsmoothed remainder, the first-order root of an elliptic operator, a short-time oscillatory parametrix and a Tauberian passage back to the spectral function. Questions about averaged higher terms are recorded, although the cited paper itself already discusses Riesz means. Several numerical coefficients in the notebook need correction.

\begin{GSitem}{1968-168.1}{E}{The spectral function of a positive elliptic realization}
\GSfield{Kind; source}{definition; \GSpage{P000273}}
\GSfield{The notes}{\((Pu,u)\geq c(u,u)\), \(\widehat P=\int\lambda\,dE_\lambda\), and \(e(x,y,\lambda)\in C^\infty(\Omega\times\Omega)\) are introduced.}
\GSfield{Statement}{Let \(P\) be a scalar elliptic differential operator of order \(m>0\) with smooth coefficients on a smooth manifold \(\Omega\), formally self-adjoint and bounded below by \(c>0\) on \(C_c^\infty\) for a fixed smooth positive density. For a self-adjoint realization \(\widehat P\geq c\), let \(E_\lambda=1_{(-\infty,\lambda]}(\widehat P)\). Its Schwartz kernel is the spectral function \(e(x,y,\lambda)\), smooth in the spatial variables locally in the interior for each finite \(\lambda\). Smoothness here does not mean continuity in the spectral parameter.}
\GSfield{Reconstruction}{The density, self-adjoint realization, scalar ellipticity and interior restriction are restored from the identified source.}
\GSfield{Fate}{The precise setup is already \citet[pp.~193--194]{hormander1968spectral}; received 7 May 1968, printed December 1968, before this exposition.}
\GSfield{Present status}{As of November 2026: established in the stated scope; this occurrence does not pose an unresolved problem.}
\GSfield{Assessment}{literal statement: correct with missing hypotheses; publication match: exact; verification: read the result; development: annotated; confidence: high}
\end{GSitem}
\begin{GSitem}{1968-168.2}{E}{The local Weyl remainder}
\GSfield{Kind; source}{theorem; \GSpage{P000273}}
\GSfield{The notes}{The diagonal expression is \(\lambda^{-n/m}e(x,x,\lambda)-(2\pi)^{-n}\int_{p^0(x,\xi)<1}d\xi=O(\lambda^{-1/m})\).}
\GSfield{Statement}{In the preceding setting, with real positive principal symbol \(p_m\), uniformly for \(x\) in each compact interior set, \[e(x,x,\lambda)=(2\pi)^{-n}\lambda^{n/m}\operatorname{vol}\{\xi:p_m(x,\xi)<1\}+O(\lambda^{(n-1)/m}).\] The cotangent measure is dual to the chosen base density.}
\GSfield{Reconstruction}{The formula is a local estimate; boundary terms and uniformity to a boundary are not asserted.}
\GSfield{Fate}{Exact match: \citet[Theorem~1.1]{hormander1968spectral}. The 1968 publication precedes Bleher\textquotesingle s talk.}
\GSfield{Present status}{As of November 2026: established in the stated scope; this occurrence does not pose an unresolved problem.}
\GSfield{Assessment}{literal statement: correct with missing hypotheses; publication match: exact; verification: read the result; development: annotated; confidence: high}
\end{GSitem}
\begin{GSitem}{1968-168.3}{E}{Off-diagonal leading term}
\GSfield{Kind; source}{claim; \GSpage{P000273}}
\GSfield{The notes}{The underlined phrase is \ru{Вне диагонали 0}.}
\GSfield{Statement}{For fixed distinct interior points, the leading \(\lambda^{n/m}\) diagonal term disappears. In the local coordinate domain of the oscillatory formula, on compact subsets away from the diagonal, \(e(x,y,\lambda)=O(\lambda^{(n-1)/m})\); in particular \(\lambda^{-n/m}e(x,y,\lambda)\to0\). The unnormalized kernel is generally not zero.}
\GSfield{Reconstruction}{Zero is read as zero leading normalized term. The notebook does not write this normalization in its short phrase.}
\GSfield{Fate}{The relevant off-diagonal parametrix and estimate are in \citet[Theorem~5.1 and (5.4), p.~215]{hormander1968spectral}.}
\GSfield{Present status}{As of November 2026: established in the stated scope; this occurrence does not pose an unresolved problem.}
\GSfield{Assessment}{literal statement: false as written; publication match: exact; verification: read the result; development: annotated; confidence: high}
\end{GSitem}
\begin{GSitem}{1968-168.4}{E}{Spherical harmonics and the spectral jump}
\GSfield{Kind; source}{construction; \GSpage{P000273}}
\GSfield{The notes}{For the Laplacian on \(S^n\), the jump is written \((C_{n+k}^{k}-C_{n+k}^{k-2})/V_n\), with large multiplicities attributed to rotations.}
\GSfield{Statement}{For the nonnegative Laplacian on the unit \(S^n\), \(n\geq1\), the eigenvalue \(\lambda_k=k(k+n-1)\) has multiplicity \[d_{n,k}=\binom{n+k}{k}-\binom{n+k-2}{k-2},\] with the second binomial zero for \(k<2\). Rotational invariance and integration over the sphere give the pointwise projector diagonal \(d_{n,k}/\operatorname{vol}(S^n)\). The upper index of the second binomial in the notebook must change.}
\GSfield{Reconstruction}{The eigenvalue and correct binomial coefficient are supplied; the transcriptions faithfully preserve the erroneous second index.}
\GSfield{Fate}{The example is analyzed in \citet[\S6]{hormander1968spectral}; the corrected formula also follows by subtracting the dimensions of homogeneous polynomials of degrees \(k\) and \(k-2\).}
\GSfield{Present status}{As of November 2026: established in the stated scope; this occurrence does not pose an unresolved problem.}
\GSfield{Assessment}{literal statement: false as written; publication match: exact; verification: read the result; development: annotated; confidence: high}
\end{GSitem}
\begin{GSitem}{1968-168.5}{E}{Why a universal smaller unsmoothed remainder is impossible}
\GSfield{Kind; source}{theorem; \GSpage{P000273}}
\GSfield{The notes}{The jump has size \(\lambda^{(n-1)/2}\); the notes say one cannot write a term \(C\lambda^{(n-1)/2}\).}
\GSfield{Statement}{For the sphere, \(d_{n,k}\sim2k^{n-1}/(n-1)!\) as \(k\to\infty\). A continuous approximation has the same limiting value on both sides of each jump, so at least one one-sided error is at least half that jump. Consequently an unsmoothed expansion with remainder \(o(\lambda^{(n-1)/2})\) is impossible, even after adding a fixed continuous term of that order. An \(O(\lambda^{(n-1)/2})\) bound remains valid.}
\GSfield{Reconstruction}{The distinction between a big-O bound and a fixed second asymptotic coefficient is essential; the Russian abbreviated sentence does not deny the valid big-O estimate.}
\GSfield{Fate}{Avakumovi\v{c}\textquotesingle s sharpness example is explicitly credited and proved in \citet[\S6]{hormander1968spectral}.}
\GSfield{Present status}{As of November 2026: established in the stated scope; this occurrence does not pose an unresolved problem.}
\GSfield{Assessment}{literal statement: correct with missing hypotheses; publication match: exact; verification: read the result; development: annotated; confidence: high}
\end{GSitem}
\begin{GSitem}{1968-168.6}{E}{Higher spectral terms after averaging}
\GSfield{Kind; source}{question; \GSpage{P000273-P000274}}
\GSfield{The notes}{The page asks for the next term of \(\lambda^{-1}\int_0^\lambda e(x,x,\nu)\,d\nu\), says this has not been done except for a square, and distinguishes obtaining subsequent terms from a formal asymptotic expansion.}
\GSfield{Statement}{For the positive elliptic \(P\) above, determine how many symbol-computable asymptotic terms are justified for the normalized Riesz means \[e^\alpha(x,x,\lambda)=\int_0^\lambda(1-\nu/\lambda)^\alpha\,d_\nu e(x,x,\nu),\quad\alpha>0.\] For \(\alpha=1\) integration by parts gives the notebook\textquotesingle s average. If the leading Weyl coefficient is \(a_0(x)\), its leading averaged coefficient is \(m a_0(x)/(m+n)\). A fixed smoothing order does not automatically yield an arbitrarily accurate expansion.}
\GSfield{Reconstruction}{The smoothing parameter and exact quantifier are editorial. The incorrect leading coefficient and the assertion of no existing work are not adopted.}
\GSfield{Fate}{Already addressed in \citet[Theorem~4.5 and pp.~215--216]{hormander1968spectral}, with remainder order \(O(\lambda^{(n-1-\alpha)/m})\) after the appropriate finite expansion. The antecedent was circulated at the November 1966 Belfer conference and translated in 1968: \citet{hormander1968riesz}.}
\GSfield{Present status}{As of November 2026: the general Riesz-mean framework was already available by 1968. Unsmoothed full expansions are obstructed by spectral jumps; no unresolved problem is inferred from this wording.}
\GSfield{Assessment}{literal statement: false as written; publication match: partial; verification: read the result; development: annotated; confidence: high; problem \hyperref[prob:GP-1968-40]{GP-1968-40}}
\end{GSitem}
\begin{GSitem}{1968-168.7}{F}{The trace-formula reading pointer}
\GSfield{Kind; source}{remark; \GSpage{P000274}}
\GSfield{The notes}{The notes recommend Dikii in UMN, dated 1957, and an earlier Gelfand work, dated 1954, on traces.}
\GSfield{Statement}{This is a bibliographical pointer to regularized spectral traces, not a recorded trace identity. No boundary conditions, potential, summation prescription or operator pair are specified; a unique theorem cannot be reconstructed from the pointer alone.}
\GSfield{Fate}{The identifiable UMN review is \citet{dikii1958traces}, published in 1958, pp.~111--143. The Gelfand pointer is too abbreviated to identify one publication without a conjectural attribution.}
\GSfield{Present status}{As of November 2026: the review is identified and its date corrected; no independent mathematical question is recoverable from the remaining pointer.}
\GSfield{Assessment}{literal statement: not determinable; publication match: partial; verification: publisher metadata checked; development: annotated; confidence: medium}
\end{GSitem}
\begin{GSitem}{1968-168.8}{E}{Reduction to a first-order pseudodifferential root}
\GSfield{Kind; source}{method; \GSpage{P000274}}
\GSfield{The notes}{\(A=\widehat P^{1/m}\), \(A=a_1+a_0+a_{-1}+\cdots\), and \(a_1=p^{1/m}\).}
\GSfield{Statement}{On a closed manifold, for a positive invertible classical scalar elliptic \(P\) of order \(m\), spectral functional calculus defines \(A=P^{1/m}\), a classical pseudodifferential operator of order one with principal symbol \(p_m^{1/m}\). Its full symbol has a homogeneous asymptotic expansion away from the zero section. On a general open manifold the corresponding assertion is used microlocally in the interior.}
\GSfield{Reconstruction}{Closedness and invertibility give a clean global version; a local statement suffices for the notes.}
\GSfield{Fate}{Already used in \citet[p.~195]{hormander1968spectral}, citing Seeley\textquotesingle s complex-power construction. It is not a new method invented in Bleher\textquotesingle s talk.}
\GSfield{Present status}{As of November 2026: established in the stated scope; this occurrence does not pose an unresolved problem.}
\GSfield{Assessment}{literal statement: correct with missing hypotheses; publication match: exact; verification: read the result; development: annotated; confidence: high}
\end{GSitem}
\begin{GSitem}{1968-168.9}{E}{Half-wave evolution and propagation of singularities}
\GSfield{Kind; source}{method; \GSpage{P000274}}
\GSfield{The notes}{\(G(t)=e^{itA}\), \((i\partial_t+A)G(t)=0\), \(G(0)=I\); the generalized Huygens principle is described as finite-speed propagation of singularities.}
\GSfield{Statement}{For self-adjoint \(A\), \(G(t)=e^{itA}\) is unitary and satisfies the displayed equation on its operator domain. When \(A\) is a real-principal-type first-order elliptic pseudodifferential operator, its short-time kernel is microlocally associated with the Hamiltonian flow, with sign determined by this evolution convention. Propagation of wavefront singularities is asserted, not compact support propagation for an arbitrary nonlocal operator.}
\GSfield{Reconstruction}{The qualification concerns singularities, which is exactly the Russian noun; finite speed of the entire support would be stronger and generally false.}
\GSfield{Fate}{The half-wave reduction and geometrical-optics construction are \citet[pp.~195--196 and \S3]{hormander1968spectral}.}
\GSfield{Present status}{As of November 2026: established in the stated scope; this occurrence does not pose an unresolved problem.}
\GSfield{Assessment}{literal statement: correct with missing hypotheses; publication match: exact; verification: read the result; development: annotated; confidence: high}
\end{GSitem}
\begin{GSitem}{1968-168.10}{E}{A short-time phase and transport amplitude}
\GSfield{Kind; source}{construction; \GSpage{P000274}}
\GSfield{The notes}{The ansatz has phase \(\psi(x,y,\xi)-t a_1(y,\xi)\), \(\psi=(x-y,\xi)+O(|x-y|^2|\xi|)\), an amplitude expansion and a smoothing remainder.}
\GSfield{Statement}{Locally choose a homogeneous phase solving the appropriate eikonal equation and construct a classical amplitude by transport equations so that the oscillatory integral solves the half-wave equation modulo a smooth kernel for sufficiently small \(|t|\). In the local normal form used here, \(\psi\) satisfies \(a_1(x,d_x\psi)=a_1(y,\xi)\); signs are adjusted if the Fourier exponential is reversed. The condition \(\psi=(x-y)\cdot\xi+O(|x-y|^2|\xi|)\) alone would not determine a valid parametrix.}
\GSfield{Reconstruction}{The omitted eikonal and transport conditions supply the substantive hypotheses of the ansatz.}
\GSfield{Fate}{Exact construction: \citet[\S3, especially the phase and transport equations]{hormander1968spectral}.}
\GSfield{Present status}{As of November 2026: established in the stated scope; this occurrence does not pose an unresolved problem.}
\GSfield{Assessment}{literal statement: correct with missing hypotheses; publication match: exact; verification: read the result; development: annotated; confidence: high}
\end{GSitem}
\begin{GSitem}{1968-168.11}{E}{Fourier cutoff, smoothed spectral density and the Tauberian step}
\GSfield{Kind; source}{method; \GSpage{P000274-P000275}}
\GSfield{The notes}{The smooth error \(E-Q\) is multiplied by \(\widehat\varrho(t)\); P000275 gives the convolution \(\int\varrho(\lambda-\mu)\,de(x,y,\mu)\) and defines \(R=\mathcal F^{-1}(q\widehat\varrho)\). A Marchenko Tauberian theorem is mentioned.}
\GSfield{Statement}{Choose \(\widehat\varrho\in C_c^\infty((-\varepsilon,\varepsilon))\). Fourier transformation of the localized smooth error is rapidly decreasing in the spectral variable. The remaining term is an oscillatory integral with amplitude \(R\). The diagonal spectral measure is positive; this positivity and a suitable cutoff permit a Tauberian estimate recovering the unsmoothed local counting bound. A displayed difference on P000275 is therefore a rapidly decreasing error, not an exact equality to zero.}
\GSfield{Reconstruction}{The order of the error, support of the cutoff and role of diagonal positivity are supplied explicitly.}
\GSfield{Fate}{The relevant estimates and Tauberian recovery are in \citet[\S4]{hormander1968spectral}; the notebook\textquotesingle s named Tauberian attribution is not needed to identify a separate theorem.}
\GSfield{Present status}{As of November 2026: established in the stated scope; this occurrence does not pose an unresolved problem.}
\GSfield{Assessment}{literal statement: correct with missing hypotheses; publication match: exact; verification: read the result; development: annotated; confidence: high}
\end{GSitem}

\GStalk{170}{1969-04-14}{I. M. Gelfand}{P000275-P000279}{Au-Yeung on a definite pair of Hermitian forms}
This short intervention is first on P000276, before Davis, although its catalogue number is larger. The February 1969 Proceedings of the AMS reference identifies Au-Yeung precisely. The field and the meaning of simultaneous reduction need restoration: the complex Hermitian theorem concerns congruence, and it does not assert that both original forms are positive. The secondary handwritten linear-algebra book reference remains unidentified.

\begin{GSitem}{1968-170.1}{E}{A pair without a common isotropic vector has a definite combination}
\GSfield{Kind; source}{theorem; \GSpage{P000276}}
\GSfield{The notes}{If two quadratic forms are not simultaneously zero, rotate their pair of values so that one form is positive definite.}
\GSfield{Statement}{Let \(A,B\) be Hermitian \(n\times n\) complex matrices. If \(x^*Ax=x^*Bx=0\) implies \(x=0\), there exist real \(a,b\) with \(aA+bB>0\). For real symmetric matrices with the hypothesis tested only on real vectors, an additional dimension qualification is needed; the unrestricted two-dimensional version is false.}
\GSfield{Reconstruction}{The complex field and exclusion of the zero vector are restored. A real counterexample is \(A=\operatorname{diag}(1,-1)\), \(B=\left(\begin{smallmatrix}0&1\\1&0\end{smallmatrix}\right)\): the forms have no common nonzero real zero, while every nonzero real combination is indefinite.}
\GSfield{Fate}{The named paper is \citet[pp.~545--548]{auyeung1969sphere}, published February 1969, before the talk. Its publisher contents and bibliographic identity were checked; the original article was not accessible. The exact complex assertion is independently readable in \citet[\S1, (1.1)]{guohighamtisseur2008arc}.}
\GSfield{Present status}{As of November 2026: established in the stated scope; this occurrence does not pose an unresolved problem.}
\GSfield{Assessment}{literal statement: correct with missing hypotheses; publication match: exact; verification: read the result; development: annotated; confidence: high}
\end{GSitem}
\begin{GSitem}{1968-170.2}{E}{The arc of normalized values and separation}
\GSfield{Kind; source}{method; \GSpage{P000276}}
\GSfield{The notes}{The value pair \((x^*Ax,x^*Bx)\) is normalized; the resulting direction runs through no more than a semicircle.}
\GSfield{Statement}{For a complex Hermitian definite pair, \(W(A+iB)=\{x^*(A+iB)x:\|x\|=1\}\) is compact and convex and excludes zero. Strict separation from zero puts its directions in a closed arc of length strictly less than \(\pi\). The separating real functional gives a positive definite combination.}
\GSfield{Reconstruction}{Compact numerical-range separation supplies the missing strictness; containment in an arbitrary closed semicircle alone would yield only semidefiniteness.}
\GSfield{Fate}{This is the geometric method associated with \citet{auyeung1969sphere}. The arc statement and its use in definite-pair algorithms are explicit in \citet[Lemma~2.1 and \S2]{guohighamtisseur2008arc}; that 2008 preprint is a later exposition and development, not the date of the theorem.}
\GSfield{Present status}{As of November 2026: established in the stated scope; this occurrence does not pose an unresolved problem.}
\GSfield{Assessment}{literal statement: correct with missing hypotheses; publication match: exact; verification: read the result; development: annotated; confidence: high}
\end{GSitem}
\begin{GSitem}{1968-170.3}{E}{Simultaneous diagonalization by congruence}
\GSfield{Kind; source}{construction; \GSpage{P000276}}
\GSfield{The notes}{The forms can be simultaneously reduced \ru{к сумме квадратов}.}
\GSfield{Statement}{Under the definite-pair hypothesis, there is an invertible complex \(S\) with \(S^*AS\) and \(S^*BS\) both real diagonal. Choose \(C=aA+bB>0\), use \(C^{-1/2}\) to make \(C\) the identity, then diagonalize the other Hermitian member by a unitary matrix. The diagonal coefficients of the original forms need not all be nonnegative.}
\GSfield{Reconstruction}{The displayed construction is an editorial verification of a standard consequence, not a claim that this proof was written in the notebook.}
\GSfield{Fate}{The precise congruence convention is also spelled out by \citet[p.~81]{auyeung1970diagonalization}. The elementary reduction explains the conclusion of the February 1969 paper without inferring positivity of both forms.}
\GSfield{Present status}{As of November 2026: established in the stated scope; this occurrence does not pose an unresolved problem.}
\GSfield{Assessment}{literal statement: correct with missing hypotheses; publication match: exact; verification: read the result; development: annotated; confidence: high}
\end{GSitem}

\GStalk{169}{1969-04-14}{Chandler Davis}{P000275-P000277}{The shell of a Hilbert-space operator}
P000276--P000277 give a compact account of \citet{davis1968shell}: an operator is replaced by a linear relation, its quadratic data are normalized to a subset of the three-dimensional ball, and norm, eigenvalues, M\"obius transformations and disk spectral sets become geometric features of that shell. The paper was published in 1968, before this lecture. P000278 already starts Buchstaber.

\begin{GSitem}{1968-169.1}{E}{Operators as linear relations}
\GSfield{Kind; source}{definition; \GSpage{P000276}}
\GSfield{The notes}{The graph is written \((y,x)\), \(y=Ax\); a more general relation is \(\alpha\subset H\oplus H\).}
\GSfield{Statement}{Let \(H\) be a complex Hilbert space. A linear relation is a linear subspace \(\alpha\subset H\oplus H\), written in output--input order \((y,x)\). A closed operator has closed graph \(\alpha=\{(Ax,x):x\in\operatorname{Dom}A\}\); a general closed relation may be multivalued and need not have dense domain.}
\GSfield{Reconstruction}{Linearity, and closedness where spectral conclusions require it, are added to the bare word set. An arbitrary subset is not the relevant notion of relation.}
\GSfield{Fate}{This is the framework of \citet[\S\S1--2]{davis1968shell}, already published in 1968.}
\GSfield{Present status}{As of November 2026: established in the stated scope; this occurrence does not pose an unresolved problem.}
\GSfield{Assessment}{literal statement: correct with missing hypotheses; publication match: exact; verification: read the result; development: annotated; confidence: high}
\end{GSitem}
\begin{GSitem}{1968-169.2}{E}{Norm ratios, numerical range and point spectrum}
\GSfield{Kind; source}{definition; \GSpage{P000276}}
\GSfield{The notes}{Three numerical data are listed: norm ratios, \(w(A)\), and \(\sigma_p(A)\). The eigenvalue quantifier has no \(x\ne0\).}
\GSfield{Statement}{For a bounded operator \(A\), \[\begin{aligned}\|A\|^2&=\sup_{x\ne0}\frac{\|Ax\|^2}{\|x\|^2},\\ W(A)&=\{\langle Ax,x\rangle/\|x\|^2:x\ne0\},\\ \sigma_p(A)&=\{\lambda:\exists x\ne0,\ Ax=\lambda x\}.\end{aligned}\] The Hilbert inner product is taken linear in its first argument so that it agrees with the notebook\textquotesingle s \(x^*Ax\).}
\GSfield{Reconstruction}{The nonzero eigenvector condition is essential; without it the displayed point spectrum would be all of the scalar field.}
\GSfield{Fate}{These classical invariants are the starting data in \citet[\S1]{davis1968shell}, not new definitions originating at this seminar.}
\GSfield{Present status}{As of November 2026: established in the stated scope; this occurrence does not pose an unresolved problem.}
\GSfield{Assessment}{literal statement: correct with missing hypotheses; publication match: exact; verification: read the result; development: annotated; confidence: high}
\end{GSitem}
\begin{GSitem}{1968-169.3}{E}{The indicatrix cone}
\GSfield{Kind; source}{construction; \GSpage{P000276}}
\GSfield{The notes}{\(\{(\|Ax\|^2,x^*Ax,\|x\|^2):x\ne0\}\subset\mathbb R\times\mathbb C\times\mathbb R\) is called a cone.}
\GSfield{Statement}{For a nonzero linear relation \(\alpha\), its indicatrix is the set of triples \((\|y\|^2,\langle y,x\rangle,\|x\|^2)\) over nonzero \((y,x)\in\alpha\). Scaling the pair scales the triple by a positive real number. Every triple satisfies \(|\langle y,x\rangle|^2\leq\|y\|^2\|x\|^2\).}
\GSfield{Reconstruction}{The relation version makes the operator formula compatible with the subsequent projective action.}
\GSfield{Fate}{Exact construction in \citet[Definition~1.1 and the preceding normalization]{davis1968shell}.}
\GSfield{Present status}{As of November 2026: established in the stated scope; this occurrence does not pose an unresolved problem.}
\GSfield{Assessment}{literal statement: correct; publication match: exact; verification: read the result; development: annotated; confidence: high}
\end{GSitem}
\begin{GSitem}{1968-169.4}{E}{The shell in the closed unit ball}
\GSfield{Kind; source}{construction; \GSpage{P000277}}
\GSfield{The notes}{The convenient section is \((\xi,h)=(2x^*Ax,(\|Ax\|^2-\|x\|^2))/(\|Ax\|^2+\|x\|^2)\), with \(|\xi|^2+h^2\leq1\).}
\GSfield{Statement}{Normalize each nonzero pair \((y,x)\in\alpha\) by \(t=\|y\|^2+\|x\|^2\): \[s(\alpha)=\{(2\langle y,x\rangle/t,(\|y\|^2-\|x\|^2)/t)\}.\] Cauchy--Schwarz gives \(s(\alpha)\subset\{(\xi,h):|\xi|^2+h^2\leq1\}\). This set need not be closed; closure must be written explicitly when used.}
\GSfield{Reconstruction}{Set braces, quantification and the distinction between the shell and its closure repair the abbreviated notation.}
\GSfield{Fate}{This is \citet[Definition~1.1]{davis1968shell}, with equivalent stereographic coordinates.}
\GSfield{Present status}{As of November 2026: established in the stated scope; this occurrence does not pose an unresolved problem.}
\GSfield{Assessment}{literal statement: correct; publication match: exact; verification: read the result; development: annotated; confidence: high}
\end{GSitem}
\begin{GSitem}{1968-169.5}{E}{Boundary points recover eigenvalues}
\GSfield{Kind; source}{theorem; \GSpage{P000277}}
\GSfield{The notes}{An eigenvector is obtained when the shell point belongs to the sphere; the eigenvalue is related by stereographic projection.}
\GSfield{Statement}{For bounded \(A\), put \(\tau(\lambda)=(2\lambda/(1+|\lambda|^2),(|\lambda|^2-1)/(1+|\lambda|^2))\). Then \(s(A)\cap S^2=\tau(\sigma_p(A))\). Equality in Cauchy--Schwarz forces \(Ax=\lambda x\), and conversely. The north pole represents the infinite value for general multivalued relations, not an eigenvalue of a bounded operator.}
\GSfield{Reconstruction}{The bounded-operator statement avoids conflating eigenvalues with approximate eigenvalues or the extra point at infinity.}
\GSfield{Fate}{Exact match with \citet[Theorem~2.2]{davis1968shell}; Theorem~2.3 separately treats the approximate point spectrum using closure.}
\GSfield{Present status}{As of November 2026: established in the stated scope; this occurrence does not pose an unresolved problem.}
\GSfield{Assessment}{literal statement: correct with missing hypotheses; publication match: exact; verification: read the result; development: annotated; confidence: high}
\end{GSitem}
\begin{GSitem}{1968-169.6}{E}{Height bounds encode an upper norm bound}
\GSfield{Kind; source}{claim; \GSpage{P000277}}
\GSfield{The notes}{If the shell does not go above \((a^2-1)/(a^2+1)\), then \(a\) is the norm.}
\GSfield{Statement}{For \(a\geq0\) and bounded \(A\), \[\|A\|\leq a\quad\Longleftrightarrow\quad h\leq\frac{a^2-1}{a^2+1}\text{ for every }(\xi,h)\in s(A).\] Equality \(\|A\|=a\) holds when \(a\) is the least such bound; the supremum need not be attained.}
\GSfield{Reconstruction}{The notebook\textquotesingle s asserted equality is replaced by the correct inequality and a separate minimality condition.}
\GSfield{Fate}{The shell formulation of the norm is in \citet[\S1 and Definition~7.1]{davis1968shell}. It also follows by solving the normalized ratio for \(\|Ax\|^2/\|x\|^2\).}
\GSfield{Present status}{As of November 2026: established in the stated scope; this occurrence does not pose an unresolved problem.}
\GSfield{Assessment}{literal statement: false as written; publication match: exact; verification: read the result; development: annotated; confidence: high}
\end{GSitem}
\begin{GSitem}{1968-169.7}{E}{Fractional linear transformations of relations}
\GSfield{Kind; source}{construction; \GSpage{P000277}}
\GSfield{The notes}{\(\mu(\lambda)=(a\lambda+b)/(c\lambda+d)\), \(ad-bc\ne0\), acts by \((y,x)\mapsto(ay+bx,cy+dx)\).}
\GSfield{Statement}{For complex \(a,b,c,d\) with \(ad-bc\ne0\), this invertible bounded block transformation sends a closed linear relation to a closed linear relation \(\mu(\alpha)\). If \(A\) is bounded and \(cA+dI\) is boundedly invertible, the transformed graph is the graph of \((aA+bI)(cA+dI)^{-1}\).}
\GSfield{Reconstruction}{The inverse condition and identity operators are made explicit; the relation construction continues to exist when an operator formula does not.}
\GSfield{Fate}{Exact relation calculus in \citet[\S5]{davis1968shell}.}
\GSfield{Present status}{As of November 2026: established in the stated scope; this occurrence does not pose an unresolved problem.}
\GSfield{Assessment}{literal statement: correct with missing hypotheses; publication match: exact; verification: read the result; development: annotated; confidence: high}
\end{GSitem}
\begin{GSitem}{1968-169.8}{E}{Möbius equivariance of the shell}
\GSfield{Kind; source}{theorem; \GSpage{P000277}}
\GSfield{The notes}{Theorem 1 states \(s(\mu(\alpha))=\varrho_\mu(s(\alpha))\).}
\GSfield{Statement}{Act on the positive Hermitian Gram matrix \(G=\left(\begin{smallmatrix}\|y\|^2&\langle y,x\rangle\\\langle x,y\rangle&\|x\|^2\end{smallmatrix}\right)\) by \(G\mapsto M G M^*\), \(M=\left(\begin{smallmatrix}a&b\\c&d\end{smallmatrix}\right)\), then divide by its trace. The resulting projective transformation \(\varrho_\mu\) of the closed ball satisfies the notebook\textquotesingle s identity and restricts on the boundary to stereographic conjugation of \(\mu\).}
\GSfield{Reconstruction}{The Gram-matrix construction defines the otherwise unnamed ball transformation; it is not an arbitrary analytic map of a three-dimensional ball.}
\GSfield{Fate}{This is \citet[Theorem~5.1]{davis1968shell}.}
\GSfield{Present status}{As of November 2026: established in the stated scope; this occurrence does not pose an unresolved problem.}
\GSfield{Assessment}{literal statement: correct with missing hypotheses; publication match: exact; verification: read the result; development: annotated; confidence: high}
\end{GSitem}
\begin{GSitem}{1968-169.9}{E}{Spectral sets and the rational functional calculus}
\GSfield{Kind; source}{definition; \GSpage{P000277}}
\GSfield{The notes}{\(X\) is spectral when every rational \(f\) with \(f(X)\subset D\) satisfies \(\|f(A)\|\leq1\).}
\GSfield{Statement}{In the standard bounded-operator formulation, a closed set \(X\subset\mathbb C\) containing \(\sigma(A)\) is a spectral set if \(\|f(A)\|\leq\sup_X|f|\) for every rational function with poles off \(X\). The unit-disk formulation is equivalent by rescaling. For general relations and subsets of the Riemann sphere, use Davis\textquotesingle s Definition~7.2, including his convention for rational transforms.}
\GSfield{Reconstruction}{The spectrum-containment and pole restrictions prevent an undefined rational calculus from being mistaken for a condition.}
\GSfield{Fate}{\citet[Definitions~7.1--7.2]{davis1968shell} makes the quantifiers precise and explicitly traces the underlying spectral-set theory to earlier work.}
\GSfield{Present status}{As of November 2026: established in the stated scope; this occurrence does not pose an unresolved problem.}
\GSfield{Assessment}{literal statement: correct with missing hypotheses; publication match: exact; verification: read the result; development: annotated; confidence: high}
\end{GSitem}
\begin{GSitem}{1968-169.10}{E}{Disk spectral sets are detected by spherical caps}
\GSfield{Kind; source}{theorem; \GSpage{P000277}}
\GSfield{The notes}{Theorem 2 says that, when \(\tau(X)\) is a spherical segment, \(X\) is spectral exactly when \(s(A)\subset\operatorname{conv}(\tau(X))\).}
\GSfield{Statement}{For a closed linear relation in Davis\textquotesingle s sense and a closed generalized disk \(X\) in the Riemann sphere (a finite disk, a half-plane with infinity, or the exterior of an open disk with infinity), \[X\text{ is a spectral set}\quad\Longleftrightarrow\quad s(\alpha)\subset\operatorname{conv}(\tau(X)).\] For bounded \(A\) and a finite disk \(X=\{z:|z-z_0|\leq r\}\), this is equivalently \(\|A-z_0I\|\leq r\). The spherical segment is a two-dimensional cap, not an arc.}
\GSfield{Reconstruction}{The disk restriction is essential. The notebook does not claim this criterion for arbitrary closed spectral sets.}
\GSfield{Fate}{Exact match: \citet[Theorem~7.2; preceding cases (i)--(iii)]{davis1968shell}.}
\GSfield{Present status}{As of November 2026: established in the stated scope; this occurrence does not pose an unresolved problem.}
\GSfield{Assessment}{literal statement: correct with missing hypotheses; publication match: exact; verification: read the result; development: annotated; confidence: high}
\end{GSitem}

\GStalk{171}{1969-04-14}{V. M. Buchstaber}{P000277-P000281}{Finiteness of smooth homotopy tori and surgery}
P000278--P000280 outline the reduction of the classification of homotopy tori to stable normal data, a finite Hurewicz fibre and surgery. The notebook conflates \(K_1\) with the Whitehead group and suppresses dimension and normal-data hypotheses. A closely matching report by Buchstaber and Mishchenko at the Moscow Mathematical Society on 22 April 1969 explicitly attributes the decisive algebraic methods to Wall and Shaneson; this is evidence of exposition of a current result, not an orphan theorem of Buchstaber. The concluding surgery step is an outline, not a proof supplied by the notes.

\begin{GSitem}{1968-171.1}{E}{Finitely many smooth manifolds of torus homotopy type}
\GSfield{Kind; source}{theorem; \GSpage{P000278}}
\GSfield{The notes}{There are finitely many smooth manifolds of homotopy type \(T^n\), up to diffeomorphism.}
\GSfield{Statement}{For each fixed \(n\geq5\), only finitely many diffeomorphism classes of closed connected smooth \(n\)-manifolds are homotopy equivalent to \(T^n=(S^1)^n\). This is a finiteness statement, not uniqueness in the smooth or PL category.}
\GSfield{Reconstruction}{Closedness and the high-dimensional range are restored; the notebook is not evidence for a four-dimensional classification.}
\GSfield{Fate}{\citet[Theorem~A]{hsiangshaneson1969fake}, communicated 6 January and published March 1969, gives PL finiteness. The passage to smooth finiteness is an additional deduction: on a fixed finite PL complex the smoothing choices and differences are controlled by finitely many cochains with finite sphere-smoothing groups, and concordant smoothings are diffeomorphic \citep[Theorems~1.1, 3.1, 4.2, 5.1]{hirsch1963obstruction}. The finiteness input is supplied by Kervaire--Milnor and the low-dimensional Cerf theorem \citep{kervairemilnor1963spheres,cerf1968gamma4}. Thus E uses these prior ingredients, not the PL theorem alone. The exact smooth statement is reported in \citet[p.~224, meeting of 22 April]{buchstabermishchenko1969tori}; that later meeting is confirmation, not evidence of prior publication. Subsequent topological uniqueness is \citet[main theorem]{hsiangwall1969tori}, received 14 May 1969. These dates distinguish prior results, the seminar exposition and later topological strengthening.}
\GSfield{Present status}{As of November 2026: established in the stated scope; this occurrence does not pose an unresolved problem.}
\GSfield{Assessment}{literal statement: correct with missing hypotheses; publication match: exact; verification: read the result; development: annotated; confidence: high; problem \hyperref[prob:GP-1968-42]{GP-1968-42}}
\end{GSitem}
\begin{GSitem}{1968-171.2}{E}{Simply connected finiteness with fixed characteristic data}
\GSfield{Kind; source}{theorem; \GSpage{P000278}}
\GSfield{The notes}{Novikov: homotopy type and Pontryagin classes determine a simply connected manifold up to finitely many possibilities.}
\GSfield{Statement}{Fix a closed simply connected smooth manifold \(M^n\), \(n\geq5\). Among closed smooth \(N^n\) with an orientation-preserving homotopy equivalence \(f:N\to M\) satisfying \(p_i(N)=f^*p_i(M)\) in rational cohomology, there are only finitely many diffeomorphism types. Neither arbitrary fundamental groups nor dimension four are covered by this assertion.}
\GSfield{Reconstruction}{The rational characteristic classes are compared under a specified homotopy equivalence. This does not say that the classes alone determine the homotopy type.}
\GSfield{Fate}{\citet[Introduction and classification theorem, pp.~366--367]{novikov1964homotopy}, received 22 March 1963 and published 1964, develops the tangential homotopy classification; its finiteness consequence with rational Pontryagin data is also stated in the author\textquotesingle s research summary \citep[Topology]{novikov2014bibliography}.}
\GSfield{Present status}{As of November 2026: established in the stated scope; this occurrence does not pose an unresolved problem.}
\GSfield{Assessment}{literal statement: correct with missing hypotheses; publication match: exact; verification: read the result; development: annotated; confidence: high}
\end{GSitem}
\begin{GSitem}{1968-171.3}{E}{Homotopy spheres as the simply connected model}
\GSfield{Kind; source}{remark; \GSpage{P000278}}
\GSfield{The notes}{Milnor\textquotesingle s sphere case is mentioned, with dimensions three and four marked unknown.}
\GSfield{Statement}{For \(n\geq5\), the oriented diffeomorphism classes of smooth homotopy \(n\)-spheres form a finite abelian group \(\Theta_n\) under connected sum, identified with the corresponding oriented h-cobordism classes. The notebook\textquotesingle s historical uncertainty in low dimensions is not extended into an assertion about their present status.}
\GSfield{Reconstruction}{Only the high-dimensional theorem needed in this talk is asserted; h-cobordism and diffeomorphism are separated before using the dimension hypothesis.}
\GSfield{Fate}{\citet[Theorems~1.1--1.2]{kervairemilnor1963spheres}, received 19 April 1962 and published in May 1963, is the precise antecedent. The earlier companion occurrence 1966-92.2 treats the seven-dimensional example; this is a related exposition, not a new unresolved sphere problem.}
\GSfield{Present status}{As of November 2026: established in the stated scope; this occurrence does not pose an unresolved problem.}
\GSfield{Assessment}{literal statement: correct with missing hypotheses; publication match: exact; verification: read the result; development: annotated; confidence: high}
\end{GSitem}
\begin{GSitem}{1968-171.4}{E}{h-cobordism}
\GSfield{Kind; source}{definition; \GSpage{P000278}}
\GSfield{The notes}{Both \(M_i^n\subset W^{n+1}\) are deformation retracts and \(\partial W=M_1\cup M_2\).}
\GSfield{Statement}{A smooth h-cobordism \((W;M_1,M_2)\) is a compact smooth cobordism whose two boundary inclusions are homotopy equivalences. For oriented manifolds the boundary convention is \(\partial W=M_1\sqcup(-M_2)\); collars make the inclusions cofibrations, so the homotopy-equivalence formulation is compatible with deformation retractions.}
\GSfield{Reconstruction}{Compactness, collars and orientation signs are made explicit.}
\GSfield{Fate}{This is the established framework in \citet[\S\S9,11]{milnor1966torsion}; the English survey appeared in 1966 and the checked Russian translation in 1967.}
\GSfield{Present status}{As of November 2026: established in the stated scope; this occurrence does not pose an unresolved problem.}
\GSfield{Assessment}{literal statement: correct with missing hypotheses; publication match: exact; verification: read the result; development: annotated; confidence: high}
\end{GSitem}
\begin{GSitem}{1968-171.5}{E}{The s-cobordism classification}
\GSfield{Kind; source}{theorem; \GSpage{P000278}}
\GSfield{The notes}{Under Smale\textquotesingle s name, an h-cobordism is determined by one boundary and an invariant \(\alpha\in\operatorname{Wh}(\pi_1M_1)\).}
\GSfield{Statement}{If \(M^n\) is closed and connected, \(n\geq5\), smooth h-cobordisms on \(M\), modulo diffeomorphism relative to the incoming boundary, are classified by their Whitehead torsion in \(\operatorname{Wh}(\pi_1M)\). In particular, torsion zero is equivalent to a product cobordism relative to that boundary.}
\GSfield{Reconstruction}{Relative equivalence, dimension and the distinction from the simply connected h-cobordism theorem are restored.}
\GSfield{Fate}{\citet[\S11, existence and uniqueness theorems]{milnor1966torsion} gives the classification. The non-simply-connected result is the s-cobordism theorem (Barden--Mazur--Stallings); attributing the full statement solely to Smale suppresses the essential extension.}
\GSfield{Present status}{As of November 2026: established in the stated scope; this occurrence does not pose an unresolved problem.}
\GSfield{Assessment}{literal statement: correct with missing hypotheses; publication match: exact; verification: read the result; development: annotated; confidence: high}
\end{GSitem}
\begin{GSitem}{1968-171.6}{E}{Stable general linear groups, K1 and Whitehead groups}
\GSfield{Kind; source}{definition; \GSpage{P000278-P000279}}
\GSfield{The notes}{The inductive general linear group is abelianized; \(\operatorname{Wh}\) is obtained by a further quotient. A noncommutative determinant is mentioned.}
\GSfield{Statement}{For a unital ring \(R\), let \(GL(R)=\varinjlim_m GL_m(R)\) under \(A\mapsto\operatorname{diag}(A,1)\), and let \(E(R)\) be the stable elementary subgroup. Then \(E(R)=[GL(R),GL(R)]\), \(K_1(R)=GL(R)/E(R)\), and \(\operatorname{Wh}(\pi)=K_1(\mathbb Z[\pi])/\langle[\pm g]:g\in\pi\rangle\). Dieudonn\'e\textquotesingle s determinant over a division ring is an analogy, not a literal determinant on an arbitrary group ring.}
\GSfield{Reconstruction}{The ring-level and group-level Whitehead terminology is disambiguated, which is necessary for the next claim.}
\GSfield{Fate}{\citet[\S1, Whitehead lemma; \S6]{milnor1966torsion} states these constructions before the talk.}
\GSfield{Present status}{As of November 2026: established in the stated scope; this occurrence does not pose an unresolved problem.}
\GSfield{Assessment}{literal statement: correct with missing hypotheses; publication match: exact; verification: read the result; development: annotated; confidence: high}
\end{GSitem}
\begin{GSitem}{1968-171.7}{E}{The Whitehead group of a free abelian group vanishes}
\GSfield{Kind; source}{claim; \GSpage{P000279}}
\GSfield{The notes}{The scan explicitly writes \(K_1(\mathbb Z+\cdots+\mathbb Z)=0\), with a gloss referring to the group ring, and attributes it to Bass.}
\GSfield{Statement}{For \(\pi=\mathbb Z^n\), \(R=\mathbb Z[\pi]=\mathbb Z[t_1^{\pm1},\ldots,t_n^{\pm1}]\), one has \(K_1(R)\simeq R^\times\simeq\{\pm1\}\times\mathbb Z^n\). Consequently \(\operatorname{Wh}(\mathbb Z^n)=0\), but \(K_1(R)\) is not zero.}
\GSfield{Reconstruction}{The zero group is the Whitehead quotient. The plus signs denote a direct sum of infinite cyclic groups, not a direct sum of coefficient rings.}
\GSfield{Fate}{Exact corrected theorem: \citet[p.~63, corollaries]{basshellerswan1964whitehead}, received 15 June 1963 and published 1964. Quotienting the surviving trivial units is the missing step.}
\GSfield{Present status}{As of November 2026: established in the stated scope; this occurrence does not pose an unresolved problem.}
\GSfield{Assessment}{literal statement: false as written; publication match: exact; verification: read the result; development: annotated; confidence: high}
\end{GSitem}
\begin{GSitem}{1968-171.8}{E}{Invertible polynomial and Laurent-polynomial matrices}
\GSfield{Kind; source}{claim; \GSpage{P000279}}
\GSfield{The notes}{For polynomial matrices, invertibility is said to be equivalent to determinant \(\pm1\); elementary reduction after increasing size is discussed.}
\GSfield{Statement}{Over the commutative polynomial ring \(\mathbb Z[t_1,\ldots,t_n]\), a square matrix is invertible exactly when its determinant is \(\pm1\). Over the group ring \(\mathbb Z[\mathbb Z^n]\), the units instead are \(\pm t_1^{a_1}\cdots t_n^{a_n}\), \(a_i\in\mathbb Z\). Stable elementary reduction detects the zero class in \(K_1\), and does not mean every invertible matrix already has determinant one.}
\GSfield{Reconstruction}{The literal polynomial assertion is correct; applying it unchanged to the preceding Laurent group ring is not.}
\GSfield{Fate}{The unit distinction and stable \(K_1\) calculation are precisely the setting of \citet[pp.~61--63]{basshellerswan1964whitehead}. The determinant criterion follows from the adjugate identity.}
\GSfield{Present status}{As of November 2026: established in the stated scope; this occurrence does not pose an unresolved problem.}
\GSfield{Assessment}{literal statement: correct with missing hypotheses; publication match: exact; verification: read the result; development: annotated; confidence: high}
\end{GSitem}
\begin{GSitem}{1968-171.9}{E}{Stabilization at sufficiently large matrix size}
\GSfield{Kind; source}{theorem; \GSpage{P000279}}
\GSfield{The notes}{Bass: it suffices to restrict to matrices of sufficiently high order.}
\GSfield{Statement}{For a commutative Noetherian ring \(R\) of finite Krull dimension, there exists a uniform \(m_0(R)\) such that every element of \(K_1(R)\) has a representative in \(GL_m(R)\) for every \(m\geq m_0(R)\). In particular this applies to \(\mathbb Z[t_1^{\pm1},\ldots,t_n^{\pm1}]\). This is a surjective stabilization statement; it does not by itself assert that every elementary equivalence is detected without further stabilization.}
\GSfield{Reconstruction}{The ring hypotheses and uniform bound are restored; the weaker representability claim is all that the terse sentence needs. No unstated optimal bound or injectivity theorem is assigned to the speaker.}
\GSfield{Fate}{The precise pre-existing result is \citet[Theorem~11.1 and Proposition~11.2(b)]{bass1964stable}. Immediately after that proposition Bass separately formulates injective stabilization as a conjecture; the two assertions must not be retroactively conflated.}
\GSfield{Present status}{As of November 2026: established in the stated scope; this occurrence does not pose an unresolved problem.}
\GSfield{Assessment}{literal statement: correct with missing hypotheses; publication match: exact; verification: read the result; development: annotated; confidence: high}
\end{GSitem}
\begin{GSitem}{1968-171.10}{E}{Stable normal data under a homotopy equivalence}
\GSfield{Kind; source}{definition; \GSpage{P000279}}
\GSfield{The notes}{An embedding \(M^n\subset\mathbb R^N\) gives \(\nu_M\), independent of the embedding stably; equivalence includes \(f^*\nu_{M_2}=\nu_{M_1}\).}
\GSfield{Statement}{For a smooth embedding, \(TM\oplus\nu_M\cong\varepsilon^N\); thus the stable class \([\nu_M]-[\varepsilon^{N-n}]=-([TM]-[\varepsilon^n])\) is intrinsic. A tangential homotopy equivalence consists of a homotopy equivalence together with a stable identification of the pulled-back tangent, or equivalently normal, bundles.}
\GSfield{Reconstruction}{Equality means stable bundle isomorphism over a specified map; it is stronger than equality of rational Pontryagin classes.}
\GSfield{Fate}{These are the normal data used by \citet[Introduction and Ch.~I]{novikov1964homotopy}.}
\GSfield{Present status}{As of November 2026: established in the stated scope; this occurrence does not pose an unresolved problem.}
\GSfield{Assessment}{literal statement: correct with missing hypotheses; publication match: exact; verification: read the result; development: annotated; confidence: high}
\end{GSitem}
\begin{GSitem}{1968-171.11}{E}{Pontryagin classes and parallelizability of homotopy tori}
\GSfield{Kind; source}{claim; \GSpage{P000279}}
\GSfield{The notes}{Adams\textquotesingle s J-functor and Novikov\textquotesingle s torus invariance are invoked; \(\nu_M=0\) is underlined.}
\GSfield{Statement}{For a smooth homotopy torus in the high-dimensional range, the rational Pontryagin classes vanish, and the classification of fake tori supplies stable triviality of the normal bundle. These are distinct statements: vanishing of rational characteristic classes alone does not prove integral stable triviality for a general manifold.}
\GSfield{Reconstruction}{The conclusion is separated from the abbreviated argument. No proof of stable triviality from rational classes alone is accepted.}
\GSfield{Fate}{\citet[\S1]{novikov1966abelian} provides the relevant free-abelian and Pontryagin-class framework. The modern classification records parallelizability explicitly in \citet[p.~2212, Remark~1]{hambleton2006tori}. The notebook\textquotesingle s unnamed Adams reference is insufficient to identify a particular J-theorem with all its hypotheses.}
\GSfield{Present status}{As of November 2026: established in the stated scope; this occurrence does not pose an unresolved problem.}
\GSfield{Assessment}{literal statement: correct with missing hypotheses; publication match: partial; verification: read the result; development: annotated; confidence: medium}
\end{GSitem}
\begin{GSitem}{1968-171.12}{E}{Infinitely many manifolds with the homotopy type of S2 times S4}
\GSfield{Kind; source}{claim; \GSpage{P000279}}
\GSfield{The notes}{\ru{на \(S^2\times S^4\) бесконечно много гладких структур}.}
\GSfield{Statement}{There are infinitely many pairwise nondiffeomorphic closed smooth manifolds homotopy equivalent to \(S^2\times S^4\), with different rational Pontryagin data. This must not be read as infinitely many smoothings of one fixed topological \(S^2\times S^4\). In dimension six that fixed topological-manifold smoothing problem has only finitely many possibilities.}
\GSfield{Reconstruction}{On is interpreted as of the homotopy type of. The corrected meaning fits the surrounding comparison of characteristic classes.}
\GSfield{Fate}{\citet[Introduction, consequences 1--2, pp.~208--209]{novikov1966abelian} explicitly distinguishes finiteness for a fixed simply connected topological manifold from infinitude within an appropriate homotopy type.}
\GSfield{Present status}{As of November 2026: established in the stated scope; this occurrence does not pose an unresolved problem.}
\GSfield{Assessment}{literal statement: false as written; publication match: exact; verification: read the result; development: annotated; confidence: high}
\end{GSitem}
\begin{GSitem}{1968-171.13}{E}{Pontryagin--Thom collapse and the degree-one fibre}
\GSfield{Kind; source}{construction; \GSpage{P000279-P000280}}
\GSfield{The notes}{A tubular neighbourhood is collapsed to obtain \(S^N\to T_{\nu_M}\), whose Hurewicz image is \(+1\).}
\GSfield{Statement}{An oriented embedding and tubular neighbourhood give a collapse \(c_M:S^N\to\operatorname{Th}(\nu_M)\). The Thom isomorphism sends the image of \([S^N]\) to \([M]\), so the Hurewicz image is the chosen generator of \(\widetilde H_N(\operatorname{Th}(\nu_M);\mathbb Z)\cong\mathbb Z\). Homotopy-equivalent bases yield equivalent Thom spaces only when the relevant normal bundles have been identified stably over that equivalence.}
\GSfield{Reconstruction}{The normal-bundle identification in the preceding page is retained; homotopy type of the base by itself is insufficient.}
\GSfield{Fate}{This is the explicit method of \citet[pp.~366--367; Ch.~I, \S3]{novikov1964homotopy}.}
\GSfield{Present status}{As of November 2026: established in the stated scope; this occurrence does not pose an unresolved problem.}
\GSfield{Assessment}{literal statement: correct with missing hypotheses; publication match: exact; verification: read the result; development: annotated; confidence: high}
\end{GSitem}
\begin{GSitem}{1968-171.14}{E}{Stable torus splitting and finiteness of H inverse of one}
\GSfield{Kind; source}{claim; \GSpage{P000280}}
\GSfield{The notes}{\(T_\nu(T^n)=\Sigma^{N-n}T^n=S^N\vee\bigvee_{i<N}S^i\); \(H^{-1}(1)\) is called a finite group and Serre is cited.}
\GSfield{Statement}{For the trivial rank \(r=N-n\geq1\) bundle, \[\operatorname{Th}(\varepsilon^r_{T^n})\simeq\Sigma^r(T^n)_+\simeq\bigvee_{j=0}^n\bigvee^{\binom nj}S^{r+j}.\] Choose \(N>2n+1\) to lie in the stable range. The Hurewicz map \(\pi_N(\operatorname{Th}(\varepsilon^{N-n}_{T^n}))\to\mathbb Z\) is projection onto the top-sphere summand, with finite kernel because the other summands are positive stable stems. Its fibre over one is a finite torsor for that kernel, not a subgroup.}
\GSfield{Reconstruction}{A disjoint basepoint, wedge multiplicities, stable-range choice and the affine nature of the fibre are supplied. The statement is restricted to the torus case, not all possible Thom spaces.}
\GSfield{Fate}{The finite-fibre mechanism is already used in \citet[p.~366]{novikov1964homotopy}; Novikov invokes Serre\textquotesingle s homotopy-finiteness theorem. The classes-of-groups framework and sphere applications are developed in \citet[Chs.~I--IV]{serre1953classes}.}
\GSfield{Present status}{As of November 2026: established in the stated scope; this occurrence does not pose an unresolved problem.}
\GSfield{Assessment}{literal statement: false as written; publication match: exact; verification: read the result; development: annotated; confidence: high}
\end{GSitem}
\begin{GSitem}{1968-171.15}{E}{From normal cobordism to h-cobordism by surgery}
\GSfield{Kind; source}{method; \GSpage{P000280}}
\GSfield{The notes}{A homotopy between collapse maps produces a film; one must arrange retractions, classify films of finite index and use obstruction theory for modifications.}
\GSfield{Statement}{A transverse homotopy of collapse maps with fixed normal data produces a normal cobordism. It need not be an h-cobordism. In dimensions at least five, surgery seeks to make the boundary inclusions homotopy equivalences; middle-dimensional surgery obstructions and the action on the structure set must be analyzed. Once an h-cobordism is obtained for fundamental group \(\mathbb Z^n\), its Whitehead torsion vanishes and the s-cobordism theorem makes it a product.}
\GSfield{Reconstruction}{The missing surgery step is identified, not filled in as though the notebook contained a complete proof.}
\GSfield{Fate}{The simply connected model is \citet[Ch.~I]{novikov1964homotopy}; the torus classification requires the non-simply-connected methods in \citet{hsiangshaneson1969fake}. The near-contemporary Buchstaber--Mishchenko abstract explicitly names Wall and Shaneson\textquotesingle s algebraic methods \citep[p.~224]{buchstabermishchenko1969tori}.}
\GSfield{Present status}{As of November 2026: established in the stated scope; this occurrence does not pose an unresolved problem.}
\GSfield{Assessment}{literal statement: correct with missing hypotheses; publication match: partial; verification: read the result; development: annotated; confidence: high; problem \hyperref[prob:GP-1968-42]{GP-1968-42}}
\end{GSitem}

\GStalk{172}{1969-04-21}{A. A. Kirillov}{P000279-P000282}{A noncompact Borel--Weil--Bott theorem}
The single mathematical page P000281 describes homogeneous line bundles on \(G/H\), harmonic \(L^2\) Dolbeault forms and discrete-series representations in a degree determined by compact and noncompact roots. The precise comparison is with the Griffiths--Schmid vanishing theorem and Schmid\textquotesingle s realization of sufficiently regular discrete series. The 1968 Schmid announcement concerns related homogeneous-cohomology work; it is not used as proof that every part of the 1969 notebook theorem had already been printed. The broader Langlands realization was settled later, and the second conjecture is named without a separate formula in the notes.

\begin{GSitem}{1968-172.1}{E}{An equal-rank homogeneous complex manifold}
\GSfield{Kind; source}{construction; \GSpage{P000281}}
\GSfield{The notes}{\(G\supset K\supset H\), \(\operatorname{rank}G=\operatorname{rank}K\); \(D=G/H\) is complex. The example is \(GL(2,\mathbb R)\), although \(G\) is called semisimple.}
\GSfield{Statement}{Let \(G\) be a connected real semisimple group with finite centre, \(K\) a maximal compact subgroup and \(H\subset K\) a compact Cartan subgroup of \(G\). Equivalently the complex ranks of \(G\) and \(K\) agree. A choice of positive roots for \((\mathfrak g_{\mathbb C},\mathfrak h_{\mathbb C})\) determines an invariant integrable complex structure on \(D=G/H\), realized as an open flag-domain orbit. A valid rank-one example is \(SL(2,\mathbb R)/SO(2)\).}
\GSfield{Reconstruction}{\(GL(2,\mathbb R)\) has a nontrivial central Lie algebra and is not semisimple. It is retained as a faithful notebook error, not silently read as \(SL_2\).}
\GSfield{Fate}{The construction is established in \citet[\S\S1--2]{griffithsschmid1969locally}; that manuscript was received 3 April 1969, before the talk, assigned to Acta volume 123 (1969), and bears the print date 23 January 1970. The 1967 thesis and January 1968 announcement are distinct earlier stages \citep{schmid1968homogeneous}.}
\GSfield{Present status}{As of November 2026: established in the stated scope; this occurrence does not pose an unresolved problem.}
\GSfield{Assessment}{literal statement: false as written; publication match: exact; verification: read the result; development: annotated; confidence: high}
\end{GSitem}
\begin{GSitem}{1968-172.2}{E}{The Weyl-group count concerns invariant structures}
\GSfield{Kind; source}{claim; \GSpage{P000281}}
\GSfield{The notes}{The number of complex structures equals the number of Weyl-group elements.}
\GSfield{Statement}{With a fixed compact Cartan and fixed identification of tangent root spaces, invariant integrable complex structures on \(G/H\) correspond to positive systems of the complex root system, hence there are \(|W|\) such marked structures. Passing to equivalence under diffeomorphisms or dropping invariance changes the counting problem.}
\GSfield{Reconstruction}{The word invariant and the marking convention are necessary; this is not a count of all complex structures on the underlying smooth manifold.}
\GSfield{Fate}{The positive-root construction is explicit in \citet[\S1]{griffithsschmid1969locally}; the Weyl group acts simply transitively on positive systems.}
\GSfield{Present status}{As of November 2026: established in the stated scope; this occurrence does not pose an unresolved problem.}
\GSfield{Assessment}{literal statement: correct with missing hypotheses; publication match: exact; verification: read the result; development: annotated; confidence: high}
\end{GSitem}
\begin{GSitem}{1968-172.3}{E}{Homogeneous line bundles from Cartan characters}
\GSfield{Kind; source}{construction; \GSpage{P000281}}
\GSfield{The notes}{\(\mathcal L_\lambda\) has fibre \(\mathbb C\); a one-dimensional representation of \(H\) determines it, with a Frobenius-reciprocity aside.}
\GSfield{Statement}{An integral character \(\chi_\lambda:H\to U(1)\) gives \(\mathcal L_\lambda=G\times_H\mathbb C_{\chi_\lambda}\to G/H\), with the holomorphic structure fixed by the chosen positive system. Sections are equivariant functions, with \(f(gh)=\chi_\lambda(h)^{-1}f(g)\) for this associated-bundle convention. Induction and Frobenius reciprocity explain the representation-theoretic interpretation; arbitrary unintegrated Lie-algebra weights do not necessarily define a bundle on the given group.}
\GSfield{Reconstruction}{Integrality and the bundle convention are stated explicitly; signs of highest weights depend on this convention.}
\GSfield{Fate}{\citet[\S\S1,4]{griffithsschmid1969locally} treats the homogeneous bundle and its invariant Hermitian geometry.}
\GSfield{Present status}{As of November 2026: established in the stated scope; this occurrence does not pose an unresolved problem.}
\GSfield{Assessment}{literal statement: correct with missing hypotheses; publication match: exact; verification: read the result; development: annotated; confidence: high}
\end{GSitem}
\begin{GSitem}{1968-172.4}{E}{Reduced L2 Dolbeault cohomology as harmonic forms}
\GSfield{Kind; source}{definition; \GSpage{P000281}}
\GSfield{The notes}{Compactly supported \((0,i)\)-forms, \(\bar\partial\), its adjoint and \(\Box\) are introduced; \(\ker\Box=\mathcal H^i(\mathcal L_\lambda)\) is called i-dimensional cohomology.}
\GSfield{Statement}{Choose invariant Hermitian metrics on \(G/H\) and \(\mathcal L_\lambda\). The complete homogeneous metric gives the closed \(L^2\) Dolbeault complex. In degree \(i\), reduced cohomology \(\ker\bar\partial_i/\overline{\operatorname{im}\bar\partial_{i-1}}\) is canonically the harmonic space \(\ker\Box_i\), where \(\Box=\bar\partial\bar\partial^*+\bar\partial^*\bar\partial\) has its natural self-adjoint realization. The superscript is cohomological degree, not vector-space dimension.}
\GSfield{Reconstruction}{The Hilbert completion is taken after the compactly supported test space; harmonic forms themselves are not required to have compact support. Reduced cohomology avoids silently assuming closed range.}
\GSfield{Fate}{\citet[pp.~293--294, before Theorem~7.8]{griffithsschmid1969locally} defines the Hilbert space of harmonic forms and uses completeness. In the noncompact case a nonzero discrete-series realization is infinite-dimensional.}
\GSfield{Present status}{As of November 2026: established in the stated scope; this occurrence does not pose an unresolved problem.}
\GSfield{Assessment}{literal statement: correct with missing hypotheses; publication match: exact; verification: read the result; development: annotated; confidence: high}
\end{GSitem}
\begin{GSitem}{1968-172.5}{E}{The compact specialization is Borel--Weil--Bott}
\GSfield{Kind; source}{theorem; \GSpage{P000281}}
\GSfield{The notes}{For compact \(G\), this is called the Borel--Weil theorem; Demazure is suggested as a reference.}
\GSfield{Statement}{For compact connected semisimple \(G\), a homogeneous holomorphic line bundle on its complete flag manifold has zero cohomology in every degree when the shifted weight \(\lambda+\rho\) is singular. For a regular shifted weight, exactly one cohomological degree can be nonzero, and it realizes the finite-dimensional irreducible representation determined by the shifted Weyl action. Dominant weights in the degree-zero convention give Borel--Weil.}
\GSfield{Reconstruction}{The compact higher-degree assertion is distinguished from the degree-zero Borel--Weil theorem; the sign convention is fixed by the bundle convention of the source.}
\GSfield{Fate}{The full higher-cohomology statement is the pre-existing Borel--Weil--Bott theorem; \citet[Theorem~6.1]{griffithsschmid1969locally} recalls it. The notebook\textquotesingle s unspecific Demazure pointer is not converted into an invented bibliographic entry.}
\GSfield{Present status}{As of November 2026: established in the stated scope; this occurrence does not pose an unresolved problem.}
\GSfield{Assessment}{literal statement: correct with missing hypotheses; publication match: exact; verification: read the result; development: annotated; confidence: high}
\end{GSitem}
\begin{GSitem}{1968-172.6}{B}{Vanishing away from one degree for sufficiently regular weights}
\GSfield{Kind; source}{theorem; \GSpage{P000281}}
\GSfield{The notes}{Theorem 1 requires \(|(\lambda,\alpha)|>b\) for every simple root, and says \(\mathcal H^i=0\) except at \(k(\lambda+\rho)\).}
\GSfield{Statement}{For the equal-rank flag domain with a fixed invariant metric, there is a constant \(b\) such that an integral weight satisfying \(|\langle\lambda,\alpha\rangle|>b\) for every root \(\alpha\) has harmonic \(L^2\) Dolbeault forms only in the degree \(q(\lambda+\rho)\). Enlarge \(b\) if necessary so that \(\lambda\) and \(\lambda+\rho\) lie in the same chamber. Conditions only on a fixed simple system with arbitrary signs do not ensure this distance from every wall.}
\GSfield{Reconstruction}{Every-root regularity replaces the insufficient literal simple-root condition. Opposite large simple-root pairings can cancel on a nonsimple positive root.}
\GSfield{Fate}{The exact all-chamber vanishing result is \citet[Theorem~7.8]{griffithsschmid1969locally}; the manuscript was received 3 April 1969 but printed 23 January 1970. Receipt alone is not evidence of public circulation. The January 1968 announcement \citep[p.~57]{schmid1968homogeneous} proves a restricted antecedent for sufficiently high powers of a bundle in a fixed antidominant chamber, not this full formulation. B records the later exact publication by Griffiths and Schmid, distinct from the irreducibility theorem.}
\GSfield{Present status}{As of November 2026: established in the stated scope; this occurrence does not pose an unresolved problem.}
\GSfield{Assessment}{literal statement: correct with missing hypotheses; publication match: exact; verification: read the result; development: annotated; confidence: high}
\end{GSitem}
\begin{GSitem}{1968-172.7}{B}{The remaining harmonic space realizes a discrete series}
\GSfield{Kind; source}{theorem; \GSpage{P000281}}
\GSfield{The notes}{In the exceptional degree, Theorem 1 identifies \(\mathcal H^{k(\lambda+\rho)}(\mathcal L_\lambda)\) with \(\pi_{\lambda+\rho}\), as a partial answer to Langlands.}
\GSfield{Statement}{For sufficiently regular integral \(\lambda\) under the preceding hypotheses, the nonzero Hilbert representation of \(G\) on \(\mathcal H^{q(\lambda+\rho)}(\mathcal L_\lambda)\) is irreducible and is the discrete-series representation with Harish-Chandra parameter \(\lambda+\rho\), with the weight convention of the chosen holomorphic bundle.}
\GSfield{Reconstruction}{The theorem is the precise first conjecture component recoverable from the notes. The full all-weights formulation in the problem index is identified from the cited literature, not quoted as notebook text.}
\GSfield{Fate}{\citet[Introduction, p.~255]{griffithsschmid1969locally} announces Schmid\textquotesingle s proof after completion of that manuscript; the full paper is \citet{schmid1971langlands}. \citet[\S3, Theorem~2]{schiffmann1971borel} states exactly the sufficiently regular result. The unrestricted harmonic realization is proved in \citet{schmid1976l2}, as explicitly confirmed in the author\textquotesingle s 1978 ICM address, published in 1980, p.~204 \citep{schmid1980representations}; the 1971 paper alone is not credited with the unrestricted theorem.}
\GSfield{Present status}{As of November 2026: this sufficiently regular realization is established; the general Langlands harmonic realization was proved by Schmid in 1976. No remaining open subproblem is inferred from this occurrence.}
\GSfield{Assessment}{literal statement: correct with missing hypotheses; publication match: exact; verification: read the result; development: annotated; confidence: high; problem \hyperref[prob:GP-1968-41]{GP-1968-41}}
\end{GSitem}
\begin{GSitem}{1968-172.8}{E}{The cohomological degree counts signed positive roots}
\GSfield{Kind; source}{claim; \GSpage{P000281}}
\GSfield{The notes}{Theorem 2 writes \(K(\mu)\) but uses \(\lambda\) on the right, counting compact roots with negative pairing and noncompact roots with positive pairing.}
\GSfield{Statement}{For a regular parameter \(\mu\), with positive roots partitioned as \(\Delta_c^+\sqcup\Delta_n^+\), \[q(\mu)=\#\{\alpha\in\Delta_c^+:\langle\mu,\alpha\rangle<0\}+\#\{\beta\in\Delta_n^+:\langle\mu,\beta\rangle>0\}.\] In the preceding realization use \(\mu=\lambda+\rho\). Counting both signs of every root would give a different, generally constant count.}
\GSfield{Reconstruction}{The variable is made uniform and positive roots are specified. These changes are faithful-error emendations because the scan agrees with the transcription.}
\GSfield{Fate}{The root-count convention already occurs in \citet[\S2, eq.~(2.1), p.~415]{griffiths1966locally}, followed by compact-quotient cohomology in \S3; it predates this talk. The precise later discrete-series convention is recorded in \citet[\S3, before Theorem~1]{schiffmann1971borel}; compare \citet[\S7]{griffithsschmid1969locally}.}
\GSfield{Present status}{As of November 2026: established in the stated scope; this occurrence does not pose an unresolved problem.}
\GSfield{Assessment}{literal statement: false as written; publication match: exact; verification: read the result; development: annotated; confidence: high}
\end{GSitem}
\begin{GSitem}{1968-172.9}{F}{The unspecified decay-at-infinity comparison}
\GSfield{Kind; source}{claim; \GSpage{P000281}}
\GSfield{The notes}{\ru{Если с убыванием на бесконечности --- будут те же когомологии}.}
\GSfield{Statement}{No exact comparison theorem is determined without a specified decay space, its topology, the closed extension of \(\bar\partial\), and the meaning of same cohomology. Harmonic \(L^2\), smooth Dolbeault, compactly supported and rapidly decreasing cohomology are different constructions and cannot simply be identified.}
\GSfield{Reconstruction}{The possible function-space interpretations are alternatives, not silently supplied hypotheses.}
\GSfield{Fate}{\citet[p.~293 and Introduction, p.~255]{griffithsschmid1969locally} explicitly treats comparison with ordinary cohomology as a separate issue; it does not validate every possible reading of this sentence. The notes leave the second Langlands conjecture similarly unnamed at formula level.}
\GSfield{Present status}{As of November 2026: the exact assertion in this sentence remains unidentified from the surviving wording; this is an evidential limit, not a claim that a definite mathematical problem is open.}
\GSfield{Assessment}{literal statement: not determinable; publication match: thematic only; verification: read the result; development: annotated; confidence: low}
\end{GSitem}

\GStalk{173}{1969-04-21}{V. M. Buchstaber}{P000280-P000284}{From finite h-cobordism classes to finite diffeomorphism classes}
The actual continuation consists of three lines on P000282: finite h-cobordism classes of homotopy tori, the question of finitely many manifolds in each class, and Morse modifications. In the high-dimensional smooth setting the second step is stronger than finiteness: vanishing of \(\operatorname{Wh}(\mathbb Z^n)\) makes every h-cobordism a product. The catalogue range extends through P000284, but P000283 explicitly begins 15 September 1969 and P000284 includes 22 September; those later-year passages are inventoried as boundary material in the report, not attributed to Buchstaber on 21 April.

\begin{GSitem}{1968-173.1}{E}{Finitely many h-cobordism classes of homotopy tori}
\GSfield{Kind; source}{claim; \GSpage{P000282}}
\GSfield{The notes}{Manifolds of the homotopy type of a torus are divided into finitely many classes by h-cobordisms.}
\GSfield{Statement}{For each fixed \(n\geq5\), the closed connected smooth manifolds homotopy equivalent to \(T^n\) comprise finitely many smooth h-cobordism classes. The dimension, smooth category and unmarked equivalence relation are essential to this reading.}
\GSfield{Reconstruction}{The same high-dimensional, closed-manifold setting is carried over from the preceding lecture.}
\GSfield{Fate}{This follows from the PL classification of \citet[Theorem~A]{hsiangshaneson1969fake} together with the finite-smoothing deduction made explicit at occurrence 1968-171.1. It is not the PL theorem alone. The exact Moscow report is \citet[p.~224]{buchstabermishchenko1969tori}, from the following day, 22 April 1969.}
\GSfield{Present status}{As of November 2026: established in the stated scope; this occurrence does not pose an unresolved problem.}
\GSfield{Assessment}{literal statement: correct with missing hypotheses; publication match: exact; verification: read the result; development: annotated; confidence: high; problem \hyperref[prob:GP-1968-42]{GP-1968-42}}
\end{GSitem}
\begin{GSitem}{1968-173.2}{E}{How many diffeomorphism types occur in one h-cobordism class?}
\GSfield{Kind; source}{problem; \GSpage{P000282}}
\GSfield{The notes}{\ru{Задача: показ., что в \(\forall\) классе конечн. число}.}
\GSfield{Statement}{Given a closed smooth homotopy \(n\)-torus, \(n\geq5\), prove that its smooth h-cobordism class contains only finitely many diffeomorphism types. In fact it contains exactly one: every h-cobordism has torsion in \(\operatorname{Wh}(\mathbb Z^n)=0\), and hence is a product by the s-cobordism theorem.}
\GSfield{Reconstruction}{Up to diffeomorphism is restored from record 171. The conclusion concerns unmarked boundary manifolds, not equality of all markings or parametrizations.}
\GSfield{Fate}{The two ingredients already predate the question: \citet[p.~63]{basshellerswan1964whitehead} and \citet[\S11]{milnor1966torsion}. Thus the question is a proof step in an exposition, not a remaining open problem.}
\GSfield{Present status}{As of November 2026: resolved in the stated high-dimensional smooth setting, with exactly one diffeomorphism type per h-cobordism class.}
\GSfield{Assessment}{literal statement: correct with missing hypotheses; publication match: exact; verification: read the result; development: annotated; confidence: high; problem \hyperref[prob:GP-1968-42]{GP-1968-42}}
\end{GSitem}
\begin{GSitem}{1968-173.3}{E}{Morse modifications of cobordisms}
\GSfield{Kind; source}{method; \GSpage{P000282}}
\GSfield{The notes}{The simplest films are classified by Morse modifications.}
\GSfield{Statement}{A relative Morse function on a compact smooth cobordism gives a handle decomposition. Handle additions, slides and cancellation are the geometric operations used to simplify the relative chain complex and expose Whitehead torsion. The notebook does not give a new finite algorithm classifying arbitrary cobordisms.}
\GSfield{Reconstruction}{The brief method label is expanded only to the standard handle-theoretic interpretation needed for the preceding problem.}
\GSfield{Fate}{This is the established method in \citet[\S\S9,11]{milnor1966torsion}.}
\GSfield{Present status}{As of November 2026: established in the stated scope; this occurrence does not pose an unresolved problem.}
\GSfield{Assessment}{literal statement: correct with missing hypotheses; publication match: partial; verification: read the result; development: annotated; confidence: high}
\end{GSitem}

\section{Developments}\label{dev:1967}
\subsection{GP-1965-08: a rational coordinate field in the first Weyl division ring}
\label{dev:weyl-coordinate}

\paragraph{Question, provenance and scope.}
The definite question in occurrence \hyperref[item:1967-145.3]{1967-145.3},
repeating 1965-80.10, is to describe the automorphisms of the fraction division
ring of the first Weyl algebra. Work throughout over \(\mathbb C\), with
\[
 A_1=\mathbb C\langle p,q\rangle/(pq-qp-1),\qquad
 D_1=\operatorname{Frac}A_1=\mathbb C(q)(p;\partial_q).
\]
All automorphisms fix \(\mathbb C\) pointwise. The notebook contrasts the
known polynomial automorphisms with the unknown division-ring case. It does
not assert that every division-ring automorphism preserves
\(K=\mathbb C(q)\), and it does not formulate the endomorphism version of the
Dixmier conjecture. Those are different questions.

The bounded problem addressed here is to describe
\[
 N_K=\{\phi\in\operatorname{Aut}_{\mathbb C}D_1:\phi(K)=K\}
 \quad\hbox{and}\quad
 U_q=\{\phi\in\operatorname{Aut}_{\mathbb C}D_1:\phi(q)=q\}.
\]
The equality \(\phi(K)=K\), rather than mere containment, matters. The
strongest result proved below is an exact description of both groups and of
their intersection with \(\operatorname{Aut}A_1\). This is a complete
verification of a known special case, not a new solution of the full
classification problem. The normalizer formula is explicitly recorded by
\citet[\S5.3.3, final Remark, printed p.~66]{dumas2011invariants}, who refers
to \citet{alevdumas1998invariants}. The formal-series embedding used below is
in \citet[\S5.2.2, pp.~59--60]{dumas2011invariants}. The proofs are included
to make the reduction independently checkable; the marker \ours\ identifies
work re-derived here, without a novelty claim.

\paragraph{What the literature does and does not supply.}
\citet[Theorem~8.10]{dixmier1968weyl} describes the polynomial automorphism
group of \(A_1\). Passing to fractions extends every such automorphism, but
also admits rational shears and inner automorphisms that need not preserve
the polynomial algebra. Dumas separates actions extending from the algebra
from general actions on its fraction field; his printed p.~66 gives the
exact coordinate-field stabilizer while explicitly leaving the full group
undetermined. The current-status audit accompanying occurrence 1967-145.3
does not identify a later complete classification. The obstruction below
explains the logical gap in this particular argument; it is not offered as
a historical explanation of why other mathematicians have not solved the
problem.

\begin{lemma}[The coordinate centralizer]\ours\ \label{lem:coordinate-centralizer}
The centralizer of \(q\) in \(D_1\) is exactly \(K=\mathbb C(q)\).
\end{lemma}
\begin{proof}
Here is a formal construction that also justifies the use of negative
powers of \(p\). Let \(\mathcal S\) consist of series
\(T=\sum_{j\leq m}a_j(q)p^j\), with \(m\in\mathbb Z\) and
\(a_j\in K\). Exponents are bounded above and can be unbounded below.
Define multiplication by the generalized Leibniz rule
\[
 (ap^i)(bp^j)=\sum_{k\geq0}\binom{i}{k}a\,\partial_q^k(b)\,p^{i+j-k},
 \qquad i,j\in\mathbb Z.
\]
For a fixed output exponent only finitely many terms contribute: if the
input orders are at most \(m,n\), a contribution to exponent \(\ell\)
requires \(i+j\geq\ell\), hence \(i\geq\ell-n\) and
\(j\geq\ell-m\). Thus there is no analytic convergence assertion.

For completeness, associativity follows coefficient by coefficient from
Leibniz' rule and the identity
\[
 \sum_{t=0}^{v}\binom jt\binom i{u+v-t}\binom{u+v-t}u
 =\binom iu\binom{i+j-u}v
 \qquad(u,v\geq0).
\]
Indeed, both sides are the coefficient of
\(a\,\partial_q^u(b)\,\partial_q^v(c)\) in the two products of
\(ap^i,bp^j,cp^k\). The identity is the generalized Vandermonde identity
after factoring out \(\binom iu\), and holds for negative as well as
nonnegative integers \(i,j\). Finite coefficient sums justify its use for
series.

Every nonzero series is invertible. If its leading term is \(a_mp^m\),
factor it as \(a_mp^m(1+R)\), where \(R\) has order at most \(-1\).
The series \(\sum_{k\geq0}(-R)^k\) is well defined coefficientwise and is
a two-sided inverse of \(1+R\). The monomial inverse is
\(p^{-m}a_m^{-1}\). Consequently \(\mathcal S\) is a division ring.
The rule \(pf=fp+f'\) embeds \(K[p;\partial_q]\) into \(\mathcal S\);
uniqueness of the polynomial expansion makes this embedding injective.
Its Ore fraction division ring therefore embeds as well.

Since \(q'=1\) and \(q^{(k)}=0\) for \(k\geq2\), the same rule gives
\[
 [p^j,q]=jp^{j-1}\quad\hbox{for every }j\in\mathbb Z,
 \qquad
 [T,q]=\sum_{j\leq m}j a_j(q)p^{j-1}.
\]
The exponents \(j-1\) are distinct. In characteristic zero the last
series vanishes only if \(a_j=0\) for every \(j\ne0\). Thus an element
of \(D_1\) commuting with \(q\) lies in \(K\). The reverse inclusion
is immediate.
\end{proof}

\begin{proposition}[All automorphisms fixing the coordinate]\ours\ 
\label{prop:rational-shear}
The map \(r\mapsto u_r\), where
\[
 u_r(q)=q,\qquad u_r(p)=p+r(q),\qquad r\in K,
\]
identifies the additive group of \(K\) with \(U_q\).
\end{proposition}
\begin{proof}
If \(\phi(q)=q\), then
\([\phi(p)-p,q]=\phi([p,q])-[p,q]=0\). The preceding lemma gives
\(\phi(p)=p+r(q)\). Conversely, replacing \(p\) by \(p+r\) and fixing
\(K\) preserves the relation \(pf-fp=f'\). It defines an automorphism of
\(K[p;\partial_q]\), with inverse obtained by replacing \(r\) by
\(-r\), and hence an automorphism of its fraction division ring.
Finally \(u_ru_s=u_{r+s}\); distinct functions have distinct images of
\(p\).
\end{proof}

\begin{lemma}[The automorphisms of the coefficient field]\ours\ 
Every \(\mathbb C\)-automorphism of \(K\) sends \(q\) to
\[
 \theta(q)=\frac{aq+b}{cq+d},\qquad ad-bc\ne0.
\]
\end{lemma}
\begin{proof}
Let its image be a nonconstant rational function \(\theta\), with
rational inverse substitution \(\psi\). A nonconstant rational function
defines a morphism of \(\mathbb P^1\) of degree equal to the common degree
of a coprime pair of homogeneous defining polynomials. Degrees multiply
under composition: composing the two homogeneous pairs gives the product
degree and has no common projective zero, because neither original pair
has one. Since \(\psi\circ\theta\) is the identity, both degrees are one.
The degree-one expressions are precisely the displayed fractions with
nonzero determinant. Conversely such a fraction has a degree-one rational
inverse and defines a field automorphism.
\end{proof}

\begin{proposition}[Normalizer of a rational coordinate field]\ours\ 
\label{prop:coordinate-normalizer}
The elements of \(N_K\) are exactly
\[
 \phi_{\theta,r}(q)=\theta(q),\qquad
 \phi_{\theta,r}(p)=\frac{1}{\theta'(q)}p+r(q),
 \quad \theta(q)=\frac{aq+b}{cq+d},\quad ad-bc\ne0,\quad r\in K.
\]
Restriction gives a split exact sequence
\[
 1\longrightarrow (K,+)\longrightarrow N_K
 \longrightarrow\operatorname{Aut}_{\mathbb C}K\longrightarrow1.
\]
With \(\sigma_\theta=\phi_{\theta,0}\), its action on rational shears is
\(\sigma_\theta u_r\sigma_\theta^{-1}=u_{\theta'\,(r\circ\theta)}\).
\end{proposition}
\begin{proof}
Restriction to \(K\) is a field automorphism, so the preceding lemma
determines \(\phi(q)=\theta(q)\). Because \(\theta'\ne0\),
\([\theta'^{-1}p,\theta]=1\). Hence
\(\phi(p)-\theta'^{-1}p\) commutes with \(\theta(q)\). Since
\(\mathbb C(\theta(q))=K\), commuting with \(\theta(q)\) is equivalent
to commuting with \(q\), and Lemma~\ref{lem:coordinate-centralizer}
gives the required rational remainder \(r\).

Conversely, substitution on \(K\), together with
\(p\mapsto\theta'^{-1}p+r\), respects the Ore relation: for \(f\in K\),
\[
 [\theta'^{-1}p+r,f(\theta(q))]
   =\theta'^{-1}\partial_q(f(\theta(q)))=f'(\theta(q)).
\]
The image of a nonzero polynomial of order \(m\) has nonzero leading
coefficient \(f_m(\theta)\theta'^{-m}\), so the resulting homomorphism
of the Ore polynomial ring is injective and extends to its fractions.
It is onto: \(q=\theta^{-1}(\phi(q))\) lies in the generated division
ring, and then \(p=\theta'(q)(\phi(p)-r(q))\) does too.

Writing \(\psi=\theta^{-1}\), an explicit inverse is
\[
 \phi_{\theta,r}^{-1}(q)=\psi(q),\qquad
 \phi_{\theta,r}^{-1}(p)=\frac{1}{\psi'(q)}p
                       -\frac{r(\psi(q))}{\psi'(q)}.
\]
The order of the coefficient and \(p\) in these expressions is essential.
The chain rule checks both compositions. The same chain rule shows that
the coefficient-field automorphism lifts with \(r=0\) give a section.
Conjugating \(u_r\) by this section yields
\(p\mapsto p+\theta'(q)r(\theta(q))\), proving the last formula.
The exact sequence uses composition of field automorphisms; if rational
maps are written as functions, substitution reverses their displayed
composition order. This convention prevents an accidental sign or order
error in a semidirect-product formula.
\end{proof}

\begin{corollary}[The polynomial boundary]\ours\ \label{cor:polynomial-boundary}
An automorphism \(\phi_{\theta,r}\) in
Proposition~\ref{prop:coordinate-normalizer} restricts to an automorphism
of \(A_1\) if and only if
\[
 \theta(q)=aq+b\quad(a\ne0),\qquad r(q)\in\mathbb C[q].
\]
In particular, a shear fixing \(q\) preserves \(A_1\) exactly when its
shear function is polynomial.
\end{corollary}
\begin{proof}
The PBW expansion in \(A_1=\mathbb C[q][p;\partial_q]\), together with
the unique formal expansion above, gives
\(A_1\cap K=\mathbb C[q]\). If \(\phi\) and \(\phi^{-1}\) preserve
\(A_1\), both \(\theta(q)\) and \(\theta^{-1}(q)\) are polynomials.
Their degrees multiply to one, so \(\theta\) is affine. Its derivative
is the nonzero constant \(a\). Subtracting \(a^{-1}p\) from
\(\phi(p)\) then shows that \(r\in A_1\cap K=\mathbb C[q]\).
Conversely affine substitution and a polynomial shear have polynomial
inverse formulas as displayed above, so they preserve \(A_1\) in both
directions.
\end{proof}

\begin{example}[A denominator changes the ambient problem]\ours\ 
The rational shear \(q\mapsto q,\ p\mapsto p+1/q\) is an automorphism
of \(D_1\), with inverse \(p\mapsto p-1/q\), and does not preserve
\(A_1\). Its hypotheses are explicit: \(q\ne0\) as an element of the
division ring, so \(q^{-1}\) exists there; the PBW intersection above
shows that it is not an element of the polynomial algebra. Similarly
\(q\mapsto-1/q,\ p\mapsto q^2p\) is a normalizer automorphism.
These examples exhibit the precise limitation of transferring the
polynomial generation theorem to the division ring.
\end{example}

\begin{remark}[Independent computation]\oursComputed\ 
The accompanying \path{code/verify_weyl.py} performs exact symbolic
checks, with no floating-point tolerances. It tests the negative-power
commutator identity in the formal Leibniz algebra, verifies both orders of
the inverse formula for several affine and nonaffine M\"obius maps and
several polynomial and rational shears, and checks the commutator by
applying the resulting first-order differential operators to independent
rational test functions. The output gives the parameter families and
number of checks. These finite checks detect ordering and sign errors;
they are not a substitute for the coefficientwise and chain-rule proofs
above. A second independent algebraic review is recorded in
\path{audit/proof_review.md}; its findings are incorporated into this
version.
\end{remark}

\paragraph{The obstruction still present.}
For a general \(\phi\in\operatorname{Aut}D_1\), put
\(Q=\phi(q)\), \(P=\phi(p)\). Necessarily \([P,Q]=1\) and the
division subring generated by \(P,Q\) is \(D_1\). The image
\(\mathbb C(Q)\) is self-centralizing, because centralizers are transported
by automorphisms. Our proof classifies the possibilities only after
\(\mathbb C(Q)=\mathbb C(q)\) is known. In general there is no reason
for that equality. The full problem requires controlling the orbit of
rational coordinate subfields, not merely their stabilizer. Moreover,
finding two elements with commutator one alone does not certify that they
generate the division ring; that further assertion cannot be inserted
without proof.

\paragraph{A precise next bounded task.}
One further research run could fix an explicit pair of rational
expressions \(P,Q\), verify \([P,Q]=1\), supply inverse expressions for
\(p,q\) to certify generation, and then determine whether its coordinate
field can be carried to \(\mathbb C(q)\) by a stated finite composition
of rational shears, Fourier exchange and inner conjugations. A finite
degree or denominator bound would make such a computation reproducible.
Success for that bounded family would give a factorization certificate
for those examples. It would not prove that the chosen transformations
generate every automorphism, nor convert a finite search with no
factorization into a counterexample. This is the exact limit of the
development undertaken here.

\section{Problem index}\label{problems:1967}
Earlier-year occurrence identifiers are retained as text; current-year identifiers link to their entries.

\subsection{GP-1967-01: Classifying noncompact groups over Hausdorff bases}\label{prob:GP-1967-01}
Construct a category of principal bundles over Hausdorff bases and universal objects that agree with ordinary bundles over completely regular bases while retaining characteristic information about a noncompact linear Lie group beyond the ordinary homotopy type of \(BG\).
\par\noindent\textit{Occurrences:} \hyperref[item:1967-132.3]{1967-132.3}, \hyperref[item:1967-132.9]{1967-132.9}, \hyperref[item:1967-134.8]{1967-134.8}, \hyperref[item:1967-135.5]{1967-135.5}, \hyperref[item:1967-135.6]{1967-135.6}.
\par\noindent\textit{Status (2026-11):} As of November 2026, the specified linear-group programme is answered by \citet[\S\S1--2 and4]{gelfandfuchs1967topology}. The more general elementary-object formulation is treated by \citet[Theorems1--2]{gelfandfuchs1968classifying}, with stronger axioms than the notebook sketch.

\subsection{GP-1967-02: Common fixed and periodic points of commuting interval maps}\label{prob:GP-1967-02}
Do two commuting continuous maps \([0,1]\to[0,1]\) necessarily share a fixed point? The accompanying differentiable statement asks for a fixed point of one map that is periodic for the other.
\par\noindent\textit{Occurrences:} \hyperref[item:1967-133.3]{1967-133.3}, \hyperref[item:1967-133.4]{1967-133.4}.
\par\noindent\textit{Status (2026-11):} As of November 2026, the common-fixed-point assertion is false: the Boyce and Huneke counterexamples circulated in1967 and appeared in print in1969. A \(C^1\) map commuting with a continuous map has the fixed/periodic point described by Schwartz1965. The latter precise formulation was checked indirectly in \citet[Introduction]{alikhanikoopaei1998commuting}. This index entry does not add the different all-continuous common-periodic-point conjecture, which the notebook does not ask.

\subsection{GP-1967-03: Metric classification of hyperbolic two-torus automorphisms}\label{prob:GP-1967-03}
When are two hyperbolic automorphisms of \(\mathbb T^2\), with Haar probability, isomorphic by an invertible measure-preserving conjugacy modulo zero?
\par\noindent\textit{Occurrences:} \hyperref[item:1967-138.2]{1967-138.2}, \hyperref[item:1967-138.10]{1967-138.10}, \hyperref[item:1967-140.17]{1967-140.17}, \hyperref[item:1967-140.18]{1967-140.18}.
\par\noindent\textit{Status (2026-11):} As of November 2026, equality of entropy is necessary and sufficient by \citet[Main Theorem]{adlerweiss1967entropy}. The paper was communicated24April1967 and publishedJune1967, before both notebook occurrences. The four-block recoding and geometric component refinement in these notes match its proof exactly.

\subsection{GP-1966-01: Algebraic classification of Anosov diffeomorphisms}\label{prob:GP-1966-01}
Is every Anosov diffeomorphism of a compact connected manifold topologically conjugate to an affine Anosov automorphism of an infranilmanifold?
\par\noindent\textit{Occurrences:} 1966-94.4, \hyperref[item:1967-140.6]{1967-140.6}.
\par\noindent\textit{Status (2026-11):} As of November 2026, the general conjecture remains open. The known infranilmanifold case is the Franks--Manning theorem; the two-torus case is proved in \citet[Theorem5.2]{sinai1968markov}. See the precise conjecture and known case in \citet[Conjecture2.3, Theorem2.4]{spatzier2016rigidity} and its still-conjectural status in \citet[\S1, Background]{damjanovic2026symmetries}.

\subsection{GP-1967-21: The unextended semisimple Gelfand--Kirillov conjecture}\label{prob:GP-1967-21}
For a complex semisimple Lie algebra \(\mathfrak g\), is \(\operatorname{Frac}U(\mathfrak g)\) isomorphic to \(D_{|\Phi^+|,\operatorname{rank}\mathfrak g}\) without extending the centre?
\par\noindent\textit{Occurrences:} \hyperref[item:1967-144.7]{1967-144.7}.
\par\noindent\textit{Status (2026-11):} November 2026: false in general by Premet (2009 preprint; 2010 publication), true for type A. The cited 2013 source leaves C and G2 open; no later exact resolution was verified. The 1969 theorem after central extension is a different assertion.

\subsection{GP-1965-08: Automorphisms of Weyl division rings}\label{prob:GP-1965-08}
Describe the automorphism groups of Weyl division rings; in the present occurrence, determine \(\operatorname{Aut}_{\mathbb C}\operatorname{Frac}A_1(\mathbb C)\), keeping polynomial algebra automorphisms and rational-subfield stabilizers distinct.
\par\noindent\textit{Occurrences:} 1965-80.10, \hyperref[item:1967-145.3]{1967-145.3}.
\par\noindent\textit{Status (2026-11):} November 2026: no complete description was found in this audit; Dumas explicitly records the full automorphism-group question as open. The known stabilizer of \(\mathbb C(q)\) consists of lifted fractional-linear maps and rational shears. Lauritzen--Thomsen's October/November 2026 theorem on equality of left/right indices is a different result and does not settle the automorphism classification.

\subsection{GP-1968-01: A universal family for the elliptic moduli problem}\label{prob:GP-1968-01}
Does an ordinary complex analytic parameter space represent all pointed elliptic families up to isomorphism, and what construction retains a universal family when this fails?
\par\noindent\textit{Occurrences:} \hyperref[item:1968-150.8]{1968-150.8}, \hyperref[item:1968-154.2]{1968-154.2}.
\par\noindent\textit{Status (2026-11):} November 2026: the unrigidified family functor is not represented by the coarse j-line. A full homology marking gives the upper-half-plane universal family; the quotient stack or Mumford moduli site retains automorphisms. These facts were already established before the seminar.

\subsection{GP-1968-20: The pointwise hyperbolic circle remainder}\label{prob:GP-1968-20}
Fix a cofinite Fuchsian group \(\Gamma<\operatorname{PSL}_2(\mathbb R)\) and \(z,w\in\mathbb H^2\), and count group elements by \(N_\Gamma(R;z,w)=\#\{\gamma\in\Gamma:d(z,\gamma w)\le R\}\). Thus stabilizer multiplicities are included. Write this count as its full spectral main term (including eigenvalues below \(1/4\)) plus \(E(X;z,w)\), where \(X\asymp e^R\). Is \(E(X;z,w)=O_{\varepsilon,\Gamma,z,w}(X^{1/2+\varepsilon})\) for every \(\varepsilon>0\)? This is an explicitly weakened and qualified reconstruction of the notebook\textquoteright s square-root-remainder line, not its exact wording.
\par\noindent\textit{Occurrences:} \hyperref[item:1968-166.2]{1968-166.2}.
\par\noindent\textit{Status (2026-11):} As of November 2026, unresolved in this general pointwise form in the searched literature. The 2025 preprint states the conjecture and proves a logarithmic improvement on the classical two-thirds bound for specified Heegner points, not the general square-root exponent \citep[Introduction and Theorem~1.1]{chatzakosetal2025hyperbolic}. The unspecified area-plus-boundary expression in the notebook is not certified as a true asymptotic expansion. The April 2026 preprint \citet[\S1.1, Theorem~1.2]{biro2026localsquare} gives a conditional local-square-mean improvement for the modular group, not a general pointwise bound.

\subsection{GP-1968-40: Higher spectral asymptotics with a specified smoothing order}\label{prob:GP-1968-40}
For a positive scalar elliptic operator with smooth coefficients, determine the asymptotic coefficients and remainder of its diagonal Riesz means of order \(\alpha\). The occurrence asks first about \(\alpha=1\), then about formal expansions; these are distinguished from a pointwise unsmoothed expansion.
\par\noindent\textit{Occurrences:} \hyperref[item:1968-168.6]{1968-168.6}.
\par\noindent\textit{Status (2026-11):} As of November 2026: solved in the appropriate finite-order smoothed sense already by \citet[Theorem~4.5 and pp.~215--216]{hormander1968spectral}; the general unsmoothed full-expansion reading is false, as the sphere in 1968-168.5 shows. No earlier exact problem identifier is supplied in the previous index.

\subsection{GP-1968-41: The Langlands harmonic realization of the discrete series}\label{prob:GP-1968-41}
For a connected real semisimple group with finite centre and compact Cartan, determine the harmonic \(L^2\) Dolbeault cohomology of the homogeneous line bundle \(\mathcal L_\lambda\) on a fixed flag domain. The notes state the sufficiently regular case. The stronger literature formulation predicts vanishing for singular \(\lambda+\rho\), and otherwise a single nonzero degree \(q(\lambda+\rho)\) realizing the corresponding discrete series.
\par\noindent\textit{Occurrences:} \hyperref[item:1968-172.7]{1968-172.7}.
\par\noindent\textit{Status (2026-11):} The notebook\textquotesingle s sufficiently regular theorem is established by Schmid\textquotesingle s work published in 1971; the 1976 paper proves the unrestricted \(L^2\) realization, as confirmed by the author\textquotesingle s ICM account, p.~204 \citep{schmid1980representations}. The all-weights formulation is an explicitly identified literature extension of this occurrence, not an additional conjecture transcribed from the notebook. The publisher metadata for both Annals papers was checked, and the exact early theorem was read in the Griffiths--Schmid paper and the 1971 Bourbaki exposition. The 1976 conclusion was checked against the author\textquotesingle s ICM account, not inferred merely from the title.

\subsection{GP-1968-42: Finiteness for smooth homotopy tori}\label{prob:GP-1968-42}
For fixed \(n\geq5\), are there finitely many diffeomorphism classes of closed connected smooth manifolds homotopy equivalent to \(T^n\)? The continuation isolates the subproblem of finiteness within one smooth h-cobordism class.
\par\noindent\textit{Occurrences:} \hyperref[item:1968-171.1]{1968-171.1}, \hyperref[item:1968-171.15]{1968-171.15}, \hyperref[item:1968-173.1]{1968-173.1}, \hyperref[item:1968-173.2]{1968-173.2}.
\par\noindent\textit{Status (2026-11):} Established. The 1969 fake-torus classification yields finitely many possibilities; the h-cobordism subproblem has the stronger answer one diffeomorphism type, by \(\operatorname{Wh}(\mathbb Z^n)=0\) and s-cobordism. Hsiang--Wall subsequently proved topological uniqueness in the high-dimensional range. None of these statements asserts smooth uniqueness or settles an unstated dimension-four problem.

\section{Corrections to the transcription and translation}\label{corr:1967}
The evidence distinguishes a transcription correction from a mathematical emendation of wording faithfully preserved in the notebook.
\subsection*{P000208 (both; faithful error): Distribution-character definition}
\noindent\textit{Transcription reads.} \(\langle\chi,\varphi\rangle=\int\varphi(g)T(g)\,dg\)\par
\noindent\textit{Critical reading.} \(\langle\chi,\varphi\rangle=\operatorname{tr}\int\varphi(g)T(g)\,dg\)\par
\noindent\textit{Evidence.} The scan omits the outer trace; its next phrase says the operator has a trace. Trace-class hypotheses must also be stated.\par
\subsection*{P000208 (both; faithful error): SO(3) Haar density}
\noindent\textit{Transcription reads.} \(p=\sin\tau/\tau\)\par
\noindent\textit{Critical reading.} \(p=(\sin\tau/\tau)^2\)\par
\noindent\textit{Evidence.} The scan has no square. The exponential Jacobian and the following \(p^{-1/2}\) character formula require the square under the half-angle convention.\par
\subsection*{P000208-P000209 (both; faithful error): Spherical Fourier transform and beta}
\noindent\textit{Transcription reads.} Ordinary surface area; transform \(\sin(\lambda\tau)/\tau\)\par
\noindent\textit{Critical reading.} \(dS/(4\pi\lambda)\) for the written transform\par
\noindent\textit{Evidence.} Direct integration with Fourier kernel exp(i<F,x>) gives 4 pi lambda sin(lambda tau)/tau for ordinary area.\par
\subsection*{P000211 (both; faithful error): Ray-space total space}
\noindent\textit{Transcription reads.} \(E=\mathbb R^\infty\)\par
\noindent\textit{Critical reading.} \(E=\mathbb R^\infty\setminus\{0\}\)\par
\noindent\textit{Evidence.} The action is not free at zero; Gelfand--Fuchs1967 explicitly removes zero.\par
\subsection*{P000211 (both; faithful error): Recovery from associated vector bundle}
\noindent\textit{Transcription reads.} \(E\) is recovered from \(E^\#\)\par
\noindent\textit{Critical reading.} \(E\) is recovered from \(E^\#\) together with its \(G\)-reduction; for \(G=GL(n)\) the full frame bundle suffices\par
\noindent\textit{Evidence.} For a proper subgroup the associated vector bundle forgets the reduction; the unqualified handwritten sentence is too strong.\par
\subsection*{P000211 (both; faithful error): Weighted cross-section}
\noindent\textit{Transcription reads.} \(\sum2^{-n}|x_n|/(1+|x_n|)=1/2\) gives a selector\par
\noindent\textit{Critical reading.} The displayed level does not meet all rays\par
\noindent\textit{Evidence.} A ray supported at coordinate \(m\ge2\) has supremum \(2^{-m}<1/2\). The formula is visible in the scan.\par
\subsection*{P000212 (en; translation): Commuting-map periodic-point remark}
\noindent\textit{Transcription reads.} of finite multiplicity\par
\noindent\textit{Critical reading.} of finite period\par
\noindent\textit{Evidence.} The context is iteration of commuting maps and Schwartz1965; multiplicity of a zero is not meant.\par
\subsection*{P000212 (both; faithful error): Three-denominator identity}
\noindent\textit{Transcription reads.} \(1/(abc)=\int_\Delta(\cdots)^{-3}\,dt_1dt_2\)\par
\noindent\textit{Critical reading.} \(1/(abc)=2\int_\Delta(\cdots)^{-3}\,dt_1dt_2\)\par
\noindent\textit{Evidence.} The missing 2! is in the scan. At a=b=c=1 the written integral is 1/2.\par
\subsection*{P000213 (both; faithful error): Paracompact sheaf cohomology}
\noindent\textit{Transcription reads.} \(H^q(X,\mathbb R)=0\) for paracompact \(X\)\par
\noindent\textit{Critical reading.} \(H^q(X,\mathcal C_X)=0\) for \(q>0\)\par
\noindent\textit{Evidence.} The notebook uses the sheaf of continuous functions; degree zero is C(X,R), not zero.\par
\subsection*{P000214 (both; ambiguity): Second-order compact dual}
\noindent\textit{Transcription reads.} \(GL(2)\) example followed by as for \(S^2\)\par
\noindent\textit{Critical reading.} \(S^2\) is the compact dual for \(SL(2,\mathbb R)/SO(2)\); the intended group switch is unstated\par
\noindent\textit{Evidence.} The scan reads second order. \(U(2)/O(2)\) is three-dimensional; the published semisimple example is \(SL(2)\).\par
\subsection*{P000217 (both; faithful error): Finite-index classifying-space claim}
\noindent\textit{Transcription reads.} I finite gives the usual classifying space\par
\noindent\textit{Critical reading.} Finite I gives finite stages; universality requires sufficiently large or infinite models\par
\noindent\textit{Evidence.} Finite Stiefel total spaces are not generally contractible. This is a mathematical qualification of the source wording.\par
\subsection*{P000219 (both; faithful error): Lorentz functional}
\noindent\textit{Transcription reads.} psi monotone\par
\noindent\textit{Critical reading.} psi increasing and concave, psi(0)=0, psi(t)>0 for t>0\par
\noindent\textit{Evidence.} Monotonicity alone does not ensure a Banach Lorentz norm; the standard hypotheses appear in Semenov1967.\par
\subsection*{P000220 (both; faithful error): Mityagin theorem}
\noindent\textit{Transcription reads.} symmetric is necessary and sufficient without qualification\par
\noindent\textit{Critical reading.} Retain the minimal/maximal, order-continuity, Fatou, or full-symmetry hypotheses of the chosen theorem\par
\noindent\textit{Evidence.} Mityagin1965 Theorems1 and3 distinguish the required function-space categories.\par
\subsection*{P000221 (both; faithful error): Hausdorff--Young}
\noindent\textit{Transcription reads.} \((|c_k|^p)^{1/p}\le\|f\|_p\)\par
\noindent\textit{Critical reading.} \(\|c\|_{\ell^{p^{\prime}}}\le\|f\|_p\), \(1\le p\le2\)\par
\noindent\textit{Evidence.} The scan contains no sum and uses p rather than its conjugate; these are mathematical emendations, not claimed OCR mistakes.\par
\subsection*{P000221 (both; ambiguity): Improving target exponent}
\noindent\textit{Transcription reads.} \(L_{1/\alpha-\Delta}\)\par
\noindent\textit{Critical reading.} \(L_{1/(\alpha-\Delta)}\) is the coherent intended reading\par
\noindent\textit{Evidence.} The scan does not group the diagonal fraction unambiguously; the \(E_\Delta\) formula gives reciprocal exponent \(\alpha-\Delta\).\par
\subsection*{P000221 (both; faithful error): Semenov fundamental-function criterion}
\noindent\textit{Transcription reads.} \(1<\ell_E\le u_E<2\) is sufficient for arbitrary symmetric \(E\)\par
\noindent\textit{Critical reading.} The unrestricted theorem is false; use dilation indices or additional structural hypotheses\par
\noindent\textit{Evidence.} Sharpley1975 pp122-125 explicitly refutes the underlying Semenov1968 representation theorem and records its consequences.\par
\subsection*{P000221 (both; faithful error): Toral eigenvalue implication}
\noindent\textit{Transcription reads.} Irrational eigenvalues imply real eigenvalues\par
\noindent\textit{Critical reading.} For A in SL(2,Z), impose |tr A|>2 for the real hyperbolic case\par
\noindent\textit{Evidence.} Finite-order integral matrices can have nonrational nonreal eigenvalues; the scan faithfully records the compressed implication.\par
\subsection*{P000221 (both; faithful error): Weighted coefficient inequality name}
\noindent\textit{Transcription reads.} Poia / \ru{Пойя}\par
\noindent\textit{Critical reading.} Paley / \ru{Пале}\par
\noindent\textit{Evidence.} Semenov1967 explicitly identifies this exact \(L^p\) weighted rearranged-coefficient specialization as Paley; the scan itself reads \ru{Пойя}.\par
\subsection*{P000222 (both; faithful error): Homeomorphic conjugacy}
\noindent\textit{Transcription reads.} U is linear for any compact abelian groups\par
\noindent\textit{Critical reading.} For ergodic toral automorphisms U is affine; U(0)=0 removes its translation\par
\noindent\textit{Evidence.} Adler--Palais1965 theorem states the affine term and ergodicity hypothesis. The general compact-group parenthesis is underspecified.\par
\subsection*{P000222 (both; faithful error): Perron vector and Markov probabilities}
\noindent\textit{Transcription reads.} \(\lambda e=\Pi^*e;\ p_{ij}=e_j\pi_{ij}/(\lambda e_i)\)\par
\noindent\textit{Critical reading.} \(\Pi r=\lambda r;\ p_{ij}=\pi_{ij}r_j/(\lambda r_i);\ w_i=l_ir_i\)\par
\noindent\textit{Evidence.} The scan clearly has a star. Right row-transition normalization requires a right eigenvector and stationary weights require both Perron vectors.\par
\subsection*{P000223 (both; faithful error): Projective descent}
\noindent\textit{Transcription reads.} The form is on projective space because it is homogeneous\par
\noindent\textit{Critical reading.} A projective differential form must be invariant under scalars and horizontal\par
\noindent\textit{Evidence.} For forms of positive degree, zero total homogeneity by itself need not imply descent. Leray works with forms on projective space.\par
\subsection*{P000224 (both; faithful error): Exponential separation on a compact manifold}
\noindent\textit{Transcription reads.} Exponential growth of ambient distance for all forward or backward iterates\par
\noindent\textit{Critical reading.} Use tangent-vector hyperbolicity, or local/intrinsic leaf distances within the valid range\par
\noindent\textit{Evidence.} Ambient distance is bounded by the diameter; Sinai1968 section1 uses induced leaf metrics for exponential estimates.\par
\subsection*{P000224 (both; faithful error): Residue cycle}
\noindent\textit{Transcription reads.} \(\gamma^{q-m}=\bigcap_{j=1}^m V_j\)\par
\noindent\textit{Critical reading.} \(\gamma^{q-m}\) is a real cycle on the complete intersection\par
\noindent\textit{Evidence.} The scan shows an equality/cycle expression; Leray1967 section1 places the cycle on the intersection. Its dimension is a homological dimension, not the complex dimension of the entire variety.\par
\subsection*{P000225 (both; faithful error): Existence of hyperbolic Lie-group automorphisms}
\noindent\textit{Transcription reads.} Exist if and only if the Lie algebra is nilpotent\par
\noindent\textit{Critical reading.} Nilpotence is not sufficient; a compatible hyperbolic lattice automorphism is required\par
\noindent\textit{Evidence.} The standard three-dimensional Heisenberg lattice forces a central derivative eigenvalue det B=plus or minus1.\par
\subsection*{P000228 (both; faithful error): Four rectangle intersections}
\noindent\textit{Transcription reads.} \(V_1\cap TV_1;\ V_2\cap TV_2;\ V_1\cap TV_2;\ V_1\cap TV_2\)\par
\noindent\textit{Critical reading.} The indexed family \(V_i\cap TV_j\), \(i,j\in\{1,2\}\)\par
\noindent\textit{Evidence.} The scan duplicates the final pair; Adler--Weiss1967 section5 uses all four pairs. The overwritten second index prevents a unique ordering-level transcription correction.\par
\subsection*{P000228 (both; ambiguity): Second entropy equation}
\noindent\textit{Transcription reads.} \(h(T)=\lambda\)\par
\noindent\textit{Critical reading.} \(h(T)=\log\lambda\), if \(\lambda\) retains the preceding expanding-eigenvalue meaning\par
\noindent\textit{Evidence.} Both lines are visible in the scan; the earlier line and Sinai1968 Theorem5.3 identify the omitted logarithm.\par
\subsection*{P000230 (both; ambiguity): circle-covering pullback}
\noindent\textit{Transcription reads.} \(U\times_U U=U\)\par
\noindent\textit{Critical reading.} Retain the arrow ambiguity; for two degree-two covers the pullback has two circular components.\par
\noindent\textit{Evidence.} Scan shows repeated U and circular coverings but no unambiguous structural-map labels; direct pullback computation.\par
\subsection*{P000233 (both; faithful error): Galois central extension}
\noindent\textit{Transcription reads.} \(\widetilde Z\) is algebraic over \(Z\), with Weyl Galois group.\par
\noindent\textit{Critical reading.} The induced extension of fraction fields is Galois; the polynomial-ring extension may ramify.\par
\noindent\textit{Evidence.} Gelfand-Kirillov1969 §§7–8 and invariant-field interpretation.\par
\subsection*{P000233 (both; faithful error): semisimple Lie field}
\noindent\textit{Transcription reads.} For semisimple groups, \(D(\mathfrak g)\cong D_{n,k}\) is almost proved.\par
\noindent\textit{Critical reading.} Retain as historical conjectural optimism; do not turn it into an established theorem.\par
\noindent\textit{Evidence.} The next notebook theorem concerns \(\widetilde U\); Premet (2010) disproves the unextended blanket assertion.\par
\subsection*{P000234 (both; faithful error): Amitsur centralizer}
\noindent\textit{Transcription reads.} \(C(x)\) is commutative.\par
\noindent\textit{Critical reading.} \(C_{A_1}(x)\) is commutative for nonscalar \(x\in A_1\).\par
\noindent\textit{Evidence.} Scalar \(x\) has centralizer \(A_1\); Dixmier (1968), \S4.2.\par
\subsection*{P000234 (both; faithful error): Euler eigenvalue}
\noindent\textit{Transcription reads.} \([pq,p^a q^b]=(a-b)p^a q^b\).\par
\noindent\textit{Critical reading.} \([pq,p^a q^b]=(b-a)p^a q^b\), when \([p,q]=1\).\par
\noindent\textit{Evidence.} The scan has \(a-b\); direct computation gives \([pq,p]=-p\) and \([pq,q]=q\).\par
\subsection*{P000234 (both; faithful error): Poisson comparison}
\noindent\textit{Transcription reads.} No connection, even with grading\par
\noindent\textit{Critical reading.} There is an associated graded Poisson algebra, but arbitrary automorphisms need not preserve the chosen filtration.\par
\noindent\textit{Evidence.} PBW and the Bernstein filtration; retain the notebook claim as a qualified historical remark.\par
\subsection*{P000234 (both; faithful error): strictly semisimple normal form}
\noindent\textit{Transcription reads.} An automorphism carries \(x\) to \(pq\).\par
\noindent\textit{Critical reading.} An automorphism carries \(x\) to \(apq+b\), with \(a\ne0\); exclude scalar \(x\).\par
\noindent\textit{Evidence.} The scan has \(pq\) without affine parameters; Dixmier (1968), Theorem 9.2 and Corollary 9.3.\par
\subsection*{P000235 (both; ambiguity): Martineau indicator limit}
\noindent\textit{Transcription reads.} \(\lim_{r\to\infty}r^{-k}\log|f(rz)|\).\par
\noindent\textit{Critical reading.} Retain the \(\lim/\limsup\) ambiguity; the radial upper indicator uses \(\limsup\) and upper-semicontinuous regularization.\par
\noindent\textit{Evidence.} The scan has an unresolved stroke above \(\lim\); ordinary radial limits can fail.\par
\subsection*{P000235 (both; faithful error): bounded disc equation}
\noindent\textit{Transcription reads.} For bounded \(f\), the equation \(\partial_z u=f\) need not have a bounded solution.\par
\noindent\textit{Critical reading.} Without an extra condition, bounded data on the disc have a bounded distributional solution; likewise for \(\bar\partial u=f\).\par
\noindent\textit{Evidence.} The scan confirms the negative claim; the Cauchy--Green operator has \(L^\infty\) norm at most four.\par
\subsection*{P000235 (both; faithful error): pseudoconvexity index}
\noindent\textit{Transcription reads.} \(a_0\); \(n-1\) positive eigenvalues.\par
\noindent\textit{Critical reading.} Strict pseudoconvexity gives the positive branch of \(a_1\); weak pseudoconvexity requires nonnegative eigenvalues.\par
\noindent\textit{Evidence.} The scan has \(a_0\); the preceding count \(n-q\) makes its positive branch dimensionally impossible.\par
\subsection*{P000235 (both; faithful error): weighted growth equivalence}
\noindent\textit{Transcription reads.} \(L^2(\phi+c\log(1+|z|^2))\) is identified with \(|f|\le C e^\phi(1+|z|^2)^k\).\par
\noindent\textit{Critical reading.} Keep the weighted \(L^2\) and pointwise assertions separate until the weight convention and local regularization are specified.\par
\noindent\textit{Evidence.} The square in the weighted norm and mean-value estimates alter the exponent; no literal same-weight equivalence follows.\par
\subsection*{P000236 (both; faithful error): boundary identity}
\noindent\textit{Transcription reads.} Displayed weighted identity without explicit adjoint boundary condition\par
\noindent\textit{Critical reading.} Require the \(\bar\partial\)-Neumann boundary condition and defining-function normalization, or divide boundary measure by \(|\nabla\rho|\).\par
\noindent\textit{Evidence.} Hormander1965 Proposition2.1.2 and equation(2.1.14).\par
\subsection*{P000237 (both; faithful error): arithmetic rationality}
\noindent\textit{Transcription reads.} A series converging near zero over \(\mathbb C\) and at every \(p\)-adic place is rational.\par
\noindent\textit{Critical reading.} Local convergence does not imply rationality; restore arithmetic coefficient restrictions and global \(p\)-adic meromorphic continuation.\par
\noindent\textit{Evidence.} The scan says convergence near zero; \(e^t\) is a counterexample. Dwork (1960), \S4, gives the genuine rationality argument.\par
\subsection*{P000237 (both; faithful error): germs and jets}
\noindent\textit{Transcription reads.} Germs and jets are weakly homotopy equivalent\par
\noindent\textit{Critical reading.} Specify the open-manifold, open invariant-relation hypotheses and homotopy of formal data.\par
\noindent\textit{Evidence.} Gromov1969 Theorem4.1.1; the general unrestricted assertion fails.\par
\subsection*{P000238 (both; faithful error): fractional-linear action}
\noindent\textit{Transcription reads.} \((az+b)/(cz+b)\).\par
\noindent\textit{Critical reading.} \((az+b)/(cz+d)\).\par
\noindent\textit{Evidence.} The scan appears to repeat \(b\); P000250--P000251 and the determinant-one matrix fix \(d\).\par
\subsection*{P000240 (both; faithful error): generic automorphisms}
\noindent\textit{Transcription reads.} Most have one automorphism \(x\mapsto-x\).\par
\noindent\textit{Critical reading.} Most pointed complex elliptic curves have one nonidentity automorphism, hence an automorphism group of order two.\par
\noindent\textit{Evidence.} Scan confirms one; the identity is also an automorphism.\par
\subsection*{P000240 (both; faithful error): rigidification terminology}
\noindent\textit{Transcription reads.} polarized (with a chosen basis of homology)\par
\noindent\textit{Critical reading.} marked by an oriented integral homology basis, with a distinguished point in genus one\par
\noindent\textit{Evidence.} The parenthetical data describe a marking; principal polarization retains the involution.\par
\subsection*{P000241 (both; faithful error): base change of a cover}
\noindent\textit{Transcription reads.} \(X_i\times_XY\to X\) is a cover.\par
\noindent\textit{Critical reading.} \(X_i\times_XY\to Y\) is the base-changed cover.\par
\noindent\textit{Evidence.} The scan visibly has target \(X\); the pullback axiom and P000245 give target \(Y\).\par
\subsection*{P000242 (both; faithful error): group cohomology characters}
\noindent\textit{Transcription reads.} \(H^2(\Gamma,\mathbb Z)\) consists of all complex one-dimensional representations.\par
\noindent\textit{Critical reading.} Valid for finite \(\Gamma\), or the required vanishing of \(H^1(\Gamma,\mathbb C)\) and \(H^2(\Gamma,\mathbb C)\); not arbitrary groups.\par
\noindent\textit{Evidence.} For \(\Gamma=\mathbb Z\), \(H^2=0\) but \(\operatorname{Hom}(\Gamma,\mathbb C^*)=\mathbb C^*\).\par
\subsection*{P000242 (both; faithful error): union of group orbits}
\noindent\textit{Transcription reads.} direct sum\par
\noindent\textit{Critical reading.} direct product for an arbitrary possibly infinite disjoint union\par
\noindent\textit{Evidence.} Scan confirms direct sum; Mumford1965 p.40 prints a product, as required by the sheaf equalizer.\par
\subsection*{P000244 (both; faithful error): self-product and multiplicities}
\noindent\textit{Transcription reads.} If \(X_1=X_2\), their product is the curve itself.\par
\noindent\textit{Critical reading.} The product over the moduli object includes every isomorphism; its base is \(\operatorname{Isom}(X_1,X_2)\).\par
\noindent\textit{Evidence.} The same page subsequently records automorphism multiplicities; Mumford1965 §3 gives the representing Isom space.\par
\subsection*{P000245 (both; provenance): opening category correction}
\noindent\textit{Transcription reads.} Earlier lectures: analytic manifolds\par
\noindent\textit{Critical reading.} Analytic spaces, including singular or nonreduced fibre products where required.\par
\noindent\textit{Evidence.} Geronimus explicitly apologizes and makes this correction in the notebook.\par
\subsection*{P000246 (both; faithful error): comparison criterion degree}
\noindent\textit{Transcription reads.} The associated sheaf of \(U\mapsto H^i(f(U),F)\) vanishes.\par
\noindent\textit{Critical reading.} Require vanishing for every \(i>0\); identify higher direct images and fix the direction convention.\par
\noindent\textit{Evidence.} Mumford (1965), p.~80, states \(R^iy_*=0\) for \(i>0\). The notebook omits the degree range.\par
\subsection*{P000246 (both; faithful error): injective acyclicity}
\noindent\textit{Transcription reads.} For an injective sheaf cohomology is zero\par
\noindent\textit{Critical reading.} Positive-degree cohomology vanishes; \(H^0\) is the sections group.\par
\noindent\textit{Evidence.} The same page states \(H^0(U,F)=F(U)\); injectivity does not force sections to vanish.\par
\subsection*{P000247 (both; ambiguity): extension neighbourhood}
\noindent\textit{Transcription reads.} Extension to a tubular neighbourhood\par
\noindent\textit{Critical reading.} Local extension of the family morphism near each point, under the modular local-isomorphism hypothesis.\par
\noindent\textit{Evidence.} The sketch does not establish a global holomorphic tubular neighbourhood; Mumford1965 §3 gives the Isom-space formulation.\par
\subsection*{P000248 (both; faithful error): jet embedding}
\noindent\textit{Transcription reads.} \(j_\mu\) is an embedding.\par
\noindent\textit{Critical reading.} It embeds the sheaf of smooth functions into the sheaf of jet sections; finite-jet evaluation at a point is not injective.\par
\noindent\textit{Evidence.} The zeroth component of the prolonged section is the function itself; Sweeney1968 ChapterI §1.\par
\subsection*{P000249 (both; faithful error): Spencer one-form sign}
\noindent\textit{Transcription reads.} \(d(\phi_I^jdx_j)-\phi_{I+e_s}^jdx_j\wedge dx_s\).\par
\noindent\textit{Critical reading.} \(d(\phi_I^jdx_j)-\phi_{I+e_s}^jdx_s\wedge dx_j\).\par
\noindent\textit{Evidence.} The scan has minus with reversed wedge order; Spencer (1969), equation (1.3.1)\(^{\prime}\), and \(D^2=0\) fix the order.\par
\subsection*{P000250 (both; faithful error): family over the coarse line}
\noindent\textit{Transcription reads.} There is a family over \(S^{\prime}=\mathbb H/\Gamma\).\par
\noindent\textit{Critical reading.} No smooth pointed elliptic family over the entire coarse \(j\)-line has \(j\)-invariant equal to its parameter.\par
\noindent\textit{Evidence.} Local square/cube branching obstructs a lift of a simple parameter; Mumford1965 p.59 and Hain Exercise26. This concerns the intended one-class-per-parameter family.\par
\subsection*{P000250 (both; faithful error): number of fixed points}
\noindent\textit{Transcription reads.} There are three fixed points\par
\noindent\textit{Critical reading.} Two exceptional \(PSL_2(\mathbb Z)\)-orbits have three drawn boundary representatives; \(-I\) fixes every point for \(SL_2(\mathbb Z)\).\par
\noindent\textit{Evidence.} The two corners are equivalent. Distinguish the effective action from \(SL_2(\mathbb Z)\).\par
\subsection*{P000251 (both; ambiguity): cohomology degree table}
\noindent\textit{Transcription reads.} Zero for \(n\ne2\); \(\mathbb Z/12\) for \(n=2\).\par
\noindent\textit{Critical reading.} \(H^0=\mathbb Z\), \(H^{2r}=\mathbb Z/12\) for \(r\ge1\), and all odd groups vanish.\par
\noindent\textit{Evidence.} The equality/congruence sign is retraced. The computation from \(C_4*_{C_2}C_6\) determines the mathematical correction.\par
\subsection*{P000251 (both; faithful error): general-field assertion}
\noindent\textit{Transcription reads.} The same answer holds over every field of characteristic different from two.\par
\noindent\textit{Critical reading.} Specify the Picard group, not integral Betti cohomology. Mumford excludes characteristics two and three; Fulton--Olsson proves \(\operatorname{Pic}=\mathbb Z/12\) over every field.\par
\noindent\textit{Evidence.} Mumford1965 p.73; Fulton–Olsson2010 Theorem1.1. The invariant and historical hypothesis are distinct.\par
\subsection*{P000252 (both; ambiguity): fibration cohomology abbreviation}
\noindent\textit{Transcription reads.} Cohomology of the base and system are triv.\par
\noindent\textit{Critical reading.} Their identification is immediate for a genuine leaf fibration; do not infer vanishing from the abbreviation.\par
\noindent\textit{Evidence.} The scan preserves the abbreviation without a decisive predicate. A nonempty connected base has \(H^0(B,\mathbb Z)=\mathbb Z\).\par
\subsection*{P000253 (both; faithful error): Coordinate-ring paragraph}
\noindent\textit{Transcription reads.} Maximal ideals are geometric points.\par
\noindent\textit{Critical reading.} Maximal ideals are closed points; geometric points include a residue-field embedding into an algebraic closure.\par
\noindent\textit{Evidence.} The Russian and scan use the same shorthand. Over a finite field a degree-d closed point has d geometric points; compare the explicit example on P000268.\par
\subsection*{P000254 (both; faithful error): Two Taylor formulas}
\noindent\textit{Transcription reads.} First derivative \(\partial_x\) multiplies \(y-x\); second \(\partial_y\) multiplies \(x-y\).\par
\noindent\textit{Critical reading.} Interchange the derivatives and read both equalities modulo \(\mathcal I_\Delta^2\).\par
\noindent\textit{Evidence.} The scan shows these derivative variables clearly. Testing the first formula on \(\varphi(x,y)=y\) gives a contradiction, while the corrected formula agrees with Taylor expansion.\par
\subsection*{P000255 (both; faithful error): Function plus differential form}
\noindent\textit{Transcription reads.} \((1,m)\)-jet: scalar function plus differential form.\par
\noindent\textit{Critical reading.} For two points, a first jet \((1,1)\) is a scalar plus one covector; for \(m+1\) points a first normal jet has \(m\) covectors.\par
\noindent\textit{Evidence.} The scan confirms \((1,m)\); the notation on P000254 instead places the jet order second. An unrecorded alternating convention cannot be supplied.\par
\subsection*{P000256 (both; ambiguity): Euler equation}
\noindent\textit{Transcription reads.} \(Au=\lambda u=0\).\par
\noindent\textit{Critical reading.} \(Au-\lambda u=0\) in the unconstrained problem; add constraint multipliers in the constrained problem.\par
\noindent\textit{Evidence.} The middle stroke is retraced in the scan, so equality versus minus is palaeographically uncertain. The variational equation determines the intended mathematical reading.\par
\subsection*{P000256 (both; provenance): Record 158 attribution}
\noindent\textit{Transcription reads.} A separate Gelfand record is listed in the catalogue.\par
\noindent\textit{Critical reading.} Retain the record but mark its independent content and attribution unresolved.\par
\noindent\textit{Evidence.} The scan contains schedule notices, a Kirillov heading, and a Weinstein heading; no Gelfand heading or identifiable intervention is visible.\par
\subsection*{P000256 (both; provenance): Weinstein identification}
\noindent\textit{Transcription reads.} Catalogue: Alan Weinstein. Scan: Prof. Weinstein (USA).\par
\noindent\textit{Critical reading.} Probably Alexander Weinstein; retain the catalogue name and flag the discrepancy.\par
\noindent\textit{Evidence.} The subject and finite-constraint method match Alexander Weinstein\textquoteright s 1963 paper. The source supplies no first name; mathematical identification is strong but personal attribution is not conclusive.\par
\subsection*{P000259 (both; faithful error): Compact-orbit sentence}
\noindent\textit{Transcription reads.} For an arbitrary compact group all orbits are in general position.\par
\noindent\textit{Critical reading.} The rho-shifted orbits entering the irreducible character formula are regular; arbitrary coadjoint orbits need not be.\par
\noindent\textit{Evidence.} The zero orbit of a nonabelian compact group is singular. Kirillov1968 Theorem2 explicitly shifts the highest weight.\par
\subsection*{P000260 (both; faithful error): Sphere Fourier integral}
\noindent\textit{Transcription reads.} \(\int\cdots=\sin(rx)/x\).\par
\noindent\textit{Critical reading.} With \(B=d\sigma/(2\pi r)\), the integral is \(2\sin(rx)/x\).\par
\noindent\textit{Evidence.} The scan confirms the coefficient is omitted. At zero the integral must equal the symplectic area \(2r\); the notebook right-hand side tends to \(r\).\par
\subsection*{P000260 (both; faithful error): Two-by-two skew determinant}
\noindent\textit{Transcription reads.} \(\det\left(\begin{smallmatrix}0&a\\-a&0\end{smallmatrix}\right)=-a^2\).\par
\noindent\textit{Critical reading.} \(\det\left(\begin{smallmatrix}0&a\\-a&0\end{smallmatrix}\right)=a^2\).\par
\noindent\textit{Evidence.} The scan has a minus sign; expansion of the determinant gives \(0-(-a^2)=a^2\).\par
\subsection*{P000264 (both; faithful error): Definition of solution}
\noindent\textit{Transcription reads.} \(\operatorname{Im}\gamma\subset A\), \(\gamma\in\Gamma(\mathcal L)\).\par
\noindent\textit{Critical reading.} \(\operatorname{Im}(j^r\gamma)\subset A\subset J^r\mathcal L\).\par
\noindent\textit{Evidence.} The scan has no visible jet operator in this line, although the preceding map is explicitly written. Source and target types require the correction.\par
\subsection*{P000264 (both; faithful error): Open condition called embedding}
\noindent\textit{Transcription reads.} Add the open condition that the map is an embedding.\par
\noindent\textit{Critical reading.} Full rank is an open first-jet immersion condition; embedding also requires global injectivity and the correct topology.\par
\noindent\textit{Evidence.} The scan clearly says embedding. A first jet cannot distinguish an immersion with self-intersections from an embedding.\par
\subsection*{P000265 (both; faithful error): Immersion-to-embedding inference}
\noindent\textit{Transcription reads.} Deform an immersion to any metric and obtain an embedding.\par
\noindent\textit{Critical reading.} The immersion conclusion and the embedding-preserving approximation theorem require distinct hypotheses.\par
\noindent\textit{Evidence.} The source changes from immersion to embedding. The change is not a formal consequence of a path of metrics.\par
\subsection*{P000267-P000268 (both; provenance): Record order and boundary}
\noindent\textit{Transcription reads.} Catalogue record163 Kazhdan precedes record164 Oleinik.\par
\noindent\textit{Critical reading.} Source order:164 Oleinik on P267, then Gelfand intervention at the top of P268, then163 Kazhdan.\par
\noindent\textit{Evidence.} The scanned headings fix this order and the scope of the intervening question.\par
\subsection*{P000268 (both; ambiguity): Finite-field defining polynomial}
\noindent\textit{Transcription reads.} An overwritten fraction is offered as the equation defining the finite field.\par
\noindent\textit{Critical reading.} Reading unresolved. Use the root-set construction \(X^{p^n}-X=0\), or a chosen irreducible polynomial of degree \(n\), as mathematical reconstruction.\par
\noindent\textit{Evidence.} The scan does not resolve the exponent hierarchy or the heavy stroke. Neither proposed transcription is adopted as an exact formula.\par
\subsection*{P000269 (both; faithful error): Fourth Euler factor}
\noindent\textit{Transcription reads.} \((1-t^4)^{(p^4-p^2)/4}\) without reciprocal.\par
\noindent\textit{Critical reading.} In the zeta Euler product the factor is \((1-t^4)^{-(p^4-p^2)/4}\).\par
\noindent\textit{Evidence.} The scan lists the fourth expression separately without a reciprocal; the preceding factors and degree-four orbit count determine its role.\par
\subsection*{P000270 (both; faithful error): Rationality inference}
\noindent\textit{Transcription reads.} Complex convergence and p-adic meromorphy imply rationality.\par
\noindent\textit{Critical reading.} Require integral coefficients and compatible radii \(Rr>1\), or global p-adic meromorphy with positive complex convergence radius.\par
\noindent\textit{Evidence.} Dwork1960 Theorem2 states the radius inequality explicitly; local convergence alone is insufficient.\par
\subsection*{P000271 (both; faithful error): Algebraic-integrality claim}
\noindent\textit{Transcription reads.} All reciprocal roots in a rational factorization with integral Taylor coefficients are algebraic integers.\par
\noindent\textit{Critical reading.} Require the factorization to be reduced.\par
\noindent\textit{Evidence.} A common nonintegral factor can be multiplied into numerator and denominator without changing the power series.\par
\subsection*{P000271 (both; faithful error): Noncompact unit disc}
\noindent\textit{Transcription reads.} The unit disc in \(\mathbb Q_p\) is not compact.\par
\noindent\textit{Critical reading.} The unit ball in \(\mathbb Q_p\) is compact. The preceding field \(\Omega=\mathbb C_p\) has noncompact unit ball.\par
\noindent\textit{Evidence.} The scan clearly reads \(\mathbb Q_p\) and not compact; this is a notebook error or field-name slip, not a mistranslation.\par
\subsection*{P000272 (both; faithful error): Last line of the schoolchild explanation}
\noindent\textit{Transcription reads.} \ru{10-адические --- деления нет} (no division).\par
\noindent\textit{Critical reading.} \(\mathbb Z_{10}\) is not a field; division by units exists, but \(5\) is not invertible.\par
\noindent\textit{Evidence.} The scan faithfully has the abbreviated phrase. Reduction modulo five rules out \(5x=1\); reduction modulo ten characterizes units. This is a qualification of the mathematics, not a mistranscription.\par
\subsection*{P000272 (both; provenance): Record166/167 boundary}
\noindent\textit{Transcription reads.} Catalogue ranges extend record166 through P273.\par
\noindent\textit{Critical reading.} Margulis166 ends in the upper part of P272; Gelfand167 starts at the underlined Gelfand heading on P272.\par
\noindent\textit{Evidence.} The scan has the two headings on one page; P273 begins a later subject.\par
\subsection*{P000273 (both; provenance): Catalogue date for record 168}
\noindent\textit{Transcription reads.} The supplied list gives 31-March-69 but labels the normalized date as absent.\par
\noindent\textit{Critical reading.} 1969-03-31.\par
\noindent\textit{Evidence.} The scan heading reads 31/III, between 10/II-1969 and 14/IV-69; the human-readable catalogue date agrees.\par
\subsection*{P000273 (both; faithful error): Isolated leading Weyl coefficient for the sphere}
\noindent\textit{Transcription reads.} \((2\pi)^{-n}\lambda^{n/2}/V_n\).\par
\noindent\textit{Critical reading.} For standard Riemannian volume, the local coefficient is \[(2\pi)^{-n}\operatorname{vol}(B^n)\lambda^{n/2}.\] The global coefficient is this times \(\operatorname{vol}(S^n)\).\par
\noindent\textit{Evidence.} The scan has the displayed denominator. The symbol \(|\xi|^2\) gives a unit cotangent ball; \(V_n\), used immediately above as sphere volume, cannot supply its reciprocal. This is a mathematical emendation.\par
\subsection*{P000273 (both; faithful error): Leading coefficient of the normalized average}
\noindent\textit{Transcription reads.} \(\lambda^{-1}\int_0^\lambda e(x,x,\nu)d\nu=a_0\lambda^{n/m}+\cdots\).\par
\noindent\textit{Critical reading.} The leading coefficient is \(\frac{m}{m+n}a_0\lambda^{n/m}\).\par
\noindent\textit{Evidence.} The scan omits the factor. Integrating the leading power \(\nu^{n/m}\) gives the factor \((1+n/m)^{-1}\); direct symbolic integration independently checks it.\par
\subsection*{P000273 (both; faithful error): Off-diagonal shorthand}
\noindent\textit{Transcription reads.} \ru{Вне диагонали 0}.\par
\noindent\textit{Critical reading.} Zero leading normalized coefficient off the diagonal, rather than zero spectral kernel.\par
\noindent\textit{Evidence.} The scan supports the words. The local oscillatory formula in \citet[Theorem~5.1]{hormander1968spectral} gives the intended meaning; the literal unnormalized assertion is false.\par
\subsection*{P000273 (both; faithful error): Sphere multiplicity numerator}
\noindent\textit{Transcription reads.} \(C_{n+k}^{k}-C_{n+k}^{k-2}\).\par
\noindent\textit{Critical reading.} \(\binom{n+k}{k}-\binom{n+k-2}{k-2}\).\par
\noindent\textit{Evidence.} The scan visibly repeats \(n+k\), confirming a faithful error. Harmonic homogeneous polynomials are the quotient of all degree-\(k\) polynomials by \(|x|^2\) times degree-\(k-2\) polynomials; compare \citet[\S6]{hormander1968spectral}.\par
\subsection*{P000274 (both; provenance): Dikii bibliography}
\noindent\textit{Transcription reads.} Dikii, UMN, 1957, on traces.\par
\noindent\textit{Critical reading.} The identified review is Dikii, UMN 13(3), 1958, pp.~111--143.\par
\noindent\textit{Evidence.} The notebook visibly writes 1957. The Math-Net publisher record for \citet{dikii1958traces} gives 1958. The abbreviated Gelfand 1954 pointer is retained as unresolved rather than silently changed.\par
\subsection*{P000276 (both; faithful error): Au-Yeung hypothesis and conclusion}
\noindent\textit{Transcription reads.} Two quadratic forms not simultaneously zero; simultaneous reduction to a sum of squares.\par
\noindent\textit{Critical reading.} Specify complex Hermitian forms, nonzero test vectors and simultaneous congruence to diagonal forms with real coefficients.\par
\noindent\textit{Evidence.} The scan writes \(x^*Ax\) and \(x^*Bx\) but omits the field and the nonzero condition. A two-dimensional real counterexample is given in item 1968-170.1. A definite combination does not make both original forms positive.\par
\subsection*{P000276 (both; faithful error): Definition of point spectrum}
\noindent\textit{Transcription reads.} \(\exists x\;Ax=\lambda x\).\par
\noindent\textit{Critical reading.} \(\exists x\ne0\;Ax=\lambda x\).\par
\noindent\textit{Evidence.} The scan has no nonzero restriction. Without it the zero vector makes every scalar an eigenvalue.\par
\subsection*{P000276 (both; provenance): Source order of records 169 and 170}
\noindent\textit{Transcription reads.} The catalogue lists Davis (169) before Gelfand (170).\par
\noindent\textit{Critical reading.} Read record 170 before record 169 in the occurrence index.\par
\noindent\textit{Evidence.} The Gelfand heading and Au-Yeung paragraph precede the Davis heading on the same scan.\par
\subsection*{P000277 (both; faithful error): Height and norm}
\noindent\textit{Transcription reads.} If the shell stays below the displayed height, then a is the norm.\par
\noindent\textit{Critical reading.} Then \(\|A\|\leq a\); equality requires the least admissible bound.\par
\noindent\textit{Evidence.} The scan faithfully states equality. Substitution of \(A=0\) and \(a=1\) gives a counterexample to that literal implication.\par
\subsection*{P000277 (en; translation): Spherical segment in Theorem 2}
\noindent\textit{Transcription reads.} A segment (spherical).\par
\noindent\textit{Critical reading.} A spherical cap, namely the stereographic image of a generalized disk.\par
\noindent\textit{Evidence.} Russian \ru{сегмент (сферич.)} and the identified Theorem~7.2 of \citet{davis1968shell} determine the dimensional meaning; an arc is not intended.\par
\subsection*{P000278 (both; provenance): Attribution of torsion classification}
\noindent\textit{Transcription reads.} Smale: an h-cobordism is classified by a boundary and an element of Wh.\par
\noindent\textit{Critical reading.} The general fundamental-group statement is the s-cobordism theorem, relative to one boundary, in boundary dimension at least five.\par
\noindent\textit{Evidence.} The scan is faithful. \citet[\S11]{milnor1966torsion} distinguishes the general torsion classification from Smale\textquotesingle s simply connected theorem.\par
\subsection*{P000279 (both; faithful error): Bass vanishing assertion}
\noindent\textit{Transcription reads.} \(K_1(\mathbb Z+\cdots+\mathbb Z)=0\) for the group ring.\par
\noindent\textit{Critical reading.} \(\operatorname{Wh}(\mathbb Z^n)=0\), while \(K_1(\mathbb Z[\mathbb Z^n])\cong\{\pm1\}\times\mathbb Z^n\).\par
\noindent\textit{Evidence.} The scan clearly has \(K_1\) and \(=0\). The corollaries on p.~63 of \citet{basshellerswan1964whitehead} identify the surviving units; this is a mathematical emendation.\par
\subsection*{P000279 (both; faithful error): Meaning of infinitely many structures}
\noindent\textit{Transcription reads.} Infinitely many smooth structures on \(S^2\times S^4\).\par
\noindent\textit{Critical reading.} Infinitely many smooth manifolds of that homotopy type, distinguished by rational Pontryagin classes.\par
\noindent\textit{Evidence.} \citet[pp.~208--209]{novikov1966abelian} states both the fixed-topological-manifold finiteness and homotopy-type infinitude. The scan\textquotesingle s shorthand conflates them.\par
\subsection*{P000279 (both; ambiguity): Polynomial determinant in the group-ring context}
\noindent\textit{Transcription reads.} Invertibility iff \(\det=\pm1\).\par
\noindent\textit{Critical reading.} This is correct for ordinary integer polynomials; for Laurent polynomials the determinant may be any \(\pm t_1^{a_1}\cdots t_n^{a_n}\).\par
\noindent\textit{Evidence.} The scan says polynomials, so the isolated sentence is correct. The preceding group ring is Laurent; the correction records the necessary change when it is applied there.\par
\subsection*{P000280 (both; faithful error): Common Thom space}
\noindent\textit{Transcription reads.} Thom spaces are homotopy equivalent because their bases are homotopy equivalent.\par
\noindent\textit{Critical reading.} Include the stable pullback identification of normal bundles written on the preceding page.\par
\noindent\textit{Evidence.} A homotopy equivalence of bases alone does not identify arbitrary vector bundles or their Thom spaces. The context supplies normal data, but the isolated implication is invalid.\par
\subsection*{P000280 (both; faithful error): Hurewicz fibre}
\noindent\textit{Transcription reads.} \(H^{-1}(1)\), finite group for all manifolds.\par
\noindent\textit{Critical reading.} In the stated stable torus setting, a finite fibre, naturally a torsor for \(\ker H\).\par
\noindent\textit{Evidence.} The scan clearly calls the fibre a group. It excludes the zero element and is not closed under addition; the general all-manifolds assertion needs additional hypotheses.\par
\subsection*{P000280 (both; faithful error): Torus Thom space}
\noindent\textit{Transcription reads.} \(T_\nu(T^n)=\Sigma^{N-n}T^n\).\par
\noindent\textit{Critical reading.} \(\operatorname{Th}(\varepsilon^{N-n}_{T^n})=\Sigma^{N-n}(T^n)_+\).\par
\noindent\textit{Evidence.} The scan has no basepoint symbol. Without a disjoint basepoint the bottom sphere summand of the Thom space is lost.\par
\subsection*{P000281 (en; translation): Cohomological degree}
\noindent\textit{Transcription reads.} i-dimensional L2 cohomology.\par
\noindent\textit{Critical reading.} L2 cohomology in degree i.\par
\noindent\textit{Evidence.} Russian \ru{\(i\)-мерные} uses the older degree terminology here. The very next sentence distinguishes finite-dimensional compact cohomology from the noncompact case.\par
\subsection*{P000281 (both; faithful error): Degree formula}
\noindent\textit{Transcription reads.} \(K(\mu)\) with pairings against \(\lambda\); roots have no positivity restriction.\par
\noindent\textit{Critical reading.} Use \(q(\mu)\) and pair with \(\mu\), counting only positive compact and positive noncompact roots.\par
\noindent\textit{Evidence.} The scan confirms the inconsistent variable and omitted positive-system qualifier. \citet[\S3]{schiffmann1971borel} gives the intended formula.\par
\subsection*{P000281 (both; faithful error): Semisimple example}
\noindent\textit{Transcription reads.} \(GL(2,\mathbb R)\), followed by real semisimple group.\par
\noindent\textit{Critical reading.} A compatible example is \(SL(2,\mathbb R)\), with compact Cartan \(SO(2)\).\par
\noindent\textit{Evidence.} The scan clearly reads GL. Its scalar matrices give a nonzero centre of the Lie algebra, so the literal example is not semisimple.\par
\subsection*{P000281 (both; faithful error): Sufficient regularity}
\noindent\textit{Transcription reads.} \(|(\lambda,\alpha)|>b\) for every simple root.\par
\noindent\textit{Critical reading.} Require separation from every root hyperplane, or specify a chamber in which the simple-root conditions imply it.\par
\noindent\textit{Evidence.} \citet[Theorem~7.8]{griffithsschmid1969locally} uses all positive roots. In type \(A_2\), pairings \(M\) and \(-M\) with simple roots vanish on their sum.\par
\subsection*{P000282-P000284 (both; provenance): End boundary of record 173}
\noindent\textit{Transcription reads.} The supplied catalogue range for Buchstaber on 21 April 1969 is P000280--P000284.\par
\noindent\textit{Critical reading.} The Buchstaber continuation ends on P000282. P000283--P000284 begin the next academic year.\par
\noindent\textit{Evidence.} P000283 has the explicit heading 15.9.69 and Vainberg; P000284 continues those remarks and introduces 22.IX with Fichera. The catalogue range overlaps later records; assigning these pages to record 173 would be a provenance error.\par
